% !TEX program = xelatex
% 编译方式：xelatex power_flow_manual.tex（需两遍以生成目录）
% 修订版 1：依据 2026-07-22 代码库（tem_write）修订，章节源文件位于 chapters1/。
\documentclass[UTF8,fontset=fandol,a4paper,11pt]{ctexart}

\usepackage[margin=2.2cm]{geometry}
\usepackage{amsmath,amssymb}
\usepackage{booktabs,longtable,array}
\usepackage{xcolor}
\usepackage{float}
\usepackage[european,straightvoltages]{circuitikz}
\usepackage{hyperref}
\hypersetup{colorlinks=true,linkcolor=black,urlcolor=blue}

% ── 统一参数表环境 ────────────────────────────────────────────────
% 六列：字段名 | 符号 | 单位 | 典型范围 | 默认值 | 说明
\emergencystretch=3em
\newcolumntype{L}[1]{>{\raggedright\arraybackslash}p{#1}}
\newenvironment{paramtable}{%
  \begin{footnotesize}
  \begin{longtable}{@{}L{3.2cm}L{1.3cm}L{1.6cm}L{2.3cm}L{1.5cm}L{4.3cm}@{}}
  \toprule
  \textbf{字段名} & \textbf{符号} & \textbf{单位} & \textbf{典型范围} & \textbf{默认值} & \textbf{说明}\\
  \midrule
  \endfirsthead
  \multicolumn{6}{@{}l}{\footnotesize（续表）}\\
  \toprule
  \textbf{字段名} & \textbf{符号} & \textbf{单位} & \textbf{典型范围} & \textbf{默认值} & \textbf{说明}\\
  \midrule
  \endhead
  \midrule
  \multicolumn{6}{r@{}}{\footnotesize（接下页）}\\
  \endfoot
  \bottomrule
  \endlastfoot
}{%
  \end{longtable}
  \end{footnotesize}
}

% 字段名排版（下划线需转义后传入；\_ 处允许换行以防超长字段名溢出版心）
\makeatletter
\newcommand{\fld}[1]{\texttt{\fld@scan#1\fld@stop}}
\def\fld@scan#1{%
  \ifx#1\fld@stop
  \else
    \ifx#1\_\_\allowbreak\else#1\fi
    \expandafter\fld@scan
  \fi}
\makeatother
% 模块元信息行：#1=求解器/类名  #2=头文件路径
\newcommand{\compmeta}[2]{\noindent\textbf{代码结构体：}\texttt{#1}\hfill
  \textbf{定义位置：}\texttt{#2}\par\smallskip}

% 可断行的源码路径/函数名（/ 与 \_ 处允许换行）
\makeatletter
\newcommand{\srcpath}[1]{\texttt{\src@scan#1\src@stop}}
\def\src@scan#1{%
  \ifx#1\src@stop
  \else
    \ifx#1\_\_\allowbreak\else
    \ifx#1/#1\allowbreak\else
    #1\fi\fi
    \expandafter\src@scan
  \fi}
\makeatother

% ── 通用理论层环境（修订版 2 新增；公共样式手册自动引入，本手册为独立导言区需手动引入）──
% theory_environments.tex —— HySim-XJTU-HRPES 模块手册"通用理论层"共享环境
% 用途：为 docs/modules/ 各手册的 theory_*.tex 理论章节提供统一排版与标记机制。
% 引入方式（二选一，不得重复引入）：
%   (a) 使用公共样式的手册：hysim_manual.sty 已自动引入，无需任何操作；
%   (b) 自带导言区的手册（power_flow、optimal_power_flow）：在导言区写
%       % theory_environments.tex —— HySim-XJTU-HRPES 模块手册"通用理论层"共享环境
% 用途：为 docs/modules/ 各手册的 theory_*.tex 理论章节提供统一排版与标记机制。
% 引入方式（二选一，不得重复引入）：
%   (a) 使用公共样式的手册：hysim_manual.sty 已自动引入，无需任何操作；
%   (b) 自带导言区的手册（power_flow、optimal_power_flow）：在导言区写
%       % theory_environments.tex —— HySim-XJTU-HRPES 模块手册"通用理论层"共享环境
% 用途：为 docs/modules/ 各手册的 theory_*.tex 理论章节提供统一排版与标记机制。
% 引入方式（二选一，不得重复引入）：
%   (a) 使用公共样式的手册：hysim_manual.sty 已自动引入，无需任何操作；
%   (b) 自带导言区的手册（power_flow、optimal_power_flow）：在导言区写
%       \input{../../_manual_common/theory_environments.tex}
% 规范见 docs/modules/THEORY_WRITING_GUIDE.md 与 docs/modules/README.md。

\RequirePackage{amsthm}
\RequirePackage[most]{tcolorbox}

% ── 定理族（中文名，按 section 编号）────────────────────────────
\theoremstyle{plain}
\newtheorem{theorem}{定理}[section]
\newtheorem{lemma}[theorem]{引理}
\newtheorem{proposition}[theorem]{命题}
\newtheorem{corollary}[theorem]{推论}
\theoremstyle{definition}
\newtheorem{definition}[theorem]{定义}
\newtheorem{example}[theorem]{例}
\theoremstyle{remark}
\newtheorem{remark}[theorem]{注记}
\renewcommand{\proofname}{证明}

% ── 理论章声明框：每个 theory_*.tex 章首必须使用 ─────────────────
% 用法：\begin{theorynote} 本章为通用数学理论……\end{theorynote}
\newtcolorbox{theorynote}{%
  enhanced, breakable, colback=blue!4, colframe=blue!45!black,
  boxrule=0.6pt, arc=1.5mm, left=2mm, right=2mm,
  title={\textbf{通用理论章说明}},
  fonttitle=\small\bfseries}

% ── 未实现/理论扩展标记框：理论内容超出当前代码实现时必须就地标记 ──
% 用法：\begin{gapnote} 本节模型在代码中尚未实现……\end{gapnote}
\newtcolorbox{gapnote}{%
  enhanced, breakable, colback=orange!6, colframe=orange!70!black,
  boxrule=0.6pt, arc=1.5mm, left=2mm, right=2mm,
  title={\textbf{未实现扩展说明}},
  fonttitle=\small\bfseries}

% ── 模型边界说明：用于明确拒绝或适用范围，不表示待实现计划 ──────
\newtcolorbox{boundarynote}{%
  enhanced, breakable, colback=gray!5, colframe=gray!65!black,
  boxrule=0.6pt, arc=1.5mm, left=2mm, right=2mm,
  title={\textbf{模型边界说明}},
  fonttitle=\small\bfseries}

% ── 实现状态标签（"与实现的对应"表专用）─────────────────────────
% \implfull{文件:函数名}  已实现（附锚点，禁止行号）
% \implpart{差异说明}     部分实现（必须说明差异）
% \implnone               未实现
\newcommand{\implfull}[1]{\textbf{已实现}\quad\srcpath{#1}}
\newcommand{\implpart}[1]{\textbf{部分实现}\quad（#1）}
\newcommand{\implnone}{\textbf{未实现}}

% 规范见 docs/modules/THEORY_WRITING_GUIDE.md 与 docs/modules/README.md。

\RequirePackage{amsthm}
\RequirePackage[most]{tcolorbox}

% ── 定理族（中文名，按 section 编号）────────────────────────────
\theoremstyle{plain}
\newtheorem{theorem}{定理}[section]
\newtheorem{lemma}[theorem]{引理}
\newtheorem{proposition}[theorem]{命题}
\newtheorem{corollary}[theorem]{推论}
\theoremstyle{definition}
\newtheorem{definition}[theorem]{定义}
\newtheorem{example}[theorem]{例}
\theoremstyle{remark}
\newtheorem{remark}[theorem]{注记}
\renewcommand{\proofname}{证明}

% ── 理论章声明框：每个 theory_*.tex 章首必须使用 ─────────────────
% 用法：\begin{theorynote} 本章为通用数学理论……\end{theorynote}
\newtcolorbox{theorynote}{%
  enhanced, breakable, colback=blue!4, colframe=blue!45!black,
  boxrule=0.6pt, arc=1.5mm, left=2mm, right=2mm,
  title={\textbf{通用理论章说明}},
  fonttitle=\small\bfseries}

% ── 未实现/理论扩展标记框：理论内容超出当前代码实现时必须就地标记 ──
% 用法：\begin{gapnote} 本节模型在代码中尚未实现……\end{gapnote}
\newtcolorbox{gapnote}{%
  enhanced, breakable, colback=orange!6, colframe=orange!70!black,
  boxrule=0.6pt, arc=1.5mm, left=2mm, right=2mm,
  title={\textbf{未实现扩展说明}},
  fonttitle=\small\bfseries}

% ── 模型边界说明：用于明确拒绝或适用范围，不表示待实现计划 ──────
\newtcolorbox{boundarynote}{%
  enhanced, breakable, colback=gray!5, colframe=gray!65!black,
  boxrule=0.6pt, arc=1.5mm, left=2mm, right=2mm,
  title={\textbf{模型边界说明}},
  fonttitle=\small\bfseries}

% ── 实现状态标签（"与实现的对应"表专用）─────────────────────────
% \implfull{文件:函数名}  已实现（附锚点，禁止行号）
% \implpart{差异说明}     部分实现（必须说明差异）
% \implnone               未实现
\newcommand{\implfull}[1]{\textbf{已实现}\quad\srcpath{#1}}
\newcommand{\implpart}[1]{\textbf{部分实现}\quad（#1）}
\newcommand{\implnone}{\textbf{未实现}}

% 规范见 docs/modules/THEORY_WRITING_GUIDE.md 与 docs/modules/README.md。

\RequirePackage{amsthm}
\RequirePackage[most]{tcolorbox}

% ── 定理族（中文名，按 section 编号）────────────────────────────
\theoremstyle{plain}
\newtheorem{theorem}{定理}[section]
\newtheorem{lemma}[theorem]{引理}
\newtheorem{proposition}[theorem]{命题}
\newtheorem{corollary}[theorem]{推论}
\theoremstyle{definition}
\newtheorem{definition}[theorem]{定义}
\newtheorem{example}[theorem]{例}
\theoremstyle{remark}
\newtheorem{remark}[theorem]{注记}
\renewcommand{\proofname}{证明}

% ── 理论章声明框：每个 theory_*.tex 章首必须使用 ─────────────────
% 用法：\begin{theorynote} 本章为通用数学理论……\end{theorynote}
\newtcolorbox{theorynote}{%
  enhanced, breakable, colback=blue!4, colframe=blue!45!black,
  boxrule=0.6pt, arc=1.5mm, left=2mm, right=2mm,
  title={\textbf{通用理论章说明}},
  fonttitle=\small\bfseries}

% ── 未实现/理论扩展标记框：理论内容超出当前代码实现时必须就地标记 ──
% 用法：\begin{gapnote} 本节模型在代码中尚未实现……\end{gapnote}
\newtcolorbox{gapnote}{%
  enhanced, breakable, colback=orange!6, colframe=orange!70!black,
  boxrule=0.6pt, arc=1.5mm, left=2mm, right=2mm,
  title={\textbf{未实现扩展说明}},
  fonttitle=\small\bfseries}

% ── 模型边界说明：用于明确拒绝或适用范围，不表示待实现计划 ──────
\newtcolorbox{boundarynote}{%
  enhanced, breakable, colback=gray!5, colframe=gray!65!black,
  boxrule=0.6pt, arc=1.5mm, left=2mm, right=2mm,
  title={\textbf{模型边界说明}},
  fonttitle=\small\bfseries}

% ── 实现状态标签（"与实现的对应"表专用）─────────────────────────
% \implfull{文件:函数名}  已实现（附锚点，禁止行号）
% \implpart{差异说明}     部分实现（必须说明差异）
% \implnone               未实现
\newcommand{\implfull}[1]{\textbf{已实现}\quad\srcpath{#1}}
\newcommand{\implpart}[1]{\textbf{部分实现}\quad（#1）}
\newcommand{\implnone}{\textbf{未实现}}


\title{\textbf{交直流混合高弹性能源电力系统仿真平台}\\[0.5em]
       {\Large 潮流计算技术手册}\\[0.3em]
       {\large 数学模型、算法流程、参数选项与验证校核}\\[0.3em]
       {\large 修订版 1（依据 2026-07-22 代码库修订）}}
\author{HySim-XJTU-HRPES（西安交通大学 HRPES 团队）}
\date{\today}

\begin{document}

\maketitle
\begin{abstract}
\noindent 本手册依据 HySim-XJTU-HRPES 平台 C++ 源代码中的潮流计算模块
（\texttt{src/\allowbreak power\_flow/} 与 \texttt{include/\allowbreak hacdcpf/\allowbreak power\_flow/} 目录）
整理，覆盖统一 Newton 混合交直流潮流核心、经典交流求解器（FDPF、线性化直流潮流、
分布式松弛、自适应孤岛）、全局化与鲁棒求解器族（非线性缩放、非单调线搜索、
LM 信赖域、伪瞬态延拓、NCP 半光滑 Newton、HELM、同伦延拓、Newton--Krylov）、
直流网络与换流器稳态建模、配电网前推回代潮流、三相相域潮流与三相混合潮流、
电压稳定连续潮流（CPF），并给出全部选项字段的默认值、结果与诊断数据结构、
以及测试验证体系。文档风格与本平台《元件模型技术手册》一致，全部公式逐条注明
源码出处（文件:函数名锚点，禁止行号）。

\smallskip
\noindent\textbf{修订版 1 说明}：本版对照 2026-07-22 代码库修订，相对初版的主要变化为——
CPF 求解器重写为默认弧长延拓（可翻越鼻点）；FDPF 与残差求值器注入计算稀疏化并支持
逐母线 ZIP；\fld{PowerFlowOptions} 新增
\fld{enforce\_converter\_physical\_limits} 硬约束选项；\fld{SolverWorkspace} 落地并被
Newton 求解器复用；求解器工厂层经新增的 \srcpath{solver\_interface.cpp} 完整接线；
HELM/同伦/Newton--Krylov/CPF 新增专项回归测试。全部受影响行号引用已逐一重新核实。
\end{abstract}

\tableofcontents
\newpage

\section{阅读指南与约定}

\subsection{数据来源}
本手册全部内容取自以下源代码与文档，公式与参数均以代码为准：
\begin{itemize}
  \item \texttt{src/power\_flow/} —— 35 个实现文件，潮流求解器全部实现；
  \item \texttt{include/hacdcpf/power\_flow/} —— 求解器头文件（含 \texttt{solvers/}、
        \texttt{globalization/}、\texttt{assembly/} 三个子目录的规范头）；
  \item \texttt{include/hacdcpf/assembly/} 与 \texttt{src/power\_flow/solver\_data.cpp} ——
        SolverData 装配与索引映射；
  \item \texttt{src/api/hacdcpf.cpp} —— 公共 API 门面（求解分发、缓存、结果回投影）；
  \item \texttt{include/hacdcpf/engine/solver/native/nle/newton\_solver.hpp} ——
        Newton 求解器规范头（\texttt{power\_flow/} 下同名头为转发层）；
  \item \texttt{docs/archive/reference/technical\_notebook/sections/03\_power\_flow.tex}、
        \texttt{05\_advanced\_pf.tex}、\texttt{12\_pf\_options\_results.tex} ——
        既有技术笔记，本手册与之交叉核对；
  \item \texttt{docs/archive/theory/multiple\_converter.md}、\texttt{docs/archive/theory/multiple\_converter\_r1.md} ——
        多换流器建模设计文档；
  \item \texttt{tests/} 与 \texttt{tools/} —— 验证与校核章节的测试矩阵来源。
\end{itemize}
文中所有“\texttt{文件:函数名}”锚点按当前代码检出处核验，禁止行号（易漂移），
但函数名与文件路径保持稳定。

\subsection{单位制约定}
平台统一采用 MW·kV·MVA 单位制配合标幺值：系统基准功率
$S_{\mathrm{base}}=100$~MVA（\texttt{Defaults::kBaseMva}），基准频率 50~Hz。
潮流方程内部一律以标幺值求解；对外结果中支路功率以 MW/Mvar 给出，
母线电压幅值以 pu、相角以度（$^\circ$）给出。三相相域潮流采用每相基准
$S_{\mathrm{base},1\phi}=S_{\mathrm{base}}/3$。ZIP 负荷权重三元组顺序为
（恒功率、恒电流、恒阻抗），默认 $(1,0,0)$ 即纯恒功率。

\subsection{编号与符号约定}
\begin{itemize}
  \item \textbf{失配约定}：全平台统一 $\mathrm{mismatch}=F^{\mathrm{spec}}-F^{\mathrm{calc}}$
        （指定值减计算值），线性系统 $J\,\Delta x=\mathrm{mismatch}$ 中
        $J=\partial F^{\mathrm{calc}}/\partial x$；对 $F=0$ 的标准 Newton 与此等价，
        符号被吸进 $J$ 的装配。
  \item \textbf{注入约定}：节点注入以\textbf{发电为正、负荷为负}；换流器端口注入为正时
        满足 $P_{ac}+P_{dc}+P_{loss}=0$。
  \item \textbf{三类 ID 纪律}：组件 \texttt{.index}（对外稳定）、向量下标（内部 0-based）、
        图节点索引（临时）不混用；对外报告必须经映射回到稳定 ID。
  \item \textbf{AC/DC 域隔离}：AC 与 DC 母线编号空间独立，跨域传递一律走
        domain-qualified 映射（\fld{ac\_bus\_id\_to\_node\_idx} /
        \fld{dc\_bus\_id\_to\_node\_idx}）。
\end{itemize}

\subsection{求解器家族总览}
\begin{center}
\footnotesize
\begin{tabular}{@{}llll@{}}
\toprule
求解器 & 入口函数 & 适用网络 & 定位\\
\midrule
统一 Newton & \srcpath{solve\_power\_flow} & 混合 AC/DC & \textbf{生产主路径}（联立 NR）\\
FDPF & \srcpath{solve\_power\_flow\_fdpf} & 纯 AC & 快速解耦（FDXB 变体）\\
线性化直流潮流 & \srcpath{solve\_ac\_linearized\_dc} & 纯 AC & $B\theta=P$ 一次求解\\
DC 网络求解器 & \srcpath{solve\_dc\_power\_flow} & 纯 DC & $V\odot GV=P$ 非线性 Newton\\
自适应孤岛 & \srcpath{solve\_power\_flow\_adaptive} & 混合 AC/DC & 分岛独立求解+死岛置零\\
分布式松弛 & \srcpath{solve\_power\_flow\_distributed\_slack} & 混合 AC/DC & 不平衡按参与因子分摊\\
HELM & \srcpath{HelmSolver::solve} & 纯 AC & 全纯嵌入，无初值依赖\\
同伦延拓 & \srcpath{HomotopyContinuationSolver} & 混合 AC/DC & 收敛兜底（鲁棒管线 Phase 4）\\
Newton--Krylov & （内嵌于 Newton 回退） & 混合 AC/DC & 病态雅可比的 GMRES 内迭代\\
配网 BFS & \srcpath{analysis::solve\_distribution\_pf} & 辐射状 AC & 前推回代+调压器外环\\
三相 NR & \srcpath{analysis::solve\_three\_phase\_nr} & 三相 AC & abc 相域极坐标 NR\\
三相混合 & \srcpath{solve\_three\_phase\_hybrid\_pf} & 三相 AC+DC & 相域 AC 与 DC 联立\\
CPF & \srcpath{CpfSolver::solve} & 混合 AC/DC & 电压稳定/负荷裕度\\
\bottomrule
\end{tabular}
\end{center}

\subsection{验证通道总览}
各求解器节末的“验证与校核”小节给出具体对照依据，总体通道如下：
\begin{itemize}
  \item \textbf{MATPOWER 算例回归}：\srcpath{external\_data/matpower/} 85 个 .m 算例
        分三档（Tier1 九算例逐案断言、Tier2 约 60 算例批收敛率、Tier3 大系统至
        case\_SyntheticUSA 约 8.2 万节点），另有 \srcpath{tools/matpower\_benchmark.py}
        与 \srcpath{tools/matpower\_pf\_compare.cpp} 的逐母线电压对照；
  \item \textbf{OpenDSS}（dss\_capi 桥）：相域潮流 IEEE13/IEEE123/Step 系列教学算例、
        配网 BFS 与 case33bw 全三相对照、调压器闭环；由
        \srcpath{HACDCPF\_ENABLE\_OPENDSS(\_COMPARE)} 门控；
  \item \textbf{GridLAB-D}：平衡馈线快照对照（运行时探测可执行文件，缺失自动 skip）；
  \item \textbf{解析解与自建算例}：两节点解析解、雅可比中心差分校验
        （\srcpath{src/power\_flow/jacobian\_check.cpp}）、独立残差复算
        （\srcpath{src/power\_flow/residual\_evaluator.cpp}，稀疏复数形式）；
  \item \textbf{能力合同与 GUI 端到端}：\srcpath{tools/validate\_case\_capabilities.py}
        （NR/FDPF/DC 三法互验）与 \srcpath{tools/gui\_api\_e2e.py}。
\end{itemize}

\subsection{现状与诚实口径声明}
以下事项在各章节均有对应标注，此处集中声明（如实）：
\begin{enumerate}
  \item \textbf{HELM、同伦延拓、CPF、Newton--Krylov 当前无生产 API 自动挂接}：
        前三者由测试/GUI 直接驱动或作为选项层暴露的能力，Newton--Krylov 默认关闭、
        仅在条件数代理越限时可选回退；四者现均已有专项冒烟回归
        （\srcpath{test\_helm\_solver.cpp}、\srcpath{test\_homotopy\_continuation.cpp}、
        \srcpath{test\_voltage\_stability.cpp}、\srcpath{test\_newton\_krylov.cpp}），
        但大系统与生产路径回归仍缺，覆盖总体仍薄（见第~11~章）。
  \item \textbf{换流器物理不等式约束（容量圆、交流/直流电流限、调制比、占空比）
        未嵌入 Newton 方程}，默认仅做事后校核并按 \texttt{ACDC-PHYS-01..05}、
        \texttt{DCDC-PHYS-01} 告警；新增 opt-in 选项
        \fld{enforce\_converter\_physical\_limits} 可将越限 Newton 根判为物理不可行
        （见第~6、10~章）；结果中的 \fld{converter\_model\_scope} 如实声明
        模型覆盖边界。
  \item \textbf{FDPF 注入计算已稀疏化}：注入按稀疏复数形式
        $S_i=V_i\,\mathrm{conj}((Y_{bus}V)_i)$ 计算（$O(\mathrm{nnz})$），
        指定注入在组件聚合负荷下采用逐母线 ZIP 系数；初版所述稠密 $O(n^2)$
        双重循环已移除。
  \item \textbf{CPF 默认启用弧长延拓}（\fld{enable\_arc\_length}=true）：
        $\lambda$ 作为增广未知量参与校正，可翻越鼻点并按
        \fld{lower\_branch\_steps} 保留若干下支点；存在任一在运 PV 母线时
        自动回退自然参数化并以 \fld{model\_scope} 如实声明；
        \fld{enable\_arc\_length}=false
        时保留初版所述自然参数化模式。电压稳定指标仍仅负荷裕度族，无 $L$ 指标。
  \item \texttt{include/hacdcpf/power\_flow/workspace.hpp} 的
        \texttt{SolverWorkspace} 已落地（\srcpath{prepare\_state} /
        \srcpath{prepare\_equations} / \srcpath{reset\_iteration}），并被
        Newton 求解器按实例复用（带递归守卫）；符号模式缓存仍靠 PatternCache。
  \item 求解器抽象层已完整接线：\srcpath{src/power\_flow/solver\_interface.cpp}
        实现 \srcpath{build\_power\_flow\_problem} 与六个工厂求解器；
        但生产门面仍直连 thread\_local \srcpath{engine::NewtonSolver}，
        不经过工厂层。
  \item \textbf{LCC 换流站稳态模型已落地}（\srcpath{src/power\_flow/lcc\_model.cpp/.hpp}）：
        \fld{SolverData::lcc\_converters}（include/hacdcpf/assembly/solver\_data.hpp:SolverData）参与统一 Newton
        残差（src/power\_flow/jacobian\_builder.cpp:build\_power\_spec）与雅可比（src/power\_flow/jacobian\_builder.cpp:evaluate\_residual\_impl），
        特征控制站锚定 DC 岛电压（src/power\_flow/dc\_solver.cpp:add（lambda）、
        src/power\_flow/newton\_solver.cpp:plan\_dc\_island\_references），求解后经 API 结果字段
        \fld{lcc\_transfers} 输出完整运行点（src/api/hacdcpf.cpp:is\_promoted（lambda））；
        混合交直流建模现有 VSC、DC/DC 与 LCC 三条路径（见第~6~章）。
\end{enumerate}

\newpage
% 第 1 章：架构总览与数据流
% 本章文件由 \input 引入，不含导言区。
\section{架构总览与数据流}\label{sec:overview}

\subsection{从工程模型到求解结果的流水线}
平台潮流计算遵循统一的五段式流水线（工程语义层 $\to$ 校验 $\to$ 投影 $\to$
装配求解 $\to$ 结果回投影），全部潮流入口共用同一条主链：

\begin{figure}[H]
\centering
\begin{tikzpicture}[font=\footnotesize,
  box/.style={draw,rounded corners,minimum height=1.0cm,align=center,inner sep=4pt},
  arr/.style={->,thick}]
  \node[box] (rich) at (0,0) {\texttt{HybridPowerSystem}\\工程语义层（rich）};
  \node[box] (valid) at (4.1,0) {四级校验\\Basic/Electrical/\\SolverReady/Strict};
  \node[box] (proj) at (8.2,0) {canonical 投影\\开关收缩·零阻抗合并\\设备聚合};
  \node[box] (data) at (12.3,0) {\texttt{SolverData} 装配\\$Y_{bus}$·$G_{dc}$·注入};
  \node[box] (solve) at (4.1,-2.3) {求解/分析\\Newton·FDPF·BFS·…};
  \node[box] (back) at (8.2,-2.3) {结果回投影\\\srcpath{unproject\_pf\_result}\\支路功率回算};
  \draw[arr] (rich) -- (valid);
  \draw[arr] (valid) -- (proj);
  \draw[arr] (proj) -- (data);
  \draw[arr] (data.south) |- (solve.west);
  \draw[arr] (solve) -- (back);
\end{tikzpicture}
\caption{潮流计算主流水线（映射 \fld{bus\_merge\_map} 随 SolverData 携带，
供结果归因回原始组件编号空间）}
\end{figure}

投影与回投影的不变量（\srcpath{docs/developer/projection\_and\_results.md} 与
\srcpath{docs/archive/reference/technical\_notebook/sections/02e\_projection\_result\_contracts.tex}，
代码实现 \srcpath{src/model/network\_utils.cpp}）：强度量（电压、相角）穿越合并
\textbf{广播}；广延量（功率、切负荷、能量）必须按参与语义\textbf{显式拆分}，
绝不盲广播；合并/死岛删除映射是复合而非覆盖；恢复等级
（strong/approximate/audit-only/unsupported）是结果契约的一部分。

\subsection{生产主调用链}
统一潮流主入口 \srcpath{hacdcpf::solve\_power\_flow}
（\srcpath{src/api/hacdcpf.cpp:solve\_power\_flow}）的执行顺序如下，每一步均注明源码位置：

\begin{enumerate}
  \item \textbf{参考母线资格校验}：\srcpath{validate\_reference\_bus\_eligibility}
        （src/api/hacdcpf.cpp:solve\_power\_flow），隔离/失电母线不得作 slack；
  \item \textbf{换流器协调预校核}（可选，默认关）：
        \srcpath{powerflow::evaluate\_converter\_coordination}（:solve\_power\_flow），
        按 DC 电压岛汇总电压源与功率平衡裕度，Error/Fatal 级问题阻断求解
        （规则目录见第~\ref{sec:hybrid}~章）；
  \item \textbf{图拓扑预检}（:solve\_power\_flow）：无任何带 slack 的合法 AC 岛时直接返回失败；
  \item \textbf{SolverData 获取与缓存}：\srcpath{get\_cached\_solver\_data}（调用点 :solve\_power\_flow，
        thread\_local 缓存，指针+存储视图+哈希签名三重匹配，:get\_cached\_solver\_data）；
        未命中时 \srcpath{make\_solver\_data} 重建
        （\srcpath{src/power\_flow/solver\_data.cpp:make\_solver\_data}）；
  \item \textbf{ZIP 权重应用}：\srcpath{apply\_zip\_weights}（\srcpath{src/api/hacdcpf.cpp:apply\_zip\_weights}）；
  \item \textbf{Newton 求解}：\srcpath{engine::NewtonSolver::solve}（调用点 :solve\_power\_flow，
        实现 \srcpath{src/power\_flow/newton\_solver.cpp:solve}）；
  \item \textbf{派生结果回算}：\srcpath{populate\_derived\_results}（\srcpath{src/api/hacdcpf.cpp:populate\_derived\_results}），
        含支路潮流 \srcpath{compute\_branch\_flows} 与换流器传输功率/损耗/限值校核；
  \item \textbf{换流器物理限值强制}（可选，默认关）：
        \fld{enforce\_converter\_physical\_limits} 开启时（\srcpath{src/api/hacdcpf.cpp:solve\_power\_flow}），
        违反换流器不等式的 Newton 根被判物理不可行（详见第~\ref{sec:hybrid}~章）；
  \item \textbf{结果回投影}：\srcpath{unproject\_pf\_result}（\srcpath{src/api/hacdcpf.cpp:unproject\_pf\_result}），
        把 canonical 解广播/归因回原始母线与支路编号空间，并恢复被合并开关/
        断路器的端子潮流。
\end{enumerate}

\subsection{公共 API 入口族}
全部潮流相关公共入口（\srcpath{src/api/hacdcpf.cpp} 与
\srcpath{include/hacdcpf/api/hacdcpf.hpp}）汇总如下；GUI 后端
（\srcpath{tests/run\_gui\_server.cpp}）的 \texttt{method} 分发与之对应
（\texttt{ac\_newton}/\allowbreak\texttt{pure\_ac}/\allowbreak\texttt{dc}/\allowbreak\texttt{fdpf}/
\texttt{adaptive}/\allowbreak\texttt{islanded}/\allowbreak\texttt{distributed\_slack}/
\texttt{three\_phase}/\texttt{three\_phase\_hybrid} 等）：

\begin{center}
\footnotesize
\begin{tabular}{@{}L{6.2cm}L{3.2cm}L{5.6cm}@{}}
\toprule
入口 & 求解器/路径 & 说明\\
\midrule
\srcpath{solve\_power\_flow} & 统一 Newton & 混合 AC/DC 主入口\\
\srcpath{powerflow::solve\_ac} & $\to$ \srcpath{solve\_power\_flow} & 薄包装（\srcpath{src/power\_flow/ac.cpp:solve\_ac}）\\
\srcpath{powerflow::solve\_hybrid} & $\to$ \srcpath{solve\_power\_flow} & 薄包装（\srcpath{src/power\_flow/hybrid.cpp}）\\
\srcpath{solve\_power\_flow\_fdpf} & \srcpath{FDPFSolver} & 快速解耦，纯 AC\\
\srcpath{solve\_ac\_linearized\_dc} & $B\theta=P$ & 线性化角度潮流（\srcpath{src/power\_flow/ac\_linearized\_pf.cpp:solve\_ac\_linearized\_dc}）\\
\srcpath{solve\_dc\_power\_flow} & \srcpath{DCSolver} & DC 网络电压 Newton（\srcpath{src/power\_flow/dc\_solver.cpp:solve}）\\
\srcpath{solve\_power\_flow\_adaptive} & \srcpath{AdaptiveSolver} & 分岛并行（\srcpath{src/power\_flow/adaptive\_solver.cpp:solve}）\\
\srcpath{solve\_power\_flow\_distributed\_slack}(\srcpath{\_full}) & \srcpath{DistributedSlackSolver} & 分布式松弛两变体\\
SolverHandle 族 & \srcpath{make\_solver\_handle} 等 & 符号分解/模式跨多次求解复用（时序仿真用）\\
\srcpath{analysis::solve\_distribution\_pf} & BFS 前推回代 & 辐射状配网（\srcpath{src/power\_flow/distribution\_power\_flow.cpp:solve\_distribution\_pf}）\\
\srcpath{analysis::solve\_three\_phase\_nr} & 相域 NR & 三相不平衡（\srcpath{src/power\_flow/three\_phase\_nr.cpp:solve\_three\_phase\_nr\_snapshot}）\\
\srcpath{solve\_three\_phase\_hybrid\_pf} & 相域 AC+DC 联立 & 矩阵级 API（\srcpath{src/power\_flow/three\_phase\_hybrid.cpp:solve\_three\_phase\_hybrid\_pf}）\\
\srcpath{CpfSolver::solve} & 连续潮流 & 电压稳定（\srcpath{src/power\_flow/voltage\_stability.cpp:solve}）\\
\bottomrule
\end{tabular}
\end{center}

\subsection{求解器抽象层与工厂}
抽象接口 \srcpath{IPowerFlowSolver}（\srcpath{include/hacdcpf/power\_flow/solver\_interface.hpp:IPowerFlowSolver}）
声明 \texttt{solve(problem)}/\texttt{name()} 两方法及 6 个子类
（\texttt{NewtonHybridSolver}/\allowbreak\texttt{FDPFSolver}/\allowbreak\texttt{DCPowerFlowSolver} 等，:NewtonHybridSolver、:FDPFSolver、:DCPowerFlowSolver、:AdaptiveHybridSolver、:ContinuationPowerFlowSolver、:NewtonKrylovSolver）；
问题描述符 \srcpath{PowerFlowProblem}（\srcpath{include/hacdcpf/power\_flow/power\_flow\_problem.hpp:PowerFlowProblem}）携带
SolverData、选项与规模摘要，工厂 \srcpath{build\_power\_flow\_problem}
（声明 \srcpath{include/hacdcpf/power\_flow/power\_flow\_problem.hpp:build\_power\_flow\_problem}，实现
\srcpath{src/power\_flow/solver\_interface.cpp:build\_power\_flow\_problem}）装配。
\srcpath{solver\_factory.hpp} 提供 \texttt{PowerFlowMethod} 枚举（含
\texttt{Continuation}）与显式 \texttt{create(PowerFlowMethod)}；六个子类的 \texttt{solve}
与工厂实现同集中于新增的 \srcpath{src/power\_flow/solver\_interface.cpp}
（:solve）。算法选择须经 \texttt{PowerFlowMethod}\textbf{显式}进行；
\texttt{PowerFlowOptions} 只配置已选方法，不再提供会被忽略的工厂重载。
如实说明：抽象层现已完整接线（\texttt{NewtonHybridSolver::\allowbreak solve}
委托 \srcpath{powerflow::NewtonSolver}，\srcpath{src/power\_flow/solver\_interface.cpp:solve}），
但生产门面 \srcpath{solve\_power\_flow} 仍不经过该抽象层，直接持有
thread\_local 的 \srcpath{engine::NewtonSolver}（src/api/hacdcpf.cpp:solve\_power\_flow；
规范头 \srcpath{include/hacdcpf/engine/solver/native/nle/newton\_solver.hpp}；
\srcpath{power\_flow/newton\_solver.hpp} 与 \srcpath{power\_flow/solvers/newton\_solver.hpp}
均为转发头）。

\subsection{零阻抗合并对潮流装配的影响}
母线合并发生在投影层（装配之前）：
\srcpath{merge\_zero\_impedance\_buses\_impl}（\srcpath{src/model/network\_utils.cpp:merge\_zero\_impedance\_buses\_impl}）
以阈值 \fld{impedance\_threshold}=$10^{-4}$（\srcpath{include/hacdcpf/projection/project\_to\_canonical.hpp:ProjectionOptions}）
合并近零阻抗支路两端母线；闭合断路器先折叠为近零阻抗支路再被合并
（src/model/network\_utils.cpp:add\_equivalent\_branch\_from\_switch）。因此 SolverData 中只剩 canonical 母线，
$Y_{bus}$ 装配时 $|z|=0$ 支路直接跳过（\srcpath{src/power\_flow/admittance\_builder.cpp:build\_admittance\_matrix}）。
合并映射 \fld{bus\_merge\_map} 随 SolverData 携带，求解后回投影恢复原始编号空间
的结果与开关端子潮流。

\subsection{缓存体系}
两层缓存避免重复装配与重复符号分解：
\begin{itemize}
  \item \textbf{API 层 SolverDataCache}（thread\_local，src/api/hacdcpf.cpp:SolverDataCache）：
        同一 \texttt{HybridPowerSystem} 对象且 \fld{loss\_model} 不变时直接复用；
  \item \textbf{求解器层 PatternCache}（规范头
        \srcpath{include/hacdcpf/engine/solver/native/nle/newton\_solver.hpp}；
        \srcpath{include/hacdcpf/power\_flow/newton\_solver.hpp} 为转发头）：
        缓存 \texttt{JacobianContext}、\texttt{JacobianPattern} 与稀疏线性求解器实例，
        键含数据指针、\fld{build\_id}（原子递增，缓存失效标记）与 Ybus/Gdc 存储地址；
        稀疏模式只在缓存失效时 \texttt{analyzePattern} 一次，之后每次迭代仅数值分解；
  \item \textbf{求解器状态工作区 SolverWorkspace}
        （\srcpath{include/hacdcpf/power\_flow/workspace.hpp:SolverWorkspace}，
        实现 \srcpath{src/power\_flow/workspace.cpp}）：已不再是空脚手架——
        \texttt{prepare\_state}/\texttt{prepare\_equations}/\texttt{reset\_iteration}
        在维度不变时保留 Eigen 分配、仅重置数值。facade 的 thread\_local
        \texttt{NewtonSolver} 跨普通求解复用同一工作区
        （成员 \fld{workspace\_}/\fld{workspace\_in\_use\_}，engine 头 :NewtonSolver；
        src/power\_flow/newton\_solver.cpp:solve 的占用守卫使同伦递归的内层求解改用局部工作区，
        避免覆盖外层状态），状态与方程向量经 \texttt{prepare\_state}（:solve）/
        \texttt{prepare\_equations}（:solve）挂接。
\end{itemize}
时序生产模拟（UC$\to$OPF$\to$PF 流水线）经 SolverHandle 族在整个时间轴上复用
同一求解器实例，拓扑仅在 \fld{in\_service} 变化时重建
（\srcpath{docs/theory/time\_series\_power\_flow\_models.md}）。

% ── 通用理论层（修订版 2 新增，专著级深度）──
% theory_fundamentals.tex —— 潮流计算 通用理论：节点方程、Newton 法与数值性态
% 本文件为理论层章节，遵循 docs/modules/THEORY_WRITING_GUIDE.md。

\section{潮流计算的数学基础}\label{sec:pf-th-fund}

\begin{theorynote}
本章为通用数学理论，按电力系统潮流计算的标准模型与数值分析文献体系撰写，覆盖：
节点功率方程从基尔霍夫电流定律与 $\pi$ 等值电路出发的一般推导（含复数变比变压器
与移相器两端口）、潮流问题的数学分类（PQ/PV/V$\theta$ 节点、方程-未知数计数、解的
存在性与多解性、可解性的压缩充分条件）、Newton--Raphson 法的一般理论（局部二次收敛定理及其证明、收敛域、
雅可比奇异的物理含义）、直角坐标下失配的精确二次结构与 Iwamoto 最优乘子法、
病态与数值性态（条件数分析、高 $R/X$ 比、重负荷/弱网、
雅可比奇异与电压崩溃的关系）、稀疏技术理论（图消元与填充、最小度/近似最小度
排序原理、LU 填充的图刻画），以及由平台求解器实测数据给出的收敛性数值算例与基准。
与平台实现（\srcpath{src/power\_flow/}）的对应关系见
本章末"与实现的对应"表，超出实现的内容以"未实现扩展说明"框就地标记。
\end{theorynote}

\subsection{记号与约定}
\begin{itemize}
  \item $\mathcal N=\{1,\dots,n\}$：AC 节点集；$n$ 为节点数；$i,j$ 为节点下标；
  \item $V_i=V_{m,i}e^{\mathrm j\theta_i}\in\mathbb C$：节点 $i$ 的复电压，
        $\theta_{ij}=\theta_i-\theta_j$；$\overline{(\cdot)}$ 表示复共轭；
        $\mathrm j$ 为虚数单位；
  \item $S_i=P_i+\mathrm jQ_i=V_i\,\overline{I_i}$：节点注入复功率（注入网络为正方向）；
  \item $Y_{bus}=G+\mathrm jB\in\mathbb C^{n\times n}$：节点导纳矩阵，
        $Y_{ij}=G_{ij}+\mathrm jB_{ij}$；
  \item 支路 $i$--$j$：串联阻抗 $z_{ij}=r_{ij}+\mathrm jx_{ij}$，串联导纳
        $y_{ij}=1/z_{ij}=g_{ij}+\mathrm jb_{ij}$，线路充电（对地）总电纳 $b_{c,ij}$，
        两端各半；
  \item $t=ae^{\mathrm j\varphi}$：变压器/移相器复数变比，$a>0$ 为变比幅值，
        $\varphi$ 为移相角；约定理想变压器位于 from 侧（节点 $i$ 侧）；
  \item $F:\mathbb R^m\to\mathbb R^m$：潮流失配（mismatch）函数，$J=DF$ 为其雅可比；
        $\lVert\cdot\rVert$ 未指明时为向量 $2$-范数及其诱导矩阵范数；
        $\kappa(J)=\lVert J\rVert\,\lVert J^{-1}\rVert$：条件数；
  \item $u$：浮点运算单位舍入（双精度 $u\approx1.1\times10^{-16}$）；
  \item 全部电气量除特别声明外取标幺值（系统基准 $S_B$，电压基准按节点额定值）。
\end{itemize}

\subsection{节点功率方程的一般推导}\label{sec:pf-th-fund-power-eq}

\subsubsection{\texorpdfstring{$\pi$}{π} 等值支路与复变比两端口}

\begin{definition}[$\pi$ 等值支路]\label{def:pf-th-fund-pi}
输电线路/电缆支路 $i$--$j$ 的集总参数 $\pi$ 等值电路由串联导纳 $y_{ij}$ 与两端
对地充电电纳 $\mathrm jb_{c,ij}/2$ 组成。记 $i$、$j$ 端注入支路的电流为
$I_{ij}$、$I_{ji}$（流入支路为正），则无变压器支路满足
\begin{equation}
\begin{bmatrix}I_{ij}\\ I_{ji}\end{bmatrix}
=
\begin{bmatrix}
y_{ij}+\mathrm jb_{c,ij}/2 & -y_{ij}\\
-y_{ij} & y_{ij}+\mathrm jb_{c,ij}/2
\end{bmatrix}
\begin{bmatrix}V_i\\ V_j\end{bmatrix}.
\label{eq:pf-th-fund-pi}
\end{equation}
该两端口是互易的（转移导纳相等）且无源，$2\times2$ 导纳矩阵对称。
\end{definition}

\begin{definition}[复变比变压器/移相器两端口]\label{def:pf-th-fund-tap}
非标准变比变压器与移相器统一建模为：$\pi$ 等值串联支路级联一台位于 from 侧
（节点 $i$ 侧）的理想复变比变压器。设理想变压器内侧（绕组侧）节点为 $k$，复数变比
$t=ae^{\mathrm j\varphi}$ 定义为
\begin{equation}
V_k=\frac{V_i}{t}.
\end{equation}
理想变压器不储能、不耗能，两侧复功率守恒：
\begin{equation}
V_i\,\overline{I_i}=V_k\,\overline{I_k'}
\quad\Longrightarrow\quad
I_i=\frac{I_k'}{\overline{t}},
\label{eq:pf-th-fund-ideal}
\end{equation}
其中 $I_k'$ 为内侧节点 $k$ 流入串联支路的电流。
\end{definition}

\begin{proposition}[含变比与移相的一般支路模型]\label{prop:pf-th-fund-tport}
定义~\ref{def:pf-th-fund-tap} 的两端口导纳参数为
\begin{equation}
\begin{bmatrix}I_i\\ I_j\end{bmatrix}
=
\begin{bmatrix}
\dfrac{y_{ij}+\mathrm jb_{c,ij}/2}{|t|^2} & -\dfrac{y_{ij}}{\overline{t}}\\[2.2ex]
-\dfrac{y_{ij}}{t} & y_{ij}+\mathrm jb_{c,ij}/2
\end{bmatrix}
\begin{bmatrix}V_i\\ V_j\end{bmatrix}.
\label{eq:pf-th-fund-tport}
\end{equation}
\end{proposition}

\begin{proof}
内侧节点 $k$ 经 $\pi$ 支路与节点 $j$ 相连，由式~\eqref{eq:pf-th-fund-pi}，
\begin{equation}
I_k'=\Bigl(y_{ij}+\mathrm j\tfrac{b_{c,ij}}{2}\Bigr)V_k-y_{ij}V_j,\qquad
I_j=-y_{ij}V_k+\Bigl(y_{ij}+\mathrm j\tfrac{b_{c,ij}}{2}\Bigr)V_j .
\end{equation}
代入 $V_k=V_i/t$ 与式~\eqref{eq:pf-th-fund-ideal} 的 $I_i=I_k'/\overline{t}$：
\begin{equation}
I_i=\frac{1}{\overline{t}}\Bigl[\Bigl(y_{ij}+\mathrm j\tfrac{b_{c,ij}}{2}\Bigr)
\frac{V_i}{t}-y_{ij}V_j\Bigr]
=\frac{y_{ij}+\mathrm jb_{c,ij}/2}{|t|^2}V_i-\frac{y_{ij}}{\overline{t}}V_j ,
\end{equation}
\begin{equation}
I_j=-\frac{y_{ij}}{t}V_i+\Bigl(y_{ij}+\mathrm j\tfrac{b_{c,ij}}{2}\Bigr)V_j .
\end{equation}
两式合并即得式~\eqref{eq:pf-th-fund-tport}。
\end{proof}

\begin{remark}[互易性、对称性与移相角的作用]\label{rem:pf-th-fund-reciprocity}
当 $\varphi=0$（实变比）时 $-\,y_{ij}/\overline{t}=-\,y_{ij}/t$，两端口互易，
$Y_{bus}$ 保持（复）对称；当 $\varphi\neq0$ 时
$Y_{ij}^{ft}\neq Y_{ij}^{tf}$，两端口在导纳参数意义下不再互易，$Y_{bus}$
失去对称性——但理想变压器仍满足功率守恒式~\eqref{eq:pf-th-fund-ideal}，
非对称恰恰是"电压相移而功率守恒"在导纳框架中的正确编码。工程上
$\varphi\neq0$ 的支路（移相器、相位角调节变压器）用于控制有功潮流的分布：
其效果是在支路两端电压相角差上叠加一个可控偏置 $\varphi$。任何假设
$Y_{bus}$ 对称的求解器或求逆算法在含移相器的网络上必须先核对适用性。
\end{remark}

\subsubsection{节点导纳矩阵与节点功率方程}

\begin{definition}[节点导纳矩阵]\label{def:pf-th-fund-ybus}
对所有在运支路按式~\eqref{eq:pf-th-fund-pi}/\eqref{eq:pf-th-fund-tport}
组装：对角元累加自导纳（含充电与节点对地并联导纳 $g_{s,i}+\mathrm jb_{s,i}$），
非对角元累加转移导纳，得 $Y_{bus}=G+\mathrm jB$。其结构性质：
(i) 稀疏——非零元个数约为 $n+2b$（$b$ 为支路数，电力网 $b\approx1.5n$）；
(ii) 无移相器时对称；（iii) 一般\textbf{不}具对角占优性（充电电纳与串联电纳
在对角元中异号抵消），因此不能直接套用对角占优类的迭代法理论。
\end{definition}

\begin{theorem}[节点功率方程]\label{thm:pf-th-fund-power-eq}
设网络无内部独立电流源，节点注入电流仅由支路与对地导纳产生，则
\begin{equation}
I_i=\sum_{j\in\mathcal N}Y_{ij}V_j,\qquad
S_i=V_i\,\overline{\Bigl(\sum_{j\in\mathcal N}Y_{ij}V_j\Bigr)},
\qquad i\in\mathcal N .
\label{eq:pf-th-fund-si}
\end{equation}
\end{theorem}

\begin{proof}
对节点 $i$ 应用基尔霍夫电流定律：注入网络的电流 $I_i$ 等于经各关联支路流出
电流之和。每条关联支路按式~\eqref{eq:pf-th-fund-pi}/\eqref{eq:pf-th-fund-tport}
的右端贡献一个关于端电压的线性式；把所有支路的端电压系数按列下标归类，
恰为 $Y_{bus}$ 第 $i$ 行与电压向量的乘积，即第一式。第二式由复功率定义
$S_i=V_i\overline{I_i}$ 直接代入。
\end{proof}

\begin{corollary}[极坐标功率方程]\label{cor:pf-th-fund-polar}
写 $V_i=V_{m,i}e^{\mathrm j\theta_i}$、$Y_{ij}=G_{ij}+\mathrm jB_{ij}$，
$\theta_{ij}=\theta_i-\theta_j$，则式~\eqref{eq:pf-th-fund-si} 的实部、虚部为
\begin{align}
P_i&=\sum_{j\in\mathcal N}V_{m,i}V_{m,j}
\bigl(G_{ij}\cos\theta_{ij}+B_{ij}\sin\theta_{ij}\bigr),\label{eq:pf-th-fund-polar-p}\\
Q_i&=\sum_{j\in\mathcal N}V_{m,i}V_{m,j}
\bigl(G_{ij}\sin\theta_{ij}-B_{ij}\cos\theta_{ij}\bigr).\label{eq:pf-th-fund-polar-q}
\end{align}
\end{corollary}

\begin{proof}
$\overline{Y_{ij}V_j}=(G_{ij}-\mathrm jB_{ij})V_{m,j}e^{-\mathrm j\theta_j}$，
乘以 $V_i=V_{m,i}e^{\mathrm j\theta_i}$ 得
$V_{m,i}V_{m,j}(G_{ij}-\mathrm jB_{ij})e^{\mathrm j\theta_{ij}}$；
展开 $e^{\mathrm j\theta_{ij}}=\cos\theta_{ij}+\mathrm j\sin\theta_{ij}$，
取实部、虚部分别即式~\eqref{eq:pf-th-fund-polar-p}、\eqref{eq:pf-th-fund-polar-q}。
\end{proof}

式~\eqref{eq:pf-th-fund-si}--\eqref{eq:pf-th-fund-polar-q} 是潮流计算的
标准模型~\cite{pff:stagg-elabiad,pff:arrillaga-watson}。其实现层转写（逐式对照
源码、ZIP 负荷注入、DC 侧 $P=V_{dc}\odot(G_{dc}V_{dc})$ 等）见本手册实现层
"装配与 SolverData"章与"统一 Newton 章"，此处不再重复。

\subsection{潮流问题的数学分类}\label{sec:pf-th-fund-classification}

\subsubsection{节点类型}

\begin{definition}[PQ、PV、V$\theta$ 节点]\label{def:pf-th-fund-bus-types}
按节点的已知量/未知量划分~\cite{pff:stott-review}：
\begin{itemize}
  \item \textbf{PQ 节点}：指定 $P_i^{\mathrm{spec}},Q_i^{\mathrm{spec}}$；
        未知量 $\theta_i,V_{m,i}$；对应负荷母线与无定压控制的电源母线；
  \item \textbf{PV 节点}：指定 $P_i^{\mathrm{spec}},V_{m,i}^{\mathrm{spec}}$；
        未知量 $\theta_i$（$Q_i$ 由解出后回算并受 $[Q_i^{\min},Q_i^{\max}]$ 校核）；
        对应具备自动电压调节的发电机母线；
  \item \textbf{V$\theta$ 节点（slack/参考节点）}：指定 $\theta_i$（通常 $0$）
        与 $V_{m,i}$；$P_i,Q_i$ 完全自由，由功率平衡残差吸收。
\end{itemize}
\end{definition}

\begin{proposition}[方程-未知数计数]\label{prop:pf-th-fund-counting}
$n$ 节点系统含 $n_{pv}$ 个 PV 节点、$1$ 个 V$\theta$ 节点，其余为 PQ 节点，
则潮流问题的未知量个数与独立方程个数均为
\begin{equation}
m=2n-2-n_{pv}.
\end{equation}
\end{proposition}

\begin{proof}
未知量：角度 $\theta_i$ 对全部非参考节点自由，共 $n-1$ 个；电压幅值 $V_{m,i}$
仅对 PQ 节点自由，共 $n-1-n_{pv}$ 个；合计 $2n-2-n_{pv}$。方程：有功平衡
$P_i^{\mathrm{spec}}-P_i^{\mathrm{calc}}=0$ 对全部非参考节点列出，共 $n-1$ 个；
无功平衡仅对 PQ 节点列出（PV 节点的 $Q$ 不平衡由发电机无功吸收、V$\theta$
节点的 $P,Q$ 均自由），共 $n-1-n_{pv}$ 个；合计 $2n-2-n_{pv}$。两者相等，
问题为方阵非线性方程组 $F(x)=0$，$F:\mathbb R^m\to\mathbb R^m$。
\end{proof}

\begin{remark}[参考节点的双重必要性]\label{rem:pf-th-fund-slack}
V$\theta$ 节点的设置有两个独立的数学理由：
(i) \emph{角度规范}：式~\eqref{eq:pf-th-fund-polar-p}--\eqref{eq:pf-th-fund-polar-q}
中 $F$ 仅通过差 $\theta_{ij}$ 依赖于角度，全体角度同时平移不改变任何失配——
角度空间带有一个 $S^1$ 对称性，不固定一个参考角则雅可比恒有一维零空间；
(ii) \emph{损耗悬空}：网络损耗在求解之前未知，$n$ 个节点的有功指定值不能
同时精确满足，必须有一个节点的 $P$（与 $Q$）退出方程组以吸收差额。
多岛系统每个连通 AC 岛各需一个参考节点；把损耗分摊到多台机组的分布式松弛
是上述单点吸收的推广，不改变计数结论。
\end{remark}

\subsubsection{解的存在性与多解性}

潮流方程是关于 $(\theta,V_m)$ 的二次型代数方程组（固定 $V_m$ 时对
$e^{\mathrm j\theta}$ 为双线性）。其解集一般不是单点：

\begin{example}[两节点系统的解析解与鼻点]\label{ex:pf-th-fund-twobus}
节点 1 为参考节点 $E\angle0$，节点 2 经纯电抗 $X$ 线路受电，负荷
$P_L+\mathrm jQ_L$。记 $V=V_{m,2}$、$\delta=-\theta_2$，则
\begin{equation}
P_L=\frac{EV}{X}\sin\delta,\qquad
Q_L=\frac{EV}{X}\cos\delta-\frac{V^2}{X}.
\end{equation}
消去 $\delta$（利用 $\sin^2+\cos^2=1$）得关于 $V^2$ 的二次方程
\begin{equation}
V^4+\bigl(2Q_LX-E^2\bigr)V^2+X^2\bigl(P_L^2+Q_L^2\bigr)=0,
\label{eq:pf-th-fund-v4}
\end{equation}
\begin{equation}
V^2=\frac{E^2-2Q_LX\pm\sqrt{\Delta}}{2},\qquad
\Delta=E^4-4E^2Q_LX-4X^2P_L^2 .
\label{eq:pf-th-fund-v2}
\end{equation}
$\Delta>0$ 且 $E^2>2Q_LX$ 时存在\textbf{两个}正根：高电压根（工程可行解）与
低电压根；$\Delta=0$ 给出鞍结分岔（鼻点）$P_L^{\max}=\sqrt{E^4-4E^2Q_LX}/(2X)$；
$\Delta<0$ 时无实解。数值上取 $E=1$、$X=0.5$、$Q_L=0$：$P_L=0.5$ 时
$V^2=(1\pm\sqrt{0.75})/2$，即 $V\approx0.966$ 或 $V\approx0.259$；
$P_L=1$ 时两根于 $V=1/\sqrt2$ 处汇合（鼻点）；$P_L>1$ 无解。
\end{example}

\begin{remark}[一般网络的多解性与不可解性]\label{rem:pf-th-fund-multiplicity}
例~\ref{ex:pf-th-fund-twobus} 的定性结论对一般网络成立：潮流方程的实解个数
通常随规模增长，且重负荷参数下可\textbf{不存在}实解（鼻点之外）。两点推论：
(i) Newton 法不收敛\textbf{不能}作为问题无解的证明——它只说明初值不在任何
解的局部收敛域内被给定迭代法捕获；(ii) 工程上关心的是高电压物理解，多解的
枚举/全解搜索属于超出常规潮流的范畴。
\end{remark}

\begin{gapnote}
平台提供 HELM（全纯嵌入）求解器
（\srcpath{src/power\_flow/helm\_solver.cpp:HelmSolver::solve}），
其方法论是沿解析延拓收敛到与平启动连通的物理解，而非枚举多解；
全解枚举（如消元法、区间分析）未实现，也无计划对应接口。
\end{gapnote}

\subsubsection{无功限值与活动集表述}

PV 节点的无功限值使潮流问题严格地说不再是纯方程组，而是带互补约束的
混合问题：
\begin{equation}
Q_g^{\min}\le Q_g(x)\le Q_g^{\max},\qquad
\text{限内时 }V_m=V^{\mathrm{set}},\ \text{触限时 }Q_g=Q_g^{\pm},
\label{eq:pf-th-fund-qlimit}
\end{equation}
三态（下限/限内/上限）可用中值互补函数
$\mathrm{mid}(V_m-V^{\mathrm{set}},\,Q_g-Q_g^{\min},\,Q_g-Q_g^{\max})=0$
统一表示~\cite{pff:hintermuller}。工程实现的通行做法是\emph{活动集}策略：
固定当前 PV/PQ 归属 $A$，解纯方程组 $F_A(x)=0$ 至收敛，在收敛点检验隐含
无功，越限则批量切换为 PQ 并钉在限值，重复至无越限；恢复侧设迟滞与单次
恢复审计以防止活动集循环。平台的活动集契约（含固定超集雅可比布局与
CHKS 光滑化的半光滑 Newton 替代）全文见
\srcpath{docs/theory/pv\_pq\_switching\_contract.md}，实现层转写见本手册
"统一 Newton 章"与"鲁棒求解器章"。

\subsection{潮流可解性的充分条件}\label{sec:pf-th-fund-solvability}

第~\ref{sec:pf-th-fund-classification}~节从两母线鼻点（例~\ref{ex:pf-th-fund-twobus}）
定性说明了潮流解可以不存在。本节给出一个\emph{定量}充分条件：当负荷相对短路容量足够
小时，存在唯一的高电压解，且不动点迭代收敛到它。设网络含单一 slack（母线 $0$，
电压 $u_0$）与 PQ 母线集 $L=\{1,\dots,n\}$，$Y_{bus}$ 按 $\{0\}$ 与 $L$ 分块，
子块 $Y_{LL}$ 非奇异（连通接地网络成立）。

\begin{proposition}[PQ 潮流的不动点形式]\label{prop:pf-th-fund-fixedpoint}
记 PQ 母线注入复功率 $s\in\mathbb C^n$（发电为正）。定义\emph{空载电压}
$w:=-Y_{LL}^{-1}Y_{L0}u_0$（$s=0$ 时的解）。则 PQ 潮流方程等价于不动点方程
\begin{equation}
u=T(u):=w+Y_{LL}^{-1}\,\overline{s\oslash u},
\label{eq:pf-th-fund-fixedpoint}
\end{equation}
其中 $\oslash$ 为逐分量除、$\overline{(\cdot)}$ 为复共轭。
\end{proposition}
\begin{proof}
PQ 母线电流注入 $i_L=Y_{LL}u+Y_{L0}u_0$；由 $s_i=u_i\overline{i_i}$ 得
$i_i=\overline{s_i/u_i}$，即 $i_L=\overline{s\oslash u}$。两式相等，左乘 $Y_{LL}^{-1}$、
移项即式~\eqref{eq:pf-th-fund-fixedpoint}。
\end{proof}

\begin{theorem}[压缩可解性条件]\label{thm:pf-th-fund-solvability}
记 $\eta:=\lVert Y_{LL}^{-1}\rVert_\infty$（行和范数）、$\sigma:=\max_i|s_i|$、
$w_{\min}:=\min_i|w_i|>0$。若
\begin{equation}
4\,\eta\,\sigma<w_{\min}^2,
\label{eq:pf-th-fund-solvcond}
\end{equation}
则映射 $T$ 在闭球 $\mathcal B=\{u:\lVert u-w\rVert_\infty\le w_{\min}/2\}$ 上是压缩自映射，
潮流方程~\eqref{eq:pf-th-fund-fixedpoint} 在 $\mathcal B$ 内有\emph{唯一}解，且 Picard
迭代 $u^{(k+1)}=T(u^{(k)})$ 从 $u^{(0)}=w$ 出发几何收敛到它。
\end{theorem}
\begin{proof}
在 $\mathcal B$ 上 $|u_i|\ge|w_i|-\tfrac{w_{\min}}2\ge\tfrac{w_{\min}}2=:\rho$。
\textbf{自映射}：对 $u\in\mathcal B$，
\begin{equation*}
\lVert T(u)-w\rVert_\infty=\bigl\lVert Y_{LL}^{-1}\overline{s\oslash u}\bigr\rVert_\infty
\le\eta\,\max_i\frac{|s_i|}{|u_i|}\le\frac{\eta\sigma}{\rho}
=\frac{2\eta\sigma}{w_{\min}}<\frac{w_{\min}}2,
\end{equation*}
末步用式~\eqref{eq:pf-th-fund-solvcond}，故 $T(u)\in\mathcal B$。
\textbf{压缩}：对 $u,u'\in\mathcal B$，逐分量
$\bigl|\overline{s_i/u_i}-\overline{s_i/u_i'}\bigr|
=|s_i|\,|u_i-u_i'|/(|u_i||u_i'|)\le(\sigma/\rho^2)\lVert u-u'\rVert_\infty$，故
\begin{equation*}
\lVert T(u)-T(u')\rVert_\infty\le\eta\,\frac{\sigma}{\rho^2}\lVert u-u'\rVert_\infty
=\frac{4\eta\sigma}{w_{\min}^2}\lVert u-u'\rVert_\infty,
\end{equation*}
压缩系数 $L=4\eta\sigma/w_{\min}^2<1$。$\mathcal B$ 是完备度量空间中的闭集，Banach
不动点定理给出唯一不动点与几何收敛。
\end{proof}

\begin{remark}[短路容量裕度的读法与两母线校核]\label{rem:pf-th-fund-solvability}
$\eta=\lVert Y_{LL}^{-1}\rVert_\infty$ 是“路径阻抗”量级、$\eta^{-1}$ 度量短路容量，
$w_{\min}$ 为最弱母线的空载电压；条件~\eqref{eq:pf-th-fund-solvcond} 即“负荷 $\times$
路径阻抗 $\ll$ 空载电压$^2$”的短路容量裕度，与前推回代的压缩条件
（式~\eqref{sec:pf-th-cls-eq:bfsc}，第~\ref{sec:pf-th-cls}~章）同源。以两母线系统
（例~\ref{ex:pf-th-fund-twobus}，$Q=0$）校核：$Y_{LL}=1/(\mathrm jX)$、$w=E$、
$\eta=X$、$\sigma=P$，条件给出 $P<E^2/(4X)$，恰为真鼻点 $P_{\max}=E^2/(2X)$ 的一半——
即该充分条件保守约 $2$ 倍，符合“压缩界不触鼻点极限”的一般预期。更锐的（计及负荷
方向与网络结构的）可解性判据见 Bolognani--Zampieri~\cite{pff:bolognani} 与
Simpson-Porco--D\"orfler--Bullo~\cite{pff:simpsonporco}。
\end{remark}

\begin{gapnote}
条件~\eqref{eq:pf-th-fund-solvcond} 作为求解前的可解性预判在平台中\textbf{未实现}：
求解器直接尝试 Newton/BFS，以收敛与否与残差水平事后判定；HELM
（\srcpath{src/power\_flow/helm\_solver.cpp:HelmSolver::solve}）沿解析延拓返回与平启动
连通的物理解，同样不输出式~\eqref{eq:pf-th-fund-solvcond} 型的先验证书。
\end{gapnote}

\subsection{Newton--Raphson 一般理论}\label{sec:pf-th-fund-newton}

\subsubsection{迭代格式与局部二次收敛}

\begin{definition}[Newton 迭代]\label{def:pf-th-fund-newton}
对 $F:\mathbb R^m\to\mathbb R^m$、$F\in C^1$，Newton 迭代为
\begin{equation}
J(x_k)\,\Delta x_k=-F(x_k),\qquad x_{k+1}=x_k+\Delta x_k .
\label{eq:pf-th-fund-newton-step}
\end{equation}
潮流语境下 $F$ 取失配 $F^{\mathrm{spec}}-F^{\mathrm{calc}}$（符号约定不影响
理论），雅可比在极坐标下分块为
$J=\begin{bmatrix}H&N\\ M&L\end{bmatrix}$，
$H=\partial P/\partial\theta$、$N=\partial P/\partial V_m$、
$M=\partial Q/\partial\theta$、$L=\partial Q/\partial V_m$
~\cite{pff:tinney-hart}。
\end{definition}

\begin{lemma}[Banach 扰动引理]\label{lem:pf-th-fund-banach}
设 $A\in\mathbb R^{m\times m}$ 非奇异，$\lVert A^{-1}\rVert\le\beta$，
$E\in\mathbb R^{m\times m}$ 满足 $\beta\lVert E\rVert<1$，则 $A+E$ 非奇异且
\begin{equation}
\bigl\lVert(A+E)^{-1}\bigr\rVert\le\frac{\beta}{1-\beta\lVert E\rVert}.
\end{equation}
\end{lemma}

\begin{proof}
写 $A+E=A(I+A^{-1}E)$。$\lVert A^{-1}E\rVert\le\beta\lVert E\rVert<1$，
Neumann 级数 $\sum_{k\ge0}(-A^{-1}E)^k$ 绝对收敛，其和为 $(I+A^{-1}E)^{-1}$，
且 $\lVert(I+A^{-1}E)^{-1}\rVert\le\sum_k(\beta\lVert E\rVert)^k
=1/(1-\beta\lVert E\rVert)$。于是 $(A+E)^{-1}=(I+A^{-1}E)^{-1}A^{-1}$，
取范数即得结论。
\end{proof}

\begin{theorem}[Newton 法的局部二次收敛]\label{thm:pf-th-fund-quadratic}
设 $F\in C^1$，$J=DF$ 在 $x^*$ 的邻域内 Lipschitz 连续：
$\lVert J(x)-J(y)\rVert\le L\lVert x-y\rVert$；
$F(x^*)=0$ 且 $J(x^*)$ 非奇异，$\lVert J(x^*)^{-1}\rVert=\beta$。
则存在 $\delta>0$，使任意 $x_0\in B(x^*,\delta)$ 出发的 Newton 序列
良定义、收敛到 $x^*$，且
\begin{equation}
\lVert e_{k+1}\rVert\le\beta L\lVert e_k\rVert^2,
\qquad e_k=x_k-x^* .
\label{eq:pf-th-fund-qconv}
\end{equation}
\end{theorem}

\begin{proof}
取 $\delta\le\min\{r,\,1/(2\beta L)\}$，其中 $r$ 为 Lipschitz 邻域半径。
归纳假设 $x_k\in B(x^*,\delta)$：则
$\lVert J(x_k)-J(x^*)\rVert\le L\delta\le1/(2\beta)$，由
引理~\ref{lem:pf-th-fund-banach}（$\varepsilon=1/(2\beta)$，$\beta\varepsilon=1/2<1$），
$J(x_k)$ 非奇异且 $\lVert J(x_k)^{-1}\rVert\le2\beta$，迭代良定义。
由 $F(x_k)-F(x^*)=\int_0^1J(x^*+te_k)e_k\,\mathrm dt$ 与 Newton 方程
$J(x_k)(x_k+\Delta x_k-x^*)=J(x_k)e_k-F(x_k)$，得误差恒等式
\begin{equation}
J(x_k)\,e_{k+1}
=\int_0^1\bigl[J(x_k)-J(x^*+te_k)\bigr]e_k\,\mathrm dt .
\end{equation}
被积函数范数 $\le L(1-t)\lVert e_k\rVert^2$，积分得
$\lVert J(x_k)e_{k+1}\rVert\le\frac{L}{2}\lVert e_k\rVert^2$，
左乘 $J(x_k)^{-1}$ 并用 $2\beta$ 界即得式~\eqref{eq:pf-th-fund-qconv}。
又因 $\beta L\lVert e_k\rVert\le\beta L\delta\le1/2$，
$\lVert e_{k+1}\rVert\le\frac12\lVert e_k\rVert$，序列保持于球内并收缩，
故收敛且为 Q-二次。
\end{proof}

\begin{remark}[潮流场景的条件核对]\label{rem:pf-th-fund-hypotheses}
定理~\ref{thm:pf-th-fund-quadratic} 的三项假设在潮流中分别对应：
(i) $C^1$+Lipschitz：式~\eqref{eq:pf-th-fund-polar-p}--\eqref{eq:pf-th-fund-polar-q}
为三角多项式，任意紧集上自动满足；(ii) $F(x^*)=0$：物理解存在性，见
第~\ref{sec:pf-th-fund-classification}~节，非平凡；(iii) $J(x^*)$ 非奇异：
在鼻点处\textbf{系统性失效}（命题~\ref{prop:pf-th-fund-fold}）。"常规工况、
平启动、$3$--$7$ 次迭代收敛"的工程经验正是二次收敛阶段
$\lVert e_{k+1}\rVert\approx c\lVert e_k\rVert^2$ 的表现：
残差小数位数近似逐次翻倍。
\end{remark}

\begin{remark}[仿射协变性与实现中的缩放]\label{rem:pf-th-fund-affine}
Newton 序列在两组变换下不变：未知量的仿射重参数化 $\tilde x=Ax+b$ 与
方程的左乘非奇异阵 $\tilde F=BF$（精确算术下逐点一致）。这一\emph{仿射协变性}
意味着：对角缩放 $D_F=\mathrm{diag}(1/\max(|F_i^{\mathrm{spec}}|,1))$、
$D_x=\mathrm{diag}(\pi,V_{m,i},\dots)$ 不改变理论迭代轨线，只改变浮点误差
的分布——其实现动机（量纲均衡、改善 $J$ 的行/列范数比）见
第~\ref{sec:pf-th-fund-conditioning}~节。反之，步长截断
（$|\Delta\theta|\le1.5$~rad 等）与线搜索回退\textbf{破坏}纯 Newton 序列，
定理~\ref{thm:pf-th-fund-quadratic} 只在进入全步接受阶段后适用。
\end{remark}

\subsubsection{收敛域与全局化}

定理~\ref{thm:pf-th-fund-quadratic} 只保证\emph{存在}收敛球，未给出其大小。
Newton 迭代的全局吸引域结构复杂（多解系统的吸引域边界可为分形），工程上
以全局化策略替代收敛域刻画~\cite{pff:nocedal-wright}：在 merit 函数
$\phi(x)=\tfrac12\lVert F(x)\rVert^2$ 上做（非单调）Armijo 线搜索、
信赖域（dogleg）、Levenberg--Marquardt 正则化
$(J^{\mathsf T}J+\lambda I)\Delta x=-J^{\mathsf T}F$，或伪瞬态延拓
$(J+I/\delta t)\Delta x=-F$。这些策略把局部方法改造为对更大初值集
"收敛或明确失败"的方法，但都不提供"解存在"的证书。

\begin{gapnote}
定量收敛判据——Kantorovich 定理（以
$h=\beta L\lVert\Delta x_0\rVert\le1/2$ 在首步数据上判定收敛并给出误差界）
与 Smale 的 $\alpha$-理论——在当前实现中\textbf{未}接线：代码不计算
$h$ 型证书，失败回退链（线搜索 $\to$ 对角正则化 $\to$ LM $\to$ PTC，
实现层见"统一 Newton 章"）是纯启发式全局化。
\end{gapnote}

\subsubsection{雅可比奇异的物理含义}

\begin{proposition}[鞍结折点处雅可比必奇异]\label{prop:pf-th-fund-fold}
设单参数潮流 $F(x,\lambda)=0$（$\lambda$ 为负荷/发电水平参数）存在光滑解弧
$(x(s),\lambda(s))$，$s$ 为弧长参数。若在 $s^*$ 处 $\lambda'(s^*)=0$ 且
$x'(s^*)\neq0$（解弧在 $\lambda$ 方向回折），则 $J^*=D_xF(x(s^*),\lambda(s^*))$
奇异，且 $x'(s^*)\in\ker J^*$。
\end{proposition}

\begin{proof}
对恒等式 $F(x(s),\lambda(s))\equiv0$ 关于 $s$ 微分：
\begin{equation}
D_xF\,x'(s)+D_\lambda F\,\lambda'(s)=0 .
\end{equation}
代入 $s=s^*$：$\lambda'(s^*)=0$ 给出 $J^*x'(s^*)=0$，而 $x'(s^*)\neq0$，
故 $J^*$ 有非平凡零空间。
\end{proof}

\begin{remark}[奇异的两类成因与工程区分]\label{rem:pf-th-fund-singularity}
(i) \emph{参数性奇异}：命题~\ref{prop:pf-th-fund-fold} 的折点，物理含义是
鼻点/静态电压稳定极限——随 $\lambda$ 逼近 $\lambda^*$，
$\sigma_{\min}(J)\to0$，$\kappa(J)\to\infty$（见
第~\ref{sec:pf-th-fund-conditioning}~节）；
(ii) \emph{结构性奇异}：无参考节点的孤岛、悬空节点形成的空行/空列、
理想回路导致的秩亏——与负荷水平无关，是建模/拓扑缺陷。两者在线性求解阶段
都表现为分解失败或巨大修正量，但处置完全不同：前者应转入连续潮流定位
极限点~\cite{pff:ajjarapu-christy}，后者应在装配期诊断剔除。实现层以
空行/空列扫描与方程闭合校核处置后者（见"与实现的对应"表）。
\end{remark}
\subsection{直角坐标表述、二次结构与最优乘子法}\label{sec:pf-th-fund-rect}

极坐标式~\eqref{eq:pf-th-fund-polar-p}--\eqref{eq:pf-th-fund-polar-q} 是平台采用的
标准形式，但改用\emph{直角坐标}可揭示潮流失配的一个精确二次代数结构，并由此导出两个
经典结论：电流注入雅可比的常量非对角块，以及沿 Newton 方向的\emph{闭式最优步长}。

\subsubsection{直角坐标功率方程}
\begin{definition}[直角坐标电压与注入电流]\label{def:pf-th-fund-rect}
记 $V_i=e_i+\mathrm jf_i$。注入电流 $I_i=\sum_jY_{ij}V_j$ 的实/虚部
\begin{equation}
a_i:=\operatorname{Re}I_i=\sum_j\bigl(G_{ij}e_j-B_{ij}f_j\bigr),\qquad
b_i:=\operatorname{Im}I_i=\sum_j\bigl(G_{ij}f_j+B_{ij}e_j\bigr)
\end{equation}
均为状态 $x=(e_1,\dots,e_n,f_1,\dots,f_n)$ 的\emph{线性}函数。
\end{definition}

\begin{proposition}[直角坐标注入功率]\label{prop:pf-th-fund-rectpq}
\begin{equation}
P_i=e_ia_i+f_ib_i,\qquad Q_i=f_ia_i-e_ib_i,
\label{eq:pf-th-fund-rectpq}
\end{equation}
均为 $x$ 的齐二次型；母线电压幅值平方 $V_i^2=e_i^2+f_i^2$ 亦为齐二次型。
\end{proposition}
\begin{proof}
$S_i=V_i\overline{I_i}=(e_i+\mathrm jf_i)(a_i-\mathrm jb_i)
=(e_ia_i+f_ib_i)+\mathrm j(f_ia_i-e_ib_i)$，取实、虚部即得。$a_i,b_i$ 线性于 $x$，
故 $P_i,Q_i$ 为 $x$ 的二次齐次多项式。
\end{proof}

\subsubsection{失配映射的精确二次结构}
\begin{theorem}[潮流失配是二次多项式映射]\label{thm:pf-th-fund-quadmap}
以直角坐标未知量 $x$（slack 的 $e,f$ 固定）装配失配 $F(x)$——PQ 母线取
$(\Delta P_i,\Delta Q_i)$、PV 母线取 $(\Delta P_i,\ (V_i^{\mathrm{set}})^2-e_i^2-f_i^2)$——
则 $F$ 是 $x$ 的二次多项式映射：存在常向量 $c_0$、常矩阵 $J_0$ 与\emph{常}对称双线性
形 $\mathcal B[\cdot,\cdot]$ 使
\begin{equation}
F(x)=c_0+J_0x+\tfrac12\mathcal B[x,x],\qquad D^3F\equiv0 .
\label{eq:pf-th-fund-quadmap}
\end{equation}
从而对任意基点 $x_k$ 与方向 $\Delta$，
\begin{equation}
F(x_k+\alpha\Delta)=F(x_k)+\alpha\,J(x_k)\Delta+\tfrac12\alpha^2\,\mathcal B[\Delta,\Delta]
\label{eq:pf-th-fund-exactstep}
\end{equation}
是\emph{精确}恒等式（无截断），其中 $J(x_k)=J_0+\mathcal B[x_k,\cdot]$。
\end{theorem}
\begin{proof}
由命题~\ref{prop:pf-th-fund-rectpq}，$F$ 的每个分量是 $x$ 的二次多项式（常数项、
线性项与二次项之和），故其三阶及以上导数恒为零，Taylor 展开在二阶精确终止，
即式~\eqref{eq:pf-th-fund-quadmap}。把 $x=x_k+\alpha\Delta$ 代入，二次映射只产生
$\alpha^0,\alpha^1,\alpha^2$ 三项，系数分别为 $F(x_k)$、$J(x_k)\Delta$、
$\tfrac12\mathcal B[\Delta,\Delta]$，即式~\eqref{eq:pf-th-fund-exactstep}。
\end{proof}

\begin{remark}[电流注入雅可比的常量非对角块]\label{rem:pf-th-fund-ci}
以\emph{电流失配} $\Delta I_i=\overline{S_i^{\mathrm{spec}}/V_i}-\sum_jY_{ij}V_j$
（实/虚部堆叠）为方程的电流注入 Newton 法~\cite{pff:dacosta} 中，网络项
$\sum_jY_{ij}V_j$ 对 $V_j$ 线性，故非对角块
$\partial(\Delta I_i^{\mathrm{re}},\Delta I_i^{\mathrm{im}})/\partial(e_j,f_j)$
（$i\ne j$）恒等于\emph{常量} $2\times2$ 实块
$\left[\begin{smallmatrix}-G_{ij}&\phantom{-}B_{ij}\\-B_{ij}&-G_{ij}\end{smallmatrix}\right]$，
与运行点无关，只有 $2\times2$ 对角块随本地负荷电流的状态依赖而变。这一“非对角块免
重算”性质是电流注入法单步更省的结构来源。
\end{remark}

\subsubsection{Iwamoto 最优乘子法}
式~\eqref{eq:pf-th-fund-exactstep} 的精确性使沿 Newton 方向的 merit 有\emph{闭式}
最小点。设 $g:=F(x_k)$，Newton 方向 $\Delta_N$ 满足 $J(x_k)\Delta_N=-g$，
记 $c:=\mathcal B[\Delta_N,\Delta_N]$，则由式~\eqref{eq:pf-th-fund-exactstep}，
$F(x_k+\alpha\Delta_N)=(1-\alpha)g+\tfrac12\alpha^2c$。

\begin{proposition}[最优乘子的三次方程]\label{prop:pf-th-fund-iwamoto}
记 $a_0=g^{\mathsf T}g$、$a_1=g^{\mathsf T}c$、$a_2=c^{\mathsf T}c$，则 merit
$\phi(\alpha)=\tfrac12\lVert F(x_k+\alpha\Delta_N)\rVert^2$ 是 $\alpha$ 的\emph{四次}
多项式，其驻点由\emph{三次}方程
\begin{equation}
a_2\,\alpha^3-3a_1\,\alpha^2+2(a_0+a_1)\,\alpha-2a_0=0
\label{eq:pf-th-fund-iwamoto}
\end{equation}
给出。落在 $(0,1]$ 内使 $\phi$ 最小的实根即\emph{最优乘子} $\mu_k$，可由 Cardano
公式闭式求得，代价可忽略~\cite{pff:iwamoto}。
\end{proposition}
\begin{proof}
$\phi(\alpha)=\tfrac12\bigl[(1-\alpha)^2a_0+(1-\alpha)\alpha^2a_1+\tfrac14\alpha^4a_2\bigr]$。
求导 $\phi'(\alpha)=\tfrac12\bigl[-2a_0(1-\alpha)+a_1(2\alpha-3\alpha^2)+a_2\alpha^3\bigr]$，
令其为零并整理即式~\eqref{eq:pf-th-fund-iwamoto}。又 $\phi'(0)=-a_0<0$
（$a_0=\lVert g\rVert^2>0$），保证在 $\alpha>0$ 侧存在使 merit 下降的根。
\end{proof}

\begin{remark}[最优乘子的全局化意义与渐近还原]\label{rem:pf-th-fund-iwamoto}
$\mu_k$ 给出沿 Newton 方向的\emph{精确}最优 merit 下降，无需回溯试探，历史上专用于
病态/发散边缘系统的稳健求解~\cite{pff:iwamoto}。近解处 $\Delta_N\to0$，
$c=\mathcal B[\Delta_N,\Delta_N]$ 为二阶小量，$a_2=\lVert c\rVert^2\to0$，
式~\eqref{eq:pf-th-fund-iwamoto} 退化为 $2(a_0+a_1)\alpha-2a_0\approx0$，
$\mu_k\to1$：方法自动还原为全步 Newton，定理~\ref{thm:pf-th-fund-quadratic} 的
二次收敛得以保留。这与极坐标下 Armijo 线搜索近解 $\alpha\to1$ 殊途同归，差别在于
二次结构给出了\emph{无试探}的闭式步长。
\end{remark}

\begin{gapnote}
平台统一 Newton 求解器采用\textbf{极坐标}功率失配
（\srcpath{src/power\_flow/ac\_kernel.cpp:evaluate\_ac\_kernel\_serial}），全局化用
Armijo/非单调线搜索、LM 与 PTC（见第~\ref{sec:pf-th-glob}~章），
\textbf{未实现}直角坐标 Newton、电流注入 Newton 与 Iwamoto 最优乘子；
式~\eqref{eq:pf-th-fund-iwamoto} 依赖的精确二次结构在极坐标下不成立（三角项使
$D^3F\ne0$）。平台内唯一的电流注入构造是三相 Newton 的定点辅助
（\srcpath{src/power\_flow/three\_phase\_nr.cpp:evaluate\_fixed\_point\_current\_injections}），
非本节的电流失配 Newton 主迭代。
\end{gapnote}
\subsection{病态与数值性态}\label{sec:pf-th-fund-conditioning}

\subsubsection{条件数与扰动界}

\begin{definition}[条件数与向后误差]\label{def:pf-th-fund-cond}
线性系统 $Jx=b$ 的条件数 $\kappa(J)=\lVert J\rVert\lVert J^{-1}\rVert$。
称 $\hat x$ 具有向后误差 $\eta$，若 $\eta$ 是使
$(J+\Delta J)\hat x=b+\Delta b$、$\lVert\Delta J\rVert\le\eta\lVert J\rVert$、
$\lVert\Delta b\rVert\le\eta\lVert b\rVert$ 成立的最小相对扰动。
\end{definition}

\begin{theorem}[线性系统扰动界]\label{thm:pf-th-fund-perturb}
设 $Jx=b$、$(J+\Delta J)(x+\Delta x)=b+\Delta b$，记
$\kappa=\kappa(J)$、$\rho=\lVert\Delta J\rVert/\lVert J\rVert$，
且 $\kappa\rho<1$，则
\begin{equation}
\frac{\lVert\Delta x\rVert}{\lVert x\rVert}
\le\frac{\kappa}{1-\kappa\rho}
\Bigl(\frac{\lVert\Delta b\rVert}{\lVert b\rVert}+\rho\Bigr).
\label{eq:pf-th-fund-perturb}
\end{equation}
\end{theorem}

\begin{proof}
两式相减：$J\Delta x=\Delta b-\Delta J\,x-\Delta J\,\Delta x$。
$J+\Delta J$ 非奇异（引理~\ref{lem:pf-th-fund-banach}，
$\lVert J^{-1}\rVert\lVert\Delta J\rVert\le\kappa\rho<1$）。左乘 $J^{-1}$
并取范数：
\begin{equation}
\lVert\Delta x\rVert
\le\lVert J^{-1}\rVert\bigl(\lVert\Delta b\rVert+\lVert\Delta J\rVert\,
\lVert x\rVert\bigr)+\lVert J^{-1}\rVert\lVert\Delta J\rVert\lVert\Delta x\rVert .
\end{equation}
移项除以 $1-\kappa\rho$，再用
$\lVert\Delta b\rVert/\lVert x\rVert
=\lVert J\rVert\,\lVert\Delta b\rVert/(\lVert J\rVert\lVert x\rVert)
\le\lVert J\rVert\,\lVert\Delta b\rVert/\lVert b\rVert$
即得式~\eqref{eq:pf-th-fund-perturb}。
\end{proof}

\begin{remark}[浮点含义与潮流量级]\label{rem:pf-th-fund-float}
部分主元 LU 的向后误差 $\eta\approx u\cdot\rho_{\mathrm{growth}}$
（增长因子，电力雅可比上实践接近 $1$）~\cite{pff:golub-vanloan}。
由式~\eqref{eq:pf-th-fund-perturb}，每步 Newton 修正的相对前向误差量级为
$u\,\kappa(J)$：$\kappa(J)\sim10^{8}$ 时双精度下修正量只剩约 8 位有效数字，
残差无法压到 $10^{-8}$ 以下；$\kappa(J)\sim10^{16}$ 时修正量完全失真。
这给出收敛容差 \fld{tol} 的一个理论下界 $\approx u\,\kappa(J)$：
把容差设到病态系统分辨率之下，表现为"残差平台期"而非继续二次收敛。
\end{remark}

\subsubsection{高 \texorpdfstring{$R/X$}{R/X} 比与解耦失效}

快速解耦潮流（FDPF）的推导假设~\cite{pff:stott-alsac}：
(i) $|G_{ij}|\ll|B_{ij}|$（即 $R/X\ll1$）；(ii) $|\theta_{ij}|$ 小，
$\cos\theta_{ij}\approx1$；(iii) $Q_i\ll B_{ii}V_i^2$。三者使
$N=\partial P/\partial V_m$ 与 $M=\partial Q/\partial\theta$ 近似可忽略，
$\theta$ 与 $V_m$ 解耦。配电网络 $R/X$ 常在 $1$ 附近甚至更大，假设 (i) 失效，
交叉块量级与对角块可比，FDPF 迭代矩阵谱半径逼近或超过 $1$，表现为线性收敛
变慢直至发散。Newton 法的渐近二次收敛\textbf{不}依赖解耦假设，但高 $R/X$
使 $P$--$V$ 强耦合、雅可比失去块对角主导，配合辐射状拓扑（无环对角的
"链式"结构）整体条件数上升；此时专用替代算法（前推回代类，利用辐射树结构
逐支路递推）在鲁棒性上优于直接 Newton~\cite{pff:stott-review}。

\subsubsection{重负荷、弱网与电压崩溃}

随负荷参数 $\lambda\to\lambda^*$（鼻点），命题~\ref{prop:pf-th-fund-fold}
给出 $\sigma_{\min}(J)\to0$；由于 $\lVert J\rVert$ 温和变化，
$\kappa(J)\to\infty$。与之平行的物理图景：短路容量低（弱网）的节点
$\partial V/\partial Q$ 灵敏度大，同样放大 $J^{-1}$ 的范数。重负荷/弱网病态
因此\emph{不是}浮点异常，而是临近静态稳定极限的准确数值反映：
雅可比的左零向量（奇异向量）同时是负荷参数变化对失配影响的方向，
这一对应是连续潮流以弧长参数化翻越鼻点的理论基础
~\cite{pff:ajjarapu-christy,pff:vancutsem-vournas,pff:kundur}。

\begin{gapnote}
当前实现未计算 $\kappa(J)$ 或其标准估计（如 Hager--Higham $1$-范数估计器、
$\sigma_{\min}$ 监控）：代码仅计算行列范数代理
\srcpath{src/power\_flow/newton\_solver.cpp:estimate\_condition\_proxy}，
并以代理越阈触发 NK-GMRES 回退（默认关闭）。鼻点附近的奇异向量分析、
按 $\sigma_{\min}$ 的裕度指标均未实现；电压稳定极限定位由连续潮流求解器
（弧长延拓，实现层见"CPF 章"）承担，而非由病态诊断推导。
\end{gapnote}

\subsection{稀疏技术理论}\label{sec:pf-th-fund-sparsity}

\subsubsection{消元图与填充}

潮流雅可比（固定 PV/PQ 活动集下）具有与 $Y_{bus}$ 图同构的对称稀疏模式：
每个 AC 节点至多四个块。对这类矩阵，稀疏 LU 的成本完全由\emph{消元顺序}
决定，其理论语言是消元图~\cite{pff:george-liu,pff:davis-suitesparse}。

\begin{definition}[矩阵图与消元图]\label{def:pf-th-fund-elimgraph}
对称模式矩阵 $A\in\mathbb R^{n\times n}$ 的图 $G(A)=(V,E)$：
$V=\{1,\dots,n\}$，$(i,j)\in E\iff a_{ij}\neq0$（$i\neq j$）。
依次消去顶点 $v_1,v_2,\dots$：消去 $v$ 时把 $v$ 的当前邻居两两连边
（已有边不重复），再删除 $v$；新产生的边称为\textbf{填充}（fill-in），
对应 LU 因子中 $A$ 原本为零位置上的非零元。
\end{definition}

\begin{proposition}[Parter 规则]\label{prop:pf-th-fund-parter}
以顶点 $v$ 为主元的消元步，在消元图中把 $v$ 的邻居集变成团（clique）：
对任意两个邻居 $i,k$，消元后 $(i,k)$ 为边。
\end{proposition}

\begin{proof}
对称高斯消元一步的 Schur 补更新为
$\tilde a_{ik}=a_{ik}-a_{iv}a_{vv}^{-1}a_{vk}$。
$i,k$ 同为 $v$ 的邻居时 $a_{iv}a_{vk}\neq0$，修正项一般非零，故
$\tilde a_{ik}\neq0$ 无论 $a_{ik}$ 是否为零——已有边保留，缺失处产生填充边。
精确相消（数值对消）是零测集事件，符号分析按一般位置处理。
\end{proof}

\begin{example}[星形网络的排序敏感性]\label{ex:pf-th-fund-star}
中心节点 $1$ 连接叶节点 $2,\dots,5$（如主 slack 汇接多条馈线）。
$A$ 的非零元 $13$ 个（$5$ 对角 $+$ $2\times4$ 非对角）。
\begin{itemize}
  \item 顺序 $(1,2,3,4,5)$：消去中心时其邻居为 $\{2,3,4,5\}$，按
        命题~\ref{prop:pf-th-fund-parter} 形成 $4$ 团，产生
        $\binom{4}{2}=6$ 条填充边；$L+U$ 非零元 $=13+2\times6=25$，
        且后续每步面对稠密团，消元工作量 $\Theta(4^2)$；
  \item 顺序 $(2,3,4,5,1)$：消去任一叶片时其唯一邻居是中心，不产生填充；
        最后消去中心。填充为零，$L+U$ 非零元保持 $13$，工作量 $\Theta(n)$。
\end{itemize}
同一矩阵，仅排序不同，填充与工作量相差一个量级——这就是排序理论要
形式化的现象。
\end{example}

\begin{proposition}[填充路径刻画]\label{prop:pf-th-fund-fillpath}
设 $F=L+U$ 为某消元顺序下的填充矩阵。对 $i>j$，$(i,j)\in G(F)$ 当且仅当
在 $G(A)$ 中存在一条从 $i$ 到 $j$ 的路径，其所有中间顶点的消元序号
小于 $j$（即小于两端序号）。
\end{proposition}

\begin{proof}[证明要点]
对消元步数归纳。基础步：原边显然满足（无中间顶点）。归纳步：消去 $v$
时由命题~\ref{prop:pf-th-fund-parter} 新产生的边 $(i,k)$ 有路径
$i$--$v$--$k$，其中 $v$ 先于 $i,k$ 消元；而既有边 $(i,j)\in G(F)$ 若在
此前各步产生，其见证路径的中间顶点均已被消去、序号小于 $j$。反向：
任给满足条件的路径，沿路径上序号最大的中间顶点归纳分裂为两条更短的
合法路径，最终归结为原边或填充边。严格归纳见 George--Liu
~\cite{pff:george-liu}。
\end{proof}

命题~\ref{prop:pf-th-fund-fillpath} 把"哪些位置会填充"完全归结为图论问题：
填充 $=$ 低序号顶点介导的可达性。排序算法的目标即选消元顺序使这种
可达性最小化。

\subsubsection{最小度与近似最小度排序}

\begin{definition}[最小度（MD）排序]\label{def:pf-th-fund-md}
贪心策略：每步在当前消元图中选度（外部度）最小的顶点消元，平局任意打破；
即 Tinney 方案 2~\cite{pff:tinney-walker}。直观依据：消去度 $d$ 的顶点至多
产生 $\binom{d}{2}$ 条填充、$d^2$ 量级算术运算，局部最小化度近似全局
最小化填充。
\end{definition}

\begin{remark}[AMD 的原理]\label{rem:pf-th-fund-amd}
精确 MD 的瓶颈在于消去顶点后邻居集的团更新与外部度重算
（最坏 $O(n\cdot nnz)$）。近似最小度（AMD）~\cite{pff:amestoy-amd} 以
\emph{商图}表示消元过程：已消元顶点被吸收为"元素"，未消元顶点的邻接
以元素集合表示，外部度用上界近似代替精确计算；配合超变量（结构相同的行
合并）与质量吸收，AMD 在实践中达到接近 MD 的填充质量而开销低一个量级以上，
是 SuiteSparse（KLU/UMFPACK 内部）的默认排序。对规则网孔类网络，嵌套分割
（递归几何/图二分）另有严格界：二维 $n\times n$ 网格上填充
$O(n^2\log n)$、分解工作量 $O(n^3)$~\cite{pff:george-nested}；电力网近似
平面图，AMD 的实际表现通常已接近嵌套分割。
\end{remark}

\subsubsection{填充分析与符号/数值分解分离}

由命题~\ref{prop:pf-th-fund-fillpath}，$L+U$ 的非零模式是 $A$ 的模式与
消元顺序的纯函数，与数值无关。因此稀疏直接法天然分两层：
\textbf{符号分解}（排序 + 模式确定，一次）与\textbf{数值分解}
（填充数值，每次 Newton 迭代）。固定 PV/PQ 活动集与网络拓扑时，雅可比
模式逐迭代不变，符号分解可跨迭代、跨求解复用——这是实现层模式缓存
（JacobianPattern/analyzePattern 一次性构建，迭代内仅数值重分解）的
理论依据。PV$\to$PQ 切换改变行/列结构时传统布局须重建模式，
固定超集布局（所有非参考母线预留 $V_m$ 列与 $Q$ 行）以常量填充冗余换取
模式不变性，其契约见 \srcpath{docs/theory/pv\_pq\_switching\_contract.md}。

\begin{gapnote}
平台\textbf{未自研}排序与块三角分解（BTF）：平衡潮流的排序委托给
KLU 的符号分析（其内部默认为 AMD，并可经环境变量切换 COLAMD/嵌套分割，
见 \srcpath{MIPSolvers/src/engine/kernel/linear\_algebra/linear\_solver.cpp:make\_default\_sparse\_solver}）；
三相混合求解器直接使用 Eigen 的 COLAMDOrdering
（\srcpath{src/power\_flow/three\_phase\_hybrid.cpp} 内 SparseLU 配置）。
本节的 MD/AMD/嵌套分割理论用于理解因子填充来源，不对应平台自有代码；
自定义排序接口亦未暴露。
\end{gapnote}

\subsection{数值算例与基准}\label{sec:pf-th-fund-numexp}

本节全部数字由 \srcpath{tools/pf\_doc\_benchmark.cpp}（模式 \texttt{nr}）在 MATPOWER
标准算例（\srcpath{data/case*.m}）上、经平台默认统一 Newton 求解器（极坐标、KLU 稀疏
分解、平启动、$\mathrm{tol}=10^{-10}$）实测生成，可随代码演进重跑复现。迭代次数与残差
轨迹为机器无关量；求解墙钟时间仅供量级参考（Apple~M4，单线程）。

\subsubsection{收敛的规模弱依赖性}
表~\ref{tab:pf-th-fund-nrconv} 给出 $9\sim300$ 母线算例的收敛数据。Newton 维数
$m=2(n-1)$（固定超集布局：非参考母线一律保留 $\theta$ 与 $V_m$ 两列，见
第~\ref{sec:pf-th-fund-sparsity}~节末）。迭代次数稳定在 $4\sim8$ 次、几乎不随规模增长，
最终残差压到 $10^{-13}$ 量级——正是注记~\ref{rem:pf-th-fund-hypotheses} 所述“平启动、
个位数次迭代收敛”的定量印证。

\begin{table}[H]\centering\footnotesize
\caption{统一 Newton 潮流在 MATPOWER 算例上的收敛数据（平启动，$\mathrm{tol}=10^{-10}$）}
\label{tab:pf-th-fund-nrconv}
\begin{tabular}{@{}lrrrrrr@{}}
\toprule
算例 & 母线数 $n$ & Newton 维数 $m$ & 迭代次数 & 最终残差 & 条件数代理 & PV$\to$PQ 切换\\
\midrule
case9   & 9   & 16  & 5 & $1.7\times10^{-14}$ & $5.9\times10^{2}$ & 0\\
case14  & 14  & 26  & 4 & $9.6\times10^{-15}$ & $1.9\times10^{3}$ & 0\\
case30  & 30  & 58  & 5 & $8.9\times10^{-15}$ & $2.0\times10^{4}$ & 0\\
case57  & 57  & 112 & 4 & $3.4\times10^{-12}$ & $2.2\times10^{4}$ & 0\\
case118 & 118 & 234 & 7 & $8.5\times10^{-14}$ & $1.3\times10^{5}$ & 6\\
case300 & 300 & 598 & 8 & $4.3\times10^{-13}$ & $1.3\times10^{9}$ & 10\\
\bottomrule
\end{tabular}
\end{table}

条件数代理列由 \srcpath{estimate\_condition\_proxy}（行/列 $\infty$-范数之积，\emph{非}真
$\kappa_2$）给出，随规模与负荷压力单调上升。case300 的代理达 $10^{9}$ 量级，但其真解
仍收敛到 $4\times10^{-13}$：若真 $\kappa_2$ 果为 $10^{9}$，由注记~\ref{rem:pf-th-fund-float}
的 $u\,\kappa$ 界残差应止步于 $\sim10^{-7}$，故该行/列范数代理\emph{显著高估}真
$\kappa_2$，只宜作趋势指标而非绝对条件数。

\subsubsection{二次收敛的逐迭代验证}
\begin{example}[case30 的二次收敛率]\label{ex:pf-th-fund-quadconv}
下表列出 case30 每次 Newton 迭代记录的最大失配范数 $r_k$（$k=0$ 为平启动初始失配）
与二次收敛率 $r_k/r_{k-1}^2$。比率在 $r_2,r_3$ 步稳定于 $\approx0.22$（有界），
即 $r_{k+1}\approx0.22\,r_k^2$——定理~\ref{thm:pf-th-fund-quadratic} 的
$\lVert e_{k+1}\rVert\le\beta L\lVert e_k\rVert^2$ 在此对应 $\beta L\approx0.2$。
末步 $r_4$ 触及 $\sim10^{-14}$ 的浮点残差地板，比率失真（$r_3^2\approx9\times10^{-19}$
已低于双精度可分辨范围），与注记~\ref{rem:pf-th-fund-float} 的残差平台期吻合。
\begin{center}\footnotesize
\begin{tabular}{@{}crl@{}}
\toprule
迭代 $k$ & 最大失配 $r_k$ & 二次率 $r_k/r_{k-1}^2$\\
\midrule
0 & $3.93\times10^{-1}$ & —\\
1 & $1.63\times10^{-2}$ & $1.06\times10^{-1}$\\
2 & $6.59\times10^{-5}$ & $2.47\times10^{-1}$\\
3 & $9.57\times10^{-10}$ & $2.20\times10^{-1}$\\
4 & $8.88\times10^{-15}$ & （残差地板）\\
\bottomrule
\end{tabular}
\end{center}
\end{example}

\subsubsection{无功限值切换对 Newton 轨迹的扰动}
\begin{example}[case300 的两段式收敛]\label{ex:pf-th-fund-qswitch}
case300 逐迭代记录的最大失配为（竖线标注一次 PV$\to$PQ 切换处）
\begin{center}\footnotesize
$9.27,\ 2.59,\ 3.54{\times}10^{-1},\ 8.45{\times}10^{-3},\ 4.65{\times}10^{-6},\
1.38{\times}10^{-12}\ \big\|\ 3.00{\times}10^{-3},\ 2.04{\times}10^{-7},\
4.26{\times}10^{-13}.$
\end{center}
轨迹呈三个阶段：(i) \emph{大修正区}（$9.27\to2.59\to0.35$，比率非二次），
对应初值远离解、超出定理~\ref{thm:pf-th-fund-quadratic} 的二次收敛邻域，
线搜索与步长限幅等全局化机制主导（第~\ref{sec:pf-th-fund-newton}~节收敛域讨论）；
(ii) \emph{二次收敛尾}（$8.45{\times}10^{-3}\to4.65{\times}10^{-6}\to1.38{\times}10^{-12}$，
残差近似逐步平方，比率 $\approx0.065$）；(iii) 竖线处一次 \textbf{无功限值 PV$\to$PQ 切换}
改变活动集、重新注入 $\sim3\times10^{-3}$ 的失配，之后再次二次收敛到 $4.3\times10^{-13}$。
这段真实轨迹同时印证三件事：远离解的全局化行为、近解的二次收敛，以及
式~\eqref{eq:pf-th-fund-qlimit} 的无功限值活动集切换如何以“重启 Newton”的形式表现——
纯 Newton 序列被活动集变更打断，正是注记~\ref{rem:pf-th-fund-affine} 所述“限值切换
破坏纯 Newton 序列”的实测样貌。
\end{example}

\subsection{与实现的对应}\label{sec:pf-th-fund-implmap}
\begin{center}\footnotesize
\begin{longtable}{@{}L{6.8cm}L{8.6cm}@{}}
\toprule
理论条目 & 实现状态\\
\midrule
\endfirsthead
\toprule
理论条目 & 实现状态\\
\midrule
\endhead
$\pi$ 等值支路组装（串联导纳 + 两端充电 $b_c/2$ + 节点/shunt 并联） &
\implfull{src/power\_flow/admittance\_builder.cpp:build\_admittance\_matrix}\\
复变比 $t=ae^{\mathrm j\varphi}$ 两端口（式~\eqref{eq:pf-th-fund-tport}，
含移相器非对称） &
\implfull{src/power\_flow/admittance\_builder.cpp:build\_admittance\_matrix}（幅值 \fld{tap}、移相角 \fld{shift\_deg} 合成复数 \fld{tap}，$Y_{ft}=-y_s/\overline{t}$、$Y_{tf}=-y_s/t$）\\
节点功率方程极坐标注入（推论~\ref{cor:pf-th-fund-polar}） &
\implfull{src/power\_flow/ac\_kernel.cpp:evaluate\_ac\_kernel\_serial}\\
PQ/PV/V$\theta$ 分类与方程-未知数闭合（命题~\ref{prop:pf-th-fund-counting}） &
\implfull{src/power\_flow/newton\_solver.cpp:build\_jacobian\_context}（闭合计数校核见同文件 scan\_empty\_rows\_cols）\\
多岛各自参考节点（注记~\ref{rem:pf-th-fund-slack}） &
\implfull{src/power\_flow/newton\_solver.cpp:build\_jacobian\_context}（每个孤立 AC 区域各自保留参考；DC 岛参考另有规划函数）\\
无功限值活动集策略（式~\eqref{eq:pf-th-fund-qlimit}） &
\implfull{src/power\_flow/newton\_solver.cpp:solve}（收敛点批量进入 + 迟滞恢复 + 单次恢复审计；契约见 docs/theory/pv\_pq\_switching\_contract.md）\\
中值互补/CHKS 光滑化的半光滑 Newton &
\implfull{src/power\_flow/newton\_solver.cpp:solve}（由 \fld{enable\_semi\_smooth\_newton} 启用，$\mu$ 延拓见 robust\_nonlinear\_options）\\
多解枚举与全解搜索（注记~\ref{rem:pf-th-fund-multiplicity}） &
\implnone（HELM 求解器 helm\_solver.cpp:HelmSolver::solve 仅返回与平启动连通的物理解，非枚举）\\
可解性的压缩充分条件（定理~\ref{thm:pf-th-fund-solvability}，
式~\eqref{eq:pf-th-fund-solvcond}） &
\implnone（无求解前可解性预判；以收敛与否事后判定）\\
Newton 迭代与局部二次收敛（定理~\ref{thm:pf-th-fund-quadratic}） &
\implfull{src/power\_flow/newton\_solver.cpp:solve}（实现层转写见本手册"统一 Newton 章"）\\
Kantorovich/$\alpha$ 定量收敛证书 &
\implnone（以线搜索/信赖域/LM/PTC 启发式全局化回退链替代）\\
折点雅可比奇异与鼻点定位（命题~\ref{prop:pf-th-fund-fold}） &
\implfull{src/power\_flow/voltage\_stability.cpp:solve}（CpfSolver 弧长延拓翻越鼻点）\\
结构性奇异诊断（注记~\ref{rem:pf-th-fund-singularity}） &
\implpart{空行/空列扫描与方程闭合校核已实现（scan\_empty\_rows\_cols）；诊断不区分参数性/结构性成因，也不输出奇异向量}\\
条件数监控与扰动界（定理~\ref{thm:pf-th-fund-perturb}） &
\implpart{仅行列范数代理 estimate\_condition\_proxy 与可选 NK-GMRES 回退；无 Hager--Higham 或 $\sigma_{\min}$ 估计}\\
仿射缩放 $D_F,D_x$（注记~\ref{rem:pf-th-fund-affine}） &
\implfull{src/power\_flow/nonlinear\_scaling.cpp:build\_nonlinear\_scaling}\\
高 $R/X$ 配网的前推回代替代 &
\implfull{src/power\_flow/distribution\_power\_flow.cpp:solve\_distribution\_pf}（辐射状配网 BFS；FDPF 本体见 fdpf\_solver.cpp:FDPFSolver::solve）\\
填充路径/消元图分析与符号-数值分解分离 &
\implfull{src/power\_flow/jacobian\_pattern.cpp:build\_jacobian\_pattern}（模式一次构建、迭代内仅数值累加；符号分析缓存见 newton\_solver.cpp:solve 的 PatternCache）\\
MD/AMD/嵌套分割排序算法 &
\implpart{平台未自研：KLU 默认 AMD（MIPSolvers 端 make\_default\_sparse\_solver），三相混合用 Eigen COLAMDOrdering；无自定义排序接口}\\
解析雅可比的差分校验工具 &
\implfull{src/power\_flow/jacobian\_check.cpp:check\_jacobian}（中心差分 $h=10^{-6}$）\\
直角坐标失配的精确二次结构与步长恒等式（定理~\ref{thm:pf-th-fund-quadmap}，
式~\eqref{eq:pf-th-fund-exactstep}） &
\implnone（主迭代用极坐标功率失配，三角项使 $D^3F\ne0$）\\
电流注入 Newton 的常量非对角块（注记~\ref{rem:pf-th-fund-ci}） &
\implpart{仅三相 Newton 含定点电流注入辅助
three\_phase\_nr.cpp:evaluate\_fixed\_point\_current\_injections；单相/正序主迭代为极坐标功率失配}\\
Iwamoto 最优乘子闭式步长（命题~\ref{prop:pf-th-fund-iwamoto}，
式~\eqref{eq:pf-th-fund-iwamoto}） &
\implnone（全局化改用 Armijo/非单调线搜索、LM、PTC，见第~\ref{sec:pf-th-glob}~章）\\
收敛性数值基准（表~\ref{tab:pf-th-fund-nrconv}、例~\ref{ex:pf-th-fund-quadconv}、
\ref{ex:pf-th-fund-qswitch}） &
\implfull{tools/pf\_doc\_benchmark.cpp}（模式 nr；调用 api/hacdcpf.hpp:solve\_power\_flow）\\
\bottomrule
\end{longtable}
\end{center}

\subsection{参考文献}
{\renewcommand{\section}[2]{}%
\begin{thebibliography}{9}\small
\bibitem{pff:stagg-elabiad} G.~W. Stagg, A.~H. El-Abiad, \emph{Computer Methods
in Power System Analysis}, McGraw-Hill, 1968.
\bibitem{pff:arrillaga-watson} J.~Arrillaga, N.~R. Watson, \emph{Computer
Modelling of Electrical Power Systems}, 2nd ed., Wiley, 2001.
\bibitem{pff:tinney-hart} W.~F. Tinney, C.~E. Hart, Power flow solution by
Newton's method, \emph{IEEE Transactions on Power Apparatus and Systems},
vol. PAS-86, no. 11, pp. 1449--1460, 1967.
\bibitem{pff:stott-alsac} B.~Stott, O.~Alsa\c{c}, Fast decoupled load flow,
\emph{IEEE Transactions on Power Apparatus and Systems}, vol. PAS-93, no. 3,
pp. 859--869, 1974.
\bibitem{pff:stott-review} B.~Stott, Review of load-flow calculation methods,
\emph{Proceedings of the IEEE}, vol. 62, no. 7, pp. 916--929, 1974.
\bibitem{pff:ortega-rheinboldt} J.~M. Ortega, W.~C. Rheinboldt,
\emph{Iterative Solution of Nonlinear Equations in Several Variables},
Academic Press, 1970.
\bibitem{pff:nocedal-wright} J.~Nocedal, S.~J. Wright, \emph{Numerical
Optimization}, 2nd ed., Springer, 2006.
\bibitem{pff:golub-vanloan} G.~H. Golub, C.~F. Van Loan, \emph{Matrix
Computations}, 4th ed., Johns Hopkins University Press, 2013.
\bibitem{pff:ajjarapu-christy} V.~Ajjarapu, C.~Christy, The continuation power
flow: a tool for steady state voltage stability analysis, \emph{IEEE
Transactions on Power Systems}, vol. 7, no. 1, pp. 416--423, 1992.
\bibitem{pff:vancutsem-vournas} T.~Van Cutsem, C.~Vournas, \emph{Voltage
Stability of Electric Power Systems}, Kluwer Academic Publishers, 1998.
\bibitem{pff:kundur} P.~Kundur, \emph{Power System Stability and Control},
McGraw-Hill, 1994.
\bibitem{pff:hintermuller} M.~Hinterm\"uller, K.~Ito, K.~Kunisch, The
primal-dual active set strategy as a semismooth Newton method, \emph{SIAM
Journal on Optimization}, vol. 13, no. 3, pp. 865--888, 2002.
\bibitem{pff:tinney-walker} W.~F. Tinney, J.~W. Walker, Direct solutions of
sparse network equations by optimally ordered triangular factorization,
\emph{Proceedings of the IEEE}, vol. 55, no. 11, pp. 1801--1809, 1967.
\bibitem{pff:george-liu} A.~George, J.~W.~H. Liu, The evolution of the minimum
degree ordering algorithm, \emph{SIAM Review}, vol. 31, no. 1, pp. 1--19, 1989.
\bibitem{pff:amestoy-amd} P.~R. Amestoy, T.~A. Davis, I.~S. Duff, An approximate
minimum degree ordering algorithm, \emph{SIAM Journal on Matrix Analysis and
Applications}, vol. 17, no. 4, pp. 886--905, 1996.
\bibitem{pff:george-nested} A.~George, Nested dissection of a regular finite
element mesh, \emph{SIAM Journal on Numerical Analysis}, vol. 10, no. 2,
pp. 345--363, 1973.
\bibitem{pff:davis-suitesparse} T.~A. Davis, \emph{Direct Methods for Sparse
Linear Systems}, SIAM, 2006.
\bibitem{pff:iwamoto} S.~Iwamoto and Y.~Tamura, A load flow calculation method
for ill-conditioned power systems, \emph{IEEE Transactions on Power Apparatus
and Systems}, vol. PAS-100, no. 4, pp. 1736--1743, 1981.
\bibitem{pff:dacosta} V.~M. da Costa, N.~Martins, and J.~L.~R. Pereira,
Developments in the Newton Raphson power flow formulation based on current
injections, \emph{IEEE Transactions on Power Systems}, vol. 14, no. 4,
pp. 1320--1326, 1999.
\bibitem{pff:bolognani} S.~Bolognani and S.~Zampieri, On the existence and linear
approximation of the power flow solution in power distribution networks,
\emph{IEEE Transactions on Power Systems}, vol. 31, no. 1, pp. 163--172, 2016.
\bibitem{pff:simpsonporco} J.~W. Simpson-Porco, F.~D\"orfler, and F.~Bullo,
Voltage collapse in complex power grids, \emph{Nature Communications}, vol. 7,
art. 10790, 2016.
\end{thebibliography}}

% theory_classic_methods.tex —— 潮流计算 通用理论：经典方法（GS/Z-bus 谱系、FDPF 解耦、直流潮流、DistFlow 与前推回代、分布式松弛与自适应孤岛）
% 本文件为理论层章节，遵循 docs/modules/THEORY_WRITING_GUIDE.md。

\section{经典潮流方法理论：不动点迭代、解耦与线性化}\label{sec:pf-th-cls}

\begin{theorynote}
本章为通用数学理论，按经典教材与原始文献体系撰写，覆盖五组主题：
(i) Gauss--Seidel 与阻抗阵法的历史定位及其作为不动点迭代的数学结构；
(ii) 快速解耦潮流（FDPF）的解耦物理原理推导、XB/BX 变体差异与收敛性分析；
(iii) 直流潮流从 AC 方程经三条假设到 $P=B\theta$ 的严格推导、线性化误差界、损耗近似，
及基于直流模型的转移分布因子（PTDF/ISF）与线路开断分布因子（LODF）；
(iv) 辐射状配电网的 DistFlow 支路潮流模型、前推回代的不动点--压缩映射收敛理论
与高 $R/X$ 弱网特性；
(v) 分布式松弛（含方程内精确加边 Newton 系统）与自适应孤岛的理论定位；
(vi) 以平台求解器实测数据给出的 FDPF、直流与前推回代收敛性数值基准。
Newton 法本体、全局化策略与电压稳定/换流器相关理论各有专章，本章不复述；
本章公式不要求逐式对应代码。与平台实现（\srcpath{src/power\_flow/}）的对应关系
见本章末"与实现的对应"表，超出实现的内容以"未实现扩展说明"框就地标记。
\end{theorynote}

\subsection{记号与约定}\label{sec:pf-th-cls-notation}
\begin{itemize}
  \item $n$：AC 母线数；母线复电压 $V_i=V_{m,i}\,e^{j\theta_i}$（极坐标），
        $\theta_{ij}=\theta_i-\theta_j$；
  \item $Y_{bus}=G+jB$：节点导纳阵，$Y_{ij}=G_{ij}+jB_{ij}$；支路 $k$（$i\!\to\!j$）
        阻抗 $z_k=r_k+jx_k$，串联导纳 $y_k=g_k+jb_k=1/z_k$（感应性支路 $b_k<0$）；
  \item $P_i,Q_i$：母线 $i$ 注入有功/无功（发电减负荷，pu）；上标
        $\mathrm{spec}/\mathrm{calc}$ 区分指定值与计算值；
        失配 $\Delta P_i=P_i^{\mathrm{spec}}-P_i^{\mathrm{calc}}$，
        $\Delta Q_i$ 同理；
  \item $B',B''\in\mathbb R^{n\times n}$：FDPF 的两个常量近似阵
        （本章取\textbf{ Laplacian 口径}：对角元为正，迭代方程写作
        $B'\Delta\theta=\Delta P/V$、$B''\Delta V=\Delta Q/V$；
        部分教材采用反号口径 $-B'$，两者仅差整体符号约定）；
  \item $\rho(\cdot)$：谱半径；$\lVert\cdot\rVert_\infty$：行和范数（向量时为最大模）；
  \item 配网一节中，$\mathcal T$ 为以根（slack）为源的辐射树，$\mathcal P(j)$ 为
        根到母线 $j$ 的唯一路径，$\mathrm{sub}(j)$ 为 $j$ 的子树（含 $j$ 自身），
        $\ell_{ij}=|I_{ij}|^2$ 为支路电流幅值平方。
\end{itemize}

\subsection{方法谱系：不动点迭代视角下的经典方法}\label{sec:pf-th-cls-lineage}

潮流计算的数字化历史几乎与数字计算机同步。Ward 与 Hale 于 1956 年给出首个
实用的导纳阵迭代法~\cite{cls:wardhale}；随后出现了以阻抗阵 $Z_{bus}=Y_{bus}^{-1}$
为核心的一类方法（Z-bus 法），其系统整理见 Stagg 与 El-Abiad~\cite{cls:stagg}；
Tinney 与 Hart 于 1967 年确立 Newton 法加稀疏排序分解的主流地位~\cite{cls:tinney}；
Stott 与 Alsa\c{c} 于 1974 年提出快速解耦潮流~\cite{cls:stott}。Stott 的
综述~\cite{cls:stottreview} 对这一谱系有经典评述。本节用统一的\emph{不动点迭代}
语言重述前 Newton 时代方法，以看清其数学结构与被取代的原因。

\begin{definition}[Y-bus Gauss--Seidel 潮流迭代]\label{sec:pf-th-cls-def-gs}
把节点方程 $I=Y_{bus}V$ 与注入约束 $S_i^{\mathrm{spec}}=V_i\overline{I_i}$
联立，对每个 PQ 母线解出
\begin{equation}
V_i^{(k+1)}=\frac{1}{Y_{ii}}\Biggl(\frac{\overline{S_i^{\mathrm{spec}}}}{\overline{V_i^{(k)}}}
-\sum_{j<i}Y_{ij}V_j^{(k+1)}-\sum_{j>i}Y_{ij}V_j^{(k)}\Biggr),
\label{sec:pf-th-cls-eq:gs}
\end{equation}
即把潮流问题写成电压向量上的不动点方程 $V=T(V)$ 并做 Gauss--Seidel 型分量迭代。
PV 母线先由 $Q$ 方程解出无功再校正幅值。Z-bus 法把式~\eqref{sec:pf-th-cls-eq:gs}
右端的邻接求和换成 $Z_{bus}$ 与注入电流的乘积，收敛性更好但每步需作用稠密
$Z_{bus}$（或其隐式因子）。
\end{definition}

式~\eqref{sec:pf-th-cls-eq:gs} 中电流项 $\overline{S}/\overline{V}$ 使 $T$ 成为
\emph{非线性}映射；其收敛性分析的第一步是看线性化的骨架——对线性方程组
$YV=I$ 的 Gauss--Seidel 迭代，下述经典结论给出充分收敛条件。

\begin{theorem}[严格对角占优线性系统的 Gauss--Seidel 收敛性]\label{sec:pf-th-cls-thm-gs}
设 $A\in\mathbb C^{n\times n}$ 严格对角占优（每行
$|a_{ii}|>\sum_{j\ne i}|a_{ij}|$），则对任意初值，求解 $Ax=b$ 的
Gauss--Seidel 迭代收敛，迭代矩阵 $T_{GS}=-(D+L)^{-1}U$ 满足 $\rho(T_{GS})<1$
（其中 $A=D+L+U$ 为对角/严格下三角/严格上三角分解）。
\end{theorem}

\begin{proof}
设 $\lambda$ 为 $T_{GS}$ 的任一特征值，$x\ne0$ 为相应特征向量：
$-(D+L)^{-1}Ux=\lambda x$，即 $M(\lambda)x=0$，其中
$M(\lambda)=\lambda(D+L)+U$。$M(\lambda)$ 的元素为
$m_{ii}=\lambda a_{ii}$，$m_{ij}=\lambda a_{ij}$（$j<i$），$m_{ij}=a_{ij}$
（$j>i$）。设 $|\lambda|\ge1$，则对每行 $i$，
\begin{equation*}
|m_{ii}|=|\lambda|\,|a_{ii}|>|\lambda|\sum_{j\ne i}|a_{ij}|
\ge \sum_{j<i}|\lambda a_{ij}|+\sum_{j>i}|a_{ij}|=\sum_{j\ne i}|m_{ij}|,
\end{equation*}
即 $M(\lambda)$ 仍严格对角占优。严格对角占优矩阵非奇异（否则由
Gershgorin 圆盘定理存在含原点的圆盘，与对角占优矛盾），故 $M(\lambda)x=0$
只有零解，与 $x\ne0$ 矛盾。因此 $T_{GS}$ 的所有特征值满足 $|\lambda|<1$，
即 $\rho(T_{GS})<1$，迭代误差 $e^{(k)}=T_{GS}^{\,k}e^{(0)}\to0$。
\end{proof}

\begin{remark}[从线性骨架到非线性潮流]\label{sec:pf-th-cls-rem-gslim}
电力网络的 $Y_{bus}$ 通常弱对角占优而非严格占优（联络线、充电电容、三绕组
变压器等都会削弱占优性），且式~\eqref{sec:pf-th-cls-eq:gs} 的非线性项
$\overline{S}/\overline{V}$ 在电压偏低时放大局部 Lipschitz 常数，使收敛域收缩、
收敛速度随系统规模恶化。Gauss--Seidel 类方法只具有\emph{线性}（几何）收敛速度，
且渐近常数接近 1；这正是 Newton 法（局部二次收敛）~\cite{cls:tinney} 与其
近似版 FDPF~\cite{cls:stott} 取而代之的根本原因。作为不动点方法的价值在于结构
透明：它把"潮流 = 求 $V=T(V)$ 的不动点"这一观点留给了后来的前推回代
（第~\ref{sec:pf-th-cls-distflow}~节）与 HELM 等方法。
\end{remark}

\begin{gapnote}
Y-bus Gauss--Seidel 与 Z-bus 阻抗阵法在平台 \srcpath{src/power\_flow/} 中
\textbf{均无实现}（求解器谱系以 Newton 核、FDPF、直流、BFS 为主，
见实现层"架构总览"章）；本节谱系内容仅作理论定位。
\end{gapnote}

\subsection{快速解耦潮流：解耦原理推导与 XB/BX 变体}\label{sec:pf-th-cls-fdpf}

\subsubsection{弱耦合假设的数学化}

极坐标 Newton 校正方程按 $(\theta,V)$ 分块（推导见本手册 Newton 理论章）：
\begin{equation}
\begin{pmatrix}\Delta P\\ \Delta Q\end{pmatrix}
=
\begin{pmatrix} H & N\\ J & L\end{pmatrix}
\begin{pmatrix}\Delta\theta\\ \Delta V\end{pmatrix},
\label{sec:pf-th-cls-eq:nrblock}
\end{equation}
其中对 $i\ne j$（记 $B_{ij}=\operatorname{Im}Y_{ij}$），
\begin{align}
H_{ij}&=\frac{\partial P_i}{\partial\theta_j}
 =V_iV_j\bigl(G_{ij}\sin\theta_{ij}-B_{ij}\cos\theta_{ij}\bigr),&
N_{ij}&=V_j\frac{\partial P_i}{\partial V_j}
 =V_iV_j\bigl(G_{ij}\cos\theta_{ij}+B_{ij}\sin\theta_{ij}\bigr),\notag\\
J_{ij}&=\frac{\partial Q_i}{\partial\theta_j}
 =-V_iV_j\bigl(G_{ij}\cos\theta_{ij}+B_{ij}\sin\theta_{ij}\bigr),&
L_{ij}&=V_j\frac{\partial Q_i}{\partial V_j}
 =V_iV_j\bigl(G_{ij}\sin\theta_{ij}-B_{ij}\cos\theta_{ij}\bigr).
\label{sec:pf-th-cls-eq:blocks}
\end{align}

\begin{definition}[FDPF 三条物理解耦假设]\label{sec:pf-th-cls-def-assump}
\textbf{(A1) 高 $X/R$}：网络支路满足 $|G_{ij}|\ll|B_{ij}|$；
\textbf{(A2) 小角度差}：运行点处 $|\theta_{ij}|$ 小，
$\sin\theta_{ij}\approx0$、$\cos\theta_{ij}\approx1$；
\textbf{(A3) 无功通量小}：对角处 $|Q_i|\ll|B_{ii}|V_i^2$
（即流经的无功相对母线短路电纳口径为小量）。
\end{definition}

在 (A1)--(A2) 下，式~\eqref{sec:pf-th-cls-eq:blocks} 的非对角元化为
\begin{equation}
H_{ij}\approx L_{ij}\approx -V_iV_jB_{ij},\qquad
N_{ij}\approx -J_{ij}\approx V_iV_jG_{ij}\approx0,
\label{sec:pf-th-cls-eq:offdiag}
\end{equation}
即 $P$--$\theta$ 与 $Q$--$V$ 两块\textbf{解耦}，交叉块 $N,J$ 消失。
对角元 $H_{ii}=-Q_i-B_{ii}V_i^2$、$L_{ii}=Q_i-B_{ii}V_i^2$ 在 (A3) 下
均 $\approx -B_{ii}V_i^2$。于是两块校正方程可统一写成
\begin{equation}
\Delta P\approx -\,\mathrm{diag}(V)\,B'\,\mathrm{diag}(V)\,\Delta\theta,
\qquad
\Delta Q\approx -\,\mathrm{diag}(V)\,B''\,\mathrm{diag}(V)\,\Delta V,
\end{equation}
其中 $B',B''$ 是由 $-B$ 派生的 Laplacian 型常量阵（对角为正）。
两端除以 $V_i$ 并把右端残留的幅值因子取为 $1$（即再叠加一条
$V\approx1$ 的\textbf{(A4)}），得 Stott--Alsa\c{c} 标准形~\cite{cls:stott}
\begin{equation}
\boxed{\;B'\,\Delta\theta=\frac{\Delta P}{V},\qquad B''\,\Delta V=\frac{\Delta Q}{V}\;}
\label{sec:pf-th-cls-eq:fdpf}
\end{equation}
两阵\textbf{常量化}，可在迭代前一次因子分解，每次迭代只做失配计算与两次
稀疏回代；收敛判据一般取
$\max(\lVert\Delta P/V\rVert_\infty,\lVert\Delta Q/V\rVert_\infty)<\varepsilon$。

\subsubsection{XB 与 BX 变体：变压器分接头归属的差异}

两条假设链在"如何处理变压器变比 $t$、线路充电与母线并联"上留下自由度，
由此产生两个经典变体~\cite{cls:stott,cls:vanamerongen}：

\begin{definition}[XB 与 BX 方案]\label{sec:pf-th-cls-def-xbbx}
\textbf{XB 方案}（Stott--Alsa\c{c} 原始版，亦称 FDXB）：
$B'$ 由纯串联电纳 $1/x$ 装配——忽略电阻、充电、母线并联，\textbf{变比取 $1$}、
相移取 $0$；$B''$ \textbf{保留}变比、充电与母线并联的影响（仍忽略电阻与相移）。
\textbf{BX 方案}：分接头归属对调——$B'$ \textbf{保留}变比（含移相的注入等效），
$B''$ 中\textbf{变比取 $1$}；其余构造相同。
\end{definition}

\begin{remark}[变体选择的经验结论]\label{sec:pf-th-cls-rem-xbbx}
van Amerongen~\cite{cls:vanamerongen} 的"通用版"FDPF 与
Monticelli--Garcia--Saavedra~\cite{cls:monticelli} 对假设链的系统推导与测试
表明：XB 方案在常规输电网（$R/X$ 小、分接头接近 1）上收敛最快；当网络
$R/X$ 偏高或分接头远离标称时，BX 方案的收敛域与稳健性明显更好。物理原因：
分接头同时影响有功与无功通道，把它放进 $B'$（BX）使 $P$--$\theta$ 子迭代
吸收了主要的对角缩放效应；而 $R/X$ 大直接破坏假设 (A1)，两个方案都会退化，
此时应回退全 Newton 或改用对 $R/X$ 不敏感的方法（配网见
第~\ref{sec:pf-th-cls-distflow}~节）。
\end{remark}

\subsubsection{收敛性：近似 Newton 的线性收敛定理}

FDPF 不再是精确 Newton，而是\textbf{常量近似雅可比的拟 Newton 迭代}
$x^{k+1}=x^k-A^{-1}F(x^k)$，其中 $A=\mathrm{diag}(-V B' V,\,-V B'' V)\big|_{V=1}$。
其收敛性由下述 Ostrowski 型定理刻画（一般形式见~\cite{cls:ortega}）。

\begin{theorem}[常量近似雅可比迭代的局部线性收敛]\label{sec:pf-th-cls-thm-fdpf}
设 $F:\mathbb R^m\to\mathbb R^m$ 在解 $x^*$ 的邻域内二阶连续可微，
$F(x^*)=0$，$J(x^*)$ 非奇异；取常量非奇异阵 $A$。若
\begin{equation}
\sigma:=\rho\bigl(I-A^{-1}J(x^*)\bigr)<1,
\end{equation}
则存在 $x^*$ 的邻域，使迭代 $x^{k+1}=x^k-A^{-1}F(x^k)$ 对其中任意初值收敛到
$x^*$，且收敛是线性的：对任意 $\varepsilon>0$ 存在范数与常数 $C$ 使
$\lVert x^{k+1}-x^*\rVert\le(\sigma+\varepsilon)\lVert x^k-x^*\rVert
+C\lVert x^k-x^*\rVert^2$。
\end{theorem}

\begin{proof}
记 $e_k=x^k-x^*$，$M=I-A^{-1}J(x^*)$。在 $x^*$ 处 Taylor 展开：
$F(x^k)=J(x^*)e_k+r_k$，其中 $\lVert r_k\rVert\le C_2\lVert e_k\rVert^2$
（$F$ 二阶连续可微）。于是
\begin{equation}
e_{k+1}=e_k-A^{-1}\bigl(J(x^*)e_k+r_k\bigr)=Me_k-A^{-1}r_k.
\label{sec:pf-th-cls-eq:errrec}
\end{equation}
由谱半径的连续性，对任意 $\varepsilon>0$ 可取 $\mathbb R^m$ 上的范数使
$\lVert M\rVert\le\sigma+\varepsilon=:\eta<1$（取 $\varepsilon<(1-\sigma)/2$）。
式~\eqref{sec:pf-th-cls-eq:errrec} 给出
$\lVert e_{k+1}\rVert\le\eta\lVert e_k\rVert+C_3\lVert e_k\rVert^2
=\bigl(\eta+C_3\lVert e_k\rVert\bigr)\lVert e_k\rVert$。
取 $\delta>0$ 使 $\eta+C_3\delta<1$；则当 $\lVert e_0\rVert<\delta$ 时归纳得
$\lVert e_k\rVert\le(\eta+C_3\delta)^k\lVert e_0\rVert\to0$，且渐近收敛因子
不超过 $\eta=\sigma+\varepsilon$。
\end{proof}

\begin{remark}[该定理对 FDPF 的实际含义]\label{sec:pf-th-cls-rem-fdpfcond}
条件 $\rho(I-A^{-1}J(x^*))<1$ 正是解耦假设 (A1)--(A4) 的\textbf{定量形式}：
假设越接近成立，$A$ 越逼近 $J(x^*)$，$\sigma$ 越接近 $0$、几何收敛越快；
$R/X$ 升高或角度差拉大时 $\sigma$ 趋近甚至超过 1，表现为迭代缓慢或发散。
与 Newton 的局部二次收敛相比，FDPF 用"线性收敛 + 每步 $O(\mathrm{nnz})$
回代 + 一次分解"换取单步成本的数量级下降，在适定输电网上的实践表现是
$2\sim5$ 次迭代达到工程容差~\cite{cls:stott,cls:monticelli}。同理，FDPF
\textbf{没有} Newton 的回退链/限值切换等全局化机制，其收敛完全依赖初值落在
定理~\ref{sec:pf-th-cls-thm-fdpf} 的邻域内。
\end{remark}

\begin{gapnote}
平台 FDPF（\srcpath{src/power\_flow/fdpf\_solver.cpp}）实现的是 \textbf{XB 方案}
（$B'$ 变比取 1、忽略电阻/充电/并联，$B''$ 保留变比与充电/并联，构型由
\srcpath{make\_bp\_flags}/\srcpath{make\_bpp\_flags} 控制）；\textbf{BX 方案未实现}。
实现层手册称该求解器为 FDXB，与定义~\ref{sec:pf-th-cls-def-xbbx} 的命名一致。
\end{gapnote}

\subsection{直流潮流：三条假设的严格推导与误差分析}\label{sec:pf-th-cls-dc}

\subsubsection{从 AC 支路方程到 $P=B\theta$}

支路 $k$（$i\!\to\!j$，串联导纳 $y_k=g_k+jb_k$，忽略充电）from 端有功潮流的
精确 AC 表达式为
\begin{equation}
P_{ij}=V_i^2 g_k-V_iV_j\bigl(g_k\cos\theta_{ij}+b_k\sin\theta_{ij}\bigr).
\label{sec:pf-th-cls-eq:acbranch}
\end{equation}

\begin{definition}[直流潮流三条假设]\label{sec:pf-th-cls-def-dc}
\textbf{(D1) 忽略电阻}：$r_k=0$，故 $g_k=0$、$b_k=-1/x_k$；
\textbf{(D2) 单位电压}：$V_i\equiv1$（标幺）；
\textbf{(D3) 小角度}：$\sin\theta_{ij}\approx\theta_{ij}$。
\end{definition}

逐条施加于式~\eqref{sec:pf-th-cls-eq:acbranch}：(D1) 消去 $g_k$ 项；
(D2) 置 $V_i=V_j=1$；得 $P_{ij}=-b_k\sin\theta_{ij}=\sin\theta_{ij}/x_k$；
(D3) 给出线性支路方程
\begin{equation}
P_{ij}=\frac{\theta_i-\theta_j}{x_k}.
\label{sec:pf-th-cls-eq:dcbranch}
\end{equation}
对全部支路求和得节点平衡 $\sum_j P_{ij}=P_i$，写成矩阵形式即
\begin{equation}
\boxed{\;B\,\theta=P\;}
\label{sec:pf-th-cls-eq:dcpf}
\end{equation}
其中 $B$ 为以 $1/x_k$ 为权的\textbf{加权图 Laplacian}：$B_{ij}=-1/x_k$（$i\ne j$，
有支路相连时），$B_{ii}=\sum_{k\ni i}1/x_k$。剔除 slack 行/列并把 slack 相角
（通常取 $0$）的贡献移到右端，得降阶方程
$B_{\mathrm{red}}\theta_{\mathrm{red}}=P_{\mathrm{red}}-B_{\mathrm{red,slack}}\theta_{\mathrm{slack}}$。
移相变压器（相移 $\varphi_k$）在式~\eqref{sec:pf-th-cls-eq:dcbranch} 中表现为
$\theta_{ij}\to\theta_{ij}-\varphi_k$，装配时等价为 from/to 端一对反向注入
$\mp b_k\varphi_k$。支路潮流回算
$P_{f,k}=(\theta_i-\theta_j-\varphi_k)/x_k$。

\subsubsection{可解性：接地 Laplacian 的正定性}

\begin{theorem}[降阶 $B$ 阵的对称正定性]\label{sec:pf-th-cls-thm-dc}
设 AC 网络连通，全部在运支路 $x_k>0$，slack 集合非空。则剔除全部 slack
行/列后的降阶阵 $B_{\mathrm{red}}$ 对称正定，直流潮流方程
\eqref{sec:pf-th-cls-eq:dcpf} 有唯一解 $\theta_{\mathrm{red}}$。
\end{theorem}

\begin{proof}
对称性由构造显然。对任意 $\theta\in\mathbb R^{n}$，二次型可写成支路平方和：
\begin{equation}
\theta^{\top}B\theta=\sum_{k=(i,j)}\frac{1}{x_k}\bigl(\theta_i-\theta_j\bigr)^2\ge0.
\label{sec:pf-th-cls-eq:quad}
\end{equation}
若 $\theta^{\top}B\theta=0$，则每条支路两端相角相等；网络连通推出
$\theta$ 为常向量。今限制到非 slack 子空间：若
$\theta_{\mathrm{red}}^{\top}B_{\mathrm{red}}\theta_{\mathrm{red}}=0$，把
$\theta_{\mathrm{red}}$ 以 slack 分量取 $0$ 补全为 $\theta$，由
式~\eqref{sec:pf-th-cls-eq:quad} 仍有 $\theta^{\top}B\theta=0$（slack 相角取 0
时含 slack 的支路项为 $(\theta_i-0)^2$，不变号），故 $\theta$ 为常向量；
而 slack 分量为 $0$，所以常数为 $0$，即 $\theta_{\mathrm{red}}=0$。
因此 $B_{\mathrm{red}}$ 正定，必非奇异，解存在唯一。
\end{proof}

\begin{remark}\label{sec:pf-th-cls-rem-dcslack}
定理~\ref{sec:pf-th-cls-thm-dc} 说明：直流潮流的适定性只需要"连通 + 至少一个
相角参考 + 正电抗"，与运行点无关——这是它作为筛选/规划工具的鲁棒性来源。
代价是信息损失：$|V|$、$Q$、网损全部不输出，slack 母线自动吸收全系统不平衡
（在无损模型中恰为 $\sum P_i^{\mathrm{spec}}$ 的反号）。
\end{remark}

\subsubsection{线性化误差分析}

保留电阻与电压幅值后，直流模型相对式~\eqref{sec:pf-th-cls-eq:acbranch} 的误差
可以逐项界定。下述命题给出 $V\equiv1$ 情形下的干净分解。

\begin{proposition}[直流潮流支路线性化误差界]\label{sec:pf-th-cls-prop-dcerr}
设支路 $k$ 阻抗 $z_k=r_k+jx_k$（$x_k>0$），$V_i=V_j=1$，记
$\theta=\theta_{ij}$，$\rho_k=r_k/x_k$。则精确有功
$P_{ij}=g_k(1-\cos\theta)+\frac{x_k}{r_k^2+x_k^2}\sin\theta$，
与直流近似 $\theta/x_k$ 的误差 $\varepsilon_k=P_{ij}-\theta/x_k$ 满足
\begin{equation}
|\varepsilon_k|
\le \frac{g_k\theta^2}{2}
+\frac{\rho_k^2}{x_k(1+\rho_k^2)}\,|\theta|
+\frac{|\theta|^3}{6x_k}.
\label{sec:pf-th-cls-eq:dcerr}
\end{equation}
三项分别对应\textbf{损耗项}（$g(1-\cos\theta)$，二阶小量但恒正、系统性偏高）、
\textbf{有效电抗项}（$1/x_{\mathrm{eff}}=x_k/(r_k^2+x_k^2)$ 与 $1/x_k$ 之差，
$\rho_k^2$ 量级）与\textbf{正弦线性化项}（三阶小量）。
\end{proposition}

\begin{proof}
由 $y_k=1/z_k=(r_k-jx_k)/(r_k^2+x_k^2)$，得 $g_k=r_k/(r_k^2+x_k^2)$、
$b_k=-x_k/(r_k^2+x_k^2)$，代入式~\eqref{sec:pf-th-cls-eq:acbranch}
（$V_i=V_j=1$）即得精确表达式。于是
\begin{equation*}
\varepsilon_k=g_k(1-\cos\theta)
+\sin\theta\Bigl(\frac{x_k}{r_k^2+x_k^2}-\frac{1}{x_k}\Bigr)
+\frac{\sin\theta-\theta}{x_k}.
\end{equation*}
逐项取绝对值：$0\le1-\cos\theta\le\theta^2/2$；
$\bigl|\frac{x_k}{r_k^2+x_k^2}-\frac{1}{x_k}\bigr|
=\frac{r_k^2}{x_k(r_k^2+x_k^2)}=\frac{\rho_k^2}{x_k(1+\rho_k^2)}$；
$|\sin\theta-\theta|\le|\theta|^3/6$（Taylor 余项的交错级数估计）。
三角不等式合并即得式~\eqref{sec:pf-th-cls-eq:dcerr}。
\end{proof}

\begin{remark}[误差的量级解读]\label{sec:pf-th-cls-rem-dcerr}
输电网典型 $\rho_k\le0.3$、$|\theta|\le0.5$~rad：有效电抗项
$\sim\rho_k^2|\theta|/x_k\le4.5\%$（相对 $|\theta|/x_k$），损耗项与正弦项更小；
这与工程文献中"直流潮流支路有功误差约百分之几"的经验一致~\cite{cls:stottreview}。
但对\textbf{高 $R/X$ 配网}（$\rho_k$ 可超 1），有效电抗项达 $O(\rho_k^2)$，
式~\eqref{sec:pf-th-cls-eq:dcerr} 失去实用意义——直流潮流不适用于配网，
与 FDPF 的失效机理同源（假设 (A1)/(D1)）。
\end{remark}

\subsubsection{损耗近似}

直流模型自身无损（$g=0$ 隐含 $\sum P_i^{\mathrm{spec}}=0$ 才可精确平衡），
工程中需要网损估计时标准的做法是\textbf{事后二次近似}：由精确支路损耗
$P_{\mathrm{loss},k}=g_k|V_i-V_j|^2=g_k\bigl(V_i^2+V_j^2-2V_iV_j\cos\theta_{ij}\bigr)$，
施加 (D2) 并展开 $\cos$ 到二阶，得
\begin{equation}
P_{\mathrm{loss},k}\approx g_k\,\theta_{ij}^2=\frac{r_k}{r_k^2+x_k^2}\,\theta_{ij}^2
\approx r_k\Bigl(\frac{\theta_{ij}}{x_k}\Bigr)^2
=r_k\,\frac{P_{f,k}^2}{1},
\label{sec:pf-th-cls-eq:loss}
\end{equation}
即"损耗 $\approx r\times$（直流潮流）$^2$"的二次口径。进一步可把
$\sum_k P_{\mathrm{loss},k}/2$ 按比例回加到负荷并重解（损耗修正直流潮流），
使 slack 吸收量逼近真实网损。

\begin{gapnote}
平台直流求解器 \srcpath{solve\_ac\_linearized\_dc} 输出仅为相角与无损支路潮流；
式~\eqref{sec:pf-th-cls-eq:loss} 的\textbf{损耗事后估计与损耗修正重解均未实现}
（结果结构 \texttt{ACLinearizedDCResult} 无损耗字段）。
\end{gapnote}

\subsection{直流灵敏度：转移分布因子与线路开断分布因子}\label{sec:pf-th-cls-ptdf}

直流模型的线性性使支路潮流对节点注入、以及对线路开断的\emph{灵敏度}具有闭式矩阵
表达。这两类因子——PTDF 与 LODF——是安全分析、市场出清与 N-1 校核的基础工具，
其价值在于\textbf{只依赖网络拓扑与电抗、与运行点无关}，可离线预计算。

\subsubsection{PTDF：注入转移分布因子}
设剔除 slack 后的降阶节点电纳阵 $B\in\mathbb R^{(n-1)\times(n-1)}$
（式~\eqref{sec:pf-th-cls-eq:dcpf}，$B=A^{\top}B_dA$），支路-节点关联阵
$A\in\mathbb R^{b\times(n-1)}$（第 $l$ 行对应支路 $l=(i\!\to\!j)$，$i$ 列为 $+1$、
$j$ 列为 $-1$，slack 列剔除），支路电纳对角阵 $B_d=\mathrm{diag}(1/x_l)$。支路有功
$f_l=(\theta_i-\theta_j)/x_l$ 即 $f=B_dA\theta$，结合 $\theta=B^{-1}P$ 得

\begin{definition}[PTDF/ISF 矩阵]\label{def:pf-th-cls-ptdf}
\begin{equation}
f=\Phi\,P,\qquad \Phi:=B_dA\,B^{-1}\in\mathbb R^{b\times(n-1)} .
\label{sec:pf-th-cls-eq:ptdf}
\end{equation}
$\Phi_{lk}$ 是“节点 $k$ 单位注入、slack 回收”引起的支路 $l$ 潮流变化，称\emph{注入
转移分布因子}（injection shift factor, ISF），亦称\emph{功率转移分布因子}（PTDF）。
\end{definition}

\begin{proposition}[点对点交易的支路灵敏度]\label{prop:pf-th-cls-ptdf-pair}
从节点 $s$ 注入、节点 $t$ 抽取的单位交易（$P_s{=}{+}1,\,P_t{=}{-}1$）引起支路 $l$ 的
潮流变化为 $\Phi_{ls}-\Phi_{lt}$；特别地，沿支路 $m=(p\!\to\!q)$ 两端的单位对流引起
支路 $l$ 的变化记 $\rho_{lm}:=\Phi_{lp}-\Phi_{lq}$（slack 端分量取零）。
\end{proposition}
（由式~\eqref{sec:pf-th-cls-eq:ptdf} 对注入的线性叠加即得。）

\subsubsection{LODF：线路开断分布因子}
\begin{theorem}[Wood--Wollenberg LODF 公式]\label{thm:pf-th-cls-lodf}
设支路 $m$ 开断前潮流为 $f_m^0$，且开断不致系统解列（$\rho_{mm}\ne1$）。则开断后
支路 $l$（$l\ne m$）的潮流增量为
\begin{equation}
\Delta f_l=\psi_{lm}\,f_m^0,\qquad
\psi_{lm}=\frac{\rho_{lm}}{1-\rho_{mm}},
\label{sec:pf-th-cls-eq:lodf}
\end{equation}
$\rho_{lm}=\Phi_{lp}-\Phi_{lq}$、$\rho_{mm}=\Phi_{mp}-\Phi_{mq}$（$m=(p\!\to\!q)$），
$\psi_{lm}$ 称\emph{线路开断分布因子}（LODF）。
\end{theorem}
\begin{proof}
用补偿定理：在\emph{完整}网络上叠加一对沿 $m$ 端点的注入（$p$ 注入 $\Delta$、$q$
抽取 $\Delta$），以其等效替代被开断支路的电流。该注入使支路 $m$ 上产生附加流
$\rho_{mm}\Delta$，使任意支路 $l$ 产生 $\rho_{lm}\Delta$。开断的自洽条件是：这对注入
所“代表”的幻象电流 $\Delta$ 恰等于叠加后支路 $m$ 的总流，即
$\Delta=f_m^0+\rho_{mm}\Delta$，解得 $\Delta=f_m^0/(1-\rho_{mm})$。于是任意
$l\ne m$ 的增量 $\Delta f_l=\rho_{lm}\Delta=\dfrac{\rho_{lm}}{1-\rho_{mm}}f_m^0$。
$\rho_{mm}\to1$ 对应开断使 $m$ 成为割边、网络解列，LODF 发散。
\end{proof}

\begin{remark}[拓扑不变性与用途]\label{rem:pf-th-cls-lodf}
$\Phi$ 与 $\psi$ 只由 $A$、$B_d$（拓扑与电抗）决定，与负荷、发电水平无关，故可离线
预计算、在线以一次矩阵-向量积完成全部注入或单一开断的支路潮流筛查——这是直流
线性化在安全约束调度、N-1 校核与市场阻塞定价中被广泛采用的根本原因；其代价是
继承直流模型的全部近似（命题~\ref{sec:pf-th-cls-prop-dcerr}）。多重开断需对
$\psi$ 作分块推广（同时开断的支路间存在耦合）。
\end{remark}

\subsection{辐射状配电网潮流：DistFlow 模型与前推回代}\label{sec:pf-th-cls-distflow}

\subsubsection{DistFlow 支路潮流模型}

配电网的结构性差异——辐射拓扑、高 $R/X$、显著电压降落——使输电网方法的前提
（(A1)、(D1)--(D3)）不再成立，催生了以支路变量为核心的模型体系。
Baran 与 Wu~\cite{cls:baranwu} 的 DistFlow 模型是其中的标准形式。

\begin{definition}[DistFlow 支路潮流方程]\label{sec:pf-th-cls-def-distflow}
设辐射树 $\mathcal T$ 的支路一律按远离根的方向定向，支路 $(i,j)$ 阻抗
$z_{ij}=r_{ij}+jx_{ij}$，记 $P_{ij},Q_{ij}$ 为从 $i$ 流向 $j$ 的有功/无功，
$\ell_{ij}=|I_{ij}|^2$，$v_i=|V_i|^2$，$p_j,q_j$ 为母线 $j$ 的净负荷
（含本地并联的 $|V|^2$ 项）。由 $V_j=V_i-z_{ij}I_{ij}$ 与
$S_{ij}=V_i\overline{I_{ij}}$ 两端取模平方并分离实虚部，得递推方程组
\begin{align}
P_{ij}-r_{ij}\ell_{ij}-p_j&=\sum_{k\in\mathcal C(j)}P_{jk},
\label{sec:pf-th-cls-eq:dfp}\\
Q_{ij}-x_{ij}\ell_{ij}-q_j&=\sum_{k\in\mathcal C(j)}Q_{jk},
\label{sec:pf-th-cls-eq:dfq}\\
v_j&=v_i-2\bigl(r_{ij}P_{ij}+x_{ij}Q_{ij}\bigr)+\bigl(r_{ij}^2+x_{ij}^2\bigr)\ell_{ij},
\label{sec:pf-th-cls-eq:dfv}\\
\ell_{ij}\,v_i&=P_{ij}^2+Q_{ij}^2,
\label{sec:pf-th-cls-eq:dfl}
\end{align}
$\mathcal C(j)$ 为 $j$ 的子节点集。方程
\eqref{sec:pf-th-cls-eq:dfp}--\eqref{sec:pf-th-cls-eq:dfq} 是含损功率平衡，
\eqref{sec:pf-th-cls-eq:dfv} 是电压降落恒等式，\eqref{sec:pf-th-cls-eq:dfl}
是支路潮流与电流的耦合约束（唯一的非线性，且为凸二阶锥形式）。
略去损耗项 $r\ell$、$x\ell$ 与 $(r^2+x^2)\ell$ 得\textbf{ LinDistFlow}
（线性 DistFlow）模型，广泛用于配网优化建模~\cite{cls:baranwu}。
\end{definition}

\begin{remark}[存在唯一性]\label{sec:pf-th-cls-rem-chiang}
Chiang 与 Baran~\cite{cls:chiangbaran} 证明：对辐射状配网，在"负荷不过分重"
的可操作条件下，潮流方程的可行解在物理有意义的电压区域内\textbf{存在且唯一}，
且可用不动点框架构造。这一结果把配网潮流与输电网多解问题区分开：辐射 +
单电源 + 适度负荷保证了良态性，是前推回代可靠工作的理论底座。
\end{remark}

\subsubsection{前推回代：作为不动点迭代的推导}

前推回代（Backward--Forward Sweep, BFS；Shirmohammadi
等~\cite{cls:shirmohammadi}）不显式组装 DistFlow 方程组，而是把辐射树结构
直接变成两遍扫描的不动点迭代。

\begin{definition}[BFS 扫描映射]\label{sec:pf-th-cls-def-bfs}
给定电压剖面 $V=(V_1,\dots,V_n)$（根 $V_0$ 固定为 slack 电压），
\textbf{回推}（叶$\to$根）按
$I_{ij}=\sum_{m\in\mathrm{sub}(j)}\overline{S_m^{\mathrm{net}}/V_m}$
汇总各支路电流；\textbf{前推}（根$\to$叶）沿树按
$V'_j=V'_i-z_{ij}I_{ij}$ 更新电压。记一遍完整扫描为映射 $V'=T(V)$，
BFS 即 Picard 迭代 $V^{(k+1)}=T(V^{(k)})$，收敛判据
$\max_j|V_j^{(k+1)}-V_j^{(k)}|<\varepsilon$。
\end{definition}

展开前推递推：$V'_j=V_0-\sum_{l\in\mathcal P(j)}z_l I_l$。由于 $I_l$ 只依赖
$\mathrm{sub}$ 内母线电压，$T$ 是良定义的显式映射；其不动点精确满足
DistFlow 方程（回推保证
\eqref{sec:pf-th-cls-eq:dfp}--\eqref{sec:pf-th-cls-eq:dfq}，
前推保证 \eqref{sec:pf-th-cls-eq:dfv}，电流定义保证
\eqref{sec:pf-th-cls-eq:dfl}）。

\begin{theorem}[BFS 的压缩收敛性（充分条件）]\label{sec:pf-th-cls-thm-bfs}
设辐射网以母线 $0$ 为 slack，$V_0=1\angle0$ 固定；负荷恒功率，
$|S_j^{\mathrm{net}}|\le\sigma_j$，记 $\sigma_\Sigma=\sum_j\sigma_j$；支路阻抗
满足 $Z_{\max}:=\max_j\sum_{l\in\mathcal P(j)}|z_l|<\infty$（根到任一母线的
路径阻抗上界）。取 $v_{\min}\in(0,1)$ 并限制电压剖面于
$\mathcal D=\{V:\ |V_j|\ge v_{\min}\ \forall j\}$。若
\begin{equation}
L:=\frac{Z_{\max}\,\sigma_\Sigma}{v_{\min}^2}<1
\qquad\text{且}\qquad
1-\frac{Z_{\max}\sigma_\Sigma}{v_{\min}}\ge v_{\min},
\label{sec:pf-th-cls-eq:bfsc}
\end{equation}
则 $T$ 在 $\mathcal D$ 上是压缩映射（$\lVert T(V)-T(W)\rVert_\infty\le
L\lVert V-W\rVert_\infty$）且 $T(\mathcal D)\subseteq\mathcal D$；
由 Banach 不动点定理，$T$ 在 $\mathcal D$ 内有\textbf{唯一}不动点，
BFS 从 $\mathcal D$ 内任意初值出发以速率 $L$ 几何收敛。
\end{theorem}

\begin{proof}
\textbf{第一步（回推的 Lipschitz 界）。}对
$I_l=\sum_{m\in\mathrm{sub}(l)}\overline{S_m/V_m}$，在 $\mathcal D$ 上
\begin{equation*}
\Bigl|\frac{\overline{S_m}}{\overline{V_m}}-\frac{\overline{S_m}}{\overline{W_m}}\Bigr|
=\frac{|S_m|\,|V_m-W_m|}{|V_m|\,|W_m|}
\le\frac{\sigma_m}{v_{\min}^2}\,\lVert V-W\rVert_\infty,
\end{equation*}
对子树求和得
$|I_l(V)-I_l(W)|\le\bigl(\sigma_\Sigma/v_{\min}^2\bigr)\lVert V-W\rVert_\infty$。
\textbf{第二步（前推传播）。}由 $T(V)_j=V_0-\sum_{l\in\mathcal P(j)}z_l I_l(V)$，
\begin{equation*}
|T(V)_j-T(W)_j|\le\sum_{l\in\mathcal P(j)}|z_l|\,
\frac{\sigma_\Sigma}{v_{\min}^2}\lVert V-W\rVert_\infty
\le L\,\lVert V-W\rVert_\infty,
\end{equation*}
对 $j$ 取上确界即压缩性。\textbf{第三步（自映射）。}对 $V\in\mathcal D$，
$|I_l(V)|\le\sigma_\Sigma/v_{\min}$，故
$|T(V)_j|\ge1-Z_{\max}\sigma_\Sigma/v_{\min}\ge v_{\min}$
（式~\eqref{sec:pf-th-cls-eq:bfsc} 第二式），即 $T(V)\in\mathcal D$。
$\mathcal D$ 是有限维空间中的闭集（完备度量空间），Banach 定理给出
唯一不动点与几何收敛 $\lVert V^{(k)}-V^*\rVert_\infty\le
\frac{L^k}{1-L}\lVert V^{(1)}-V^{(0)}\rVert_\infty$。
\end{proof}

\begin{remark}[条件 \eqref{sec:pf-th-cls-eq:bfsc} 的物理读法]\label{sec:pf-th-cls-rem-bfsc}
$Z_{\max}\sigma_\Sigma$ 是"路径阻抗 $\times$ 总负荷"——即\textbf{短路容量相对
负荷规模的裕度}：条件要求馈线足够"强"（短路容量大、负荷轻、电压不贴近崩塌点
$v_{\min}$ 过小）。这与工程直觉一致：重载长馈线末端电压低，BFS 迭代变慢甚至
失效；也解释了为何实现必须用拓扑 fail-close 与残差监控兜底，而不能依赖先验
条件检验（$L<1$ 只是充分条件，实践中 $L$ 略超 1 仍常收敛）。
恒阻抗/恒电流负荷使 $T$ 更接近线性映射（注入电流对 $V$ 的敏感度降低），
压缩条件更宽松——这就是 ZIP 负荷有助于 BFS 收敛的理论原因。
\end{remark}

\begin{remark}[高 $R/X$ 弱网特性与方法选择]\label{sec:pf-th-cls-rem-rx}
配网 $R/X$ 常大于 1：(i) FDPF/直流的假设 (A1)/(D1) 被破坏
（命题~\ref{sec:pf-th-cls-prop-dcerr} 的有效电抗项达 $O(\rho_k^2)$）；
(ii) $P$--$V$、$Q$--$\theta$ 交叉耦合增强，解耦迭代发散；
(iii) 电压降落项 $2(rP+xQ)$ 中电阻贡献占主导，$|V|$ 对有功潮流敏感。
BFS 与 DistFlow 不显式使用任何弱耦合假设（式
\eqref{sec:pf-th-cls-eq:dfp}--\eqref{sec:pf-th-cls-eq:dfl} 是精确恒等式），
对 $R/X$ 天然免疫，配合辐射拓扑的 $O(n)$ 扫描复杂度，成为配网正序潮流的
标准选择；环网/弱环网则需补偿法~\cite{cls:shirmohammadi} 或三相 Newton。
\end{remark}

\begin{gapnote}
平台实现的是\textbf{精确 AC 支路模型的 BFS}
（\srcpath{run\_bfs\_sweep}，含变比折算与 ZIP 电压依赖），收敛点与
DistFlow 方程组等价，但\textbf{没有}独立的 DistFlow 方程组求解器；
LinDistFlow 仅以 MIP 建模层形式存在于 power\_models 模块
（\srcpath{src/power\_models/lindistflow\_builder.cpp:solve\_lindistflow}），
不属于潮流求解路径。定理~\ref{sec:pf-th-cls-thm-bfs} 的压缩条件
\eqref{sec:pf-th-cls-eq:bfsc} 亦\textbf{未}作为先验检验实现——实现以辐射/连通
fail-close 检查加 $\max|\Delta V|$ 收敛判据代替。
\end{gapnote}

\subsection{分布式松弛与自适应孤岛的理论定位}\label{sec:pf-th-cls-slack}

\subsubsection{松弛母线的自由度账与物理意义}

$n$ 母线 AC 潮流每节点有四个量 $P,Q,|V|,\theta$、两条规范方程（注入定义），
因此每母线必须指定两个量；但全系统的\textbf{网损在求解前未知}，若所有母线的
$P$ 都事先钉死，功率平衡 $\sum P_i=\text{损耗}$ 一般无解。松弛（slack）母线的
本质是\textbf{开放其 $P,Q$ 两个自由度}以吸收先验未知的不平衡（主要是网损），
同时提供相角参考——$\theta$ 只在连通分量内有相对意义，必须钉住一个零点
（自由度账：$n$ 个 PQ/PV 母线贡献 $2n-2$ 个失配方程，对应 $\theta$ 与 $|V|$
的 $2n-2$ 个未知量，slack 的 $\theta,|V|$ 钉值）。

\begin{definition}[分布式松弛与参与因子]\label{sec:pf-th-cls-def-dslack}
把单一 slack 吸收的不平衡 $\Delta P_\Sigma$ 按归一化\textbf{参与因子}
$\alpha_i\ge0$、$\sum_i\alpha_i=1$ 分摊到一组可调电源母线
$\mathcal G$：$P_{g,i}\leftarrow P_{g,i}+\alpha_i\Delta P_\Sigma$。
因子的常见生成口径：按容量（$\alpha_i\propto P_{g,i}^{\max}$）、
按下垂系数（$\alpha_i\propto1/R_i$，对应一次调频的稳态份额）或均分。
\end{definition}

\begin{remark}[slack 分摊的物理意义与两种数学口径]\label{sec:pf-th-cls-rem-dslack}
单 slack 模型隐含"一台无穷大机组瞬时平衡全网"，在网损分摊、AGC 与经济调度的
衔接上不真实；分布式松弛把不平衡按参与因子映射到多台机组，模拟调速器
稳态分配。\emph{数学上}有两种口径：\textbf{方程内口径}把
$P_{g,i}=P_{g,i}^0+\alpha_i\,\pi$（$\pi$ 为增广未知的不平衡量）直接嵌入
Newton 失配方程同步求解，结果自洽（分摊后网损变化已被迭代吸收）；
\textbf{事后口径}先用单 slack 解出基准潮流，再把 slack 吸收量按 $\alpha$ 外部分摊。
事后口径忽略分摊引起的网损二阶变化（$\Delta P_\Sigma$ 改变潮流分布进而改变
损耗），属于一阶近似；要自洽需把分摊值写回后\textbf{重解}直至残余失配收敛。
限值处理（机组 $P_{\min}/P_{\max}$ 越限钉限再分摊剩余）使口径进一步成为
带投影的迭代。
\end{remark}

\subsubsection{方程内精确分布式松弛的加边 Newton 系统}\label{sec:pf-th-cls-augslack}
注记~\ref{sec:pf-th-cls-rem-dslack} 的“事后口径”只一阶近似分摊引起的网损变化。
\emph{方程内}口径把分摊量作为增广未知量与潮流方程联立，自洽吸收网损的二阶效应。

\begin{proposition}[加边分布式松弛系统]\label{prop:pf-th-cls-augslack}
引入标量未知 $\pi$（待分摊的总不平衡）；参与母线 $i\in\mathcal G$ 的有功注入取
$P_i^{\mathrm{spec}}(\pi)=P_i^0+\alpha_i\pi$（$\sum_{i\in\mathcal G}\alpha_i=1$）。
保留一个\emph{仅}提供角度与电压基准的参考母线（$\theta_{\mathrm{ref}}=0$、
$V_{\mathrm{ref}}$ 固定、其无功自由），并\textbf{恢复}参考母线的有功平衡方程纳入方程组。
则未知量 $(\theta_{\ne\mathrm{ref}},V_{m,\mathrm{PQ}},\pi)$ 与方程（全部母线有功平衡
$+$ PQ 母线无功平衡）个数相等，其 Newton 系统为\emph{加边}结构
\begin{equation}
\begin{bmatrix} J_{\mathrm{pf}} & \boldsymbol\alpha\\[2pt]
\boldsymbol c^{\top} & 0\end{bmatrix}
\begin{bmatrix}\Delta x\\ \Delta\pi\end{bmatrix}
=-\begin{bmatrix}F_{\mathrm{pf}}\\ F_{\mathrm{ref}}\end{bmatrix},
\label{sec:pf-th-cls-eq:augslack}
\end{equation}
其中 $J_{\mathrm{pf}}$ 为常规潮流雅可比，$\boldsymbol\alpha$ 为
$\partial F_{\mathrm{pf}}/\partial\pi$ 的参与因子列，$\boldsymbol c^{\top}$ 为新增
参考母线有功平衡对 $x$ 的偏导行。
\end{proposition}

\begin{remark}[自洽性与加边求解]\label{rem:pf-th-cls-augslack}
式~\eqref{sec:pf-th-cls-eq:augslack} 的解 $\pi^\star$ 是收敛点处各参与母线共同吸收的
总不平衡（对应全网损耗份额）；网损随分摊变化的二阶效应已在迭代内消化，正是事后
口径所缺。加边系统只比基潮流多一行一列，用块消元（Schur 补/矩阵求逆引理）可在\emph{一次}
基潮流分解上求解——$\Delta x=-J_{\mathrm{pf}}^{-1}(F_{\mathrm{pf}}+\boldsymbol\alpha\Delta\pi)$
代入底行得关于标量 $\Delta\pi$ 的方程，$\Delta\pi=-(F_{\mathrm{ref}}-\boldsymbol c^{\top}
J_{\mathrm{pf}}^{-1}F_{\mathrm{pf}})/(\boldsymbol c^{\top}J_{\mathrm{pf}}^{-1}\boldsymbol\alpha)$，
仅需两次以 $J_{\mathrm{pf}}$ 的回代，无需如“重复完整重解”那样多次重构与再分解。
多参考（多岛）时 $\pi$ 逐岛设立、$\alpha$ 逐岛归一。
\end{remark}

\begin{gapnote}
式~\eqref{sec:pf-th-cls-eq:augslack} 的方程内加边分布式松弛在平台中\textbf{未实现}
（实现为事后分摊与重复重解，见本节末“与实现的对应”表及其函数命名说明）。
\end{gapnote}

\subsubsection{自适应孤岛：连通分量、参考角与死岛}

拓扑分裂（故障开断、计划孤岛）使一个求解对象变成多个电气解耦的子问题。
其代数判据是图 Laplacian 的秩亏结构。

\begin{lemma}[Laplacian 零空间的维数 = 连通分量数]\label{sec:pf-th-cls-lem-lap}
设无向图有 $n$ 个顶点、$c$ 个连通分量，边权 $w_k>0$，$B$ 为其加权
Laplacian（$B_{ij}=-w_k$、$B_{ii}=\sum_{k\ni i}w_k$）。则
$\dim\ker B=c$，$\ker B$ 由各分量的示性向量张成。
\end{lemma}

\begin{proof}
由式~\eqref{sec:pf-th-cls-eq:quad} 的二次型展开
$x^{\top}Bx=\sum_{k=(i,j)}w_k(x_i-x_j)^2\ge0$：$Bx=0\Rightarrow x^{\top}Bx=0$
$\Rightarrow$ 每条边两端分量相等 $\Rightarrow$ $x$ 在每个连通分量上取常值。
反之每个分量上的常值向量显然在 $\ker B$ 中。故 $\ker B$ 恰为
$\{\text{各分量示性向量的线性组合}\}$，维数为 $c$。
\end{proof}

\begin{remark}[引理~\ref{sec:pf-th-cls-lem-lap} 对潮流求解的含义]\label{sec:pf-th-cls-rem-island}
$c$ 个孤岛的网络，其潮流雅可比（以及直流 $B$ 阵）天然有 $c$ 重零特征值：
每个孤岛都需要\textbf{各自}的相角参考与功率平衡机制，混用单一 slack 的
整体方程组必然奇异。\emph{自适应孤岛}的理论定位因此是：先分岛（图的连通分量
分解），把原问题分解为 $c$ 个独立的适定子问题，活岛各自指定/自选 slack 求解；
\textbf{无源岛}（无任何注入设备）在潮流框架内\textbf{无可行解}——其负荷无法满足、
电压无支撑，物理上对应失压停运，诚实的计算处理是作为\emph{死岛}以零电压
哨兵标记并排除出求解，而非强行给出无意义数值。这与定理
\ref{sec:pf-th-cls-thm-dc} 的适定性条件（连通 + 参考 + 正电抗）互为表里。
\end{remark}

\begin{gapnote}
方程内口径的精确分布式松弛（$\pi$ 作为增广未知量进入 Newton 失配）在平台中
\textbf{未实现}：\srcpath{DistributedSlackSolver::solve\_simplified} 是事后分摊
（一次基准 Newton + \srcpath{compute\_total\_participating\_mismatch\_pu} 复算
分摊），\srcpath{solve\_full\_jacobian} 是"重调度 + 反复完整重解 + 钉限重分摊"
的迭代逼近（兼容名 \texttt{full\_jacobian} 不表示增广 Jacobian）。
\end{gapnote}

\subsection{数值算例与基准}\label{sec:pf-th-cls-numexp}

本节数字由 \srcpath{tools/pf\_doc\_benchmark.cpp}（模式 \texttt{fdpf}/\texttt{dc}/\texttt{bfs}）
在 MATPOWER 算例上经平台公共求解 API 实测生成（平启动，Apple~M4，单线程），
迭代次数与误差为机器无关量。

\subsubsection{FDPF 与 Newton 的迭代数对比}
表~\ref{tab:pf-th-cls-fdpf} 给出统一 Newton 与 FDPF（XB 方案）达到 $\mathrm{tol}=10^{-8}$
所需迭代数。FDPF 迭代数（$11\sim44$）显著多于 Newton（$3\sim8$），且随系统解耦质量
（$R/X$、角度差）波动——这正是定理~\ref{sec:pf-th-cls-thm-fdpf} 的\emph{线性}收敛与
Newton \emph{二次}收敛之别；FDPF 以更多次但每次仅两回代（常量 $B'/B''$ 一次分解）
换取单步成本的数量级下降。
\begin{table}[H]\centering\footnotesize
\caption{FDPF 与统一 Newton 的迭代数（平启动，$\mathrm{tol}=10^{-8}$）}
\label{tab:pf-th-cls-fdpf}
\begin{tabular}{@{}lrrr@{}}
\toprule
算例 & 母线数 $n$ & Newton 迭代 & FDPF(XB) 迭代\\
\midrule
case9   & 9   & 5 & 11\\
case14  & 14  & 3 & 26\\
case30  & 30  & 4 & 27\\
case57  & 57  & 4 & 39\\
case118 & 118 & 7 & 15\\
case300 & 300 & 8 & 44\\
\bottomrule
\end{tabular}
\end{table}

\subsubsection{直流潮流的支路有功误差}
表~\ref{tab:pf-th-cls-dc} 给出直流线性化支路有功与全交流解的偏差（逐支路 from 端有功）。
平均误差为 $0.4\sim11$~MW（相对各算例总负荷约 $0.05\%\sim0.5\%$），但重载支路的最大误差
可达数十至数百 MW，随规模/负荷压力上升——与命题~\ref{sec:pf-th-cls-prop-dcerr} 的误差
界（损耗项 $\propto\theta^2$、正弦线性化项 $\propto\theta^3$）一致：角度差越大，线性化
误差越显著。直流适合把握整体潮流格局，不宜用于重载支路的精确定量。
\begin{table}[H]\centering\footnotesize
\caption{直流线性化相对全交流的支路有功误差（MW）}
\label{tab:pf-th-cls-dc}
\begin{tabular}{@{}lrrrrr@{}}
\toprule
算例 & 母线数 & 支路数 & 最大误差 & 平均误差 & 均方根误差\\
\midrule
case9   & 9   & 9   & $4.6$   & $1.25$ & $1.95$\\
case14  & 14  & 20  & $9.0$   & $1.25$ & $2.42$\\
case30  & 30  & 41  & $1.7$   & $0.39$ & $0.53$\\
case57  & 57  & 80  & $9.6$   & $1.57$ & $2.45$\\
case118 & 118 & 186 & $59.4$  & $3.60$ & $7.46$\\
case300 & 300 & 411 & $409.5$ & $10.6$ & $29.8$\\
\bottomrule
\end{tabular}
\end{table}

\subsubsection{前推回代在配电馈线上的表现}
表~\ref{tab:pf-th-cls-bfs} 在四个辐射配电算例上对比 BFS 与统一 Newton。两法最低母线电压
一致到 $10^{-9}$（交叉校核），BFS 迭代数（$7\sim9$）略多于 Newton（$4\sim5$），但每次
仅 $O(n)$ 次树上扫描、无雅可比装配与分解；馈线末端电压低至 $0.87\sim0.93$~pu，体现
高 $R/X$ 长馈线的显著压降。BFS 对 $R/X$ 天然免疫（式~\eqref{sec:pf-th-cls-eq:dfp}--%
\eqref{sec:pf-th-cls-eq:dfl} 为精确恒等式），是注记~\ref{sec:pf-th-cls-rem-rx} 所述配网
首选方法的实测印证。
\begin{table}[H]\centering\footnotesize
\caption{前推回代（BFS）与统一 Newton 在辐射配电算例上的对比（$\mathrm{tol}=10^{-8}$）}
\label{tab:pf-th-cls-bfs}
\begin{tabular}{@{}lrrrrrr@{}}
\toprule
算例 & 母线数 & BFS 迭代 & BFS $\min|V|$ & Newton 迭代 & Newton $\min|V|$ & $\min|V|$ 差\\
\midrule
case33bw & 33  & 8 & $0.9131$ & 4 & $0.9131$ & $2\times10^{-9}$\\
case69   & 69  & 8 & $0.9092$ & 5 & $0.9092$ & $4\times10^{-10}$\\
case85   & 85  & 9 & $0.8739$ & 5 & $0.8739$ & $5\times10^{-10}$\\
case141  & 141 & 7 & $0.9279$ & 4 & $0.9279$ & $1\times10^{-10}$\\
\bottomrule
\end{tabular}
\end{table}

\subsection{与实现的对应}\label{sec:pf-th-cls-implmap}
\begin{center}\footnotesize
\begin{longtable}{@{}L{6.8cm}L{8.6cm}@{}}
\toprule
理论条目 & 实现状态\\
\midrule
\endfirsthead
\toprule
理论条目 & 实现状态\\
\midrule
\endhead
Y-bus Gauss--Seidel 不动点迭代（定义~\ref{sec:pf-th-cls-def-gs}） &
\implnone\\
Z-bus 阻抗阵法 &
\implnone\\
FDPF 主迭代 $B'\Delta\theta=\Delta P/V$、$B''\Delta V=\Delta Q/V$
（式~\eqref{sec:pf-th-cls-eq:fdpf}） &
\implfull{src/power\_flow/fdpf\_solver.cpp:FDPFSolver::solve}（失配装配见同文件
compute\_power\_injections、compute\_specified\_injections）\\
$B'/B''$ 的 XB 口径常量阵装配（定义~\ref{sec:pf-th-cls-def-xbbx}） &
\implfull{src/power\_flow/admittance\_builder.cpp:build\_susceptance\_matrix}（构型标志
\srcpath{include/hacdcpf/power\_flow/assembly/admittance\_builder.hpp:make\_bp\_flags}、\srcpath{make\_bpp\_flags}）\\
FDPF 常量阵一次分解 + 线性收敛结构
（定理~\ref{sec:pf-th-cls-thm-fdpf}） &
\implpart{Eigen SparseLU 一次性分解（FDPFSolver::solve）；无全局化与 Q 限值切换，
收敛域外的行为无保障}\\
FDPF BX 变体 &
\implnone\\
直流潮流 $B\theta=P$ 降阶求解
（式~\eqref{sec:pf-th-cls-eq:dcpf}，定理~\ref{sec:pf-th-cls-thm-dc}） &
\implfull{src/power\_flow/ac\_linearized\_pf.cpp:solve\_ac\_linearized\_dc}\\
直流潮流移相变压器注入等效 &
\implfull{src/power\_flow/ac\_linearized\_pf.cpp:solve\_ac\_linearized\_dc}（同函数内
$\pm b_k\varphi_k$ 注入装配）\\
直流损耗二次近似与损耗修正重解（式~\eqref{sec:pf-th-cls-eq:loss}） &
\implnone\\
DistFlow 支路潮流方程组
（式~\eqref{sec:pf-th-cls-eq:dfp}--\eqref{sec:pf-th-cls-eq:dfl}） &
\implpart{无独立方程组求解器；辐射支路关系由精确 AC 模型的 BFS 扫描实现，
收敛点与 DistFlow 等价}\\
LinDistFlow 线性化模型 &
\implpart{潮流路径内无；MIP 建模层
src/power\_models/lindistflow\_builder.cpp:solve\_lindistflow（属 power\_models 模块）}\\
前推回代扫描映射（定义~\ref{sec:pf-th-cls-def-bfs}） &
\implfull{src/power\_flow/distribution\_power\_flow.cpp:run\_bfs\_sweep}（BFS 树
build\_bfs\_tree，支路关系 branch\_series\_current、branch\_child\_voltage\_from\_parent，
ZIP 电压依赖 voltage\_dependent\_net\_demand）\\
辐射/单连通前提的 fail-close 检查 &
\implfull{src/power\_flow/distribution\_power\_flow.cpp:solve\_distribution\_pf\_ac\_snapshot}\\
BFS 压缩收敛先验检验（定理~\ref{sec:pf-th-cls-thm-bfs} 的 $L<1$） &
\implpart{无先验条件检验；以 $\max|\Delta V|$ 收敛判据与拓扑 fail-close 兜底}\\
高 $R/X$ 网络的自动方法切换 &
\implnone（方法由用户显式选择，无基于 $R/X$ 的自动分派）\\
参与因子三口径生成（定义~\ref{sec:pf-th-cls-def-dslack}） &
\implfull{src/power\_flow/distributed\_slack\_solver.cpp:create\_participation\_factors}（含
sanitize\_slack\_cfg 校验与 normalize\_factors 归一）\\
事后分摊与带限值重解两口径
（注记~\ref{sec:pf-th-cls-rem-dslack}） &
\implfull{src/power\_flow/distributed\_slack\_solver.cpp:solve\_simplified、solve\_full\_jacobian}（限值重分摊
distribute\_with\_limits，失配复算 compute\_total\_participating\_mismatch\_pu）\\
方程内精确分布式松弛（增广未知量 $\pi$） &
\implnone\\
孤岛分解与每岛参考/slack 自选
（引理~\ref{sec:pf-th-cls-lem-lap} 的算法化） &
\implfull{src/power\_flow/island\_detector.cpp:detect\_islands、extract\_island\_subsystem；自选摇摆
src/power\_flow/adaptive\_solver.cpp:auto\_select\_swing\_bus}\\
死岛零电压哨兵（注记~\ref{sec:pf-th-cls-rem-island}） &
\implfull{src/power\_flow/adaptive\_solver.cpp:set\_dead\_island\_zero}（调度入口
AdaptiveSolver::solve）\\
PTDF/ISF 与 LODF（定义~\ref{def:pf-th-cls-ptdf}、定理~\ref{thm:pf-th-cls-lodf}） &
\implpart{仅市场 N-1 安全校核实现：src/market/market\_simulation.cpp:build\_lodf\_model
构造 PTDF 与 $\psi_{lm}=\rho_{lm}/(1-\rho_{mm})$（lodf\_value）；power\_flow 直流求解器不产出因子}\\
方程内精确分布式松弛加边系统（命题~\ref{prop:pf-th-cls-augslack}，
式~\eqref{sec:pf-th-cls-eq:augslack}） &
\implnone（实现为事后分摊 solve\_simplified 与重复完整重解 solve\_full\_jacobian，见上）\\
FDPF/直流/BFS 收敛性数值基准（表~\ref{tab:pf-th-cls-fdpf}、\ref{tab:pf-th-cls-dc}、
\ref{tab:pf-th-cls-bfs}） &
\implfull{tools/pf\_doc\_benchmark.cpp}（模式 fdpf/dc/bfs）\\
\bottomrule
\end{longtable}
\end{center}

\subsection{参考文献}
{\renewcommand{\section}[2]{}%
\begin{thebibliography}{9}\small
\bibitem{cls:wardhale} J.~B. Ward and H.~W. Hale, Digital computer solution of
power-flow problems, \emph{AIEE Transactions (Power Apparatus and Systems)},
vol. 75, 1956.
\bibitem{cls:stagg} G.~W. Stagg and A.~H. El-Abiad, \emph{Computer Methods in
Power System Analysis}, McGraw-Hill, 1968.
\bibitem{cls:tinney} W.~F. Tinney and C.~E. Hart, Power flow solution by
Newton's method, \emph{IEEE Transactions on Power Apparatus and Systems},
vol. PAS-86, no. 11, 1967.
\bibitem{cls:stott} B.~Stott and O.~Alsa\c{c}, Fast decoupled load flow,
\emph{IEEE Transactions on Power Apparatus and Systems}, vol. PAS-93, no. 3, 1974.
\bibitem{cls:stottreview} B.~Stott, Review of load-flow calculation methods,
\emph{Proceedings of the IEEE}, vol. 62, no. 7, 1974.
\bibitem{cls:vanamerongen} R.~A.~M. van Amerongen, A general-purpose version of
the fast decoupled load flow, \emph{IEEE Transactions on Power Systems},
vol. 4, no. 2, 1989.
\bibitem{cls:monticelli} A.~Monticelli, A.~Garcia and O.~R. Saavedra, Fast
decoupled load flow: hypothesis, derivations, and testing, \emph{IEEE
Transactions on Power Systems}, vol. 5, no. 4, 1990.
\bibitem{cls:ortega} J.~M. Ortega and W.~C. Rheinboldt, \emph{Iterative Solution
of Nonlinear Equations in Several Variables}, Academic Press, 1970.
\bibitem{cls:arrillaga} J.~Arrillaga and N.~R. Watson, \emph{Computer Modelling
of Electrical Power Systems}, 2nd ed., Wiley, 2001.
\bibitem{cls:baranwu} M.~E. Baran and F.~F. Wu, Network reconfiguration in
distribution systems for loss reduction and load balancing, \emph{IEEE
Transactions on Power Delivery}, vol. 4, no. 2, 1989.
\bibitem{cls:shirmohammadi} D.~Shirmohammadi, H.~W. Hong, A.~Semlyen and
G.~X. Luo, A compensation-based power flow method for weakly meshed distribution
and transmission networks, \emph{IEEE Transactions on Power Systems},
vol. 3, no. 2, 1988.
\bibitem{cls:chiangbaran} H.-D. Chiang and M.~E. Baran, On the existence and
uniqueness of load flow solution for radial distribution power networks,
\emph{IEEE Transactions on Circuits and Systems}, vol. 37, no. 3, 1990.
\bibitem{cls:grainger} J.~J. Grainger and W.~D. Stevenson, \emph{Power System
Analysis}, McGraw-Hill, 1994.
\bibitem{cls:kundur} P.~Kundur, \emph{Power System Stability and Control},
McGraw-Hill, 1994.
\bibitem{cls:woodwollenberg} A.~J. Wood, B.~F. Wollenberg and G.~B. Shebl\'e,
\emph{Power Generation, Operation, and Control}, 3rd ed., Wiley, 2014.
\end{thebibliography}}

% theory_globalization_advanced.tex —— 潮流计算 通用理论：全局化与高级求解方法
% 本文件为理论层章节，遵循 docs/modules/THEORY_WRITING_GUIDE.md。

\section{全局化与高级求解方法理论}\label{sec:pf-th-glob}

\begin{theorynote}
本章为通用数学理论，按数值最优化与非线性方程组求解的标准文献体系撰写，覆盖：
线搜索（Armijo/Wolfe 条件、Zoutendijk 全局收敛定理）、信赖域子问题与柯西点、
Levenberg--Marquardt 与信赖域的等价关系、GLL 非单调线搜索、NCP/半光滑 Newton
理论（Fischer--Burmeister 函数、Clarke 广义 Jacobian、超线性收敛）与 PV/PQ 切换
的统一表述、HELM 全纯嵌入的复分析基础（幂级数、收敛半径、Padé 逼近与 Stahl
定理）、同伦延拓（概率一正则性、路径跟踪、转折点）、Newton--Krylov/JFNK
与非精确 Newton 的收敛理论，以及不动点迭代的 Anderson 加速（与 GMRES 的等价）；
并以平台求解器实测数据给出高级求解器收敛与可解性边界的数值基准。本章公式不要求逐式对应代码；与平台实现
（\srcpath{src/power\_flow/} 与 \srcpath{include/hacdcpf/power\_flow/}）的对应关系
见本章末"与实现的对应"表，超出实现的内容以"未实现扩展说明"框就地标记。
实现层的组件级转写见本手册"全局化与鲁棒求解器族"章。
\end{theorynote}

\subsection{记号与约定}\label{sec:pf-th-glob-notation}
\begin{itemize}
  \item $F:\mathbb R^n\to\mathbb R^n$：潮流失配向量（P/Q 平衡、直流功率平衡、
        换流器局部方程等的堆叠）；$x=(\theta,V_m,V_{dc},\ldots)$ 为堆叠状态；
        $J(x)=DF(x)$ 为 Jacobian；
  \item $\phi(x)=\tfrac12\lVert F(x)\rVert_2^2$：最小二乘 merit 函数，
        $g(x)=\nabla\phi(x)=J(x)^{\mathsf T}F(x)$；
  \item $d_k$ 或 $\Delta x_k$：第 $k$ 次迭代的搜索方向；$\alpha_k\in(0,1]$：步长；
        $\beta\in(0,1)$：回溯折减因子；$c_1,c_2$：Armijo/Wolfe 常数；
  \item $\Delta_k$：信赖域半径；$m_k(p)=\phi_k+g_k^{\mathsf T}p+\tfrac12
        p^{\mathsf T}B_k p$（$B_k=J_k^{\mathsf T}J_k$）：二次模型；
        $\rho_k=\mathrm{ared}_k/\mathrm{pred}_k$：增益比；
  \item $\lambda$ 在本章按上下文两用：信赖域/LM 小节中为 LM 阻尼参数，
        同伦小节中为同伦参数；$\mu\ge0$ 专用于 NCP 光滑化参数；
  \item $s\in\mathbb C$：HELM 嵌入参数；$V_i(s)$：母线 $i$ 电压的全纯嵌入；
        $[L/M]_f$：函数 $f$ 的 $(L,M)$ 阶 Padé 逼近；
  \item $\eta_k\in[0,1)$：非精确 Newton 强制项（forcing term）；
        $\mathcal K_m(A,v)$：$m$ 维 Krylov 子空间；
  \item $a\perp b$ 表示互补性 $a\ge0,\ b\ge0,\ ab=0$；
        $\partial F(x)$ 表示 Clarke 广义 Jacobian。
\end{itemize}

\subsection{线搜索与信赖域的全局化框架}\label{sec:pf-th-glob-framework}

潮流 Newton 法的局部二次收敛只在解的邻域内成立；全局化机制回答的问题是：
从任意初值出发，如何修正纯 Newton 步使迭代在全局意义下有收敛保证。
两类标准机制——线搜索与信赖域——都作用在 merit 函数 $\phi$ 上。

\begin{definition}[Armijo 充分下降条件]\label{sec:pf-th-glob-def-armijo}
给定下降方向 $d_k$（即 $g_k^{\mathsf T}d_k<0$）与常数 $c_1\in(0,1)$，
步长 $\alpha_k>0$ 满足 Armijo 条件，若
\begin{equation}
\phi(x_k+\alpha_k d_k)\le \phi(x_k)+c_1\alpha_k\,g_k^{\mathsf T}d_k .
\end{equation}
回溯线搜索从 $\alpha=1$ 出发按 $\alpha\leftarrow\beta\alpha$ 折减，
直到条件成立。由 Taylor 展开，只要 $d_k$ 是下降方向，充分小的 $\alpha$
必满足该条件，故回溯在有限步内终止。
\end{definition}

\begin{definition}[Wolfe 条件]\label{sec:pf-th-glob-def-wolfe}
步长 $\alpha_k$ 满足 Wolfe 条件，若 Armijo 条件成立且曲率条件
\begin{equation}
g(x_k+\alpha_k d_k)^{\mathsf T}d_k\ge c_2\,g_k^{\mathsf T}d_k,
\qquad 0<c_1<c_2<1
\end{equation}
成立。曲率条件排除"步长过小"的接受：它要求新点处方向导数的下降速率
已被足够多地"用掉"。
\end{definition}

\begin{theorem}[Zoutendijk 全局收敛定理]\label{sec:pf-th-glob-thm-zoutendijk}
设 $\phi$ 下有界，$g=\nabla\phi$ 在水平集
$\{x:\phi(x)\le\phi(x_0)\}$ 上 Lipschitz 连续（常数 $L$），
迭代 $x_{k+1}=x_k+\alpha_k d_k$ 的步长满足 Wolfe 条件。则
\begin{equation}
\sum_{k=0}^{\infty}\cos^2\theta_k\,\lVert g_k\rVert^2<\infty,
\qquad
\cos\theta_k=\frac{-g_k^{\mathsf T}d_k}{\lVert g_k\rVert\,\lVert d_k\rVert}.
\end{equation}
特别地，若方向满足\emph{夹角条件} $\cos\theta_k\ge\delta>0$，
则 $\lim_{k\to\infty}\lVert g_k\rVert=0$。
\end{theorem}

\begin{proof}
记 $s_k=\alpha_k d_k$。由曲率条件，
\[
(g_{k+1}-g_k)^{\mathsf T}s_k\ge(c_2-1)\,g_k^{\mathsf T}s_k .
\]
由 Lipschitz 连续性，
\[
(g_{k+1}-g_k)^{\mathsf T}s_k\le\lVert g_{k+1}-g_k\rVert\,\lVert s_k\rVert
\le L\,\alpha_k\lVert d_k\rVert\,\lVert s_k\rVert
=L\,\alpha_k\lVert d_k\rVert^2 .
\]
两式合并得
$\alpha_k\ge\dfrac{(c_2-1)g_k^{\mathsf T}d_k}{L\lVert d_k\rVert^2}
=\dfrac{(1-c_2)\cos\theta_k\lVert g_k\rVert}{L\lVert d_k\rVert}$，
其中用到 $-g_k^{\mathsf T}d_k=\cos\theta_k\lVert g_k\rVert\lVert d_k\rVert$。
代入 Armijo 条件：
\[
\phi_{k+1}\le\phi_k-c_1\alpha_k\,\cos\theta_k\lVert g_k\rVert\lVert d_k\rVert
\le\phi_k-\frac{c_1(1-c_2)}{L}\cos^2\theta_k\lVert g_k\rVert^2 .
\]
递推并对 $k$ 求和，由 $\phi$ 下有界，
$\frac{c_1(1-c_2)}{L}\sum_k\cos^2\theta_k\lVert g_k\rVert^2
\le\phi_0-\lim_k\phi_k<\infty$。
若 $\cos\theta_k\ge\delta>0$，则
$\delta^2\sum_k\lVert g_k\rVert^2<\infty$，故 $\lVert g_k\rVert\to0$。
\end{proof}

\begin{remark}[对非线性方程组的含义]\label{sec:pf-th-glob-rem-stationary}
Zoutendijk 定理保证的是 $g=J^{\mathsf T}F\to0$，即 $\phi$ 的驻点，
而不是 $F=0$。当 $J$ 亏秩时可能存在 $F\ne0$ 而 $J^{\mathsf T}F=0$ 的
\emph{非解驻点}（merit 函数的局部极小）。这是以最小二乘 merit 全局化
非线性方程组的固有局限；实践中以发散探测、重启动与回退链兜底。
Newton 方向 $d_k=-J_k^{-1}F_k$ 在 $J_k$ 非奇异且条件数有界时满足
夹角条件：$g_k^{\mathsf T}d_k=-F_k^{\mathsf T}F_k=-2\phi_k<0$。
\end{remark}

\begin{definition}[信赖域子问题与柯西点]\label{sec:pf-th-glob-def-tr}
信赖域方法在第 $k$ 步求解约束子问题
\begin{equation}
\min_{p\in\mathbb R^n}\ m_k(p)=\phi_k+g_k^{\mathsf T}p
+\tfrac12 p^{\mathsf T}B_k p,
\qquad \lVert p\rVert\le\Delta_k,
\end{equation}
$B_k=J_k^{\mathsf T}J_k$（Gauss--Newton 模型）。沿负梯度方向的约束极小点
称为\emph{柯西点}：
\begin{equation}
p_k^{C}=-\tau_k\frac{\Delta_k}{\lVert g_k\rVert}\,g_k,
\qquad
\tau_k=\begin{cases}1,&\lVert g_k\rVert^3\le\Delta_k\,g_k^{\mathsf T}B_k g_k
\ \text{或}\ g_k^{\mathsf T}B_k g_k\le0,\\[2pt]
\dfrac{\lVert g_k\rVert^3}{\Delta_k\,g_k^{\mathsf T}B_k g_k},&\text{其余}.\end{cases}
\end{equation}
即无约束最速下降极小点 $\bar\tau=\lVert g_k\rVert^3/(g_k^{\mathsf T}B_k g_k)$
对应的点未出界时取其本身，否则截断到边界。
\end{definition}

\begin{lemma}[柯西点的充分下降]\label{sec:pf-th-glob-lem-cauchy}
柯西点满足
\begin{equation}
m_k(p_k^{C})-m_k(0)\le-\tfrac12\lVert g_k\rVert\,
\min\Bigl(\Delta_k,\ \frac{\lVert g_k\rVert}{\lVert B_k\rVert}\Bigr).
\end{equation}
\end{lemma}

\begin{proof}
沿 $p(\tau)=-\tau g_k/\lVert g_k\rVert$（$\tau\in(0,\Delta_k]$）有
\[
m_k(p(\tau))-m_k(0)=-\tau\lVert g_k\rVert
+\frac{\tau^2}{2\lVert g_k\rVert^2}\,g_k^{\mathsf T}B_k g_k .
\]
\textbf{情形 1}：$g_k^{\mathsf T}B_k g_k\le0$。二次项非正，
$\tau_k=1$，故
$m_k-m_k(0)\le-\Delta_k\lVert g_k\rVert\le-\tfrac12\lVert g_k\rVert
\min(\Delta_k,\lVert g_k\rVert/\lVert B_k\rVert)$。

\textbf{情形 2}：$g_k^{\mathsf T}B_k g_k>0$ 且
$\bar\tau=\lVert g_k\rVert^3/(g_k^{\mathsf T}B_k g_k)\le\Delta_k$。
取 $\tau=\bar\tau$ 得
$m_k-m_k(0)=-\tfrac12\lVert g_k\rVert^4/(g_k^{\mathsf T}B_k g_k)
=-\tfrac12\lVert g_k\rVert\bar\tau$。
由 $g_k^{\mathsf T}B_k g_k\le\lVert B_k\rVert\lVert g_k\rVert^2$ 有
$\bar\tau\ge\lVert g_k\rVert/\lVert B_k\rVert$；结合
$\bar\tau\le\Delta_k$ 知
$\min(\Delta_k,\lVert g_k\rVert/\lVert B_k\rVert)\le\bar\tau$，结论成立。

\textbf{情形 3}：$g_k^{\mathsf T}B_k g_k>0$ 且 $\bar\tau>\Delta_k$。
取 $\tau=\Delta_k$，此时 $\Delta_k^2 g_k^{\mathsf T}B_k g_k/
(2\lVert g_k\rVert^2)<\tfrac12\Delta_k\lVert g_k\rVert$，故
$m_k-m_k(0)<-\tfrac12\Delta_k\lVert g_k\rVert
\le-\tfrac12\lVert g_k\rVert\min(\Delta_k,\lVert g_k\rVert/\lVert B_k\rVert)$。
\end{proof}

\begin{theorem}[信赖域全局收敛（陈述）]\label{sec:pf-th-glob-thm-tr}
设迭代只接受满足 $\mathrm{pred}_k=m_k(0)-m_k(p_k)$ 不小于柯西点下降之固定
比例（如 $\tfrac14$）的试探步，并按增益比
$\rho_k=\mathrm{ared}_k/\mathrm{pred}_k$ 调节半径
（$\rho_k$ 小则收缩并拒绝，$\rho_k$ 接近 1 则扩张）。在 $\phi$ 下有界、
$g$ Lipschitz、$\lVert B_k\rVert$ 一致有界的假设下，
$\liminf_{k\to\infty}\lVert g_k\rVert=0$；若再有 $\eta>0$ 的接受阈值
与聚点处 Jacobian 非奇异，则整个序列收敛且渐近取纯 Newton 步。
\end{theorem}

\begin{proof}[证明思路]
若 $\sum\Delta_k$ 在某一拒绝循环中保持有正下界以外的行为，则对足够小
半径，Taylor 余项使 $\rho_k\to1$，必出现接受步；于是无限次接受的子列上
$\mathrm{ared}_k\ge\eta\,\mathrm{pred}_k
\ge\frac{\eta}{4}\cdot\tfrac12\lVert g_k\rVert
\min(\Delta_k,\lVert g_k\rVert/\lVert B_k\rVert)$。
由 $\phi_k$ 单调下降且下有界，$\mathrm{ared}_k\to0$，从而
$\lVert g_k\rVert\min(\Delta_k,\cdot)\to0$；若 $\lVert g_k\rVert$ 有正下界
则 $\Delta_k\to0$，与"半径充分小必接受"矛盾。完整论证见
Nocedal \& Wright~\cite{pfglob-nw} 第 4 章。
\end{proof}

\begin{remark}[dogleg 折线步]\label{sec:pf-th-glob-rem-dogleg}
精确求解信赖域子问题代价高；Powell 的 dogleg
方法~\cite{pfglob-powell,pfglob-ds} 以折线近似解轨迹：当 Gauss--Newton 步
$p^{GN}=-(J^{\mathsf T}J)^{-1}g$ 在界内时取之；否则沿"柯西点 $\to$ Newton
步"的折线走到边界。折线包含柯西点，故引理~\ref{sec:pf-th-glob-lem-cauchy}
的下降估计成立，全局收敛性随之保证。
\end{remark}

\begin{theorem}[Levenberg--Marquardt 与信赖域的对应]\label{sec:pf-th-glob-thm-lmtr}
设 $\lambda>0$ 使 $B+\lambda I=J^{\mathsf T}J+\lambda I$ 正定，LM 步
\begin{equation}
p^{LM}(\lambda)=-\bigl(J^{\mathsf T}J+\lambda I\bigr)^{-1}J^{\mathsf T}F
\end{equation}
恰好是信赖域子问题在半径 $\Delta=\lVert p^{LM}(\lambda)\rVert$ 下的全局解；
反之，信赖域解在 $\lVert p^*\rVert=\Delta$（边界情形）时必具有 LM 形式。
\end{theorem}

\begin{proof}
信赖域子问题的最优性条件（Nocedal \&
Wright~\cite{pfglob-nw} 定理 4.1；Moré~\cite{pfglob-more}）：
$p^*$ 为全局解当且仅当存在 $\lambda^*\ge0$ 使
\begin{equation}
(B+\lambda^* I)p^*=-g,\qquad B+\lambda^* I\succeq0,\qquad
\lambda^*\bigl(\Delta-\lVert p^*\rVert\bigr)=0.
\end{equation}
对给定 $\lambda>0$，$B+\lambda I$ 正定，$p^{LM}(\lambda)$ 由第一式唯一
确定，互补条件以 $\Delta=\lVert p^{LM}(\lambda)\rVert$ 平凡成立，故
$p^{LM}(\lambda)$ 是该半径下的全局解。反向由互补条件：边界解
（$\lVert p^*\rVert=\Delta$）要求 $\lambda^*\ge0$ 且满足第一式，即
LM 形式。
\end{proof}

\begin{remark}\label{sec:pf-th-glob-rem-lmtr}
定理~\ref{sec:pf-th-glob-thm-lmtr} 表明 LM 阻尼与信赖域半径是同一族正则化
步长的两种参数化：$\lambda$ 单调地反比于有效半径
（$\lambda\to0$ 时 $p^{LM}\to p^{GN}$，$\lambda\to\infty$ 时步长趋于零且
方向趋于最速下降）。因此实现中以增益比 $\rho$ 调节 $\lambda$
（接受且 $\rho$ 大则 $\lambda/3$，拒绝则 $10\lambda$）与以 $\rho$ 调节
$\Delta$ 在理论上等价，差别只在参数化与数值细节。LM 的历史出处为
Levenberg~\cite{pfglob-levenberg} 与 Marquardt~\cite{pfglob-marquardt}。
\end{remark}

\begin{remark}[伪瞬态延拓（PTC）]\label{sec:pf-th-glob-rem-ptc}
PTC 将方程 $F(x)=0$ 视为伪动力系统 $\dot x=-F(x)$ 的稳态，以隐式 Euler
离散的 Newton 修正求解 $(J+\tfrac1{\delta t}I)\Delta x=-F$：
$\delta t\to0$ 时为强阻尼的显式 Euler 小步，$\delta t\to\infty$ 时退化为
纯 Newton 步。其收敛性分析（稳态吸引域内局部指数收敛到稳态）见
Kelley \& Keyes~\cite{pfglob-kelley-keyes}。SER（Self-Equalising Residual）
更新 $\delta t\leftarrow\delta t\,(r_{k-1}/r_k)^\gamma$ 是实践上有效的
启发式调度，其理论性质没有与 Zoutendijk 定理同等严格的结果。
\end{remark}

\begin{gapnote}
精确信赖域子问题求解（Moré--Sorensen 迭代）与大规模情形常用的
Steihaug--Toint 截断共轭梯度法在当前代码中均未实现；信赖域分支采用
dogleg 折线近似（经 \srcpath{include/hacdcpf/engine/kernel/globalization/globalization.hpp}
转发到 MIPSolvers 的 \srcpath{dogleg\_step}）。
此外实现只检查 Armijo 充分下降，未检查 Wolfe 曲率条件——对非线性方程组
merit 这是常见简化，但严格意义上
定理~\ref{sec:pf-th-glob-thm-zoutendijk} 的假设因此不被完整满足。
\end{gapnote}

\subsection{非单调线搜索（GLL）}\label{sec:pf-th-glob-gll}

单调 Armijo 搜索要求 merit 每步下降；当迭代序列需穿过 merit 面的
"窄谷"时，这会导致过小的步长与停滞。Grippo--Lampariello--Lucidi
（GLL）~\cite{pfglob-gll} 放松了单调性。

\begin{definition}[GLL 非单调 Armijo 条件]\label{sec:pf-th-glob-def-gll}
取窗口 $M\ge1$，定义非单调参考值
\begin{equation}
\phi_k^{R}=\max_{0\le j\le\min(k,M-1)}\phi_{k-j},
\end{equation}
接受条件为
\begin{equation}
\phi(x_k+\alpha_k d_k)\le\phi_k^{R}
+c_1\alpha_k\,g_k^{\mathsf T}d_k .
\end{equation}
$M=1$ 时退化为单调 Armijo 搜索；$M>1$ 时容许 merit 在窗口内暂时上升，
只要相对最近 $M$ 步的最坏值仍有充分下降。
\end{definition}

\begin{theorem}[GLL 全局收敛性（论证思路）]\label{sec:pf-th-glob-thm-gll}
在定理~\ref{sec:pf-th-glob-thm-zoutendijk} 的光滑性假设与夹角条件
$\cos\theta_k\ge\delta>0$ 下，采用 GLL 接受准则的迭代满足
$\liminf_k\lVert g_k\rVert=0$。
\end{theorem}

\begin{proof}[证明思路]
记参考下标 $\ell(k)$ 为窗口内取得最大值者：
$\phi_{\ell(k)}=\phi_k^{R}$。论证分三步（细节见
Grippo 等~\cite{pfglob-gll}）。
\textbf{(i)} 参考序列单调：由接受条件，
$\phi_{k+1}\le\phi_{\ell(k)}+c_1\alpha_k g_k^{\mathsf T}d_k
\le\phi_{\ell(k)}$，故下一段窗口的最大值不超过上一段，
$\{\phi_{\ell(k)}\}$ 非增且下有界，收敛到某 $\phi^R$。
\textbf{(ii)} 参考子列充分下降：沿参考下标构成的子列应用与
Zoutendijk 定理相同的估计，得
$\sum\cos^2\theta_{\ell(k)}\lVert g_{\ell(k)}\rVert^2<\infty$，
故 $\lVert g_{\ell(k)}\rVert\to0$。
\textbf{(iii)} 全序列由参考子列控制：每个 $k$ 距其参考下标不超过
$M$ 步；窗口内每步的回溯把 merit 增长控制在参考值的有限邻域内，
方向与步长的一致有界性再把 $g_k$ 与 $g_{\ell(k)}$ 的差控制在
$o(1)$，从而 $\liminf\lVert g_k\rVert=0$。
\end{proof}

\begin{remark}\label{sec:pf-th-glob-rem-gll}
窗口 $M$ 是鲁棒性与速度的工程折中：$M$ 大则更宽容、可能接受更多全
Newton 步（减少回溯次数与 Jacobian 重分解），但非单调震荡的幅度也更
大；$M$ 小则接近单调行为。潮流实践中 $M\in[3,10]$ 是常见取值。
\end{remark}

\subsection{Anderson 加速}\label{sec:pf-th-glob-anderson}

前推回代、FDPF 与 HELM 胚迭代等\emph{不动点}方法只具线性收敛
（第~\ref{sec:pf-th-cls}~章）。Anderson 加速~\cite{pfglob-anderson} 利用历史残差
外推，常能显著提升其收敛速度，且不需要 Jacobian。

\begin{definition}[Anderson 加速 AA($m$)]\label{def:pf-th-glob-anderson}
设不动点迭代 $x\mapsto g(x)$、残差 $f(x)=g(x)-x$。取记忆深度 $m$，第 $k$ 步令
$m_k=\min(m,k)$，用最近 $m_k{+}1$ 个残差解约束最小二乘
\begin{equation}
\alpha^{(k)}=\arg\min_{\sum_i\alpha_i=1}
\Bigl\lVert\sum_{i=0}^{m_k}\alpha_i\,f_{k-m_k+i}\Bigr\rVert_2,
\end{equation}
再取（阻尼因子 $\beta\in(0,1]$，$\beta{=}1$ 为无阻尼）
\begin{equation}
x_{k+1}=(1-\beta)\sum_{i=0}^{m_k}\alpha_i^{(k)}x_{k-m_k+i}
+\beta\sum_{i=0}^{m_k}\alpha_i^{(k)}g(x_{k-m_k+i}).
\end{equation}
$m=0$ 退化为（阻尼）Picard 迭代。约束最小二乘可等价化为对残差差分
$\Delta f_i=f_{i+1}-f_i$ 的无约束最小二乘，规模 $n\times m_k$，代价低。
\end{definition}

\begin{theorem}[线性情形与 GMRES 的等价（Walker--Ni）]\label{thm:pf-th-glob-anderson-gmres}
设 $g(x)=Ax+b$ 为仿射映射，不动点即解 $(I-A)x=b$。则无阻尼、全记忆（$m=\infty$）的
Anderson 加速在不停滞前提下，其第 $k$ 步迭代 $x_k$ 与 GMRES 施于 $(I-A)x=b$ 的第 $k$
步迭代重合，残差范数逐步相等 $\lVert f_k\rVert=\lVert r_k^{\mathrm{GMRES}}\rVert$
~\cite{pfglob-walkerni}。
\end{theorem}
\begin{proof}[证明要点]
无约束形式下 AA 更新为 $x_{k+1}=x_k+\beta f_k-(\Delta X_k+\beta\Delta F_k)\gamma^{(k)}$，
$\gamma^{(k)}=\arg\min_\gamma\lVert f_k-\Delta F_k\gamma\rVert$，其中
$\Delta X_k,\Delta F_k$ 收集历史差分 $\Delta x_i,\Delta f_i$。仿射映射使
$\Delta f_i=(A-I)\Delta x_i$，故 $\Delta F_k=(A-I)\Delta X_k$；最小二乘
$\min_\gamma\lVert f_k-\Delta F_k\gamma\rVert$ 是在 Krylov 仿射空间上极小化
$(I-A)$-残差，正是 GMRES 的最小残差性质（命题~\ref{sec:pf-th-glob-prop-gmres}）。
逐步对照生成子空间即得残差相等；完整对应见 Walker--Ni~\cite{pfglob-walkerni} 定理 2.2。
\end{proof}

\begin{remark}[非线性收敛与工程意义]\label{rem:pf-th-glob-anderson}
对压缩不动点 $g$（Lipschitz 常数 $\kappa<1$），AA($m$) 局部 r-线性收敛，渐近因子不劣于
Picard，实践中常显著更快~\cite{pfglob-tothkelley}。定理
\ref{thm:pf-th-glob-anderson-gmres} 揭示其本质：AA 是不动点迭代上的“类 GMRES”外推，
以 $O(m)$ 存储与一次小最小二乘换取超线性加速——对配网 BFS、FDPF 这类线性收敛且单步
廉价的方法尤为契合。$m$ 过大时 $\Delta F_k$ 病态，需列主元 QR 或 Tikhonov 正则与条件
数门限兜底。
\end{remark}

\begin{gapnote}
Anderson 加速在平台中\textbf{未实现}：不动点类求解器（\srcpath{run\_bfs\_sweep}、
\srcpath{FDPFSolver::solve}）按朴素 Picard/拟 Newton 迭代运行，HELM 用固定阶 Padé 而非
残差外推；平台的加速手段是 Newton 类（二次收敛）与 GMRES 内迭代
（\srcpath{src/power\_flow/newton\_krylov.cpp:gmres\_solve}），非不动点外推。
\end{gapnote}

\subsection{互补问题与半光滑 Newton 方法}\label{sec:pf-th-glob-ncp}

发电机无功限值（PV$\to$PQ 切换）与换流器限流都是"控制器在若干运行分支
间切换"的问题。统一它们的数学对象是互补问题。

\begin{definition}[互补问题与 NCP 函数]\label{sec:pf-th-glob-def-ncp}
标量互补条件写作 $0\le a\perp b\ge0$，即 $a\ge0,\ b\ge0,\ ab=0$。
函数 $\varphi:\mathbb R^2\to\mathbb R$ 称为 \emph{NCP 函数}，若
\begin{equation}
\varphi(a,b)=0\iff a\ge0,\ b\ge0,\ ab=0 .
\end{equation}
把逐分量互补条件堆叠进 $F(x)=0$，即用 NCP 函数把混合不等式-等式系统
化为单个非光滑方程组。
\end{definition}

\begin{proposition}[Fischer--Burmeister 函数]\label{sec:pf-th-glob-prop-fb}
Fischer--Burmeister（FB）函数~\cite{pfglob-fischer}
\begin{equation}
\Phi_{FB}(a,b)=\sqrt{a^2+b^2}-a-b
\end{equation}
是 NCP 函数。
\end{proposition}

\begin{proof}
$\sqrt{a^2+b^2}\ge\max(|a|,|b|)$。
（$\Leftarrow$）若 $a,b\ge0$ 且 $ab=0$，不妨 $b=0$，则
$\sqrt{a^2}-a-0=a-a=0$。
（$\Rightarrow$）设 $\Phi_{FB}(a,b)=0$，即 $\sqrt{a^2+b^2}=a+b$。
右端须非负，且 $a+b\ge\max(|a|,|b|)\ge|a|$ 给出 $b\ge0$，对称地
$a\ge0$。平方得 $a^2+b^2=a^2+2ab+b^2$，即 $ab=0$。
\end{proof}

\begin{proposition}[FB 的光滑性结构]\label{sec:pf-th-glob-prop-fbsmooth}
$\Phi_{FB}$ 在 $(0,0)$ 以外处处连续可微，
\begin{equation}
\nabla\Phi_{FB}(a,b)=\Bigl(\frac{a}{\sqrt{a^2+b^2}}-1,\
\frac{b}{\sqrt{a^2+b^2}}-1\Bigr);
\end{equation}
在原点不可微，但 $\Psi(a,b)=\tfrac12\Phi_{FB}(a,b)^2$ 在全平面连续
可微且 $\nabla\Psi(0,0)=0$（原点外 $\nabla\Psi=\Phi_{FB}\nabla\Phi_{FB}$，
其模被 $\sqrt2\,|\Phi_{FB}|$ 控制，随 $(a,b)\to0$ 趋于零）。
\end{proposition}

\begin{definition}[Clarke 广义 Jacobian 与半光滑性]\label{sec:pf-th-glob-def-clarke}
局部 Lipschitz 函数 $F:\mathbb R^n\to\mathbb R^m$ 在 $x$ 处的
B-次微分 $\partial_B F(x)$ 为所有"可微点列 Jacobian 的极限"之集合，
Clarke 广义 Jacobian 为其凸包
$\partial F(x)=\mathrm{co}\,\partial_B F(x)$。
$F$ 在 $x$ \emph{半光滑}，若对任意 $h\to0$ 与
$V\in\partial F(x+h)$ 有
\begin{equation}
F(x+h)-F(x)-Vh=o(\lVert h\rVert);
\end{equation}
若加强为 $O(\lVert h\rVert^2)$，则称\emph{强半光滑}。
分片光滑函数（如 $\min/\max/\Phi_{FB}$）皆强半光滑
（Facchinei \& Pang~\cite{pfglob-fp} 第 7 章）。
\end{definition}

\begin{example}[FB 在原点的广义 Jacobian]\label{sec:pf-th-glob-ex-fb}
$\Phi_{FB}$ 在原点的 B-次微分是单位圆周平移后的点集
$\{(\frac{u}{r}-1,\frac{v}{r}-1):u^2+v^2=r^2,\ r>0\}$ 的闭包
（以 $1/\sqrt2$ 为半径的圆），其 Clarke 广义 Jacobian 为对应圆盘。
实现对偶激活点 $(a,b)=(0,0)$ 选取对称元素
$(1/\sqrt2-1,\ 1/\sqrt2-1)$——即广义 Jacobian 的"中心"元，
避免对某一侧分支的偏向。
\end{example}

\begin{theorem}[半光滑 Newton 的局部超线性收敛（Qi--Sun）]\label{sec:pf-th-glob-thm-qisun}
设 $F$ 半光滑，$x^*$ 满足 $F(x^*)=0$，且 $\partial F(x^*)$ 中所有元素
非奇异（BD-正则性）。则广义 Newton 迭代
\begin{equation}
x_{k+1}=x_k-V_k^{-1}F(x_k),\qquad V_k\in\partial F(x_k)
\end{equation}
在 $x^*$ 的邻域内适定并 q-超线性收敛；若 $F$ 强半光滑，收敛是 q-二次的
（Qi \& Sun~\cite{pfglob-qisun}）。
\end{theorem}

\begin{proof}[证明思路]
BD-正则性与广义 Jacobian 的上半连续性给出 $x^*$ 邻域内
$\lVert V^{-1}\rVert$ 的一致界 $C$。于是
\[
x_{k+1}-x^*=V_k^{-1}\bigl[F(x^*)-F(x_k)-V_k(x^*-x_k)\bigr],
\]
括号内由半光滑性为 $o(\lVert x_k-x^*\rVert)$，故
$\lVert x_{k+1}-x^*\rVert=o(\lVert x_k-x^*\rVert)$，即超线性；
强半光滑时余项为 $O(\lVert x_k-x^*\rVert^2)$，得二次收敛。
邻域适定性由同一致界保证 $V_k$ 非奇异。
\end{proof}

\begin{remark}[BD-正则性不可由单点证书替代]\label{sec:pf-th-glob-rem-bd}
定理~\ref{sec:pf-th-glob-thm-qisun} 的假设要求解点邻域内
B-次微分的\emph{所有}元素非奇异，这比"所选广义 Jacobian 元非奇异"
强。运行中实测的速率商
$\lVert F_{k+1}\rVert/\lVert F_k\rVert$ 只是局部速率的经验证据，
不能反推 BD-正则性成立；退化分支交点（例如两个限值同时临界）处
所选广义 Jacobian 可能病态，实现以后向误差证书拒绝该步并回退全
稀疏 LU。
\end{remark}

\subsubsection{PV/PQ 切换的统一表述：中位数互补方程}

\begin{proposition}[中位数 NCP 方程的零点集]\label{sec:pf-th-glob-prop-median}
设母线电压偏差 $e=V_m-V_m^{set}$，无功裕度
$\underline q=Q_g-Q_g^{\min}$、$\bar q=Q_g-Q_g^{\max}$
（$\underline q>\bar q$）。方程
\begin{equation}
H(e,\underline q,\bar q)
=\max\bigl(\min(e,\underline q),\ \bar q\bigr)=0
\end{equation}
的零点集恰为以下三个物理分支的并：
\begin{equation}
\text{(i)}\ e=0,\ Q_g^{\min}\le Q_g\le Q_g^{\max};\qquad
\text{(ii)}\ Q_g=Q_g^{\min},\ e\ge0;\qquad
\text{(iii)}\ Q_g=Q_g^{\max},\ e\le0 .
\end{equation}
即"限值内为 PV（电压控制）"、触下/上限时转为 PQ 且电压向正确方向
松弛。
\end{proposition}

\begin{proof}
记 $u=\min(e,\underline q)$，则 $H=\max(u,\bar q)=0$ 当且仅当
$u\le0$，$\bar q\le0$，且两者至少其一为零。
\textbf{(i)} 若 $\bar q<0$（未触上限），须 $u=0$，即
$\min(e,\underline q)=0$：$e=0$ 且 $\underline q\ge0$，或
$\underline q=0$ 且 $e\ge0$（并入情形 (ii) 的边界）。
\textbf{(ii)} 若 $\underline q=0$（触下限），$u=\min(e,0)=0$
当且仅当 $e\ge0$，而 $\bar q=\underline q-(Q_g^{\max}-Q_g^{\min})<0$
自动满足。
\textbf{(iii)} 若 $\bar q=0$（触上限），须
$u=\min(e,\underline q)\le0$；此时 $\underline q>0$，
故须 $e\le0$。
三种情形之外：$\bar q<0$、$\underline q>0$、$e\ne0$ 时
$u=\min(e,\underline q)\ne0$ 且若 $e<0$ 则 $u<0$ 使 $H<0$，
$e>0$ 时 $u>0$ 使 $H>0$，均非零。
\end{proof}

\begin{remark}[与原始-对偶积极集的关系]\label{sec:pf-th-glob-rem-activeset}
把不等式约束的积极集显式枚举（PV$\to$PQ 批量切换加有界恢复审计）是
同一数学对象的另一算法化路径：对固定积极集解光滑系统、在收敛点处
按违反度更新积极集，其有限识别性质与半光滑 Newton 等价地建立在
互补结构上（Hintermüller--Ito--Kunisch~\cite{pfglob-hik}）。两者的
差别在工程侧：积极集切换改变方程布局与稀疏模式（需重做符号分析），
NCP 方程保持固定布局但引入非光滑点。
\end{remark}

\begin{definition}[CHKS 与 Kanzow 光滑化]\label{sec:pf-th-glob-def-smoothing}
CHKS 光滑化（Chen--Harker~\cite{pfglob-chenharker} 与
Kanzow~\cite{pfglob-kanzow} 的非内点延拓构造）以
\begin{equation}
\max\nolimits_\mu(a,b)=\tfrac12\bigl(a+b+\sqrt{(a-b)^2+4\mu^2}\bigr),
\qquad \mu>0
\end{equation}
替代 $\max$（$\min$ 对偶地），Kanzow 光滑 FB 为
\begin{equation}
\Phi_\mu(a,b)=\sqrt{a^2+b^2+2\mu^2}-a-b .
\end{equation}
两者在 $\mu>0$ 时光滑，$\mu\to0$ 时逐点（实际一致）收敛到对应非光滑
算子。
\end{definition}

\begin{proposition}[光滑化的一致逼近阶]\label{sec:pf-th-glob-prop-smootherr}
对所有 $(a,b)\in\mathbb R^2$ 与 $\mu>0$，
\begin{equation}
\bigl|\Phi_\mu(a,b)-\Phi_{FB}(a,b)\bigr|\le\sqrt2\,\mu .
\end{equation}
\end{proposition}

\begin{proof}
记 $r=\sqrt{a^2+b^2}$，则
$\Phi_\mu-\Phi_{FB}=\sqrt{r^2+2\mu^2}-r
=\dfrac{2\mu^2}{\sqrt{r^2+2\mu^2}+r}\le\dfrac{2\mu^2}{\sqrt2\,\mu}
=\sqrt2\,\mu$。
\end{proof}

\begin{remark}[光滑延拓的收敛逻辑]\label{sec:pf-th-glob-rem-smoothcont}
以 $\mu_k\downarrow0$ 的调度依次解光滑方程 $F_{\mu_k}(x)=0$、每个
$\mu$ 层以前一层解为初值，是标准的非内点延拓
（Kanzow~\cite{pfglob-kanzow}）：
命题~\ref{sec:pf-th-glob-prop-smootherr} 保证 $F_{\mu_k}\to F$ 一致，
若极限方程的解满足 BD-正则性，则足够小的 $\mu$ 层落在下一层的
Newton 吸引域内，整条延拓链以有限次逐层 Newton 求解完成。工程上
以两阶段调度（粗阶段慢退火、近收敛快退火）平衡鲁棒性与迭代数；
最终证书必须以 $\mu=0$ 的精确 NCP 残差重估，防止把光滑代理的收敛
误报为精确互补性。
\end{remark}

\subsection{全纯嵌入负荷流方法（HELM）的数学基础}\label{sec:pf-th-glob-helm}

HELM（Trias~\cite{pfglob-trias}）把潮流方程嵌入复参数 $s$ 的全纯族，
用解析延拓代替迭代。

\begin{definition}[全纯嵌入与幂级数]\label{sec:pf-th-glob-def-helm}
对 PQ 母线 $i$、PV 母线 $j$ 与平衡母线 $k$，嵌入取
\begin{align}
\sum_{\ell}Y_{i\ell}\,V_\ell(s)
&=s\,\frac{S_i^{*}}{V_i^{*}(s^{*})}-s\,Y_i^{sh}\,V_i(s),\\
\sum_{\ell}Y_{j\ell}\,V_\ell(s)
&=s\,\frac{P_j-\mathrm j Q_j(s)}{V_j^{*}(s^{*})},\qquad
V_j(s)V_j^{*}(s^{*})=W_j+s\bigl(|V_j^{set}|^2-W_j\bigr),\\
V_k(s)&=V_k^{spec},
\end{align}
其中 $W_j=|V_j(0)|^2$。$s=0$ 对应无注入\emph{胚}（germ）系统，
其解 $V(0)$ 由带平衡母线恒等行的齐次线性系统给出；$s=1$ 还原原
潮流方程。假设 $s\mapsto V(s)$ 在 $s=0$ 邻域全纯，展开
$V_i(s)=\sum_{n\ge0}c_i^{(n)}s^n$。
\end{definition}

\begin{proposition}[系数递推]\label{sec:pf-th-glob-prop-helmrec}
若胚系统矩阵非奇异，则各阶系数由一系列具有\emph{同一}系数矩阵的
线性系统递推确定：定义逆电压级数 $W_i(s)=1/V_i(s)=\sum_n
w_i^{(n)}s^n$，Cauchy 乘积 $V_i(s)W_i(s)=1$ 给出卷积递推
$\sum_{m=0}^{n}c_i^{(m)}w_i^{(n-m)}=\delta_{n0}$；
代入嵌入方程逐阶比较 $s^n$ 的系数，第 $n$ 阶右端只含
$c^{(0)},\ldots,c^{(n-1)}$（及 $w^{(n-1)}$）。
故全部系数经一次矩阵分解加多次回代求得。
\end{proposition}

\begin{proof}[证明思路]
嵌入方程除 $1/V_i(s^{*})$ 项外对 $V(s)$ 是线性的；$s=0$ 处
$V_i(0)\ne0$（胚非奇异）保证 $W_i(s)$ 同邻域全纯。代入幂级数并
比较系数是标准操作；线性项贡献第 $n$ 阶未知量
$c^{(n)}$ 的胚矩阵作用，含 $s$ 因子的项把注入右端推迟一阶，
故第 $n$ 阶系统闭式可解。严格性由隐函数定理的复形式保证
（胚处 Jacobian 即胚矩阵，非奇异）。
\end{proof}

\begin{theorem}[Cauchy--Hadamard 收敛半径]\label{sec:pf-th-glob-thm-ch}
幂级数 $\sum_n c^{(n)}s^n$ 的收敛半径
\begin{equation}
R=\Bigl(\limsup_{n\to\infty}\lVert c^{(n)}\rVert^{1/n}\Bigr)^{-1}
\end{equation}
等于 $s=0$ 到 $V(s)$ 最近奇点的距离；级数在 $|s|<R$ 内收敛到嵌入解，
在 $|s|>R$ 发散。
\end{theorem}

\begin{remark}[奇点的物理含义]\label{sec:pf-th-glob-rem-helmsing}
潮流方程的解支在鞍节点（saddle-node，电压失稳临界点）处两支合并，
映射 $s\mapsto V(s)$ 在该处出现平方根型分支点。因此 $R$ 度量的是
沿嵌入路径到最近鞍节点的"复距离"：重载系统分支点靠近 $s=1$，
$R<1$ 时原级数无法在 $s=1$ 处求值——这正是需要解析延拓的原因。
\end{remark}

\begin{definition}[Padé 逼近]\label{sec:pf-th-glob-def-pade}
函数 $f$ 的 $(L,M)$ 阶 Padé 逼近是满足
\begin{equation}
f(z)-\frac{P_L(z)}{Q_M(z)}=O(z^{L+M+1}),\qquad Q_M(0)=1
\end{equation}
的有理函数 $[L/M]_f=P_L/Q_M$（$\deg P_L\le L,\ \deg Q_M\le M$）。
其系数由幂级数系数的线性方程组（Toeplitz 结构）确定；
Padé 逼近的极点自动逼近 $f$ 的奇点，因而可在收敛圆外继续
$f$——这是有理逼近相对幂级数的本质优势
（Baker \& Graves-Morris~\cite{pfglob-baker}）。
\end{definition}

\begin{theorem}[Stahl 定理（陈述）]\label{sec:pf-th-glob-thm-stahl}
设 $f$ 在 $z=0$ 解析，在 $\mathbb C$ 中只有有限多个分支点型奇点。
则对角（及近对角）Padé 序列 $[n/n]_f$ 在 $\mathbb C$ 去掉一个
对数容量最小的割集系统后的紧子集上\emph{依容量收敛}到 $f$
（Stahl~\cite{pfglob-stahl}）；分母零点的极限分布集中在该极值割集上。
\end{theorem}

\begin{remark}[对 HELM 的操作性含义]\label{sec:pf-th-glob-rem-stahlops}
Stahl 定理把"Padé 是否可信"化为奇点结构问题：若嵌入解支在实轴段
$[0,1]$ 上无奇点，则（在有限分支点的假设下）近对角 Padé 在 $s=1$
处收敛到物理解；反之，若 $[0,1]$ 内存在鞍节点分支点，Padé 部分
极点被吸引到该处，逼近出现震荡/发散。实现据此把"Padé 逼近不收敛"
作为潮流不可行（或解支不可达）的实用判据。严格地说，依容量收敛
容许零容量例外集，且"发散必因不可行"的逆否口径依赖奇点结构假设，
故该判据是强证据而非纯数学证书；其实现层边界见本手册
"全局化与鲁棒求解器族"章的 HELM 小节。
\end{remark}

\begin{gapnote}
Stahl 定理的完整假设核验（嵌入映射奇点结构的辨识、极值割集计算）
与多解分支的系统性枚举在当前代码中均未实现；HELM 求解器采用固定的
对角 Padé 矩阵法与相邻逼近差的启发式收敛判据
（\srcpath{src/power\_flow/helm\_solver.cpp:pade\_at\_one}），
未做容量意义下的误差分析。
\end{gapnote}

\subsection{同伦延拓理论}\label{sec:pf-th-glob-homotopy}

\begin{definition}[同伦映射]\label{sec:pf-th-glob-def-homotopy}
设 $F(x)=0$ 为目标方程，$G(x)=0$ 为有已知解 $x_0$ 的平凡方程。
同伦是满足 $H(x,0)=G(x)$、$H(x,1)=F(x)$ 的连续映射
$H:\mathbb R^n\times[0,1]\to\mathbb R^n$。线性凸同伦
$H(x,\lambda)=(1-\lambda)G(x)+\lambda F(x)$ 是最简单的构造；
对潮流，等价做法是按 $\lambda$ 缩放负荷/注入：
$H(x,\lambda)=F_\lambda(x)$，其中 $\lambda=0$ 为零注入平启动系统，
$\lambda=1$ 为目标工况。
\end{definition}

\begin{theorem}[概率一正则性（参数化 Sard 定理，陈述）]\label{sec:pf-th-glob-thm-sard}
对带参数 $a\in\mathbb R^p$ 的光滑同伦 $H(x,\lambda;a)$，
若 $0$ 是 $(x,\lambda,a)$ 联合映射的正则值，则对 Lebesgue 几乎
所有 $a$，$0$ 是 $H(\cdot,\cdot;a)$ 的正则值，从而
$H^{-1}(0)$ 是 $\mathbb R^n\times[0,1]$ 中的一维光滑流形——
其连通分支或为闭曲线，或为两端落在 $\lambda\in\{0,1\}$ 截口上的
光滑路径；从 $(x_0,0)$ 出发的路径或以有限弧长到达 $\lambda=1$
截口（得到解），或发散至无穷。
（Allgower \& Georg~\cite{pfglob-allgower}。）
\end{theorem}

\begin{remark}[证明思想]\label{sec:pf-th-glob-rem-sard}
核心是 Sard 定理的参数化形式：临界值集测度为零，故"路径横截
$\lambda$ 截口"对几乎一切初值/参数成立。这把延拓法从"碰运气的
启发式"提升为有概率一保证的方法；剩下的失败模式只有路径无界
（目标问题无解或解支断裂）与数值跟踪丢失。
\end{remark}

\begin{definition}[预测-校正路径跟踪]\label{sec:pf-th-glob-def-pc}
增广系统 $H(x,\lambda)=0$ 的一维解流形以弧长 $s$ 参数化。
预测步沿切向 $(\dot x,\dot\lambda)$（由
$[D_xH\ \ D_\lambda H](\dot x,\dot\lambda)^{\mathsf T}=0$ 与归一化
条件确定）取 Euler 步；校正步在垂直于切向的超平面上以 Newton 法
回到路径（Newton 校正子的 Jacobian 为加边矩阵，非奇异性由路径
正则性保证）。
\end{definition}

\begin{remark}[转折点（turning points）]\label{sec:pf-th-glob-rem-turning}
路径可能相对 $\lambda$ 折叠：在转折点处 $D_xH$ 奇异、
$d\lambda/ds=0$，"固定 $\lambda$ 增量+逐点重解"的朴素策略在该处
失败（试探点无收敛解）。弧长参数化使增广 Jacobian 在普通转折点
仍非奇异，路径可平滑地"越过"折叠；鞍节点恰恰是潮流嵌入路径最
典型的转折点。按弧长连续编号的路径上，奇点指标（$D_xH$ 负
特征值计数）在简单转折点两侧差一，这也是 CPF 电压稳定分析的
理论入口（其工具化属本手册 CPF 章的主题）。
\end{remark}

\begin{remark}[步长自适应]\label{sec:pf-th-glob-rem-homstep}
实现采用的简化策略——成功则步长加倍（封顶）、失败则折半（低于
下界判失败）——是预测-校正框架去掉切向预测后的退化形式：以
上一收敛解直接作为下一 $\lambda$ 层的 Newton 初值（纯校正）。
在每层 Newton 收敛且层间距离小时等价于零阶预测。
\end{remark}

\begin{gapnote}
弧长参数化、加边 Newton 校正与转折点穿越在当前同伦实现中均未实现：
\srcpath{src/power\_flow/homotopy\_continuation.cpp:HomotopyContinuationSolver::solve}
按 $\lambda$ 单调递增逐层调用完整混合潮流求解
（\srcpath{detail::scale\_system\_by\_lambda} 缩放注入），
失败仅折半步长，路径折叠会被判为同伦失败而非穿越。
\end{gapnote}

\subsection{Newton--Krylov 与非精确 Newton 理论}\label{sec:pf-th-glob-nk}

大规模系统的 Newton 步 $J_k s_k=-F_k$ 不必精确求解；
Newton--Krylov 方法以 Krylov 子空间迭代近似求解（JFNK 中矩阵-向量
积 $Jv$ 可用差分 $[F(x+\epsilon v)-F(x)]/\epsilon$ 免 Jacobian 装配地
近似），其收敛理论由非精确 Newton 框架给出。

\begin{definition}[非精确 Newton 条件]\label{sec:pf-th-glob-def-inexact}
第 $k$ 步接受满足
\begin{equation}
\lVert F_k+J_k s_k\rVert\le\eta_k\lVert F_k\rVert,
\qquad \eta_k\in[0,1)
\end{equation}
的试探步 $s_k$，$\eta_k$ 称为强制项（forcing term）。
$\eta_k=0$ 退化为精确 Newton。
\end{definition}

\begin{theorem}[非精确 Newton 局部收敛（Dembo--Eisenstat--Steihaug）]\label{sec:pf-th-glob-thm-des}
设 $F$ 在解 $x^*$ 邻域连续可微，$J(x^*)$ 非奇异，$x_0$ 充分接近
$x^*$。则：
\begin{enumerate}
  \item 若 $\eta_k\le\bar\eta<1$，迭代局部 q-线性收敛；
  \item $x_k\to x^*$ q-超线性当且仅当
        $\lVert F_k+J_ks_k\rVert=o(\lVert F_k\rVert)$
       （特别地 $\eta_k\to0$ 是充分条件）；
  \item 若 $J$ Lipschitz 且
        $\lVert F_k+J_ks_k\rVert=O(\lVert F_k\rVert^2)$
        （如 $\eta_k=O(\lVert F_k\rVert)$），收敛为 q-二次。
\end{enumerate}
（Dembo, Eisenstat \& Steihaug~\cite{pfglob-des}；选择策略见
Eisenstat \& Walker~\cite{pfglob-ew}。）
\end{theorem}

\begin{proof}[证明思路]
记 $r_k=F_k+J_ks_k$ 为线性化残差，则
$x_{k+1}-x^*=x_k-x^*-J_k^{-1}(F_k-r_k)$。
对 $F(x^*)=0$ 作 Taylor 展开
$F_k=J_k(x_k-x^*)+O(\lVert x_k-x^*\rVert^2)$（$J$ Lipschitz），
得
\[
x_{k+1}-x^*
=J_k^{-1}r_k+O(\lVert x_k-x^*\rVert^2),
\]
故
$\lVert x_{k+1}-x^*\rVert
\le C\bigl(\eta_k\lVert F_k\rVert+\lVert x_k-x^*\rVert^2\bigr)$。
由 $\lVert F_k\rVert=O(\lVert x_k-x^*\rVert)$（$J(x^*)$ 非奇异
给出双边等价），三条结论依次成立。常数 $C$ 依赖
$\lVert J(x^*)^{-1}\rVert$ 与局部 Lipschitz 常数。
\end{proof}

\begin{remark}[强制项的自适应选择]\label{sec:pf-th-glob-rem-forcing}
远离解时取小 $\eta$ 是浪费（线性化模型本身误差大），近解时必须
$\eta\to0$ 以保超线性。Eisenstat--Walker~\cite{pfglob-ew} 给出两类
标准调度：$\eta_k=\lVert\phi_k-\phi_{k-1}^{\mathrm{pred}}\rVert/
\phi_{k-1}$ 型（反映模型保真度）与
$\eta_k=\gamma(\lVert F_k\rVert/\lVert F_{k-1}\rVert)^\alpha$ 型
（反映收敛速率），并配安全回退。定理~\ref{sec:pf-th-glob-thm-des}
只约束渐近行为，调度影响的是实际迭代数与内迭代总成本。
\end{remark}

\begin{proposition}[GMRES 的最小残差性质]\label{sec:pf-th-glob-prop-gmres}
GMRES~\cite{pfglob-saad-schultz} 第 $m$ 步在仿射空间
$x_0+\mathcal K_m(A,r_0)$ 中取使 $\lVert b-Ax\rVert_2$ 最小的
$x_m$。精确算术下 $m=n$ 必终止于真解；若 Arnoldi 过程在第 $j$ 步
满足 $h_{j+1,j}=0$（良性中断），则当前子空间已含精确解。
\end{proposition}

\begin{proof}
Arnoldi 关系 $AV_m=V_{m+1}\bar H_m$（$\bar H_m$ 为
$(m{+}1)\times m$ 上 Hessenberg 矩阵）与 $V_m$ 列正交性给出：对
$x=x_0+V_my$，
\[
\lVert b-Ax\rVert=\lVert r_0-AV_my\rVert
=\bigl\lVert\beta e_1-\bar H_m y\bigr\rVert,
\]
右端是小规模最小二乘问题，以 Givens 旋转在线更新，其解给出最小
残差。$m=n$ 时 $\mathcal K_n=\mathbb R^n$，最小残差为 0；
$h_{j+1,j}=0$ 时 $\mathcal K_j$ 是 $A$-不变子空间且含 $r_0$，
真解的残差校正落在其中。
\end{proof}

\begin{remark}[重启与预条件]\label{sec:pf-th-glob-rem-gmresrestart}
存储与正交化成本随 $m$ 线性增长，实用中取 GMRES($m$)：
每 $m$ 步以当前解重启。重启破坏有限终止性，可能停滞，故预条件
几乎总是必要：左预条件以 $P^{-1}A$ 替代 $A$ 构造 Krylov 子空间，
$P$ 应逼近 $A$ 且易解。对本平台的混合 AC/DC 系统，按实现头文件
契约注释，统一 Newton 的 Jacobian 呈块上三角结构
（AC 块 $A$、DC 块 $D$、耦合块 $B$、左下为零），此时"Schur 块
预条件子"实为精确块分解 $P=J$：预处理算子为恒等，GMRES 理论上
一步收敛；其价值不在 Krylov 加速，而在于分别分解 $A$、$D$ 以
降低填充与分解成本（Kelley~\cite{pfglob-kelley} 第 3、6 章的一般
讨论）。
\end{remark}

\begin{gapnote}
定理~\ref{sec:pf-th-glob-thm-des} 意义下的自适应强制项在当前实现中
未接线：GMRES 内迭代的相对残差容差取固定配置值
\fld{gmres\_tol}（\srcpath{src/power\_flow/newton\_krylov.cpp:gmres\_solve}），
NK 步的接受另加宽松阈值与有限性检查
（\srcpath{src/power\_flow/newton\_krylov.cpp:newton\_krylov\_step}），
因此外层只享有结论 (i) 的线性收敛口径；免 Jacobian 的差分矩阵-向量
积（严格 JFNK）亦未采用，矩阵-向量积直接作用于已装配的 Jacobian。
\end{gapnote}

\subsection{数值算例与基准}\label{sec:pf-th-glob-numexp}

本节数字由 \srcpath{tools/pf\_doc\_benchmark.cpp}（模式 \texttt{robust}/\texttt{stress}）
在 MATPOWER 算例上经平台公共 API 实测（平启动，Apple~M4）。

\subsubsection{高级求解器在标准算例上的收敛}
表~\ref{tab:pf-th-glob-robust} 给出平启动下\emph{关闭}全局化的纯 Newton、同伦延拓、HELM
与 Newton--Krylov 的最终残差。四者在 $9\sim57$ 母线算例上均达到 $10^{-8}$ 容差；
\textbf{HELM 在 case118/case300 上止步于 $\sim10^{-7}$}（固定对角 Padé 阶的精度上限，未达
$10^{-8}$，表中标 $\dagger$），与注记~\ref{sec:pf-th-glob-rem-stahlops} 及实现层 HELM
边界一致。纯 Newton 在标准算例平启动下即收敛，说明全局化阶梯的价值在于\emph{更难}的
工况（见下）而非常规潮流。
\begin{table}[H]\centering\footnotesize
\caption{高级求解器在标准算例上的最终残差（平启动，$\mathrm{tol}=10^{-8}$）}
\label{tab:pf-th-glob-robust}
\begin{tabular}{@{}lrllll@{}}
\toprule
算例 & 母线数 & 纯 Newton & 同伦延拓 & HELM & Newton--Krylov\\
\midrule
case9   & 9   & $1.7\times10^{-14}$ & $3.4\times10^{-9}$  & $8.2\times10^{-10}$      & $1.7\times10^{-14}$\\
case14  & 14  & $1.3\times10^{-10}$ & $5.8\times10^{-11}$ & $1.6\times10^{-10}$      & $1.3\times10^{-10}$\\
case30  & 30  & $9.6\times10^{-10}$ & $1.4\times10^{-12}$ & $3.5\times10^{-11}$      & $9.6\times10^{-10}$\\
case57  & 57  & $3.4\times10^{-12}$ & $9.9\times10^{-9}$  & $1.6\times10^{-10}$      & $3.4\times10^{-12}$\\
case118 & 118 & $8.5\times10^{-14}$ & $8.5\times10^{-14}$ & $8.9\times10^{-8}\,\dagger$ & $8.5\times10^{-14}$\\
case300 & 300 & $4.3\times10^{-13}$ & $6.2\times10^{-13}$ & $2.5\times10^{-7}\,\dagger$ & $4.3\times10^{-13}$\\
\bottomrule
\end{tabular}
\end{table}

\subsubsection{负荷加载与可解性边界}
表~\ref{tab:pf-th-glob-stress} 把算例负荷按因子 $\gamma$ 整体放大（发电不变、slack 吸收），
记录纯 Newton、带鲁棒升级的 Newton 与同伦三者的收敛性。三者在每个 $\gamma$ 处\emph{一致}
收敛或一致失败（故合并为一列）：case9 的边界落在 $\gamma\in(2.2,2.5)$、case118 在
$(1.5,1.8)$。这印证：越过该边界的失败是\textbf{物理无解}（鼻点之外），而非求解器鲁棒性
不足——对应同伦理论中路径“发散至无穷”的失败模式（定理~\ref{sec:pf-th-glob-thm-sard}）
与可解性充分条件越界（定理~\ref{thm:pf-th-fund-solvability}）。全局化阶梯能扩大\emph{可
达}解集，却不能制造不存在的解。
\begin{table}[H]\centering\footnotesize
\caption{负荷加载因子 $\gamma$ 下的收敛性（纯 Newton、鲁棒 Newton、同伦三者一致）}
\label{tab:pf-th-glob-stress}
\begin{tabular}{@{}lccccccccc@{}}
\toprule
$\gamma$ & 1.0 & 1.2 & 1.5 & 1.8 & 2.0 & 2.2 & 2.5 & 2.8 & 3.0\\
\midrule
case9   & Y & Y & Y & Y & Y & Y & N & N & N\\
case118 & Y & Y & Y & N & N & N & N & N & N\\
\bottomrule
\end{tabular}
\end{table}

\subsection{与实现的对应}\label{sec:pf-th-glob-map}
\begin{center}\footnotesize
\begin{longtable}{@{}L{6.8cm}L{8.6cm}@{}}
\toprule
理论条目 & 实现状态\\
\midrule
\endfirsthead
\toprule
理论条目 & 实现状态\\
\midrule
\endhead
Armijo 充分下降 + 回溯 &
\implfull{include/hacdcpf/power\_flow/globalization/nonmonotone\_linesearch.hpp:NonmonotoneLineSearch::search}（$M{=}1$ 时即单调 Armijo；主循环挂接见 src/power\_flow/newton\_solver.cpp:NewtonSolver::solve）\\
Wolfe 曲率条件 & \implnone\\
Zoutendijk 全局收敛（定理~\ref{sec:pf-th-glob-thm-zoutendijk} 的渐近保证） &
\implpart{实现以有限回溯上限 + 停滞检测 + 五级失败回退链兜底（RobustNonlinearOptions 停滞与回退字段），无渐近收敛性证书}\\
信赖域子问题、柯西点与 dogleg（定义~\ref{sec:pf-th-glob-def-tr}、
引理~\ref{sec:pf-th-glob-lem-cauchy}） &
\implfull{include/hacdcpf/engine/kernel/globalization/globalization.hpp:dogleg\_step}（转发到 MIPSolvers；增益比半径调节同文件 trust\_region\_update，调用侧为 newton\_solver.cpp 的 TrustRegion 分支）\\
精确信赖域子问题（Moré--Sorensen）、Steihaug--Toint CG & \implnone\\
LM 正规方程 $(J^{\mathsf T}J+\lambda I)\Delta x=-J^{\mathsf T}F$
与增益比调 $\lambda$（定理~\ref{sec:pf-th-glob-thm-lmtr} 的参数化） &
\implpart{LM 恢复步内联实现于 src/power\_flow/newton\_solver.cpp:NewtonSolver::solve（线搜索耗尽后至多 4 次 $\lambda\times10$ 尝试）；头文件 LMController::update 是预留控制器，未被主循环消费}\\
PTC 与 SER 伪时间步（注记~\ref{sec:pf-th-glob-rem-ptc}） &
\implfull{include/hacdcpf/power\_flow/globalization/lm\_trust\_region.hpp:PtcSerController::update}（PTC 分支装配 $J+I/\delta t$，挂接 NewtonSolver::solve）\\
GLL 非单调窗口参考值（定义~\ref{sec:pf-th-glob-def-gll}） &
\implfull{include/hacdcpf/power\_flow/globalization/nonmonotone\_linesearch.hpp:NonmonotoneLineSearch}（push\_merit/reference\_merit，窗口默认 $M=5$）\\
FB 函数与广义 Jacobian 元素
（命题~\ref{sec:pf-th-glob-prop-fb}、例~\ref{sec:pf-th-glob-ex-fb}） &
\implfull{include/hacdcpf/power\_flow/ncp\_functions.hpp:fischer\_burmeister}（原点对称元见同文件 fischer\_burmeister\_jacobian）\\
中位数 NCP 统一 PV/PQ（命题~\ref{sec:pf-th-glob-prop-median}） &
\implfull{include/hacdcpf/power\_flow/ncp\_functions.hpp:pv\_pq\_ncp}（广义导数 pv\_pq\_ncp\_jacobian，并列处取相邻单侧斜率平均）\\
CHKS/Kanzow 光滑化与 $\mu$ 延拓（定义~\ref{sec:pf-th-glob-def-smoothing}） &
\implfull{include/hacdcpf/power\_flow/ncp\_functions.hpp:smooth\_fb}（pv\_pq\_smooth\_ncp；两阶段退火调度在 NewtonSolver::solve 内，参数 RobustNonlinearOptions::enable\_smooth\_ncp 系列）\\
半光滑 Newton 超线性/二次收敛（定理~\ref{sec:pf-th-glob-thm-qisun}） &
\implpart{求解器按广义 Newton 步运行并上报实测速率商；BD-正则性不做事先验证，病态所选块以后向误差证书拒绝并回退全稀疏 LU（见 vsc\_limit\_ncp 契约）}\\
VSC 限流的互补表述（同一 NCP 框架的设备级实例） &
\implfull{src/power\_flow/vsc\_limit\_ncp.cpp:evaluate\_vsc\_limit\_ncp}（median\_ncp；限流圆盘与优先级分支见 vsc\_limit\_command）\\
HELM 嵌入与系数递推（定义~\ref{sec:pf-th-glob-def-helm}、
命题~\ref{sec:pf-th-glob-prop-helmrec}） &
\implfull{src/power\_flow/helm\_solver.cpp:HelmSolver::solve}（胚矩阵 build\_germ\_matrix、$|V|^2$ 约束系数 magnitude\_squared\_coefficient）\\
Padé 解析延拓（定义~\ref{sec:pf-th-glob-def-pade}） &
\implfull{src/power\_flow/helm\_solver.cpp:pade\_at\_one}（对角 Padé 矩阵法在 $s=1$ 求值）\\
Stahl 定理层面的收敛性分析（定理~\ref{sec:pf-th-glob-thm-stahl}） &
\implpart{实现以相邻 Padé 逼近差为启发式收敛判据、以发散作不可行性实用判据；无容量/割集层面的理论核验}\\
同伦映射与注入缩放（定义~\ref{sec:pf-th-glob-def-homotopy}） &
\implfull{src/power\_flow/homotopy\_continuation.cpp:detail::scale\_system\_by\_lambda}（逐层驱动 HomotopyContinuationSolver::solve）\\
概率一正则性与预测-校正跟踪
（定理~\ref{sec:pf-th-glob-thm-sard}、定义~\ref{sec:pf-th-glob-def-pc}） &
\implpart{实现为零阶预测（上层解直接作初值）+ 逐层完整重解 + 步长成败折半/加倍（HomotopyState）；无正则性检验}\\
弧长参数化与转折点穿越（注记~\ref{sec:pf-th-glob-rem-turning}） & \implnone\\
非精确 Newton 条件与 DES 收敛理论
（定义~\ref{sec:pf-th-glob-def-inexact}、
定理~\ref{sec:pf-th-glob-thm-des}） &
\implpart{内层残差容差为固定 \fld{gmres\_tol} 而非自适应强制项（newton\_krylov.cpp:newton\_krylov\_step），外层仅享线性收敛口径}\\
GMRES($m$)：MGS Arnoldi + Givens 最小二乘
（命题~\ref{sec:pf-th-glob-prop-gmres}） &
\implfull{src/power\_flow/newton\_krylov.cpp:gmres\_solve}（重启、真残差复核、良性中断处理齐全）\\
Schur 块预条件（注记~\ref{sec:pf-th-glob-rem-gmresrestart}） &
\implfull{src/power\_flow/newton\_krylov.cpp:SchurBlockPreconditioner::build}（apply 实现两段块回代；块结构契约见 newton\_krylov.hpp 类注释）\\
Anderson 加速与 GMRES 等价（定义~\ref{def:pf-th-glob-anderson}、
定理~\ref{thm:pf-th-glob-anderson-gmres}） &
\implnone（不动点求解器用朴素 Picard/拟 Newton，无残差外推；加速手段为 Newton/GMRES）\\
高级求解器与可解性边界数值基准（表~\ref{tab:pf-th-glob-robust}、
\ref{tab:pf-th-glob-stress}） &
\implfull{tools/pf\_doc\_benchmark.cpp}（模式 robust/stress）\\
\bottomrule
\end{longtable}
\end{center}

\subsection{参考文献}
{\renewcommand{\section}[2]{}%
\begin{thebibliography}{99}\small
\bibitem{pfglob-nw} J.~Nocedal and S.~J. Wright, \emph{Numerical
Optimization}, 2nd ed., Springer, 2006.
\bibitem{pfglob-ds} J.~E. Dennis and R.~B. Schnabel, \emph{Numerical
Methods for Unconstrained Optimization and Nonlinear Equations}, SIAM,
1996.
\bibitem{pfglob-powell} M.~J.~D. Powell, A hybrid method for nonlinear
equations, in \emph{Numerical Methods for Nonlinear Algebraic Equations}
(P.~Rabinowitz, ed.), Gordon and Breach, 1970.
\bibitem{pfglob-levenberg} K.~Levenberg, A method for the solution of
certain non-linear problems in least squares, \emph{Quarterly of Applied
Mathematics}, 2, 1944.
\bibitem{pfglob-marquardt} D.~W. Marquardt, An algorithm for least-squares
estimation of nonlinear parameters, \emph{SIAM Journal on Applied
Mathematics}, 11(2), 1963.
\bibitem{pfglob-more} J.~J. Mor\'e, The Levenberg--Marquardt algorithm:
implementation and theory, in \emph{Numerical Analysis}, Lecture Notes in
Mathematics 630, Springer, 1978.
\bibitem{pfglob-gll} L.~Grippo, F.~Lampariello and S.~Lucidi, A nonmonotone
line search technique for Newton's method, \emph{SIAM Journal on Numerical
Analysis}, 23(4), 1986.
\bibitem{pfglob-fp} F.~Facchinei and J.-S. Pang, \emph{Finite-Dimensional
Variational Inequalities and Complementarity Problems}, Vol.~I, Springer,
2003.
\bibitem{pfglob-fischer} A.~Fischer, A special Newton-type optimization
method, \emph{Optimization}, 24, 1992.
\bibitem{pfglob-qisun} L.~Qi and J.~Sun, A nonsmooth version of Newton's
method, \emph{Mathematical Programming}, 58, 1993.
\bibitem{pfglob-kanzow} C.~Kanzow, Some noninterior continuation methods
for linear complementarity problems, \emph{SIAM Journal on Matrix Analysis
and Applications}, 17(4), 1996.
\bibitem{pfglob-chenharker} B.~Chen and P.~T. Harker, A non-interior-point
continuation method for linear complementarity problems, \emph{SIAM Journal
on Matrix Analysis and Applications}, 14(4), 1993.
\bibitem{pfglob-hik} M.~Hinterm\"uller, K.~Ito and K.~Kunisch, The
primal-dual active set strategy as a semismooth Newton method,
\emph{SIAM Journal on Optimization}, 13(3), 2002.
\bibitem{pfglob-trias} A.~Trias, The holomorphic embedding load flow
method, in \emph{Proceedings of the IEEE PES General Meeting}, San Diego,
2012.
\bibitem{pfglob-stahl} H.~Stahl, On the convergence of generalized Pad\'e
approximants, \emph{Constructive Approximation}, 5, 1989.
\bibitem{pfglob-baker} G.~A. Baker and P.~Graves-Morris, \emph{Pad\'e
Approximants}, 2nd ed., Cambridge University Press, 1996.
\bibitem{pfglob-allgower} E.~L. Allgower and K.~Georg, \emph{Numerical
Continuation Methods: An Introduction}, Springer, 1990.
\bibitem{pfglob-des} R.~S. Dembo, S.~C. Eisenstat and T.~Steihaug, Inexact
Newton methods, \emph{SIAM Journal on Numerical Analysis}, 19(2), 1982.
\bibitem{pfglob-ew} S.~C. Eisenstat and H.~F. Walker, Choosing the forcing
terms in an inexact Newton method, \emph{SIAM Journal on Scientific
Computing}, 17(1), 1996.
\bibitem{pfglob-saad-schultz} Y.~Saad and M.~H. Schultz, GMRES: a
generalized minimal residual algorithm for solving nonsymmetric linear
systems, \emph{SIAM Journal on Scientific and Statistical Computing},
7(3), 1986.
\bibitem{pfglob-kelley} C.~T. Kelley, \emph{Iterative Methods for Linear
and Nonlinear Equations}, SIAM, 1995.
\bibitem{pfglob-kelley-keyes} C.~T. Kelley and D.~E. Keyes, Convergence
analysis of pseudo-transient continuation, \emph{SIAM Journal on Numerical
Analysis}, 35(2), 1998.
\bibitem{pfglob-anderson} D.~G. Anderson, Iterative procedures for nonlinear
integral equations, \emph{Journal of the ACM}, 12(4), 1965.
\bibitem{pfglob-walkerni} H.~F. Walker and P.~Ni, Anderson acceleration for
fixed-point iterations, \emph{SIAM Journal on Numerical Analysis}, 49(4), 2011.
\bibitem{pfglob-tothkelley} A.~Toth and C.~T. Kelley, Convergence analysis for
Anderson acceleration, \emph{SIAM Journal on Numerical Analysis}, 53(2), 2015.
\end{thebibliography}}

% theory_voltage_stability_converters.tex —— 潮流计算 通用理论：电压稳定、连续潮流与换流器稳态模型
% 本文件为理论层章节，遵循 docs/modules/THEORY_WRITING_GUIDE.md。

\section{电压稳定、连续潮流与换流器稳态理论}\label{sec:pf-th-vs}

\begin{theorynote}
本章为通用数学理论，按电压稳定与高压直流输电领域的标准模型与文献体系撰写，
覆盖四组主题：(i) 电压稳定的分岔理论基础——鞍结分岔（SNB）与鼻点的分岔本质、
负荷裕度的数学定义、$L$ 指标与模态分析/参与因子体系、戴维南等值下的最大功率传输与阻抗匹配判据；(ii) 连续潮流（CPF）的
参数化理论——自然/弧长/伪弧长参数化、增广系统、切向量预测与校正、步长控制的
理论依据；(iii) 换流器稳态理论——VSC 注入模型与控制模式自由度分析、容量圆与
调制比约束的几何意义、损耗模型推导，LCC 准稳态方程，DC-DC 变换器平均模型；
(iv) AC/DC 接口方程的统一表述与 DC 岛可解性条件；(v) 以平台弧长连续潮流实测数据与两母线解析鼻点对照的数值基准。
本章公式不要求逐式对应代码；与平台实现（\srcpath{src/power\_flow/}）的对应关系
见本章末"与实现的对应"表，超出实现的内容以"未实现扩展说明"框就地标记。
\end{theorynote}

\subsection{记号与约定}\label{sec:pf-th-vs-notation}
\begin{itemize}
  \item $x\in\mathbb{R}^{n}$：潮流状态向量（AC 相角 $\theta$、PQ 母线电压幅值
        $V_m$、DC 母线电压 $v_{dc}$ 的堆叠）；$\lambda\in\mathbb{R}$：延拓参数
        （负荷水平因子）；$F(x,\lambda)=0$：参数化潮流方程，
        $J=F_x=\partial F/\partial x$ 为潮流雅可比；
  \item $y=(x,\lambda)\in\mathbb{R}^{n+1}$：增广状态；$s$：弧长参数；$t$：
        归一化切向量；上标 $0$ 表示基态（$\lambda=0$）取值；
  \item $v,w$：奇异雅可比 $J$ 的右/左零向量（$Jv=0$，$w^{\top}J=0$）；
  \item VSC 换流器 $k$ 连接 AC 母线 $i$ 与 DC 母线 $j$：$P_{ac},Q_{ac}$ 为 AC
        端口注入（注入网络为正），$P_{dc}$ 为 DC 端口注入，$P_{loss}\ge0$ 为损耗，
        $m$ 为调制比；$S_{base}$ 为系统基准功率；
  \item LCC 换流站：$E_0$ 为换流变漏抗后的理想空载阀侧线电压，$E_t$ 为带载阀侧
        母线线电压，$X_c$ 为每桥每相换相电抗，$n_b$ 为每极串联 6 脉动桥数，
        $\alpha$ 为触发角，$\gamma$ 为关断角，$\mu$ 为换相重叠角，$I_d$、$U_d$
        为直流电流与电压；
  \item DC-DC 变换器 $m$：$V_{in},V_{out}$ 为端口电压，$D$ 为占空比，$\eta$ 为
        效率；$G_{dc}=[g_{ij}]$ 为 DC 网络电导 Laplacian（$g=1/r_{pu}$）。
  与手册其余部分冲突时以本章声明为准。
\end{itemize}

\subsection{电压稳定理论基础：鞍结分岔与负荷裕度}\label{sec:pf-th-vs-snb}

\begin{definition}[参数化潮流方程与解支]
设系统注入按方向 $d=(\Delta P,\Delta Q)$ 仿射增长，
\begin{equation}
F(x,\lambda)=F_0(x)+\lambda\,d(x)=0,\qquad F:\mathbb{R}^{n}\times\mathbb{R}\to\mathbb{R}^{n},
\end{equation}
$\lambda=0$ 对应基态运行点 $x_0$（$F(x_0,0)=0$）。满足 $F(x,\lambda)=0$ 的
点集在 $(x,\lambda)$ 空间中形成\emph{解支}；其在指定母线上的投影
$(\lambda,V_i(\lambda))$ 即 P--V 鼻曲线。
\end{definition}

\begin{definition}[正则点与鞍结分岔（折点）]
解支上的点 $(x^*,\lambda^*)$ 称为\emph{正则点}，若 $J=F_x(x^*,\lambda^*)$ 非奇异；
由隐函数定理，正则点附近解支可局部写成光滑函数 $x(\lambda)$。若 $J$ 奇异且
秩亏一（$\mathrm{rank}\,J=n-1$），并满足下文定理~\ref{thm:pf-th-vs-snb} 的横截
条件，则该点为\emph{鞍结分岔点}（saddle-node bifurcation，亦称折点/鼻点）：
两条解支在此汇合并消失，$\lambda$ 在该点取局部极值 $\lambda^*$。
\end{definition}

\begin{lemma}[鼻点处雅可比奇异]\label{lem:pf-th-vs-singular}
设解支可光滑地以弧长 $s$ 参数化：$(x(s),\lambda(s))$，$\|(\dot x,\dot\lambda)\|=1$。
若 $\lambda(s)$ 在 $s=s^*$ 处取局部极值（鼻点），则 $F_x(x(s^*),\lambda(s^*))$
奇异。
\end{lemma}
\begin{proof}
对 $F(x(s),\lambda(s))\equiv0$ 关于 $s$ 求导得
\begin{equation}
F_x\,\dot x+F_\lambda\,\dot\lambda=0.
\end{equation}
极值处 $\dot\lambda(s^*)=0$，故 $F_x\,\dot x(s^*)=0$。由弧长参数化的归一化，
$\|(\dot x,\dot\lambda)\|=1$ 且 $\dot\lambda=0$ 蕴含 $\dot x(s^*)\neq0$，
即 $F_x$ 有非平凡零空间，奇异。
\end{proof}

引理~\ref{lem:pf-th-vs-singular} 给出"鼻点 $\Rightarrow$ 奇异"的方向；其逆命题
需要横截条件排除退化情形，这正是下述核心定理的内容。

\begin{theorem}[二次折点的刻画]\label{thm:pf-th-vs-snb}
设 $F$ 光滑，$F(x^*,\lambda^*)=0$，$J=F_x(x^*,\lambda^*)$ 满足
$\mathrm{rank}\,J=n-1$，右/左零空间分别由 $v,w$ 张成。若横截条件
\begin{equation}
w^{\top}F_\lambda\neq0,\qquad a:=w^{\top}F_{xx}[v,v]\neq0
\end{equation}
成立，则解集在 $(x^*,\lambda^*)$ 附近是一条光滑曲线，可参数化为
$(x(\xi),\lambda(\xi))$（$\xi$ 为沿 $v$ 方向的局部坐标），且
\begin{equation}
\lambda(\xi)=\lambda^*-\frac{a}{2\,w^{\top}F_\lambda}\,\xi^{2}+O(\xi^{3}),\qquad
\dot\lambda(0)=0,\quad \ddot\lambda(0)\neq0,
\end{equation}
即 $\lambda$ 在该点经历二次折点（鞍结分岔），解支在 $\lambda<\lambda^*$ 侧有两支
（上支/下支），$\lambda>\lambda^*$ 侧无解。
\end{theorem}
\begin{proof}
作分解 $x=x^*+\xi v+z$，其中 $z\in Z$，$Z$ 为 $\mathrm{span}\{v\}$ 的任一补子空间
（如取 $Z=\{v\}^{\perp}$，维数 $n-1$）。考察映射
\begin{equation}
H(z,\lambda;\xi)=F(x^*+\xi v+z,\ \lambda)=0,
\end{equation}
其对 $(z,\lambda)$ 的导数为 $[J|_{Z}\ \ F_\lambda]$：这是 $n\times n$ 矩阵，
其值域包含 $\mathrm{range}\,J$（维数 $n-1$）与 $F_\lambda$；由
$w^{\top}F_\lambda\neq0$ 知 $F_\lambda\notin\mathrm{range}\,J$（否则存在 $u$
使 $Ju=F_\lambda$，则 $w^{\top}F_\lambda=w^{\top}Ju=0$，矛盾），故该矩阵秩为
$n$，非奇异。由隐函数定理，存在光滑的 $z(\xi)$、$\lambda(\xi)$ 使
$H(z(\xi),\lambda(\xi);\xi)\equiv0$，$z(0)=0$，$\lambda(0)=\lambda^*$。

对恒等式 $F(x^*+\xi v+z(\xi),\lambda(\xi))\equiv0$ 关于 $\xi$ 求导：
\begin{equation}
J\bigl(v+z'(0)\bigr)+F_\lambda\,\lambda'(0)=0.
\end{equation}
左乘 $w^{\top}$：$w^{\top}J=0$ 与 $w^{\top}F_\lambda\neq0$ 给出
$\lambda'(0)=0$；代回得 $Jz'(0)=-Jv=0$，而 $z'(0)\in Z$ 与 $Z\cap\ker J=\{0\}$
给出 $z'(0)=0$。再求二阶导并代入 $z'(0)=0$、$\lambda'(0)=0$：
\begin{equation}
J\,z''(0)+F_{xx}[v,v]+F_\lambda\,\lambda''(0)=0,
\end{equation}
左乘 $w^{\top}$ 得
\begin{equation}
\lambda''(0)=-\frac{w^{\top}F_{xx}[v,v]}{w^{\top}F_\lambda}=-\frac{a}{w^{\top}F_\lambda}\neq0.
\end{equation}
于是 $\lambda(\xi)=\lambda^*+\tfrac12\lambda''(0)\xi^2+O(\xi^3)$，为二次折点。
\end{proof}

\begin{remark}[鼻曲线上下支的稳定性解读]
定理~\ref{thm:pf-th-vs-snb} 中 $\lambda''(0)$ 的符号决定折点开口方向；常规
负荷增长下 $\lambda^*>0$ 为局部最大值，上支（高电压支）为小信号稳定运行支，
下支（低电压支）不稳定（详见 Kundur~\cite{pfvs:kundur} 与
Van Cutsem--Vournas~\cite{pfvs:vancutsem} 的小扰动分析）。工程负荷裕度只关心
上支终点，CPF 翻越鼻点进入下支是\emph{数值追踪}行为，不对应可持续运行方式。
\end{remark}

\begin{example}[两母线系统的解析鼻点]\label{ex:pf-th-vs-twobus}
送端为无穷大母线 $E\angle0$，经纯电抗 $jX$ 向受端母线（电压 $V\angle\delta$）
输送恒功率负荷 $(P,Q)=(\lambda P_0,\lambda Q_0)$。潮流方程
\begin{equation}
P=\frac{EV\sin\delta}{X},\qquad Q=\frac{EV\cos\delta-V^{2}}{X}.
\end{equation}
消去 $\delta$（两式平方相加），令 $u=V^{2}$：
\begin{equation}
u^{2}+\bigl(2QX-E^{2}\bigr)u+X^{2}\bigl(P^{2}+Q^{2}\bigr)=0,
\end{equation}
\begin{equation}
u=\frac{E^{2}-2QX\pm\sqrt{\bigl(E^{2}-2QX\bigr)^{2}-4X^{2}\bigl(P^{2}+Q^{2}\bigr)}}{2},
\end{equation}
$+$ 号对应上支。实解存在的充要条件是判别式非负；判别式 $=0$ 即鼻点。特取
$Q_0=0$：鼻点条件 $E^{4}=4X^{2}\lambda^{2}P_0^{2}$，得
\begin{equation}
\lambda^{*}=\frac{E^{2}}{2XP_0},\qquad
V_{\mathrm{nose}}=\frac{E}{\sqrt2},\qquad
P_{\max}=\lambda^{*}P_0=\frac{E^{2}}{2X}.
\end{equation}
此时鼻点电压恰为 $E/\sqrt2\approx0.707$~pu——输电极限不仅取决于线路参数，
还直接压低临界电压。该解析结果是 CPF 回归算例（两母线电压稳定算例）的理论
期望值。
\end{example}

\begin{definition}[负荷裕度]\label{def:pf-th-vs-margin}
沿增长方向 $d$，系统的\emph{负荷裕度}为基态到鼻点的传输能力增量：
\begin{equation}
\lambda^{*}=\sup\bigl\{\lambda\ge0:\ \exists\,x,\ F(x,\lambda)=0,\ (x,\lambda)
\text{ 与基态同支}\bigr\},
\end{equation}
\begin{equation}
P_{\mathrm{margin}}=P_{\mathrm{total}}(\lambda^{*})-P_{\mathrm{total}}(0)
\quad(\mathrm{MW}),\qquad
P_{\mathrm{total}}(\lambda)=\sum_i\bigl(P_{L,i}^{0}+\lambda\,\Delta P_{L,i}\bigr).
\end{equation}
$\lambda^{*}$ 由分岔点给出（连续负荷模型下即 SNB 点）；当计及发电机无功上限、
离散分接头等不等式时，崩溃点可由限值诱导分岔（LIB）提前出现，负荷裕度取各
机制的最小值。
\end{definition}

\begin{remark}[裕度对方向的依赖]
负荷裕度是\emph{方向依赖}量：同一系统沿不同 $d$ 的 $\lambda^{*}$ 一般不同。
工程报告必须同时声明增长方向（比例增长、单母线增长、区域增长等），否则裕度
数字没有意义。
\end{remark}

\subsubsection*{$L$ 指标（Kessel--Glavitsch）}
把母线划分为发电机母线集 $G$ 与负荷母线集 $L$，节点方程分块：
\begin{equation}
\begin{bmatrix}I_L\\ I_G\end{bmatrix}
=
\begin{bmatrix}Y_{LL} & Y_{LG}\\ Y_{GL} & Y_{GG}\end{bmatrix}
\begin{bmatrix}V_L\\ V_G\end{bmatrix},
\qquad
V_L^{0}:=-Y_{LL}^{-1}Y_{LG}\,V_G,
\end{equation}
$V_L^{0}$ 为全部负荷电流置零时负荷母线的等效开路电压。

\begin{definition}[$L$ 指标~\cite{pfvs:kessel}]
对每个负荷母线 $j\in L$，
\begin{equation}
L_j=\Bigl|\,1-\frac{V_j^{0}}{V_j}\,\Bigr|,\qquad
L=\max_{j\in L}L_j\in[0,1].
\end{equation}
\end{definition}
$L_j=0$ 对应空载（$V_j=V_j^{0}$），$L\to1$ 对应 $V_j$ 远低于等效源电压、
逼近电压崩溃；$L$ 是逐母线弱节点排序指标。在恒功率负荷、发电机母线电压
幅值恒定的假设下，可以证明鼻点处 $L=1$；一般情形下 $L$ 是单调逼近 $1$ 的
启发式裕度度量。

\begin{gapnote}
$L$ 指标、模态分析/参与因子与 Q--V 曲线在当前平台 CPF 中\textbf{均未实现}：
指标族仅含负荷裕度类标量（$\lambda_{\max}$、$P_{\max}$、$P_{\mathrm{margin}}$
与监控母线电压，实现层连续潮流章"结果结构与电压稳定指标"节如实声明）。
上述两小节为完整理论体系所需而保留，属未实现扩展。
\end{gapnote}

\subsubsection*{模态分析与参与因子（Gao--Morison--Kundur）}
把潮流雅可比按 $(\theta,V)$ 分块，并令 $\Delta P=0$（有功已平衡的扰动口径）：
\begin{equation}
\begin{bmatrix}\Delta P\\ \Delta Q\end{bmatrix}
=
\begin{bmatrix}J_{P\theta} & J_{PV}\\ J_{Q\theta} & J_{QV}\end{bmatrix}
\begin{bmatrix}\Delta\theta\\ \Delta V\end{bmatrix}.
\end{equation}

\begin{proposition}[降阶 Q--V 雅可比~\cite{pfvs:gao}]\label{prop:pf-th-vs-modal}
设 $J_{P\theta}$ 非奇异，则 $\Delta P=0$ 时
\begin{equation}
\Delta Q=J_R\,\Delta V,\qquad
J_R=J_{QV}-J_{Q\theta}J_{P\theta}^{-1}J_{PV},
\end{equation}
$J_R$ 即潮流雅可比关于 $\theta$ 的 Schur 补。$J$ 奇异蕴含 $J_R$ 奇异
（$\det J=\det J_{P\theta}\cdot\det J_R$），故电压崩溃点可用 $J_R$ 的最小模
特征值刻画。
\end{proposition}
\begin{proof}
由 $\Delta P=0$ 得 $\Delta\theta=-J_{P\theta}^{-1}J_{PV}\Delta V$，代入
$\Delta Q$ 行即得 $J_R$。行列式恒等式为分块 Gauss 消元的标准结论。
\end{proof}

设 $J_R=\Xi\Lambda\Xi^{-1}$，$\Xi$ 的列为右特征向量 $\xi_k$，左特征向量
$\eta_k=\Xi^{-1}$ 的行。模态分解 $\Delta V=\sum_k c_k\xi_k$ 给出
$\Delta Q=\sum_k c_k\lambda_k\xi_k$：特征值 $\lambda_k$ 度量第 $k$ 个
电压模态的 Q--V 灵敏度，$\lambda_k$ 小（$>0$）表示弱模态，$\lambda_k<0$ 表示
电压不稳定（无功支撑增加反而压低电压）。母线 $i$ 对模态 $k$ 的\emph{参与因子}
为 $p_{ik}=\xi_{ik}\eta_{ki}$，满足 $\sum_i p_{ik}=1$，用于弱节点排序。

\subsection{戴维南等值与最大功率传输判据}\label{sec:pf-th-vs-thevenin}

例~\ref{ex:pf-th-vs-twobus} 的两母线鼻点可推广到\emph{任意}负荷母线：把其余网络对该
母线作戴维南等值 $(E_{th},Z_{th})$，则该母线的静态电压稳定极限由\emph{阻抗匹配}刻画。

\begin{theorem}[最大功率传输与临界电压]\label{thm:pf-th-vs-thevenin}
负荷母线经戴维南等值 $E_{th}\angle0$、$Z_{th}\angle\theta_{th}$ 馈电，负荷阻抗
$Z_L\angle\varphi$（$\varphi$ 为功率因数角，$\tan\varphi=Q/P$）。沿固定 $\varphi$ 增大
负荷（即减小 $Z_L$），则传输有功在\emph{阻抗模匹配} $Z_L=Z_{th}$ 处取最大，且该临界点
\begin{equation}
V_{\mathrm{crit}}=\frac{E_{th}}{\sqrt{2\bigl(1+\cos(\theta_{th}-\varphi)\bigr)}},\qquad
P_{\max}=\frac{E_{th}^2}{2Z_{th}}\cdot\frac{\cos\varphi}{1+\cos(\theta_{th}-\varphi)} .
\label{eq:pf-th-vs-thevenin}
\end{equation}
\end{theorem}
\begin{proof}
回路电流 $I=E_{th}/|Z_{th}+Z_L|$，其中
$|Z_{th}+Z_L|^2=Z_{th}^2+Z_L^2+2Z_{th}Z_L\cos(\theta_{th}-\varphi)$。负荷有功
\begin{equation*}
P=I^2Z_L\cos\varphi
=\frac{E_{th}^2\,Z_L\cos\varphi}{Z_{th}^2+Z_L^2+2Z_{th}Z_L\cos(\theta_{th}-\varphi)} .
\end{equation*}
对 $Z_L$ 求导，分子为 $E_{th}^2\cos\varphi\,(Z_{th}^2-Z_L^2)$，令其为零得
$Z_L=Z_{th}$（阻抗匹配）。代回：分母
$=2Z_{th}^2\bigl(1+\cos(\theta_{th}-\varphi)\bigr)$，即得 $P_{\max}$；负荷端电压
$V=IZ_L=E_{th}Z_{th}/|Z_{th}+Z_L|$，匹配点处
$|Z_{th}+Z_L|=Z_{th}\sqrt{2(1+\cos(\theta_{th}-\varphi))}$，即 $V_{\mathrm{crit}}$。
\end{proof}

\begin{remark}[与鼻点、$L$ 指标、在线监视的联系]\label{rem:pf-th-vs-thevenin}
纯电抗、单位功率因数（$\theta_{th}=\pi/2$、$\varphi=0$）时
$\cos(\theta_{th}-\varphi)=0$，式~\eqref{eq:pf-th-vs-thevenin} 化为
$V_{\mathrm{crit}}=E_{th}/\sqrt2$、$P_{\max}=E_{th}^2/(2Z_{th})$——正是
例~\ref{ex:pf-th-vs-twobus} 的两母线鼻点。判据 $Z_L=Z_{th}$ 即“负荷阻抗逼近戴维南
阻抗”，是基于本地相量在线辨识 $Z_{th}$ 的阻抗裕度指标~\cite{pfvs:vu} 与 $L$ 指标
（第~\ref{sec:pf-th-vs-snb}~节）的共同物理内核。其局限：戴维南等值随其余网络运行点
缓变，需在线更新；多负荷同时增长时各母线 $Z_{th}$ 耦合变化，单母线判据仅局部有效。
\end{remark}

\begin{gapnote}
基于戴维南等值的阻抗匹配裕度指标在平台 CPF 中\textbf{未实现}（指标族仅含 $\lambda$-裕度类，
见本章末对应表与实现层 CPF 章）；本判据作为电压稳定的解析内核与 CPF 数值追踪互补。
\end{gapnote}

\subsection{连续潮流（CPF）理论}\label{sec:pf-th-vs-cpf}

\subsubsection{问题提法与参数化}
CPF 的目标是数值追踪解支 $F(x,\lambda)=0$ 经过鼻点的全过程。核心困难由
引理~\ref{lem:pf-th-vs-singular} 给出：鼻点处 $F_x$ 奇异，\emph{自然参数化}
（固定 $\lambda$ 递增、以 $x$ 为未知量用 Newton 求解）在鼻点附近必然恶化——
Newton 迭代面对病态雅可比，且 $\lambda>\lambda^*$ 时无解可收。

\begin{definition}[三种参数化]\label{def:pf-th-vs-param}
设 $y=(x,\lambda)\in\mathbb{R}^{n+1}$，$G(y)=F(x,\lambda)$。
\begin{itemize}
  \item \emph{自然参数化}：以 $\lambda$ 为参数，逐步固定 $\lambda$ 求解 $x$；
        在折点失效；
  \item \emph{弧长参数化}：以解支弧长 $s$ 为参数，
        $\mathrm{d}s^{2}=\|\mathrm{d}x\|^{2}+\mathrm{d}\lambda^{2}$，
        补充方程 $\|y-y_{\mathrm{curr}}\|^{2}=s^{2}$；
  \item \emph{伪弧长（Keller）参数化}~\cite{pfvs:keller}：以当前点切向量 $t$
        张成的超平面局部参数化，补充方程
        \begin{equation}
        N(y)=t^{\top}\bigl(y-y_{\mathrm{pred}}\bigr)=0,
        \end{equation}
        其中 $y_{\mathrm{pred}}=y_{\mathrm{curr}}+s\,t$ 为预测点。伪弧长是弧长
        约束的一阶（线性）版本，保证增广系统线性代数结构简单。
    \end{itemize}
\end{definition}

\subsubsection{切向量预测}
\begin{proposition}[切向量的确定]\label{prop:pf-th-vs-tangent}
设 $y$ 为解支上的正则点（$\mathrm{rank}[F_x\ F_\lambda]=n$），则解支在该点的
单位切向量 $t$ 由
\begin{equation}
[F_x\ \ F_\lambda]\,t=0,\qquad \|t\|=1,\qquad t^{\top}t_{\mathrm{prev}}>0
\end{equation}
唯一确定（符号条件保证沿支定向一致）。实践中以相邻两解点的归一化割线
$t=(y_{\mathrm{curr}}-y_{\mathrm{prev}})/\|y_{\mathrm{curr}}-y_{\mathrm{prev}}\|$
近似切向量（割线预测）。
\end{proposition}
\begin{proof}
$[F_x\ F_\lambda]$ 为 $n\times(n+1)$ 满秩矩阵，零空间恰为一维，故单位切向量
仅差符号；定向条件取与 $t_{\mathrm{prev}}$ 同号者。
\end{proof}

\begin{remark}[折点处切向量的形态]
在二次折点，定理~\ref{thm:pf-th-vs-snb} 的证明给出 $\dot\lambda=0$、
$\dot x\parallel v$：切向量退化为 $t\propto(v,0)$，即"$\lambda$ 不动、$x$ 沿
$J$ 的零向量方向走"。这正是鼻曲线在鼻点处切线竖直（$\mathrm{d}\lambda/
\mathrm{d}V=0$）的向量表述，也是自然参数化失效的几何原因。
\end{remark}

\subsubsection{增广校正与增广雅可比的非奇异性}
伪弧长校正在预测点处求解增广系统
\begin{equation}
\Phi(y)=
\begin{bmatrix}
F(x,\lambda)\\[2pt] t^{\top}(y-y_{\mathrm{pred}})
\end{bmatrix}=0,
\qquad
\Phi_y=\begin{bmatrix}F_x & F_\lambda\\ t_x^{\top} & t_\lambda\end{bmatrix}.
\end{equation}
CPF 方法的全部价值浓缩在下述定理中：增广化把折点从奇异点变回正则点。

\begin{theorem}[增广雅可比在二次折点非奇异]\label{thm:pf-th-vs-augmented}
设 $(x^*,\lambda^*)$ 为满足定理~\ref{thm:pf-th-vs-snb} 条件的二次折点，$t$ 为
该点切向量。则增广雅可比 $\Phi_y$ 在 $(x^*,\lambda^*)$ 非奇异。
\end{theorem}
\begin{proof}
先定 $t$ 的形态。切向量条件 $F_x t_x+F_\lambda t_\lambda=0$ 左乘 $w^{\top}$ 得
$t_\lambda\,w^{\top}F_\lambda=0$，由横截条件 $t_\lambda=0$；于是 $F_xt_x=0$，
$t_x=c\,v$，归一化后 $t=(v/\|v\|,\,0)$。

设 $\Phi_y[u;\ \mu]=0$，即
\begin{equation}
Ju+\mu F_\lambda=0,\qquad t^{\top}\begin{bmatrix}u\\ \mu\end{bmatrix}=0.
\end{equation}
第一式左乘 $w^{\top}$：$\mu\,w^{\top}F_\lambda=0\Rightarrow\mu=0$；于是
$Ju=0\Rightarrow u=c'v$；代入第二式：$t^{\top}[u;0]=c'\,v^{\top}v/\|v\|
=c'\|v\|=0\Rightarrow c'=0$。故零空间平凡，$\Phi_y$ 非奇异。
\end{proof}

\begin{remark}[校正的收敛性与实现含义]
定理~\ref{thm:pf-th-vs-augmented} 表明伪弧长校正在折点处仍是\emph{正则}
Newton 问题：预测点足够接近解支时 Newton 校正保持二次收敛（标准 Newton--
Kantorovich 框架），可以数值上"翻越"鼻点并采样下支。代价是每步需求解
$(n+1)\times(n+1)$ 增广系统；Ajjarapu--Christy~\cite{pfvs:ajjarapu} 的原始
CPF 与 CPFLOW~\cite{pfvs:chiang} 均利用增广矩阵相对 $J$ 只多一行一列的
加边结构（bordered matrix）做块消元以保持稀疏效率。
\end{remark}

\subsubsection{步长控制的理论依据}
弧长步长 $s$ 的自适应控制遵循两条原则：
\begin{itemize}
  \item \emph{曲率原则}：解支曲率大（鼻点附近）时切向量预测的一阶误差
        $O(s^{2}\kappa)$（$\kappa$ 为曲率）放大，需减小 $s$；平直段可放大 $s$。
        成败倍率（校正成功 $s\leftarrow\min(\gamma_{+}s,s_{\max})$，失败
        $s\leftarrow\gamma_{-}s$ 重试，$s<s_{\min}$ 终止）是其最简实现；
  \item \emph{收敛域原则}：Newton 校正对初值的吸引域有限，$s$ 过大时预测点
        落在吸引域外，收缩重试相当于在解支上做信任域回退。
\end{itemize}
步长上下限 $[s_{\min},s_{\max}]$ 同时承担鼻点定位精度的角色：报告的
$\lambda_{\max}$ 是采样点中的最大值，分辨率受 $s_{\max}$ 控制，需要精确鼻点
时应收紧步长或做局部加密。

\begin{gapnote}
平台 CPF 实现了割线切向量+伪弧长超平面校正的完整框架，但两处与理论全配置
有差距：(i) 翻越鼻点后仅保留固定个数（默认 2）的下支采样点即主动停止，
非完整下支追踪；(ii) 校正器使用有限差分稠密雅可比 + 稠密 LU，未实现
定理~\ref{thm:pf-th-vs-augmented} 加边结构的稀疏块消元，大规模系统下校正
代价为 $O(n^{3})$ 级。另：存在任一在运 PV 母线时求解器自动回退自然参数化
（实现层连续潮流章），此时定理~\ref{thm:pf-th-vs-augmented} 的保障不再适用，
只能从上支逼近鼻点。
\end{gapnote}

\subsection{VSC 换流器通用稳态理论}\label{sec:pf-th-vs-vsc}

\subsubsection{注入模型与功率守恒}
\begin{definition}[VSC 稳态注入模型]\label{def:pf-th-vs-vsc}
稳态潮流层面，VSC 换流器 $k$ 抽象为连接 AC 母线 $i$ 与 DC 母线 $j$ 的
双端口耦合设备，端口注入 $(P_{ac},Q_{ac},P_{dc})$（注入网络为正）满足
\emph{功率守恒}
\begin{equation}
P_{ac}+P_{dc}+P_{loss}=0,\qquad P_{loss}\ge0,
\end{equation}
换流器内部开关过程全部平均化，不显式出现。AC 网络只看到注入
$(P_{ac},Q_{ac})$，DC 网络只看到注入 $P_{dc}$，两域通过守恒方程耦合。
\end{definition}

\subsubsection{控制模式、自由度与方程闭合原则}
\begin{proposition}[换流器方程闭合的计数判据]\label{prop:pf-th-vs-closure}
每台双端口 VSC 向网络引入三个端口量 $(P_{ac},Q_{ac},P_{dc})$；功率守恒（损耗
为端口量的确定函数）提供一条方程，闭合 $P_{dc}$。因此每台 VSC 的控制模式必须
额外交由\emph{恰好两条}独立等式约束：\emph{有功通道一条}（定 $P_{ac}$ /
定 $P_{dc}$ / 控 $V_{dc}$ 刚性或下垂——此时 $P_{ac}$ 释放为自由平衡量，相应
方程转移到 DC 侧），\emph{无功/电压通道一条}（定 $Q_{ac}$ 或控 $V_{ac}$）。
少于两条则方程组欠定（自由度未被吸收，岛奇异）；多于两条则过约束（整定冲突，
无解）。
\end{proposition}
\begin{remark}
命题~\ref{prop:pf-th-vs-closure} 是\emph{必要性计数判据}而非充分条件：计数
闭合不保证物理解存在（可能无可行运行点），也不排除病态。平台以"所有自由
注入必须有闭合方程"为设计原则，并以规则目录形式落地为求解前预校核。
\end{remark}

用布尔变量记各通道占用：$z^{P_{ac}}$（AC 有功硬给定）、$z^{P_{dc}}$（DC 有功
硬给定）、$z^{V_{dc}}$（DC 电压刚性控制）、$z^{Q}$、$z^{V_{ac}}$、$z^{pf}$
（功率因数控制），闭合原则的互斥形式为
\begin{equation}
z^{P_{ac}}+z^{P_{dc}}+z^{V_{dc}}\le1,\qquad
z^{Q}+z^{V_{ac}}+z^{pf}\le1.
\end{equation}
DC 侧构网（$z^{V_{dc}}=1$）时有功通道被占用，$P_{ac}$、$P_{dc}$ 均不得再硬
给定（调度值只能作初值）；双侧同时构网（AC 侧定相角+幅值且 DC 侧定电压）
对普通两端口换流器缺少承担有功不平衡的自由度，属于需要内部储能等额外机制
的高级模式。

VSC 稳态控制的七类典型分类（$\delta_s{+}V_s$ / $P_s{+}V_s$ / $P_s{+}Q_s$ /
$U_{dc}{+}Q_s$ / $U_{dc}{+}V_s$ / droop $U_{dc}{+}V_s$ / droop $U_{dc}{+}Q_s$）
是上述通道规则的枚举展开：每类模式恰好占用两个通道各一次。

\subsubsection{DC 电压下垂律}
多换流器共联 DC 网络时，DC 电压控制必须允许多源并联。标准形式为电压下垂
\begin{equation}
P_{dc}=P^{0}+K\bigl(V_{dc}^{set}-V_{dc}\bigr),\qquad K>0,
\end{equation}
电压偏低时注入增加，符号正确。平台采用\emph{平方电压下垂}变体
\begin{equation}
P_{tr}=k_{vdc}\bigl(v_{dc}^{2}-v_{dc,set}^{2}\bigr),\qquad
P_{ac}=P_{tr}-P_{loss}(P_{tr},v_{dc}).
\end{equation}
\begin{remark}[$v^{2}$ 下垂的动机]
恒电阻 DC 支路功率 $p_{ij}=g_{ij}v_i(v_i-v_j)$ 在标称点附近
$v_i\approx v_j\approx1$ 有 $p_{ij}\approx g_{ij}(v_i-v_j)\approx
\tfrac{g_{ij}}{2}(u_i-u_j)$（$u=v^{2}$）：网络方程对平方电压近似仿射。
以 $u$ 为下垂变量使换流器律与网络律在同一变量上线性相容，小信号增益
$\mathrm{d}P/\mathrm{d}v=2k_{vdc}v$ 在标称点即 $2k_{vdc}$，量纲与调度口径
（MW/pu$^{2}$）直接可读。
\end{remark}

\subsubsection{损耗模型推导}
VSC 稳态损耗的工程口径有两类：
\begin{itemize}
  \item \emph{恒定效率（Linear）}：由 $\eta=P_{out}/P_{in}$ 直接得
        \begin{equation}
        P_{loss}=(1-\eta)\,|p|,
        \end{equation}
        $p$ 为传输功率。损耗与电流解耦，适合规划级估算；
  \item \emph{电流相关二次模型（CurrentBased）}：计及损耗的物理组成——
        固定项（控制、辅助电源、空载）$a$；开关损耗与阀正向压降近似正比于
        电流 $bI$；导通（欧姆）损耗正比于 $cI^{2}$：
        \begin{equation}
        P_{loss}=a+bI+cI^{2},\qquad I=\frac{|p|}{v_{dc}},
        \end{equation}
        电流取 DC 侧口径（$I=|p|/v_{dc}$，低压重载时损耗显著上升，与
        $\eta$ 恒定的 Linear 模型定性不同）。
\end{itemize}
任何损耗模型必须满足 $P_{loss}\ge0$；模型给出负损耗属于数据错误而非数值
现象。双向运行时效率口径按功率方向分段（注入约定下以 $|p|$ 对称化）。

\subsubsection{容量圆、电流限与调制比约束的几何意义}
VSC 的物理不等式约束在运行平面上有明确几何形态：
\begin{itemize}
  \item \emph{容量圆}：换流阀与内电抗的电流定额把 AC 侧运行点限制在圆盘
        \begin{equation}
        P_{ac}^{2}+Q_{ac}^{2}\le S_{rated}^{2},
        \end{equation}
        几何上为 $P$--$Q$ 平面中以原点为中心的圆；与有功盒约束
        $P\in[P_{\min},P_{\max}]$、无功盒约束求交即得可行域；
  \item \emph{AC 电流限}：$I_{ac}=S_{ac}/(\sqrt3\,V_{ll})\le I_{ac}^{\max}$
        （名值；标幺口径 $I_{ac}=S_{ac}/V_m$），本质是容量圆随端电压收缩的
        动态版本；
  \item \emph{DC 电流限}：$I_{dc}=|P_{dc}|/v_{dc}\le I_{dc}^{\max}$，
        $v_{dc}$ 低时同功率下电流放大，是低压 DC 网络的关键校核；
  \item \emph{调制比窗口}：PWM 合成要求
        \begin{equation}
        V_{ac}^{model}=K_m\,m\,V_{dc},\qquad m\in[m_{\min},m_{\max}],
        \end{equation}
        即给定 $V_{dc}$ 时可合成的 AC 电压幅值被限制在线段
        $[K_m m_{\min}V_{dc},\,K_m m_{\max}V_{dc}]$ 上；$V_{dc}$ 过低时窗口
        无法覆盖所需 $V_{ac}$（几何上窗口平移出可行区），等价可行性校核为
        $m=V_{ac}^{model}/(K_mV_{dc})\in[m_{\min},m_{\max}]$。
\end{itemize}

\begin{gapnote}
上述不等式在平台默认路径下\textbf{不嵌入 Newton 方程}：容量圆、AC/DC 电流限、
调制比窗口与 DC-DC 占空比窗口均在收敛后逐台校核并以告警输出
（\texttt{ACDC-PHYS-01..05}、\texttt{DCDC-PHYS-01}），opt-in 开关可把越限
Newton 根整体判为物理不可行（硬门卫），但方程内部不做再调度。唯一的方程
内嵌例外是 VSC 交流电流限的半光滑 NCP 路径（min--max 中值 NCP 函数及其
Chen--Harker--Kanzow--Smale 平滑，见~\cite{pfvs:facchinei}），为 opt-in
且仅覆盖部分控制模式；完整的容量圆/调制比方程内嵌处理属未实现扩展。
\end{gapnote}

\subsection{LCC 换流器准稳态理论}\label{sec:pf-th-vs-lcc}

\subsubsection{六脉动桥的平均直流电压}
LCC（电网换相）换流器以晶闸管阀按电网电压过零点自然换相。6 脉动桥在一个
周期内由 6 段 $60^{\circ}$ 的线电压片段拼成直流电压波形；设阀侧线电压有效值
$E_0$，触发延迟角 $\alpha$，不计换相重叠时平均直流电压为片段积分
\begin{equation}
U_{d0}=\frac{1}{\pi/3}\int_{-\pi/6+\alpha}^{\pi/6+\alpha}\sqrt2\,E_0\cos\theta\,
\mathrm{d}\theta
=\frac{3\sqrt2}{\pi}E_0\cos\alpha.
\end{equation}
定义理想空载直流电压 $U_{d0}^{(0)}=\tfrac{3\sqrt2}{\pi}E_0$，则
$U_d=U_{d0}^{(0)}\cos\alpha$。$\alpha$ 是整流侧唯一控制量：
$\alpha<90^{\circ}$ 时 $U_d>0$（整流），$\alpha>90^{\circ}$ 时反号（逆变）。

\subsubsection{换相重叠与外特性}
换相期间两个阀同时导通，重叠角 $\mu$ 内直流电压被两相电压平均值拉低，
损失面积正比于换相电抗与直流电流。标准结果（Kimbark~\cite{pfvs:kimbark}、
Arrillaga~\cite{pfvs:arrillaga}）为
\begin{equation}
\cos\alpha-\cos(\alpha+\mu)=\frac{\sqrt2\,X_c I_d}{E_0},
\end{equation}
\begin{equation}
U_d=U_{d0}^{(0)}\cos\alpha-\frac{3}{\pi}X_c I_d,
\end{equation}
重叠的净效应等效为一个\emph{不消耗有功}的换相电阻 $R_c=\tfrac{3}{\pi}X_c$
（压降代表交流侧无功消耗，而非热损耗）。每极串联 $n_b$ 个桥时电压、压降
同倍放大：$U_{d0}=\tfrac{3\sqrt2}{\pi}n_bE_0$，$U_d=U_{d0}\cos\alpha-
\tfrac{3}{\pi}n_bX_cI_d$。

逆变侧以越前角 $\beta=\pi-\alpha$ 与\emph{关断角} $\gamma=\beta-\mu$ 描述，
$\gamma$ 是阀恢复正向阻断能力所需的最低裕度，逆变站正常控制方式为定关断角
（CEA，$\gamma=\gamma_{\min}$）以防换相失败；外特性写为
\begin{equation}
U_d=U_{d0}\cos\gamma-\frac{3}{\pi}n_bX_cI_d+n_b\,dU_v,
\end{equation}
$dU_v$ 为阀正向压降项。整流侧对应方程以 $\alpha$ 描述，$dU_v$ 取负号。

\subsubsection{无功消耗与功率因数}
基波电流因触发延迟与换相重叠而滞后阀侧电压，换流器\emph{恒吸收无功}：
\begin{equation}
Q=P\tan\varphi,\qquad
\cos\varphi\approx\frac{U_d}{U_{d0,t}},\qquad
U_{d0,t}=\frac{3\sqrt2}{\pi}n_bE_t,
\end{equation}
功率因数用带载阀侧电压 $E_t$（区别于经分接头决定的空载量 $E_0$）评估。
LCC 站的无功需求约为其传输有功的 $40\%\sim60\%$，必须就地补偿，这是
LCC-HVDC 系统设计的核心约束之一~\cite{pfvs:arrillaga}。

\begin{remark}[控制模式与电流单向性]
准稳态控制模式：整流侧定电流（CC，正常方式，$\alpha$ 为自由量）、定功率
（CP）、定 $\alpha$（备用）；逆变侧定 $\gamma$（CEA，正常方式）、定电流、
定功率。晶闸管阀单向导电决定 $I_d$ 不可反号：功率反转靠电压极性反转实现。
$I_d$ 触及边界 $[0,I_{rated}]$ 时原控制律不再保持，站退化为恒流特性
（$\gamma$ 保持在 $\gamma_{\min}$ 之上）。分接头作为慢速离散自由度在
Newton 方程外（外环）调整，不是潮流状态量。
\end{remark}

\subsection{DC-DC 变换器平均模型}\label{sec:pf-th-vs-dcdc}

\subsubsection{CCM 平均模型推导}
潮流级 DC-DC 模型基于开关周期平均与电感伏秒平衡：连续导通模式（CCM）下，
稳态电感电流周期回复要求一个开关周期内电感两端伏秒积为零。

\begin{example}[三种基本拓扑的电压律推导]
设占空比 $D$（导通占比），开关周期 $T$。
\begin{itemize}
  \item \emph{Buck}：导通段电感承受 $V_{in}-V_{out}$，关断段承受 $-V_{out}$，
        伏秒平衡 $(V_{in}-V_{out})DT=V_{out}(1-D)T$ 给出
        $V_{out}=DV_{in}$；
  \item \emph{Boost}：导通段 $V_{in}$，关断段 $V_{in}-V_{out}<0$，
        $V_{in}DT+(V_{in}-V_{out})(1-D)T=0$ 给出
        $V_{out}=V_{in}/(1-D)$；
  \item \emph{Buck-Boost}（非反相口径）：$V_{in}DT=V_{out}(1-D)T$ 给出
        $V_{out}/V_{in}=D/(1-D)$；
  \item \emph{隔离型}：经变比 $n$ 的变压器耦合，$V_{out}=nM(D)V_{in}$，
        $M(D)$ 为调制增益。
\end{itemize}
\end{example}

稳态潮流关心的是反问题——由端口电压反演占空比并做可行性校核：
\begin{equation}
D=\begin{cases}
V_{out}/V_{in} & \text{Buck}\\
1-V_{in}/V_{out} & \text{Boost}\\
V_{out}/(V_{in}+V_{out}) & \text{Buck-Boost}\\
V_{out}/(nV_{in}) & \text{隔离型（}M\text{ 口径）}
\end{cases},
\qquad D\in[D_{\min},D_{\max}].
\end{equation}

\subsubsection{功率守恒、效率与欧姆损耗}
端口注入约定下 $P_{in}^{port}+P_{out}^{port}+P_{loss}=0$。效率按功率方向
分段施加：正向（$P_{out}^{ref}\ge0$）时 $P_{in}=P_{out}^{ref}/\eta$，反向时
$P_{in}=\eta\,P_{out}^{ref}$；等值电阻产生的 $I^{2}R$ 损耗
\begin{equation}
P_{loss}^{I^2R}=r_{eq}\,\frac{P_{in}^{2}}{v_{dc,in}^{2}},\qquad
P_{out}=P_{out}^{ref}-P_{loss}^{I^2R}.
\end{equation}

\begin{remark}[平均模型的适用边界]
上述公式默认 CCM、稳态平均，忽略开关压降、电感电阻、死区、轻载 DCM、控制
环动态与热约束。因此它们适用于潮流级可行性校核与稳态控制方程，不等同于
开关级精确仿真；电流约束 $I_{in/out}=|P_{in/out}|/V_{in/out}\le I^{\max}$
与单向器件的功率方向约束属不等式校核范畴。
\end{remark}

\subsection{AC/DC 接口方程的统一表述}\label{sec:pf-th-vs-unified}

\subsubsection{统一方程组结构}
混合 AC/DC 潮流把两域未知量放入同一方程组联立求解。状态与方程结构：
\begin{equation}
x=\begin{bmatrix}\theta_{\text{非 slack}}\\ V_{m,\mathrm{PQ}}\\ v_{dc,\text{非参考}}\end{bmatrix},
\qquad
F(x)=\begin{bmatrix}
P^{spec}(x)-P^{calc}(x)\\
Q^{spec}(x)-Q^{calc}(x)\\
p_{dc}^{spec}(x)-v_{dc}\odot\bigl(G_{dc}v_{dc}\bigr)
\end{bmatrix}=0,
\end{equation}
其中换流器贡献全部经 $spec$ 向量进入：AC 侧 $P^{spec}\mathrel{+}=P_{ac}$、
$Q^{spec}\mathrel{+}=Q_{ac}$，DC 侧 $p_{dc}^{spec}\mathrel{+}=P_{dc}$；
VSC 注入、LCC 站注入与 DC-DC 端口注入按各自模式在当前迭代点求值。雅可比
呈块结构
\begin{equation}
J=\begin{bmatrix}
J_{AC} & J_{AC\leftarrow DC}\\
J_{DC\leftarrow AC} & J_{DC}
\end{bmatrix},
\end{equation}
对角块内换流器自洽导数（如 $\partial p_{dc}^{spec}/\partial v_{dc}$、CEA 站
外特性导数）是 DC 方程 Newton 收敛的必要条件；交叉块
$J_{AC\leftarrow DC}$、$J_{DC\leftarrow AC}$ 是否装配是精度/稳健性取舍。

\subsubsection{DC 岛可解性：电压参考的必要性}
\begin{theorem}[DC 岛电压参考条件]\label{thm:pf-th-vs-dcref}
设某 DC 电压岛连通、仅含恒电阻支路（电导 Laplacian $L=G_{dc}$）。
(i) 若岛内无任何电压成形机制（无钉住母线、无下垂源、无特征控制换流站），
则 DC 子方程在无载点 $v=c\,\mathbf 1$ 处雅可比奇异：$J_{dc}=cL$；
(ii) 钉住岛内任一节点电压（删除对应行列）后，接地 Laplacian $L_r$ 在连通
假设下正定，DC 子块局部正则。
\end{theorem}
\begin{proof}
(i) Laplacian 行和为零：$L\mathbf1=0$。DC 方程
$F(v)=v\odot(Lv)-p^{spec}$ 的雅可比为
$J_{dc}=\mathrm{diag}(Lv)+L\,\mathrm{diag}(v)$；在 $v=c\mathbf1$ 处
$Lv=0$，故 $J_{dc}=cL$，$\mathbf1\in\ker J_{dc}$，奇异。
(ii) 删去参考节点后的接地 Laplacian $L_r$ 满足能量恒等式：对任意 $u$
（参考分量置零），
\begin{equation}
u^{\top}L_r\,u=\sum_{(i,j)}g_{ij}(u_i-u_j)^{2}\ge0,
\end{equation}
等号成立要求所有相连节点分量相等；连通性归纳至参考节点分量 $0$，故
$u=0$，$L_r$ 正定、非奇异。
\end{proof}

\begin{remark}[下垂源与参与因子的正则化作用]
下垂源注入 $P=P^{0}+K(V^{set}-V)$ 为雅可比对角贡献正项 $K$：岛矩阵由
$L$ 变为 $L+\mathrm{diag}(K_e)$，无需钉住任何节点即正定（弱对角占优加严格
正对角）。多下垂源并联要求总增益 $K_{total}=\sum_eK_e>0$；多刚性源位于同一
电气节点且整定不一致（$V_1^{set}\neq V_2^{set}$）时方程直接矛盾；刚性电压
控制与参与因子功率分摊（$P_e=P_e^{0}+\alpha_e\lambda_d$，$\sum_e\alpha_e=1$）
叠加占用同一有功通道，构成过约束——这些一致性条件与
命题~\ref{prop:pf-th-vs-closure} 的计数判据同源。
\end{remark}

\begin{remark}[方程闭合计数的系统级形式]
设 AC 非 slack 母线 $n_p$ 个、PQ 母线 $n_q$ 个、DC 非参考母线 $n_{dc}$ 个，
则网络方程共 $n_p+n_q+n_{dc}$ 条，与状态维数相等；每台换流器/换流站按
命题~\ref{prop:pf-th-vs-closure} 自带两条控制等式（或以 slack/参考身份换取
网络行列的剔除），增广模式（如 AC 电压成形母线的增广方程、NCP 限值块）
同步增广状态与方程。系统级闭合等价于：每个 DC 岛恰有一个电压成形机制
（定理~\ref{thm:pf-th-vs-dcref}），且每台设备的自由注入均被控制等式吸收。
\end{remark}

\subsection{数值算例与基准}\label{sec:pf-th-vs-numexp}

\srcpath{tools/pf\_doc\_benchmark.cpp}（模式 \texttt{cpf}）在例~\ref{ex:pf-th-vs-twobus} 的
两母线系统（$E=1$、$X=0.5$~pu、负荷 $50+\mathrm j20$~MVA、基准 100~MVA）上运行平台弧长
连续潮流（\srcpath{src/power\_flow/voltage\_stability.cpp:CpfSolver}），与解析鼻点对照。

\begin{example}[两母线 CPF 与解析鼻点]\label{ex:pf-th-vs-cpfnum}
比例加载 $P(\lambda)=(1{+}\lambda)\cdot0.5$、$Q(\lambda)=(1{+}\lambda)\cdot0.2$（pu）。
解析上判别式 $\Delta=1-0.4s-0.25s^2=0$（$s=1{+}\lambda$）给出 $s^*=1.3541$，故
$\lambda^*=0.3541$、$P_{\max}=0.677$~pu$=67.70$~MW、
$V_{\mathrm{nose}}=\sqrt{(1-2Q_LX)/2}=0.6038$~pu。平台弧长 CPF（默认翻越鼻点）实测：
\begin{center}\footnotesize
\begin{tabular}{@{}lccc@{}}
\toprule
量 & 解析 & CPF 实测 & 相对误差\\
\midrule
$\lambda_{\max}$      & $0.3541$ & $0.3536$ & $0.14\%$\\
$P_{\max}$ (MW)       & $67.70$  & $67.68$  & $0.03\%$\\
$V_{\mathrm{nose}}$ (pu) & $0.6038$ & $0.6130$ & $1.5\%$\\
\bottomrule
\end{tabular}
\end{center}
$\lambda_{\max}$ 与 $P_{\max}$ 吻合到 $0.1\%$ 量级，验证弧长校正在鼻点邻域仍为正则
Newton 问题（定理~\ref{thm:pf-th-vs-augmented}）；$V_{\mathrm{nose}}$ 的 $1.5\%$ 偏差
源于鼻点采样步长（$\Delta\lambda\!\sim\!0.02$）——CPF 报告的是最接近鼻尖的\emph{采样}点，
其电压略高于精确鼻尖 $E/\!\sqrt2$。求解器以 $22$ 个追踪点翻越鼻点并采样下支，
$\mathtt{nose\_found}$、$\mathtt{arc\_length\_used}$ 均为真，与弧长连续潮流的设计一致。
\end{example}

\subsection{与实现的对应}\label{sec:pf-th-vs-impl}
\begin{center}\footnotesize
\begin{longtable}{@{}L{6.8cm}L{8.6cm}@{}}
\toprule
理论条目 & 实现状态\\
\midrule
\endfirsthead
\toprule
理论条目 & 实现状态\\
\midrule
\endhead
仿射负荷方向 $P(\lambda)=P^{0}+\lambda\Delta P$ 与方向构造 &
\implfull{src/power\_flow/voltage\_stability.cpp:apply\_direction}（比例/指定母线两工厂见 include/hacdcpf/power\_flow/voltage\_stability.hpp:CpfDirection）\\
自然参数化（固定 $\lambda$ 递增 + NR 校正） &
\implfull{src/power\_flow/voltage\_stability.cpp:solve}（legacy 模式与首个延拓点；有在运 PV 母线时自动回退此口径并如实声明）\\
弧长增广系统（$\lambda$ 为未知量、状态布局与打包） &
\implfull{src/power\_flow/voltage\_stability.cpp:make\_state\_layout}（pack/unpack 见同文件 pack\_augmented\_state、unpack\_augmented\_state）\\
割线切向量预测与符号一致性校正 &
\implfull{src/power\_flow/voltage\_stability.cpp:solve}（定方向预测可经 use\_secant\_predictor=false 关闭）\\
伪弧长超平面校正 $t^{\top}(y-y_{pred})=0$ &
\implpart{src/power\_flow/voltage\_stability.cpp:arc\_length\_correct 实现超平面校正，但用有限差分稠密雅可比 + FullPivLU + 回溯线搜索，未做加边块消元与稀疏解析雅可比}\\
成败倍率步长控制（$s_{\min}/s_{\max}/\gamma_\pm$） &
\implfull{src/power\_flow/voltage\_stability.cpp:solve}（step\_grow/step\_shrink/step\_min/step\_max）\\
鼻点判定（$\lambda$ 切向符号反转）与下支采样 &
\implpart{翻越检测已实现（solve，置 nose\_found），但仅保留 lower\_branch\_steps（默认 2）个下支点即停，非完整下支追踪；相邻采样点间不做鼻点加密二分}\\
负荷裕度指标族（$\lambda^*$、$P_{\max}$、$P_{margin}$） &
\implfull{src/power\_flow/voltage\_stability.cpp:compute\_vsi}（汇总口径取 trace 中 $\lambda$ 最大点）\\
$L$ 指标、模态分析/参与因子、Q--V 曲线 & \implnone\\
VSC 注入模型与功率守恒 $P_{ac}+P_{dc}+P_{loss}=0$ &
\implfull{src/power\_flow/converter\_model.cpp:converter\_ac\_injection}（DC 侧见同文件 converter\_dc\_injection）\\
Linear/CurrentBased 损耗模型及导数 &
\implfull{src/power\_flow/converter\_model.cpp:converter\_loss}（导数 converter\_loss\_jacobian；$I=|p|/\max(v_{dc},0.1)$ 带下限钳制）\\
$v^{2}$ 下垂律 $P_{tr}=k_{vdc}(v_{dc}^{2}-v_{dc,set}^{2})$ &
\implfull{src/power\_flow/converter\_model.cpp:converter\_ac\_injection}（交叉导数 converter\_ac\_jacobian\_vdc、converter\_dc\_jacobian\_vdc）\\
控制模式枚举与七类分类映射 &
\implfull{include/hacdcpf/model/enums/converter\_enums.hpp:ConverterMode}（六枚举；七类分类 ACDCControlMode 同文件；角色判定 include/hacdcpf/model/device\_control\_role.hpp:resolve\_device\_control\_role）\\
方程闭合原则（计数判据、互斥规则） &
\implpart{以规则目录形式落地为求解前预校核 src/power\_flow/converter\_coordination.cpp:evaluate\_converter\_coordination（ACDC-GFM/DCISLAND 等系列），非严格自由度计数；PQ$\leftrightarrow$VDC\_Q 切换见 src/power\_flow/newton\_solver.cpp:apply\_converter\_mode\_switching}\\
容量圆/电流限/调制比窗口不等式 &
\implpart{默认收敛后逐台校核仅告警（ACDC-PHYS-01..05，src/api/hacdcpf.cpp:populate\_derived\_results）；opt-in 硬门卫 enforce\_converter\_physical\_limits 判根不可行；VSC 交流电流限另有方程内嵌半光滑 NCP 路径 src/power\_flow/vsc\_limit\_ncp.cpp:evaluate\_vsc\_limit\_ncp（opt-in enable\_limit\_ncp，限 PQ/VDC\_Q/带阻抗 GFM 模式）}\\
LCC 六脉动桥外特性与 $R_c=\frac{3}{\pi}X_c$ 压降 &
\implfull{include/hacdcpf/power\_flow/lcc\_model.hpp:lcc\_operating\_point}（实现 src/power\_flow/lcc\_model.cpp；整流/逆变、CEA/CC/CP/定 $\alpha$ 分派）\\
LCC 无功 $Q=P\tan\varphi$、$\cos\varphi\approx U_d/U_{d0,t}$ &
\implfull{include/hacdcpf/power\_flow/lcc\_model.hpp:lcc\_ac\_injection}（$E_t$/$E_0$ 双电压口径与实现一致）\\
LCC 电流单向钳制 $[0,I_{rated}]$ 与恒流退化 &
\implfull{include/hacdcpf/power\_flow/lcc\_model.hpp:lcc\_operating\_point}（id\_at\_limit 与 [LCC-PHYS-03] 诊断）\\
分接头的 Newton 外环地位 &
\implfull{include/hacdcpf/power\_flow/lcc\_model.hpp}（头文件注释约定 R 卡外环包络；分接头非 Newton 状态量）\\
DC-DC 平均模型（Power/Droop/Voltage 三模式、效率分段、$I^2R$ 损耗） &
\implfull{src/power\_flow/converter\_model.cpp:dcdc\_power\_transfer}（雅可比 dcdc\_jacobian\_vdc）\\
占空比反演与可行性窗 $D\in[d_{\min},d_{\max}]$ &
\implfull{src/power\_flow/converter\_model.cpp:dcdc\_duty\_ratio}（Generic 拓扑不评估；校核为事后告警 DCDC-PHYS-01）\\
统一联立方程组与换流器 spec 注入 &
\implfull{src/power\_flow/jacobian\_builder.cpp:build\_power\_spec}（DC 注入装配 src/power\_flow/pf\_injection\_assembly.cpp:assemble\_dc\_injections；残差/雅可比 src/power\_flow/jacobian\_builder.cpp:evaluate\_residual\_impl）\\
DC 岛电压参考必要性（定理~\ref{thm:pf-th-vs-dcref}） &
\implfull{src/power\_flow/newton\_solver.cpp:plan\_dc\_island\_references}（岛识别 compute\_dc\_components；无参考自动提升 PQ 换流器或钉母线并告警）\\
多源协调（下垂并联 $K_{total}>0$、主从、参与因子 $\sum\alpha=1$） &
\implfull{src/power\_flow/newton\_solver.cpp:apply\_participation\_factor\_sharing}（规则校核 converter\_coordination.cpp 的 DCISLAND-*/PARTICIPATION 系列）\\
AC$\leftrightarrow$DC 交叉雅可比块 &
\implpart{src/power\_flow/jacobian\_builder.cpp:evaluate\_residual\_impl 已实现交叉块，但默认关闭（opt-in enable\_coupled\_jacobian）；DC 块内自洽导数与 LCC 外特性导数恒常装配}\\
戴维南最大功率/阻抗匹配判据（定理~\ref{thm:pf-th-vs-thevenin}，
式~\eqref{eq:pf-th-vs-thevenin}） &
\implnone（指标族仅 $\lambda$-裕度类；电压稳定极限由弧长 CPF 追踪，非阻抗匹配指标）\\
两母线 CPF 与解析鼻点数值基准（例~\ref{ex:pf-th-vs-cpfnum}） &
\implfull{tools/pf\_doc\_benchmark.cpp}（模式 cpf；调用 voltage\_stability.cpp:CpfSolver::solve）\\
\bottomrule
\end{longtable}
\end{center}

\subsection{参考文献}
{\renewcommand{\section}[2]{}%
\begin{thebibliography}{9}\small
\bibitem{pfvs:kundur} P.~Kundur, \emph{Power System Stability and Control},
McGraw-Hill, 1994.
\bibitem{pfvs:vancutsem} T.~Van Cutsem, C.~Vournas, \emph{Voltage Stability of
Electric Power Systems}, Kluwer Academic Publishers, 1998.
\bibitem{pfvs:ajjarapu} V.~Ajjarapu, C.~Christy, ``The continuation power flow:
a tool for steady state voltage stability analysis,'' \emph{IEEE Transactions on
Power Systems}, vol.~7, no.~1, pp.~416--423, 1992.
\bibitem{pfvs:chiang} H.-D. Chiang, A.~J. Flueck, K.~S. Shah, N.~Balu,
``CPFLOW: a practical tool for tracing power system steady-state stationary
behavior due to load and generation variations,'' \emph{IEEE Transactions on
Power Systems}, vol.~10, no.~2, pp.~623--634, 1995.
\bibitem{pfvs:gao} B.~Gao, G.~K. Morison, P.~Kundur, ``Voltage stability
evaluation using modal analysis,'' \emph{IEEE Transactions on Power Systems},
vol.~7, no.~4, pp.~1529--1542, 1992.
\bibitem{pfvs:kessel} P.~Kessel, H.~Glavitsch, ``Estimating the voltage
stability of a power system,'' \emph{IEEE Transactions on Power Delivery},
vol.~1, no.~3, pp.~346--354, 1986.
\bibitem{pfvs:keller} H.~B. Keller, ``Numerical solution of bifurcation and
nonlinear eigenvalue problems,'' in P.~H. Rabinowitz (ed.), \emph{Applications
of Bifurcation Theory}, Academic Press, pp.~359--384, 1977.
\bibitem{pfvs:kimbark} E.~W. Kimbark, \emph{Direct Current Transmission},
vol.~I, Wiley-Interscience, 1971.
\bibitem{pfvs:arrillaga} J.~Arrillaga, \emph{High Voltage Direct Current
Transmission}, 2nd ed., IEE Power Engineering Series, 1998.
\bibitem{pfvs:facchinei} F.~Facchinei, J.-S. Pang, \emph{Finite-Dimensional
Variational Inequalities and Complementarity Problems}, Springer, 2003.
\bibitem{pfvs:vu} K.~Vu, M.~M. Begovic, D.~Novosel, and M.~M. Saha, Use of local
measurements to estimate voltage-stability margin, \emph{IEEE Transactions on
Power Systems}, vol.~14, no.~3, pp.~1029--1035, 1999.
\end{thebibliography}}

% ── 实现转写层（源码等价）──
% 第 2 章：SolverData 装配与网络矩阵
% 本章文件由 \input 引入，不含导言区。
\section{SolverData 装配与网络矩阵}\label{sec:assembly}

\subsection{SolverData 结构}

\compmeta{SolverData}{\srcpath{include/hacdcpf/assembly/solver\_data.hpp}}

\subsubsection*{功能说明}
\texttt{SolverData} 是潮流求解的\textbf{扁平 POD 数据底座}：rich 工程模型经
canonical 投影后，全部网络、设备、预聚合注入与矩阵以 0-based 向量形式集中存放，
求解器不再接触工程语义层。它同时携带投影回溯映射与精确 Newton 的开关量，
是"装配一次、多求解器复用"的核心（FDPF、DC、Newton、CPF 校正器均以其为输入）。

\subsubsection*{字段分组}
（\srcpath{include/hacdcpf/assembly/solver\_data.hpp:SolverData}）
\begin{itemize}
  \item \textbf{网络}：\fld{ac\_buses}、\fld{ac\_branches}、\fld{dc\_buses}、
        \fld{dc\_branches}、\fld{converters}（VSC）、\fld{dcdc\_converters}、
        \fld{energy\_routers}、\fld{switches}；
  \item \textbf{设备}：AC 侧 \fld{generators}、\fld{loads}、\fld{flexible\_loads}、
        \fld{static\_generators}、\fld{renewable\_gens}、\fld{pv\_systems}、
        \fld{storage\_units}；DC 侧 \fld{dc\_loads}、\fld{dc\_storage}、
        \fld{dc\_static\_generators}、\fld{dc\_pv\_arrays}；另有 \fld{shunts}、
        \fld{charging\_stations}、\fld{chargers}、\fld{external\_grids}、
        \fld{vpps}、\fld{microgrids}、\fld{mobile\_storage}；
  \item \textbf{标量与 ZIP}：\fld{base\_mva}=100、\fld{loss\_model}=Linear、
        \fld{zip\_pw[3]}/\fld{zip\_qw[3]}=\{1,0,0\}；
  \item \textbf{负荷聚合缓存}：\fld{has\_component\_loads}、\fld{pd\_pu}/\fld{qd\_pu}、
        每母线 ZIP 权重 \fld{bus\_zip\_zp}/\fld{\_ip}/\fld{\_pp}/\fld{\_zq}/\fld{\_iq}/\fld{\_pq}；
  \item \textbf{精确 Newton 开关}：\fld{enable\_coupled\_jacobian}、
        \fld{enable\_augmented\_equations}、\fld{enable\_semi\_smooth\_newton}、\fld{ncp\_mu}；
  \item \textbf{预聚合注入与矩阵}：\fld{pg}、\fld{qg}（Eigen 向量）、
        \fld{ybus}（稀疏复数阵）、\fld{gdc}（稀疏实数阵）、\fld{build\_id}（原子递增缓存键）；
  \item \textbf{投影回溯}：\fld{bus\_merge\_map}、\fld{projection\_certificate}、
        \fld{projection\_report}（\texttt{std::optional}）。
\end{itemize}
配套映射类型 \srcpath{BusIndexMap}/\srcpath{BranchIndexMap}/\srcpath{SystemIndexMap}
见 \srcpath{include/hacdcpf/assembly/index\_map.hpp:BusIndexMap}。

\subsection{装配流程}
入口 \srcpath{make\_solver\_data(const HybridPowerSystem\&, LossModelType)}
（\srcpath{src/power\_flow/solver\_data.cpp:make\_solver\_data}），步骤：
\begin{enumerate}
  \item rich$\to$canonical 投影 \srcpath{projection::RichToCanonicalOperator::apply}
        （:268--269，含开关/零阻抗合并）；
  \item \srcpath{materialize\_dc\_storage} 与逐字段 move；\texttt{StaticGeneratorDC}
        转换为 \texttt{StaticGenerator}（:make\_solver\_data）；
  \item \fld{build\_id} 取自原子计数器（:266,319），作为全部缓存的失效标记；
  \item \textbf{ExternalGrid 处理}（:make\_solver\_data）：把母线 \fld{vm\_pu}/\fld{va\_deg}
        写为外部电网整定，并将该母线 PQ 提升为 SLACK；
  \item \textbf{成网换流器提升} \srcpath{apply\_acpv\_voltage\_control}（:apply\_acpv\_voltage\_control）：
        按 \srcpath{resolve\_device\_control\_role}，AC\_GRID\_FORMING 换流器的交流母线
        提升为 SLACK 并钉住 \fld{v\_ac\_set\_pu}/\fld{v\_ac\_angle\_set\_deg}；
        AC\_PV / DC\_V\_DROOP\_AC\_V 模式提升为 PV；已有 SLACK 不动；
  \item \srcpath{rebuild\_matrices}（:rebuild\_matrices）$=$ 发电聚合 $+$ 负荷聚合 $+$
        $Y_{bus}$ 构建 $+$ $G_{dc}$ 构建。
\end{enumerate}

\subsection{注入聚合（指定功率）}
\subsubsection*{发电聚合}
\srcpath{aggregate\_generation}（src/power\_flow/solver\_data.cpp:aggregate\_generation）把下列注入聚进
\fld{pg}/\fld{qg}（除以 \fld{base\_mva} 转 pu）：发电机、
\texttt{StaticGenerator}（$p\cdot$\fld{scaling}）、可再生能源、
PV（经功率曲线 \srcpath{compute\_pv\_power\_mw}，见第~\ref{sec:distribution}~章）、
储能（放电为正）、VPP（PCC 交换功率）、并网微网（\fld{p\_exchange} 出口为正）、
非在途移动储能。

\subsubsection*{负荷聚合与 ZIP 归一化}
\srcpath{aggregate\_load\_demand}（src/power\_flow/solver\_data.cpp:aggregate\_load\_demand）：
Load 表为空（且无充电站）时 \fld{has\_component\_loads}=false，回退母线级
\fld{pd\_mw}/\fld{qd\_mvar}；否则以母线级需求为底（恒定功率，ZIP 种子 $(1,0,0)$），
叠加 $p_{\mathrm{mw}}\cdot$\fld{scaling} 的分量负荷，并按 $|P|/|Q|$ 加权平均
归一化出每母线 ZIP 系数（:160--168, 193--211）。充电站恒功率叠加
（kW$\to$MW，:173--188）。潮流方程中的指定注入为
（\srcpath{src/power\_flow/jacobian\_builder.cpp:build\_power\_spec}）：
\begin{align*}
P_i^{\mathrm{spec}} &= P_{g,i}-p_{d,i}\big(w_{p0}+w_{p1}V_i+w_{p2}V_i^2\big)\\
Q_i^{\mathrm{spec}} &= Q_{g,i}-q_{d,i}\big(w_{q0}+w_{q1}V_i+w_{q2}V_i^2\big)
\end{align*}
权重顺序 $w_0$=恒功率、$w_1$=恒电流、$w_2$=恒阻抗。VSC 交流注入经
\srcpath{converter\_ac\_injection} 加入 spec（src/power\_flow/jacobian\_builder.cpp:build\_power\_spec）。

\subsection{交流导纳阵 $Y_{bus}$}
\srcpath{build\_admittance\_matrix(data)}（\srcpath{src/power\_flow/admittance\_builder.cpp:build\_admittance\_matrix}），
π 型支路 + 复变比（MATPOWER \texttt{makeYbus} 约定）。对在役支路 $i\to j$
（\fld{from\_bus}/\fld{to\_bus} 为 1-based，代码中 $-1$）：
\begin{align*}
y_s &= \frac{1}{r+\mathrm{j}x},\qquad
y_{tt} = y_s+\mathrm{j}\frac{b}{2},\qquad
\hat{t} = |\mathrm{tap}|\,\mathrm{e}^{\mathrm{j}\,\mathrm{shift}\cdot\pi/180}\\
Y_{ff} &= \frac{y_{tt}}{|\hat{t}|^2},\qquad
Y_{ft} = -\frac{y_s}{\hat{t}^{*}},\qquad
Y_{tf} = -\frac{y_s}{\hat{t}}
\end{align*}
对角元累加 $Y_{ii}\leftarrow Y_{ii}+Y_{ff}$、$Y_{jj}\leftarrow Y_{jj}+y_{tt}$，
非对角元 $Y_{ij}=Y_{ft}$、$Y_{ji}=Y_{tf}$（:build\_admittance\_matrix）。装配层为 fail-close 口径：
$|z|$ 非有限或 $\le\texttt{kMinSeriesImpedancePu}=10^{-12}$ 即抛
\texttt{std::runtime\_error}（:50--55，阈值 :15）；\fld{tap} 非有限或
$\le 10^{-12}$ 同样抛出（:build\_admittance\_matrix）——零阻抗支路须在投影层先行合并，
装配层不做静默回退。母线对地并联（$V=1.0$~pu 功率口径，
MATPOWER 惯例）直接入对角（:build\_admittance\_matrix）：
\[
Y_{ii}\leftarrow Y_{ii}+\frac{g_s+\mathrm{j}b_s}{S_{\mathrm{base}}}
\]
独立 \texttt{Shunt} 表同样入对角；可投切并联按
\fld{bs\_per\_step}$\times$\fld{current\_step} 取值（:build\_admittance\_matrix）。
装配用三元组列表 + \texttt{setFromTriplets}。

\subsection{直流电导阵 $G_{dc}$}
\srcpath{build\_dc\_conductance(data)}（src/power\_flow/admittance\_builder.cpp:build\_dc\_conductance）：
正对角 Laplacian。对在运 DC 支路取（$r_{pu}$ 非有限或 $\le 10^{-12}$ 即抛
\texttt{std::runtime\_error}，fail-close，:128--132）
\[
g=\frac{1}{r_{pu}},\qquad G_{ii}\leftarrow G_{ii}+g,\qquad G_{ij}=G_{ji}=-g
\]
（:build\_dc\_conductance）。DC 节点功率方程为 $\mathbf p_{dc}=\mathbf v_{dc}\odot(G_{dc}\mathbf v_{dc})$
（src/power\_flow/jacobian\_builder.cpp:evaluate\_residual\_impl）。

\subsection{FDPF 电纳阵 $B'/B''$}
\srcpath{build\_susceptance\_matrix}（src/power\_flow/admittance\_builder.cpp:build\_susceptance\_matrix）取
$-\mathrm{Im}(Y^{\mathrm{mod}}_{ij})$，构型由 \texttt{BpBuildFlags} 控制
（\fld{zero\_resistance}/\fld{zero\_charging}/\fld{unit\_tap}/\fld{zero\_phase\_shift}/%
\fld{add\_bus\_shunts}/\fld{min\_x\_pu}，头文件
\srcpath{include/hacdcpf/assembly/ybus\_builder.hpp:BpBuildFlags}）：
\begin{itemize}
  \item $B'$（\srcpath{make\_bp\_flags}）：\textbf{FDXB} 方案——$R=0$、忽略充电、
        变比幅值取 1、相移取 0、不含母线并联；$|x|<$\fld{min\_x\_pu} 的支路跳过；
  \item $B''$（\srcpath{make\_bpp\_flags}）：\textbf{FDBX} 风格——保留 $R$、充电与变比
        （即 $-\mathrm{Im}(Y)$ 全量）、相移取 0，并加母线并联 $-B_s/S_{\mathrm{base}}$。
\end{itemize}

\subsection{DC 注入装配与残差独立复算}
\srcpath{assemble\_dc\_injections(data, vm, va, vdc, pdc\_spec)}
（\srcpath{src/power\_flow/pf\_injection\_assembly.cpp:assemble\_dc\_injections}），净注入约定发电为正：
DC 负荷为负（Load 表空时回退母线 \fld{pd\_mw}，:25--37）、DC 储能放电为正、
静态发电机 $p\cdot$\fld{scaling}、PV 阵列 \fld{p\_set}、VSC 经
\srcpath{converter\_dc\_injection}、DC/DC（入端 $-P_{in}$、出端 $+P_{out}$，
含 $I^2R$ 损耗，:79--88）、能量路由器 DC 端口。DC 参考选择
\srcpath{first\_dc\_slack\_or\_default}（\srcpath{include/hacdcpf/power\_flow/pf\_utils.hpp:first\_dc\_slack\_or\_default}）。

\srcpath{evaluate\_power\_flow\_residual}（\srcpath{src/power\_flow/residual\_evaluator.cpp:evaluate\_power\_flow\_residual}）
提供独立于 Newton 热路径的残差复算，返回 \{ac, dc, full\} 三块 $\ell_\infty$ 范数，
供诊断与测试交叉校验。交流注入采用\textbf{稀疏复数形式}
\srcpath{compute\_ac\_power\_injections}（:compute\_ac\_power\_injections）：$\mathbf V=\mathrm{polar}(V_m,\theta)$，
$\mathbf I=Y_{bus}\mathbf V$ 为一次稀疏矩阵-向量乘，$S_i=V_i\,\overline{I_i}$
（:23--32，调用点 :72），复杂度 $O(\mathrm{nnz}(Y_{bus}))$，不再把 $Y_{bus}$
稠密化后做 $O(n^2)$ 双重循环（数学上与逐元素 $g\cos\theta+b\sin\theta$ 求和等价）。
其母线分类用
\srcpath{classify\_ac\_buses}（\srcpath{include/hacdcpf/power\_flow/pf\_utils.hpp:classify\_ac\_buses}，额外 SLACK 归 PV），
与 Newton 上下文（额外 SLACK 整体剔除）略有差异。

\subsection{支路功率回算}
求解收敛后 \srcpath{compute\_branch\_flows}（\srcpath{src/power\_flow/branch\_flow.cpp:compute\_branch\_flows}）
按 π 等值 + 复变比逐支路回算两端功率（:compute\_branch\_flows）：
\begin{align*}
S_f &= V_f\,\overline{\big(y_{ff}V_f+y_{ft}V_t\big)}\cdot S_{\mathrm{base}}\\
S_t &= V_t\,\overline{\big(y_{tf}V_f+y_{tt}V_t\big)}\cdot S_{\mathrm{base}}
\end{align*}
输出 \texttt{BranchFlow\{pf\_mw, qf\_mvar, pt\_mw, qt\_mvar\}}（MW/Mvar），
按 \fld{data.ac\_branches} 顺序；停运/零阻抗支路保持 0。注意：本文件名为
"branch\_flow" 但\textbf{不是} Baran--Wu 支路潮流法；配网前推回代见
第~\ref{sec:distribution}~章。

\subsection{验证与校核}
装配层的验证以端到端方式覆盖：85 个 MATPOWER 算例的电压对照直接检验 $Y_{bus}$
装配（含分接头/移相/充电的正确性）；\srcpath{tools/transformer\_tap\_crosscheck.py}
对 12 组双节点变压器配置做 pandapower/MATPOWER/本平台三向对照；ZIP 聚合与
换流器提升逻辑由 \srcpath{tests/test\_all\_components\_pf.cpp}（29 类组件逐类禁用、
解变化量阈值断言）与 \srcpath{tests/test\_dc\_island\_reference.cpp} 覆盖；
投影幂等性、合并映射复合与结果恢复由投影蜕变测试与
\srcpath{tools/gui\_api\_e2e.py}（Bus 1 P/Q 闭合）保证。

% 第 3 章：统一 Newton-Raphson 潮流核心
% 本章文件由 \input 引入，不含导言区。
\section{统一 Newton--Raphson 潮流核心}\label{sec:newton}

\subsection{功能定位与入口}

\compmeta{engine::NewtonSolver}{\srcpath{include/hacdcpf/engine/solver/native/nle/newton\_solver.hpp}}

统一 Newton 求解器是全平台潮流的\textbf{生产主路径}：AC（$\theta,V_m$）、
DC（$V_{dc}$）与换流器耦合在同一方程组、同一雅可比中联立求解
（\textbf{不是} AC/DC 主从交替迭代）。入口签名
（include/hacdcpf/power\_flow/newton\_solver.hpp；实现 \srcpath{src/power\_flow/newton\_solver.cpp:solve}）：
\begin{quote}\footnotesize
\texttt{PowerFlowResult solve(const SolverData\& data, const PowerFlowOptions\& opt,\\
\hspace*{2em}const InitialState* init = nullptr) const;}
\end{quote}
可变成员 \texttt{PatternCache cache\_} 缓存雅可比上下文、稀疏模式与线性求解器实例
（见第~\ref{sec:overview}~章缓存体系）。每个求解器实例另复用一个
\texttt{SolverWorkspace}（include/hacdcpf/power\_flow/newton\_solver.hpp，接口
\srcpath{include/hacdcpf/power\_flow/workspace.hpp:SolverWorkspace}）：\texttt{vm}/\texttt{va}/\texttt{vdc}
经 \srcpath{prepare\_state} 按平启动重置、\texttt{mismatch}/\texttt{dx} 经
\srcpath{prepare\_equations} 按方程数重置，求解体内均为工作区成员引用
（src/power\_flow/newton\_solver.cpp:solve），而非每次求解在栈上新建；递归守卫
\texttt{workspace\_in\_use\_} 使同伦延拓的内层求解改用局部工作区，避免覆写外层状态
（:solve）。

\subsection{状态变量与节点类型}
未知量排序（\srcpath{build\_jacobian\_context}，src/power\_flow/newton\_solver.cpp:build\_jacobian\_context）：
\[
x=\big[\Delta\theta_{\,n_p};\ \Delta V_{m,\,n_q};\ \Delta V_{dc,\,n_{dc}}\big]^{\mathsf T},
\qquad n_{\mathrm{var}}=n_p+n_q+n_{dc}
\]
行/列映射（:build\_jacobian\_context）：非 slack 母线 $k$ 的有功行=角度列；PQ 母线的无功行=
$V_m$ 列（偏移 $n_p$）；非 slack DC 母线的 $P_{dc}$ 行=$V_{dc}$ 列（偏移 $n_p+n_q$）。
节点角色：
\begin{itemize}
  \item \textbf{Slack}：首个 SLACK 为主参考；\textbf{额外 SLACK 同样整体剔除}
        （$\theta$、$V_m$ 全固定），每个孤立 AC 区域各自保留参考（:build\_jacobian\_context）；
  \item \textbf{PV}：仅 P 方程、$\theta$ 未知，$V_m$ 固定；
  \item \textbf{PQ}：P、Q 方程与 $\theta$、$V_m$ 皆未知；
  \item \textbf{DC}：\texttt{DC\_V} 母线为 DC 参考（钉住）；\texttt{DC\_ISOLATED}
        剔除出方程组（:build\_jacobian\_context）；其余为 P 型 DC 母线。DC 参考由
        \srcpath{plan\_dc\_island\_references} 逐 DC 孤岛规划（:plan\_dc\_island\_references，
        详见第~\ref{sec:hybrid}~章）。
\end{itemize}
两种精确建模扩展（Direction 2/3）：\texttt{VDC\_VAC} 换流器交流母线在增强方程模式下
强制改 PQ、Q 行替换为 $V_m=V_{set}$（:build\_jacobian\_context）；半光滑 Newton 模式下
PV 全转 PQ 并记录 NCP 限值（:build\_jacobian\_context，见第~\ref{sec:robust}~章）。

\subsection{失配方程}
约定 $\mathrm{mismatch}=F^{\mathrm{spec}}-F^{\mathrm{calc}}$，解
$J\,\Delta x=\mathrm{mismatch}$，$J=\partial F^{\mathrm{calc}}/\partial x$
（注释见 \srcpath{src/power\_flow/jacobian\_builder.cpp:evaluate\_residual\_impl}）。
AC 节点注入（\srcpath{src/power\_flow/ac\_kernel.cpp:evaluate\_ac\_kernel\_serial}，
$\theta_{ij}=\theta_i-\theta_j$）：
\begin{align*}
P_i^{\mathrm{calc}}&=\sum_{j}V_iV_j\big(G_{ij}\cos\theta_{ij}+B_{ij}\sin\theta_{ij}\big)\\
Q_i^{\mathrm{calc}}&=\sum_{j}V_iV_j\big(G_{ij}\sin\theta_{ij}-B_{ij}\cos\theta_{ij}\big)
\end{align*}
指定注入含 ZIP 负荷（见第~\ref{sec:assembly}~章）；VSC 交流注入加入 spec。
DC 节点：$\mathbf p_{dc}^{\mathrm{calc}}=\mathbf v_{dc}\odot(G_{dc}\mathbf v_{dc})$，
失配行 $=p_{dc}^{\mathrm{spec}}-p_{dc}^{\mathrm{calc}}$（src/power\_flow/jacobian\_builder.cpp:evaluate\_residual\_impl）。
装配入口 \srcpath{evaluate\_residual\_and\_jacobian}（:evaluate\_residual\_and\_jacobian）与
\srcpath{evaluate\_residual\_only}（:evaluate\_residual\_only）共享 \srcpath{evaluate\_residual\_impl}（:evaluate\_residual\_impl），
返回 $\lVert\mathrm{mismatch}\rVert_\infty$（:evaluate\_residual\_impl）。

\subsection{雅可比结构}
$J=\begin{bmatrix}H&N\\ M&L\end{bmatrix}$ 为极坐标非归一化（MATPOWER）形式，
增量为 $\Delta V$ 本身（非 $\Delta V/V$）：
$H=\partial P/\partial\theta$、$N=\partial P/\partial V$、$M=\partial Q/\partial\theta$、
$L=\partial Q/\partial V$。

非对角元（$i\neq j$，AC kernel 累加，src/power\_flow/ac\_kernel.cpp:evaluate\_ac\_kernel\_serial）：
\begin{align*}
H_{ij}&\mathrel{+}=V_iV_j\big(G_{ij}\sin\theta_{ij}-B_{ij}\cos\theta_{ij}\big) &
M_{ij}&\mathrel{+}=-V_iV_j\big(G_{ij}\cos\theta_{ij}+B_{ij}\sin\theta_{ij}\big)\\
N_{ij}&\mathrel{+}=V_i\big(G_{ij}\cos\theta_{ij}+B_{ij}\sin\theta_{ij}\big) &
L_{ij}&\mathrel{+}=V_i\big(G_{ij}\sin\theta_{ij}-B_{ij}\cos\theta_{ij}\big)
\end{align*}
对角元（src/power\_flow/jacobian\_builder.cpp:evaluate\_residual\_impl），含 ZIP 负荷导数，
$\underline V=$ \fld{min\_vm\_pu}（默认 $10^{-8}$）：
\begin{align*}
H_{ii}&=-Q_i^{\mathrm{calc}}-B_{ii}V_i^2 &
M_{ii}&=P_i^{\mathrm{calc}}-G_{ii}V_i^2\\
N_{ii}&=\frac{P_i^{\mathrm{calc}}}{\max(|V_i|,\underline V)}+G_{ii}V_i+\big(p_dw_{p1}+2V_ip_dw_{p2}\big)\\
L_{ii}&=\frac{Q_i^{\mathrm{calc}}}{\max(|V_i|,\underline V)}-B_{ii}V_i+\big(q_dw_{q1}+2V_iq_dw_{q2}\big)
\end{align*}
DC 块：对角 $p_{dc,i}^{\mathrm{linear}}+g_{ii}v_{dc,i}$、非对角 $g_{ij}v_{dc,i}$
（:evaluate\_residual\_impl）。换流器导数分三档注入（详见第~\ref{sec:hybrid}~章）：DC 自洽导数
始终注入（:evaluate\_residual\_impl）；AC$\leftrightarrow$DC 交叉耦合块仅在
\fld{enable\_coupled\_jacobian} 时注入（:evaluate\_residual\_impl）；AC 构网换流器的能量导管
无条件生效（:evaluate\_residual\_impl）。

\subsection{稀疏模式与求值内核}
\begin{itemize}
  \item \srcpath{build\_jacobian\_pattern}（\srcpath{src/power\_flow/jacobian\_pattern.cpp:build\_jacobian\_pattern}）：
        从 $Y_{bus}$ 稀疏结构生成 $4\times nnz(Y)$ 个 AC 块槽位 + DC 块 + 耦合槽；
        以 \fld{entry\_to\_nz} 哈希把 (row,col)$\to$nz 下标写入
        \texttt{ACEntry\{p\_va\_nz, p\_vm\_nz, q\_va\_nz, q\_vm\_nz, g, b\}}（:build\_jacobian\_pattern）；
        模式固定后每次迭代仅 \texttt{valuePtr()} 清零再累加
        （src/power\_flow/jacobian\_builder.cpp:evaluate\_residual\_impl），不再分配；
  \item \srcpath{evaluate\_ac\_kernel\_serial}/\srcpath{\_parallel}
        （src/power\_flow/ac\_kernel.cpp:evaluate\_ac\_kernel\_serial）：遍历 \fld{ac\_entries} 同时累加
        $P^{\mathrm{calc}}$/$Q^{\mathrm{calc}}$ 与（可选）雅可比值；并行版每线程
        独立累加器再归约，条目数 $<1024$ 或单线程回退串行（:effective\_ac\_threads，
        \srcpath{util::ThreadPool::global()}），线程数由 \fld{ac\_eval\_threads} 控制；
  \item 一致性校验工具 \srcpath{check\_jacobian}（\srcpath{src/power\_flow/jacobian\_check.cpp:check\_jacobian}）：
        中心差分 $h=10^{-6}$、容差 $10^{-5}$。
\end{itemize}

\subsection{线性求解}
抽象 \srcpath{SparseLinearSolver}\{analyze\_pattern, factorize, solve, backend\_name\}
（\srcpath{include/hacdcpf/engine/kernel/linear\_algebra/linear\_solver.hpp}，转发 MIPSolvers）。
默认后端优先级 \textbf{KLU $>$ UMFPACK $>$ MKL Pardiso $>$ SuperLU $>$ Eigen SparseLU}
（\srcpath{MIPSolvers/src/engine/kernel/linear\_algebra/linear\_solver.cpp:make\_default\_sparse\_solver}；
CMake 选项 \srcpath{HACDCPF\_USE\_SUITESPARSE}，缺失回退 Eigen SparseLU）。
稀疏模式仅在缓存失效时 \texttt{analyzePattern} 一次；每次 Newton 迭代重新数值分解
（\srcpath{solve\_linear}，src/power\_flow/newton\_solver.cpp:solve\_linear（lambda）；计数
\fld{factorization\_calls}/\fld{linear\_solve\_calls}）。坏条件数
（代理 $\kappa=(\max/\min$ 行范数$)\times(\max/\min$ 列范数$)$，:estimate\_condition\_proxy）
可触发 NK-GMRES 回退（:solve\_linear\_nk（lambda），默认关闭，见第~\ref{sec:robust}~章）。

\subsection{更新、步长截断与收敛}
\begin{itemize}
  \item \textbf{步长截断} \srcpath{clip\_step}（:clip\_step）：
        $|\Delta\theta|\le 1.5$~rad、$|\Delta V_m|\le 0.5$~pu、$|\Delta V_{dc}|\le 0.5$~pu；
        $V_m,V_{dc}$ 下限 \texttt{kMinVm}$=0.05$（:kMinVm，:NewtonSolver::solve）；
  \item \textbf{收敛判据}：缩放后 $\lVert F\rVert_\infty<\texttt{opt.tol}$；
        \fld{tol}=$10^{-8}$、\fld{max\_iter}=50（\texttt{Defaults::kPFTol}/\texttt{kPFMaxIter}，
        \srcpath{include/hacdcpf/model/defaults.hpp:Defaults}）；NaN/Inf 立即判失败（src/power\_flow/newton\_solver.cpp:solve）；
  \item \textbf{诚实收敛}：最终以末态残差重新判定 \fld{out.converged}（:solve）；
        DC 电压越限仅告警（:solve）；
  \item \textbf{缩放}（默认全开）：$D_F=\mathrm{diag}(1/\max(|F_i^{\mathrm{spec}}|,1))$、
        $D_x=\mathrm{diag}(\pi,V_i,V_{dc,i})$，解 $\hat J\Delta\hat x=\hat F$，
        $\hat J=D_FJD_x^{-1}$（\srcpath{src/power\_flow/nonlinear\_scaling.cpp:build\_nonlinear\_scaling}，
        钳制 $[10^{-6},10^6]$；详见第~\ref{sec:robust}~章）。
\end{itemize}

\subsection{PV$\leftrightarrow$PQ 外环}
MATPOWER 风格无功限值切换（src/power\_flow/newton\_solver.cpp:solve）：外环上限 30 次；
内层收敛后检查隐含无功 $Q_g^{\mathrm{implied}}=Q^{\mathrm{calc}}+Q^{\mathrm{load}}-Q^{\mathrm{conv}}$
越限（迟滞 \fld{pv\_q\_hysteresis\_pu}=0.01），越限则转 PQ 并把 \fld{qg} 钉在限值；
全部限内后尝试恢复 PV（\fld{pv\_recover\_vm\_tol\_pu}=0.01）。切换\textbf{只在外环}
发起，迭代内不做切换：内层 Newton 收敛后统一检查越限（:solve）；内层未收敛时
仅当残差 $\le 0.1$ 才允许按当前 $Q$ 估计切换重试（:solve，守卫避免大残差下
按不可靠 $Q$ 值在 PV/PQ 间来回振荡）。
换流器 PQ$\leftrightarrow$VDC\_Q 模式切换在残差 $<10^{-6}$ 后按迟滞评估
（高/低阈值 0.03/0.01 pu、持续 2 次迭代，:apply\_converter\_mode\_switching，:NewtonSolver::solve；
逻辑见第~\ref{sec:hybrid}~章）。

\subsection{失败回退链与全局化}
默认全局化策略为线搜索（\fld{globalization}=LineSearch），回退链依次为
（src/power\_flow/newton\_solver.cpp:solve）：
\begin{enumerate}
  \item 非单调 Armijo-GLL 线搜索（\srcpath{run\_line\_search}，
        src/power\_flow/newton\_solver.cpp:run\_line\_search（lambda））：回溯 $\alpha\leftarrow\beta\alpha$、
        上限 \fld{max\_line\_search\_steps}=10 次试探（:run\_line\_search（lambda）），接受条件
        $\phi(x+\alpha\Delta x)\le\varphi_k^{\max}+c_1\,\alpha\,\nabla\phi(x)^{\mathsf T}\Delta x$
        （\srcpath{include/hacdcpf/power\_flow/nonmonotone\_linesearch.hpp}）；窗口 $M=5$、$c_1=10^{-4}$、
        $\beta=0.5$（\fld{nonmonotone\_window}/\fld{armijo\_c}/\fld{line\_search\_beta}，
        include/hacdcpf/power\_flow/robust\_nonlinear\_options.hpp），merit 每轮压入窗口
        （src/power\_flow/newton\_solver.cpp:solve）；\fld{enable\_nonmonotone\_linesearch}=false
        时以 $M=1$ 构造、退化为单调 Armijo（:run\_line\_search（lambda））；
  \item 对角正则化 $J+\lambda I$（$\lambda_0=10^{-6}$、增长 $\times10$、最多 8 步），
        每个正则化方向再经同一线搜索判定接受（:solve）；
  \item 自动恢复调度（\fld{enable\_auto\_fallback\_scheduling} 默认\textbf{开}，
        include/hacdcpf/power\_flow/robust\_nonlinear\_options.hpp）：LM 恢复
        $(J^{\mathsf T}J+\lambda I)\Delta x=-J^{\mathsf T}F$，最多 4 次尝试、
        $\lambda$ 自 \fld{lm\_lambda0} 起 $\times10$ 增至 \fld{lm\_lambda\_max}
        （src/power\_flow/newton\_solver.cpp:solve）；仍失败则以 $J+I/\delta t$ 做 PTC 恢复并
        无条件接受（:solve）；全部失败记 \fld{inner\_failed}（:solve）。
\end{enumerate}
另两种全局化策略：信赖域 dogleg（半径 1.0$\to$10）与伪瞬态延拓
（$J+I/\delta t$，SER 时间步由 \fld{ptc\_dt0} 族驱动，详见
第~\ref{sec:robust:lmptc}~节），
由 \fld{opt.globalization} 三选一。完整算法伪代码见技术笔记
\srcpath{docs/archive/reference/technical\_notebook/sections/03\_power\_flow.tex} 的 Algorithm PF-1；
五阶段鲁棒管线的组件细节见本手册第~\ref{sec:robust}~章。

\subsection{结构闭合检查与结果}
首次装配后扫描空行/空列（\srcpath{scan\_empty\_rows\_cols}，:scan\_empty\_rows\_cols），
做 $n_{eq}=n_{var}$ 闭合校核，结果入 \fld{diagnostics.equation\_closure\_*}。
输出 \texttt{PowerFlowResult\{vm, va, vdc, converged, iterations, residual,
profiling, diagnostics, branch\_flows, vsc\_transfers, ...\}}
（\srcpath{include/hacdcpf/power\_flow/power\_flow\_result.hpp:PowerFlowResult}），字段全集见第~\ref{sec:options}~章。
迭代计数 \fld{out.iterations} = 外层 + 内层 + 1（src/power\_flow/newton\_solver.cpp:solve）。

\subsection{验证与校核}
\begin{itemize}
  \item \textbf{MATPOWER 回归}：\srcpath{tests/test\_matpower\_cases.cpp} 三档
        85 算例；Tier1 九算例逐案断言（收敛 + $|V|\in[0.5,1.5]$、$|\delta|<60^\circ$、
        残差 $<10^{-4}$），Tier2 批收敛率 $\ge 85\%$，Tier3 大系统至约 8.2 万节点；
        \srcpath{tests/test\_powerflow\_performance.cpp} 用独立残差复算交叉验证；
  \item \textbf{雅可比对差分}：\srcpath{tests/test\_unified\_pf\_upgrade.cpp}
        解析雅可比 vs 中心差分 max\_rel\_error $<10^{-3}$（差分 $\varepsilon=10^{-6}$）；
  \item \textbf{缩放回归}：\srcpath{tests/test\_scaling\_regression.cpp}
        （case9/14/30/57/118 默认缩放开收敛、关缩放回退仍收敛）；
  \item \textbf{基准性能}：技术笔记 03 节 84 算例基准表（Apple M4 单线程）：
        NR 对 8.2 万节点 1.75~s/17 次迭代，实际标度约 $O(N^{1.2})$；
  \item 已知边界：PV$\leftrightarrow$PQ 为离散切换外环（半光滑 NCP 为可选替代）；
        换流器不等式约束仅事后校核（见第~\ref{sec:hybrid}~章诚实口径）。
\end{itemize}

% 第 4 章：经典交流求解器族
% 本章文件由 \input 引入，不含导言区。
\section{经典交流求解器族}\label{sec:classic}

本章覆盖四个经典交流求解器：FDPF 快速解耦、线性化直流潮流、分布式松弛与
自适应孤岛。前两者直接以 \texttt{SolverData} 为输入，复用第~\ref{sec:assembly}~章的
母线分类与 $B'/B''$ 电纳阵装配；后两者以 \texttt{HybridPowerSystem} 为输入，内部各自
重建 SolverData 并委托统一 Newton 核心（第~\ref{sec:newton}~章）完成数值求解，自身
只负责参与因子分摊与分岛调度。四者的公共入口在公共 API 族表（第~\ref{sec:overview}~章）
中列出。

\subsection{快速解耦潮流 FDPF（FDXB 变体）}

\compmeta{powerflow::FDPFSolver}{\srcpath{include/hacdcpf/power\_flow/solvers/fdpf\_solver.hpp}}

\subsubsection*{功能说明与入口}
经典 FDXB 方案快速解耦潮流，\textbf{纯 AC}：不读取 DC 网络与换流器，
\fld{vdc} 输出为空。入口签名（include/hacdcpf/power\_flow/fdpf\_solver.hpp；实现
\srcpath{src/power\_flow/fdpf\_solver.cpp:solve}）：
\begin{quote}\footnotesize
\texttt{PowerFlowResult solve(const SolverData\& data, const PowerFlowOptions\& opt) const;}
\end{quote}
\srcpath{include/hacdcpf/power\_flow/fdpf\_solver.hpp} 是 3 行转发头，规范头在
\texttt{solvers/} 子目录。公共 API \srcpath{solve\_power\_flow\_fdpf}
（\srcpath{src/api/hacdcpf.cpp:solve\_power\_flow\_fdpf}）经 SolverDataCache（见第~\ref{sec:overview}~章）
取得装配数据、施加 ZIP 权重后调用本求解器，再统一做支路潮流回算与结果回投影
（:solve\_power\_flow\_fdpf）。

\subsubsection*{数学模型与算法}
母线分类复用 \srcpath{classify\_ac\_buses}
（\srcpath{include/hacdcpf/power\_flow/pf\_utils.hpp:classify\_ac\_buses}）：首个 SLACK 为参考
（\srcpath{first\_slack\_or\_default}），\textbf{额外 SLACK 按 PV 处理}——
与统一 Newton"额外 SLACK 整体剔除"的约定不同（见第~\ref{sec:newton}~章）。
状态初值取母线记录 \fld{vm\_pu}/\fld{va\_deg}（src/power\_flow/fdpf\_solver.cpp:solve）；
PV 母线电压在 Q 步中不更新（Q 步仅作用 PQ 集合，:solve），即 PV 幅值钉在
母线记录初值，\textbf{不读取发电机 \fld{vg} 整定}。

常量阵 $B'/B''$ 由 \srcpath{build\_susceptance\_matrix} 装配
（实现 \srcpath{src/power\_flow/admittance\_builder.cpp:build\_susceptance\_matrix}，
取 $-\mathrm{Im}(Y^{\mathrm{mod}}_{ij})$），构型由 \texttt{BpBuildFlags} 控制
（\srcpath{include/hacdcpf/assembly/ybus\_builder.hpp:BpBuildFlags}，详见第~\ref{sec:assembly}~章）：
$B'$ 置 $R=0$、忽略充电、变比幅值取 1、相移取 0、不含母线并联（\srcpath{make\_bp\_flags}），
并跳过 $|x|<\texttt{min\_x\_pu}$ 的支路；$B''$ 保留 $R$、充电与变比，相移取 0，
加母线并联 $-B_s/S_{\mathrm{base}}$（\srcpath{make\_bpp\_flags}）。
两阵分别降阶至非 slack 全集与 PQ 集合并用 Eigen SparseLU \textbf{一次性分解}
（src/power\_flow/fdpf\_solver.cpp:solve），分解失败直接返回 \fld{converged}=false（:solve）。

每次迭代（:solve）先以\textbf{稀疏复数形式}计算注入
（\srcpath{compute\_power\_injections}）：$V_i=\mathrm{polar}(V_{m,i},\theta_i)$，
$\mathbf I=Y_{bus}\mathbf V$（稀疏矩阵--向量乘），$S_i=V_i\,\overline{I_i}$，
\textbf{不再生成稠密} $Y_{bus}$：
\begin{align*}
P_i^{\mathrm{calc}}&=\sum_{j}V_iV_j\big(G_{ij}\cos\theta_{ij}+B_{ij}\sin\theta_{ij}\big)\\
Q_i^{\mathrm{calc}}&=\sum_{j}V_iV_j\big(G_{ij}\sin\theta_{ij}-B_{ij}\cos\theta_{ij}\big)
\end{align*}
（与第~\ref{sec:newton}~章失配方程同式）。指定注入由
\srcpath{compute\_specified\_injections}（:compute\_specified\_injections）计算，含 ZIP 电压项
（权重顺序恒功率/恒电流/恒阻抗，见第~\ref{sec:assembly}~章）：\fld{has\_component\_loads}
为真时改用\textbf{逐母线} ZIP 系数（\fld{bus\_zip\_pp}/\fld{\_ip}/\fld{\_zp} 与
\fld{bus\_zip\_pq}/\fld{\_iq}/\fld{\_zq}，配 \fld{pd\_pu}/\fld{qd\_pu}，:compute\_specified\_injections；
由 \srcpath{aggregate\_load\_demand} 装配，\srcpath{src/power\_flow/solver\_data.cpp:aggregate\_load\_demand}），
否则沿用母线级 \fld{pd\_mw}/\fld{qd\_mvar} 加全局 \fld{zip\_pw}/\fld{zip\_qw}
（src/power\_flow/fdpf\_solver.cpp:compute\_specified\_injections）。如实说明：旧版 FDPF 只用全局 ZIP 权重，逐母线口径是本次的行为修正，
使 FDPF 与统一 Newton 的负荷模型对齐。
失配按 FDPF 惯例除以幅值：
\[
\Delta P_i/V_i=\big(P_i^{\mathrm{spec}}-P_i^{\mathrm{calc}}\big)/V_i,\qquad
\Delta Q_i/V_i=\big(Q_i^{\mathrm{spec}}-Q_i^{\mathrm{calc}}\big)/V_i
\]
（src/power\_flow/fdpf\_solver.cpp:solve）。收敛判据为
$\max\big(\lVert\Delta P/V\rVert_\infty,\lVert\Delta Q/V\rVert_\infty\big)<\texttt{tol}$
（:solve）。校正分两步单序执行（每迭代 P 步、Q 步各一次）：
\[
B'\,\Delta\theta=\Delta P/V\;(:solve),\qquad
B''\,\Delta V=\Delta Q/V\;(:solve)
\]
如实说明三处简化与缺陷：
\begin{itemize}
  \item \textbf{无全局化与限值逻辑}：无步长截断、无回退链（对照第~\ref{sec:newton}~章），
        无 PV$\leftrightarrow$PQ 无功限值切换，发电机 Q 限值完全不参与；
  \item \textbf{注入计算已稀疏化}（:compute\_power\_injections）：复杂度 $O(\mathrm{nnz}(Y_{bus}))$，
        旧版的 $Y_{bus}$ 转稠密 $+O(n^2)$ 双重循环已移除（主文档诚实口径相应条目
        随之过时，以本节为准）；
  \item 未收敛时复用同一稀疏注入内核与指定注入函数复算最终残差写入 \fld{residual}
        （跳过 slack 的 P 行与非 PQ 母线的 Q 行，:solve）。
\end{itemize}

\subsubsection*{求解参数}
\begin{paramtable}
\fld{tol} & $\varepsilon$ & pu & $10^{-6}\sim10^{-10}$ & $10^{-8}$ &
  收敛容差（$\Delta P/V$、$\Delta Q/V$ 的 $\ell_\infty$ 范数），
  \texttt{Defaults::kPFTol}（\srcpath{include/hacdcpf/model/defaults.hpp:Defaults}）\\
\fld{fdpf\_max\_iter} & $k_{\max}$ & 次 & 50--5000 & 1000 &
  最大迭代数（\srcpath{include/hacdcpf/power\_flow/power\_flow\_options.hpp:PowerFlowOptions}；
  读于 src/power\_flow/fdpf\_solver.cpp:solve）\\
\fld{robust\_nonlinear.min\_branch\_x\_pu} & $x_{\min}$ & pu & --- & $10^{-20}$ &
  $B'$ 装配的小电抗支路剔除阈值（真实头
  \srcpath{include/hacdcpf/power\_flow/globalization/robust\_nonlinear\_options.hpp:RobustNonlinearOptions}；
  传参于 src/power\_flow/fdpf\_solver.cpp:solve；注意它\textbf{覆盖} \srcpath{make\_bp\_flags}
  自带的 $10^{-6}$ 缺省）\\
\end{paramtable}
全局 ZIP 权重 \fld{zip\_pw}/\fld{zip\_qw} 与 \fld{loss\_model} 由 API 层在装配侧读取
（\srcpath{apply\_zip\_weights}，src/api/hacdcpf.cpp:get\_cached\_solver\_data，调用于 :solve\_power\_flow\_fdpf）；
求解器核在 \fld{has\_component\_loads} 为真时改读逐母线 ZIP 系数
（见上文算法小节），否则回退全局权重。

\subsubsection*{验证与校核}
\begin{itemize}
  \item \textbf{冒烟}：\srcpath{tests/test\_hacdcpf.cpp:TEST\_CASE} 两节点算例断言收敛；
  \item \textbf{稀疏注入与逐母线负荷回归}：\srcpath{tests/test\_helm\_solver.cpp}
        （注册于 \srcpath{tests/CMakeLists.txt}）——「FDPF sparse injection path」
        两节点算例与 Newton 逐母线 $|V|/\theta$ 对照，容差 $2\times10^{-4}$（:make\_two\_bus\_hybrid\_dc\_case）；
        「FDPF uses aggregated component loads」挂显式 \texttt{Load} 组件触发
        \fld{has\_component\_loads} 逐母线 ZIP 路径，与 Newton 对照容差
        $3\times10^{-4}$（:make\_two\_bus\_hybrid\_dc\_case）；
  \item \textbf{与 Newton 对照}：\srcpath{tests/test\_matpower\_cases.cpp:TEST\_CASE}
        case30 下两法均收敛时逐母线 $|V|$ 对照，容差 \texttt{WithinAbs}$=10^{-3}$；
        任一不收敛则 \texttt{SKIP}——\textbf{FDPF 无 Tier2/3 批量回归}，覆盖薄；
  \item 能力合同工具 \srcpath{tools/validate\_case\_capabilities.py} 做
        NR/FDPF/DC 三法互验（见第~\ref{sec:validation}~章）。
\end{itemize}

\subsection{线性化直流潮流（$B\theta=P$）}

\compmeta{powerflow::solve\_ac\_linearized\_dc}{\srcpath{include/hacdcpf/power\_flow/ac\_linearized\_pf.hpp}}

\subsubsection*{功能说明与入口}
经典 DC 潮流：$|V|\equiv1$、忽略电阻与无功，角度方程一次线性求解。入口为自由函数
（include/hacdcpf/power\_flow/ac\_linearized\_pf.hpp:solve\_ac\_linearized\_dc；实现 \srcpath{src/power\_flow/ac\_linearized\_pf.cpp:solve\_ac\_linearized\_dc}）：
\begin{quote}\footnotesize
\texttt{ACLinearizedDCResult solve\_ac\_linearized\_dc(const SolverData\& data);}
\end{quote}
结果结构 \texttt{ACLinearizedDCResult\{va（弧度）, pf\_mw, success\}}
（include/hacdcpf/power\_flow/ac\_linearized\_pf.hpp:ACLinearizedDCResult）。公共 API \srcpath{solve\_ac\_dc\_power\_flow}
（\srcpath{src/api/hacdcpf.cpp:solve\_ac\_dc\_power\_flow}）负责回投影相角与支路功率到原始编号空间
（:solve\_ac\_dc\_power\_flow）。如实说明：GUI 将该入口标为 \texttt{hybrid\_linearized}，但实现
\textbf{只解 AC 相角}，不求解 DC 电压、无功与网损。

\subsubsection*{数学模型与算法}
slack 取首个 SLACK 母线，缺省回退 bus 0（src/power\_flow/ac\_linearized\_pf.cpp:solve\_ac\_linearized\_dc）。对每条在运且
$|x_{pu}|\ge10^{-20}$ 的支路（:solve\_ac\_linearized\_dc）取
\[
b_k=\frac{1}{x_{pu}\cdot t_k}\qquad(t_k=0\text{ 按 }1.0\text{ 处理})
\]
（:solve\_ac\_linearized\_dc），装配 Laplacian $B_{ii}\mathrel{+}=b_k$、$B_{ij}=-b_k$
（:solve\_ac\_linearized\_dc）。移相变压器以注入等效：from 端 $+b_k(-\varphi_k)$、to 端
$-b_k(-\varphi_k)$（:solve\_ac\_linearized\_dc，$\varphi_k$ 为 \fld{shift\_deg} 化弧度）。代码先定义
$P^{\mathrm{shift}}$，再严格按注释“$B\theta+P^{\mathrm{shift}}=P^{\mathrm{spec}}$”
装配右端：
\[
P_i=P_{g,i}-P_{d,i}^{\mathrm{pu}}-P_i^{\mathrm{shift}}
\]
（:solve\_ac\_linearized\_dc；当 \fld{pd\_pu} 维数不匹配时才回退为
$P_{d,i}^{\mathrm{pu}}=\fld{pd\_mw}/\fld{base\_mva}$，恒功率口径，
\textbf{不含 ZIP 电压项}）。剔除 slack 行/列得降阶方程（:solve\_ac\_linearized\_dc），右端计入
slack 相角：
\[
B_{\mathrm{red}}\,\theta_{\mathrm{red}}
=\mathbf P_{\mathrm{red}}-\mathbf B_{\mathrm{red,slack}}\,\theta_{\mathrm{slack}}
\]
（:solve\_ac\_linearized\_dc），Eigen SparseLU 一次求解（:solve\_ac\_linearized\_dc）；分解或回代失败时
\fld{success}=false 静默返回。解出全状态后，代码只在非参考节点计算
$\max_i|(B\theta)_i-P_i|$ 并写入 \fld{residual\_pu}（:solve\_ac\_linearized\_dc），不检查参考节点
平衡。支路有功回算
\[
P_{f,k}=b_k\big(\theta_i-\theta_j-\varphi_k\big)\cdot S_{\mathrm{base}}
\]
（:solve\_ac\_linearized\_dc，MW）。近似边界如实声明：电压幅值、无功、网损均不输出；
$r_{pu}$、充电、并联全部忽略；多 slack 时仅首个被剔除，其余按普通母线计入
$B$（与 Newton/FDPF 的分类约定均不同）。

\subsubsection*{求解参数}
本求解器\textbf{不读取任何} \texttt{PowerFlowOptions} 字段（函数签名无 \texttt{opt}，
include/hacdcpf/power\_flow/ac\_linearized\_pf.hpp:solve\_ac\_linearized\_dc），全部阈值为硬编码常数：
\begin{paramtable}
（内部常数）$|x_{pu}|$ 下限 & --- & pu & --- & $10^{-20}$ &
  低于此值的支路视为零阻抗跳过（src/power\_flow/ac\_linearized\_pf.cpp:solve\_ac\_linearized\_dc）\\
（内部常数）$|t|$ 下限 & --- & pu & --- & $10^{-12}$ &
  低于此值变比按 1.0 处理（src/power\_flow/ac\_linearized\_pf.cpp:solve\_ac\_linearized\_dc）\\
\end{paramtable}

\subsubsection*{验证与校核}
\begin{itemize}
  \item \textbf{冒烟}：\srcpath{tests/test\_hacdcpf.cpp:TEST\_CASE} 两节点算例断言
        相角向量非空；
  \item \textbf{投影空间回归}：\srcpath{tests/test\_sppt\_metamorphic.cpp:TEST\_CASE}
        断言 \fld{va} 尺寸=原始母线数、\fld{pf\_mw} 尺寸=原始支路数、被合并支路
        功率保持 0；
  \item 覆盖薄弱如实说明：无对 MATPOWER \texttt{rundcpf} 的批量数值对照，
        本求解器的正确性主要靠 $B\theta=P$ 公式简单性与上述结构断言背书。
\end{itemize}

\subsection{分布式松弛}

\compmeta{powerflow::DistributedSlackSolver}{\srcpath{include/hacdcpf/power\_flow/solvers/distributed\_slack\_solver.hpp}}

\subsubsection*{功能说明与入口}
把 slack 吸收的不平衡功率按参与因子分摊到多台电源。两个入口
（\srcpath{include/hacdcpf/power\_flow/distributed\_slack\_solver.hpp}；实现
\srcpath{src/power\_flow/distributed\_slack\_solver.cpp:solve\_simplified}）：
\begin{quote}\footnotesize
\texttt{DistributedSlackResult solve\_simplified(const HybridPowerSystem\&, const DistributedSlack\&,\\
\hspace*{2em}const PowerFlowOptions\&) const;}\\
\texttt{DistributedSlackResult solve\_full\_jacobian(...同名参数...) const;}
\end{quote}
配套工厂 \srcpath{create\_participation\_factors}（:create\_participation\_factors）
按 \texttt{capacity}/\texttt{droop}/\texttt{equal} 三种方式生成配置结构
\texttt{DistributedSlack}（四字段定义见
\srcpath{include/hacdcpf/power\_flow/power\_flow\_options.hpp:DistributedSlack}，语义列于本节末）。
公共 API \srcpath{solve\_power\_flow\_distributed\_slack}(\srcpath{\_full})
（\srcpath{src/api/hacdcpf.cpp:solve\_power\_flow\_distributed\_slack}）。

\subsubsection*{数学模型与算法}
\textbf{参与因子}（src/power\_flow/distributed\_slack\_solver.cpp:create\_participation\_factors）：默认参与母线为有实际发电（$P_g>0$）的 SLACK/PV
在运母线（:create\_participation\_factors）。发电聚合与参与列表均使用 AC 母线稳定
\fld{index}，不把 ID 当作 vector 位置；显式列表过滤不在运、不存在和重复 ID。
空集抛 \texttt{invalid\_argument}（:create\_participation\_factors）；原始权重
\[
\alpha_i^{0}=\max(P_{g,i},10^{-4})\;(\text{capacity}),\qquad
\alpha_i^{0}=1/R_i\;(\text{droop},\ R_i>0\text{ 否则取容量}),\qquad
\alpha_i^{0}=1\;(\text{equal})
\]
（:create\_participation\_factors），统一归一化 $\sum_i\alpha_i=1$（:normalize\_factors, :create\_participation\_factors）；单机分摊上限
与该母线机组群的调节裕度挂钩（三种模式同一口径），
$\overline P_i=\max\big(|P_{\min,i}^{\Sigma}-P_{g,i}^{\mathrm{sched}}|,\,
|P_{\max,i}^{\Sigma}-P_{g,i}^{\mathrm{sched}}|\big)/S_{\mathrm{base}}$
（:create\_participation\_factors），写入 \fld{max\_participation\_p}（:create\_participation\_factors）。\srcpath{sanitize\_slack\_cfg}
（:sanitize\_slack\_cfg）：参与列表为空时自动按 capacity 生成；因子数不匹配时全置 1 再归一；
\fld{reference\_bus}=0 时取首个在运 SLACK 母线（无则抛
\texttt{invalid\_argument}，:sanitize\_slack\_cfg），并校验其确为在运 AC 母线（:sanitize\_slack\_cfg）。

\textbf{简化变体} \srcpath{solve\_simplified}（:solve\_simplified）：先做换流器协调预校核
（可阻断，:solve\_simplified），随后 \srcpath{make\_solver\_data} $+$ 统一 Newton 求基准解
（:solve\_simplified，电压经 \srcpath{unproject\_bus\_vector} 回投影，:solve\_simplified）；
基准解收敛后计算\textbf{全部 AC 母线}上的总不平衡
\[
\Delta P_{\Sigma}=\sum_{i\in\text{全部 AC 母线}}\big(P_i^{\mathrm{calc}}-P_i^{\mathrm{spec}}\big)
\]
（\srcpath{compute\_total\_participating\_mismatch\_pu}；遍历 Ybus 非零元，按
$P_i=\sum_{j:Y_{ij}\ne0}V_iV_j[G_{ij}\cos(\theta_i-\theta_j)+B_{ij}\sin(\theta_i-\theta_j)]$
复算，时间 $O(n+\mathrm{nnz}(Y))$、附加向量空间 $O(n)$，不构造稠密 Ybus；
列主序遍历保留各行按 $j$ 递增的累加顺序；
含换流器 AC 注入 \srcpath{converter\_ac\_injection}，:compute\_total\_participating\_mismatch\_pu）。求和取全部已解 AC
母线而非仅参与母线：非 slack 方程在数值容差内相消，剩余即全体 slack 母线提供的
有功；只取参与母线会在物理 slack 不在参与集合时静默返回零（函数内注释，
src/power\_flow/distributed\_slack\_solver.cpp:compute\_total\_participating\_mismatch\_pu）。按
$\Delta P_i=\alpha_i\,\Delta P_{\Sigma}$ 分摊输出（:solve\_simplified）。

AUD-092/096 的注册回归为 \srcpath{tests/test\_review\_regressions.cpp} 与
\srcpath{tests/test\_advanced\_pf.cpp}；可复现局部性能对照为
\srcpath{tools/review\_performance\_benchmark.py}，实测范围与误差见
\srcpath{docs/testing/module\_code\_audit.md}，不把 mismatch 耗时解释为整次 PF 耗时。

\textbf{限值变体} \srcpath{solve\_full\_jacobian}：每轮完整 Newton 求解后，用
\srcpath{distribute\_with\_limits} 迭代裁剪重分摊——越限母线钉在发电机
$P_{\min}/P_{\max}$ 与 \fld{max\_participation\_p} 的交集边界并计入
\fld{hit\_limits}，再把修正写回发电机并重解；剩余失配小于
$\max(10^{-10},\fld{tol})$ pu 时结束，最多 20 次完整求解。

诚实口径（重要）：
\begin{itemize}
  \item 简化变体只做一次 Newton，\fld{distributed\_slack\_p} 是事后分摊；限值变体
        会重调度并反复重解。兼容名称 \texttt{full\_jacobian} 不表示把参与因子作为
        单个增广 Jacobian 未知量；
  \item 失配复算复用统一 AC 注入装配，覆盖组件负荷、ZIP 电压项、VSC/LCC 与
        energy-router 注入；网络功率仍由 Ybus 致密复算，复杂度为 $O(n^2)$；
  \item \fld{reference\_bus} 会实际改写角度参考；输出
        \fld{distributed\_slack\_p} 的单位是系统 \fld{base\_mva} 基准上的 pu，
        不是 MW。
\end{itemize}

\subsubsection*{求解参数}
\begin{paramtable}
\fld{tol} & $\varepsilon$ & pu & $10^{-6}\sim10^{-10}$ & $10^{-8}$ &
  基准 Newton 收敛容差（转发，见第~\ref{sec:newton}~章）\\
\fld{max\_iter} & $k_{\max}$ & 次 & 20--300 & 50 &
  基准 Newton 迭代上限（转发；\texttt{Defaults::kPFMaxIter}，
  \srcpath{include/hacdcpf/model/defaults.hpp:Defaults}）\\
\fld{enable\_converter\_coordination\_check} & --- & 布尔 & --- & false &
  换流器协调预校核开关，阻断级问题直接返回失败
  （src/power\_flow/distributed\_slack\_solver.cpp:solve\_simplified）\\
\fld{loss\_model} & --- & 枚举 & --- & Linear &
  换流器损耗模型（\srcpath{make\_solver\_data} 装配与注入复算共用，:compute\_total\_participating\_mismatch\_pu, :solve\_simplified）\\
\end{paramtable}
\texttt{DistributedSlack} 配置字段（\fld{participating\_buses}、
\fld{participation\_factors}、\fld{reference\_bus}、\fld{max\_participation\_p}）
不属于 \texttt{PowerFlowOptions}，语义见上文算法小节。

\subsubsection*{验证与校核}
\srcpath{tests/test\_advanced\_pf.cpp}（标签 \texttt{[distributed\_slack]}）：
\begin{itemize}
  \item \textbf{因子构造}：capacity/equal 因子归一化误差 $\le10^{-12}$（:TEST\_CASE），
        droop 因子随下垂系数严格递减（:TEST\_CASE）；
  \item \textbf{求解}：3 节点两变体收敛且电压在 $(0.85,1.15)$（:TEST\_CASE）；
        限值 $10^{-12}$ 强制触发 \fld{hit\_limits}（:TEST\_CASE）；
  \item \textbf{MATPOWER 回归}：case9 capacity、case14 两变体一致收敛、case30 droop
        收敛（:TEST\_CASE）；集成用例 case14 下 slack 电压保持 $(0.95,1.10)$
        （:TEST\_CASE）。
\end{itemize}

\subsection{自适应孤岛}

\compmeta{powerflow::AdaptiveSolver}{\srcpath{include/hacdcpf/power\_flow/solvers/adaptive\_solver.hpp}}

\subsubsection*{功能说明与入口}
面向拓扑可能分裂的系统：先分岛，活岛各自独立求解，死岛电压置零。入口签名
（include/hacdcpf/power\_flow/adaptive\_solver.hpp；实现 \srcpath{src/power\_flow/adaptive\_solver.cpp:solve}）：
\begin{quote}\footnotesize
\texttt{AdaptiveSolveResult solve(const HybridPowerSystem\& sys, const PowerFlowOptions\& opt,\\
\hspace*{2em}const std::unordered\_map<int, ReactiveLimit>\& q\_limits = \{\}) const;}
\end{quote}
如实说明：\texttt{q\_limits} 形参已接线——\srcpath{apply\_reactive\_limit\_overrides}
（src/power\_flow/adaptive\_solver.cpp:apply\_reactive\_limit\_overrides）先校验
\fld{base\_mva} 为正、限值有限且 \fld{qmin}$\le$\fld{qmax}，再把限值按
\fld{base\_mva} 折算改写该母线首台在运发电机的 \fld{qmin\_mvar}/\fld{qmax\_mvar}
（同母线其余机组清零）；限值非法或引用的 AC 母线无在运发电机时抛
\texttt{std::runtime\_error}。岛描述结构 \texttt{IslandInfo}
（\srcpath{include/hacdcpf/model/ac\_components.hpp:IslandInfo}）。公共 API
\srcpath{solve\_power\_flow\_adaptive}（\srcpath{src/api/hacdcpf.cpp:solve\_power\_flow\_adaptive}）在收敛后
于全系统 SolverData 上补算支路潮流（:solve\_power\_flow\_adaptive）；
\srcpath{solve\_power\_flow\_islanded}（:solve\_power\_flow\_islanded）为薄包装。

\subsubsection*{数学模型与算法}
\textbf{分岛} \srcpath{detect\_islands}
（\srcpath{src/power\_flow/island\_detector.cpp:detect\_islands}）：AC/DC 母线合并为一张图，
停运与 \texttt{ISOLATED}/\texttt{DC\_ISOLATED} 母线剔除（:detect\_islands）；真实母线编号
映射到数组位置，编号\textbf{不要求连续}（:ac\_node（lambda））。边集涵盖 AC 支路、闭合开关与
断路器、2W/3W 变压器、DC 支路、闭合 DC 断路器、VSC、DCDC 与能量路由器端口链
（:detect\_islands），DFS 求连通分量（:detect\_islands）。活岛判据 \fld{has\_generators} 只看
\textbf{真实在运注入}（发电机 $P_g>10^{-9}$ 或 \fld{is\_slack}、换流器
$|P_{set}|$、静态/可再生/PV/储能/VPP、外部电网，:detect\_islands），明确不以 SLACK
母线类型作代理（注释 :detect\_islands）。如实说明：\textbf{每个连通分量都生成}
\texttt{IslandInfo}（含纯 DC 岛，:detect\_islands）；无源岛（\fld{has\_generators}=false）
不删除，而是在 \texttt{AdaptiveSolver::solve} 中按死岛处理、电压置零
（\srcpath{src/power\_flow/adaptive\_solver.cpp:solve}，
\srcpath{set\_dead\_island\_zero}）。

\textbf{求解流程}（src/power\_flow/adaptive\_solver.cpp:solve）：
\begin{enumerate}
  \item 初值取母线记录（:solve），调用 \srcpath{detect\_islands}（:solve）；
  \item \textbf{单岛快路径}：不做子系统抽取，直接全系统 Newton，失败则平启动重试
        （:solve）；
  \item \textbf{多岛}：死岛（\fld{has\_generators}=false）立即置零
        （:solve，\srcpath{set\_dead\_island\_zero}）；活岛缺 AC slack 且
        \fld{enable\_auto\_swing\_selection} 时自动选摇摆母线
        （\srcpath{auto\_select\_swing\_bus}：SLACK $\to$ 最大 $P_g$ $\to$
        PV $\to$ 首母线；覆盖在抽取时提升为 SLACK 且 $\theta=0$，
        src/power\_flow/island\_detector.cpp:extract\_island\_subsystem）；
  \item 每岛 \srcpath{extract\_island\_subsystem} 重编号抽取（src/power\_flow/island\_detector.cpp:extract\_island\_subsystem，
        含开关/变压器/换流器与全部组件表），独立 Newton $+$ 平启动重试
        （src/power\_flow/adaptive\_solver.cpp:solve）；
  \item \textbf{并行}：可解岛 $>1$ 时经 \srcpath{util::ThreadPool::global()} 并发，
        每岛强制 \fld{ac\_eval\_threads}=1 防嵌套并行死锁（:solve）；单可解岛保留
        雅可比级并行（:solve）；
  \item 结果按排序母线序映射回全局向量（:map\_island\_result\_back）；
        \fld{converged} 取各岛合取、\fld{iterations} 求和、\fld{residual} 取 max
        （:solve）。
\end{enumerate}

\subsubsection*{求解参数}
\begin{paramtable}
\fld{enable\_auto\_swing\_selection} & --- & 布尔 & --- & true &
  无 slack 活岛自动选摇摆母线（src/power\_flow/adaptive\_solver.cpp:solve）\\
\fld{ac\_eval\_threads} & $n_{thr}$ & 线程 & 0（自动） & 0 &
  单可解岛的雅可比并行线程数；多岛时内部强制为 1（:solve）\\
\fld{tol}、\fld{max\_iter} & --- & --- & --- & $10^{-8}$/50 &
  转发各岛 Newton（见第~\ref{sec:newton}~章）\\
\fld{loss\_model} & --- & 枚举 & --- & Linear &
  各岛 \srcpath{make\_solver\_data} 装配损耗模型（:merge\_island\_diagnostics, :merge\_island\_diagnostics）\\
\end{paramtable}

\subsubsection*{验证与校核}
\srcpath{tests/test\_advanced\_pf.cpp}（标签 \texttt{[island]}/\texttt{[adaptive]}）：
\begin{itemize}
  \item \textbf{分岛检测}：9 母线网络断 1/2 条支路得 2/3 岛（:build\_9bus\_island\_network）；
        非连续编号回归（12 母线三馈线，断言恰为 3 岛且自适应求解后全部
        $V>0.8$，:TEST\_CASE("Island detection: non-contiguous bus IDs (case\_2\_1\_backup) → 3 islands")）；case14/case118 单岛（:TEST\_CASE）；
  \item \textbf{孤岛求解}：2 岛下 DG 岛存活、死岛电压近零、活岛
        $V\in(0.7,1.2)$（:TEST\_CASE）；混合 AC/DC 微网"并网--孤岛--重并网"
        全生命周期（:TEST\_CASE）；网状微网选择性孤岛 4 岛与互济合并
        （:TEST\_CASE）；
  \item \textbf{自适应求解}：14 母线 AC/DC 系统死母线 $|V_m|\le10^{-12}$
        （:TEST\_CASE）；case9 收敛且残差 $<10^{-4}$（:TEST\_CASE）；
        两岛集成用例功率平衡（:TEST\_CASE）；另有两节点冒烟
        （\srcpath{tests/test\_hacdcpf.cpp:TEST\_CASE}）。
\end{itemize}

% 第 5 章：全局化与鲁棒求解器族
% 本章文件由 \input 引入，不含导言区。
\section{全局化与鲁棒求解器族}\label{sec:robust}

本章给出统一 Newton 核心（第~\ref{sec:newton}~章）所用全局化组件与鲁棒回退
求解器的实现细节：非线性缩放、非单调（GLL）线搜索、LM 信赖域与伪瞬态延拓、
NCP 半光滑 Newton、HELM 全纯嵌入、同伦延拓与 Newton--Krylov 内迭代，以及把
它们组织为五阶段管线的 \texttt{RobustNonlinearOptions}。第~\ref{sec:newton}~章
已概述缩放、线搜索与失败回退链的主流程，本章逐组件给出数学模型、参数与验证依据。

\subsection{非线性缩放（$D_F$、$D_x$ 对角缩放）}\label{sec:robust:scaling}

\compmeta{powerflow::NonlinearScaling}{\srcpath{include/hacdcpf/power\_flow/globalization/nonlinear\_scaling.hpp}}
\paragraph{功能说明} 对 Newton 方程做行/列对角平衡，使缩放后残差与步长均为
无量纲 $O(1)$ 量，收敛判据与病态探测都作用在缩放系统上。类只存两个向量并
施加逐元运算（\srcpath{include/hacdcpf/power\_flow/globalization/nonlinear\_scaling.hpp}），构造由自由函数
\srcpath{build\_nonlinear\_scaling} 完成（\srcpath{src/power\_flow/nonlinear\_scaling.cpp:build\_nonlinear\_scaling}）；
三个开关 \fld{enable\_residual\_scaling}/\fld{enable\_variable\_scaling}/
\fld{enable\_jacobian\_row\_col\_equilibration} 默认全开
（\srcpath{include/hacdcpf/power\_flow/globalization/robust\_nonlinear\_options.hpp}）。
\paragraph{数学模型与算法} 缩放系统（头文件契约注释
\srcpath{include/hacdcpf/power\_flow/globalization/nonlinear\_scaling.hpp}）：
\[
\hat F=D_F F,\qquad \hat J=D_F\,J\,D_x^{-1},\qquad
\Delta\hat x=D_x\,\Delta x,\qquad \Delta x=D_x^{-1}\Delta\hat x
\]
行缩放 $D_F$（src/power\_flow/nonlinear\_scaling.cpp:build\_nonlinear\_scaling）：P/Q/DC 三类方程行分别取
$D_{F,i}=1/\max(|F_i^{\mathrm{spec}}|,1)$，经 \srcpath{clamp\_inv\_scale}（:clamp\_inv\_scale）
钳制到 $[\,$\fld{min\_scale},\fld{max\_scale}$\,]$。列缩放 $D_x$（:build\_nonlinear\_scaling）：
$\theta$ 列取 $\pi$；$V_m$ 列取当前 $V_m$（$>10^{-3}$ 时，否则取 1，:build\_nonlinear\_scaling）；
$V_{dc}$ 列同规则（:build\_nonlinear\_scaling）。施加方式：残差逐元乘、雅可比遍历非零元行乘列除、
物理步长逐元除恢复（头 :apply\_residual\_scaling、:apply\_jacobian\_scaling、:unscale\_step）。主循环挂接
（\srcpath{src/power\_flow/newton\_solver.cpp:NewtonSolver::solve}）：每次迭代重建缩放；
\fld{enable\_residual\_scaling} 开时以缩放残差的 $\ell_\infty$ 范数判定收敛
（:NewtonSolver::solve）；仅当 equilibration 开\textbf{且}至少一个缩放开关开时才替换
线性系统（:NewtonSolver::solve）。条件数代理
$\kappa=(\max/\min$ 行$\ell_1$范数$)\times(\max/\min$ 列$\ell_1$范数$)$
（\srcpath{estimate\_condition\_proxy}，:estimate\_condition\_proxy）在
\fld{enable\_condition\_monitor} 或 NK 回退开启时每轮写入
\fld{profiling.condition\_estimate}（:NewtonSolver::solve）。
\paragraph{参数表}
\begin{paramtable}
\fld{enable\_residual\_scaling} & — & — & on/off & true & 行缩放 $D_F$ 开关；同时决定收敛判据用缩放残差\\
\fld{enable\_variable\_scaling} & — & — & on/off & true & 列缩放 $D_x$ 开关\\
\fld{enable\_jacobian\_row\_col\_equilibration} & — & — & on/off & true & 是否把 $\hat J=D_FJD_x^{-1}$ 送入分解\\
\fld{enable\_condition\_monitor} & — & — & on/off & true & 每轮计算条件数代理（NK 触发也依赖它）\\
\fld{min\_scale} & $s_{\min}$ & — & $10^{-8}$--$10^{-4}$ & $10^{-6}$ & 缩放因子下限（防除零）\\
\fld{max\_scale} & $s_{\max}$ & — & $10^{4}$--$10^{8}$ & $10^{6}$ & 缩放因子上限\\
\fld{min\_vm\_pu} & $\underline V$ & pu & $10^{-9}$--$10^{-6}$ & $10^{-8}$ & 雅可比 $\partial/\partial V$ 除法下限（\srcpath{src/power\_flow/jacobian\_builder.cpp:build\_power\_spec}）\\
\end{paramtable}
\paragraph{验证与校核} \srcpath{tests/test\_scaling\_regression.cpp}：
case9/14/30/57/118 默认缩放开全部收敛且 \fld{scaled\_residual\_norm}/
\fld{raw\_residual\_norm} 有限（:load\_case）；条件监控开启时
\fld{condition\_estimate} 为正有限（:TEST\_CASE）；三开关全关的 case14 回退仍收敛
（:TEST\_CASE）。\srcpath{equilibrate\_nonlinear\_scaling} 另提供四遍 Jacobian
行/列最大范数均衡实验 API；$\hat J=D_FJD_x^{-1}$ 的回代步与原物理 Newton 步在
构造的三阶病态矩阵上误差小于 $10^{-9}$。该实验均衡\textbf{未接入生产 Newton 主循环}：
合并试验虽保持单步代数等价，却改变浮点分解与全局化轨迹，完整回归曾引发
VSC/GFM/NCP/MATPOWER 收敛退化，故生产路径仍只消费
\srcpath{build\_nonlinear\_scaling} 的名义 $D_F,D_x$。这一区分防止把单步恒等式
误写成端到端收敛认证。

\subsection{非单调线搜索（GLL）}\label{sec:robust:gll}

\compmeta{powerflow::NonmonotoneLineSearch}{\srcpath{include/hacdcpf/power\_flow/globalization/nonmonotone\_linesearch.hpp}}
\paragraph{功能说明} Grippo--Lampariello--Lucidi（1986）非单调 Armijo 搜索，
header-only 实现（\srcpath{src/power\_flow/nonmonotone\_linesearch.cpp} 仅为
空壳翻译单元）。在 Newton 主循环中它维护最近 $M$ 轮 merit 历史、给出非单调
参考值，并对每个试探步长直接做 Armijo-GLL 接受判定
（\srcpath{NonmonotoneLineSearch::search}）。
\paragraph{数学模型与算法} Merit 函数 $\phi(x)=\tfrac12\lVert F(x)\rVert_2^2$，
滑动窗口参考 $\varphi_k^{\max}=\max_{0\le j\le M-1}\phi_{k-j}$，接受条件
（\srcpath{include/hacdcpf/power\_flow/globalization/nonmonotone\_linesearch.hpp:NonmonotoneLineSearch}）：
\[
\phi(x+\alpha\,\Delta x)\le \varphi_k^{\max}+c_1\,\alpha\,\nabla\phi(x)^{\mathsf T}\Delta x,
\qquad \alpha\leftarrow\beta\alpha
\]
$M=1$ 退化为标准单调 Armijo。主循环中的实际形态
（src/power\_flow/newton\_solver.cpp:NewtonSolver::solve 构造、\srcpath{run\_line\_search}（lambda））：
merit 每轮压入窗口（:NewtonSolver::solve）；回溯 $\alpha\leftarrow\beta\alpha$、
试探上限 \fld{max\_line\_search\_steps}$=10$ 次（:run\_line\_search（lambda）, :run\_line\_search（lambda））；
\fld{enable\_nonmonotone\_linesearch}=false 时改用 $M=1$ 的单调搜索实例
（:run\_line\_search（lambda））。如实说明：生产回溯次数由 \fld{max\_line\_search\_steps}
（默认 10）控制；折减系数即 \fld{line\_search\_beta}（默认 0.5，经构造函数传入，
:run\_line\_search（lambda））；类自带的 \fld{max\_line\_search\_trials}$=12$
（include/hacdcpf/power\_flow/globalization/robust\_nonlinear\_options.hpp）当前未被主循环消费，为保留字段。
\paragraph{参数表}
\begin{paramtable}
\fld{enable\_nonmonotone\_linesearch} & — & — & on/off & true & GLL 参考窗口与非单调接受开关（关则退化为 $M=1$ 单调 Armijo）\\
\fld{nonmonotone\_window} & $M$ & — & 3--10 & 5 & 滑动窗口长度（\srcpath{include/hacdcpf/power\_flow/nonmonotone\_linesearch.hpp}）\\
\fld{armijo\_c} & $c_1$ & — & $10^{-5}$--$10^{-3}$ & $10^{-4}$ & Armijo 充分下降常数\\
\fld{line\_search\_beta} & $\beta$ & — & 0.3--0.7 & 0.5 & 回溯折减系数（主循环经构造函数消费，src/power\_flow/newton\_solver.cpp:NewtonSolver::solve）\\
\fld{max\_line\_search\_trials} & — & 次 & 8--20 & 12 & 类内试探上限（当前未被主循环消费，保留字段）\\
\end{paramtable}
\paragraph{验证与校核} 无独立单元测试；由
\srcpath{tests/test\_scaling\_regression.cpp} 与
\srcpath{tests/test\_unified\_pf\_upgrade.cpp} 的全局化对比用例间接覆盖
（默认配置下 GLL 开启，全部基线用例收敛）。

\subsection{LM 信赖域与伪瞬态延拓（PTC）}\label{sec:robust:lmptc}

\compmeta{powerflow::LMController / PtcSerController}{\srcpath{include/hacdcpf/power\_flow/globalization/lm\_trust\_region.hpp}}
\paragraph{功能说明} 两个 header-only 自适应控制器
（\srcpath{src/power\_flow/lm\_trust\_region.cpp} 为空壳翻译单元）：
\texttt{LMController} 按增益比 $\rho$ 调节 LM 阻尼 $\lambda$
（$\rho\ge0.75\Rightarrow\lambda/3$，$0.1\le\rho<0.75$ 保持，
$\rho<0.1$ 或 pred$\le0\Rightarrow\lambda\times10$，
\srcpath{include/hacdcpf/power\_flow/lm\_trust\_region.hpp}）；\texttt{PtcSerController} 按 SER 公式
$\delta t\leftarrow\mathrm{clamp}(\delta t(r_{k-1}/r_k)^\gamma,
\delta t_{\min},\delta t_{\max})$ 调节伪时间步（\srcpath{include/hacdcpf/power\_flow/globalization/lm\_trust\_region.hpp:PtcSerController}）。
如实说明：\texttt{LMController} 当前\textbf{未被} Newton 主循环调用，属预留接口；
\texttt{PtcSerController} 已接线——PTC 全局化分支以 \fld{ptc\_dt0}/\fld{ptc\_dt\_min}/
\fld{ptc\_dt\_max}/\fld{ptc\_gamma} 构造该控制器，用其 \texttt{dt()} 装配
$J+I/\delta t$ 并以 \texttt{update()} 逐步更新伪时间步
（src/power\_flow/newton\_solver.cpp:NewtonSolver::solve）。
\paragraph{数学模型与算法}
\begin{itemize}
  \item \textbf{信赖域 dogleg}（\fld{globalization}=TrustRegion，
        src/power\_flow/newton\_solver.cpp:NewtonSolver::solve）：Cauchy 步长
        $\alpha_c=\lVert g\rVert^2/\lVert Jg\rVert^2$（$g=J^{\mathsf T}F$），
        dogleg 路径在 Cauchy 点与 Newton 步间插值并截断到半径 $\Delta$
        （\srcpath{include/hacdcpf/engine/kernel/globalization/globalization.hpp}）；
        增益比 $\rho=\mathrm{ared}/\mathrm{pred}$：$\rho>0.75$ 且触界时
        $\Delta\leftarrow\min(2\Delta,\Delta_{\max})$，$\rho<0.25$ 时
        $\Delta\leftarrow\max(0.25\Delta,10^{-10})$，$\rho>0$ 才接受
        （同文件 :trust\_region\_update；调用 :NewtonSolver::solve）。
  \item \textbf{PTC}（\fld{globalization}=PseudoTransient，
        src/power\_flow/newton\_solver.cpp:NewtonSolver::solve）：
        求解 $(J+I/\delta t)\Delta x=-F$，步\textbf{无条件接受}；
        SER 更新 $\delta t\leftarrow\delta t\cdot(r_{k-1}/r_k)^\gamma$，
        残差 $\le10^{-30}$ 时改乘 $2$（\texttt{PtcSerController::update}
        内硬编码），并钳位到 $[\delta t_{\min},\delta t_{\max}]$
        （\srcpath{include/hacdcpf/power\_flow/globalization/lm\_trust\_region.hpp:PtcSerController}）。
        初值 $\delta t_0$、指数 $\gamma$ 与上下限分别取
        \fld{ptc\_dt0}$=0.1$、\fld{ptc\_gamma}$=0.7$、
        \fld{ptc\_dt\_min}$=10^{-4}$、\fld{ptc\_dt\_max}$=10^{6}$
        （\srcpath{include/hacdcpf/power\_flow/globalization/robust\_nonlinear\_options.hpp:RobustNonlinearOptions}）。
  \item \textbf{LM/PTC 自动恢复}（\fld{enable\_auto\_fallback\_scheduling}，
        默认\textbf{开}，\srcpath{include/hacdcpf/power\_flow/globalization/robust\_nonlinear\_options.hpp}；
        src/power\_flow/newton\_solver.cpp:NewtonSolver::solve）：线搜索与正则化全部失败后，LM 恢复解
        $(J^{\mathsf T}J+\lambda I)\Delta x=-J^{\mathsf T}F$（:NewtonSolver::solve），
        最多 4 次尝试、$\lambda$ 自 \fld{lm\_lambda0} 起 $\times10$ 增至
        \fld{lm\_lambda\_max}（:reduce\_ncp\_mu（lambda）, :reduce\_ncp\_mu（lambda））；仍失败则以 $J+I/\delta t$ 做
        PTC 恢复并无条件接受（:NewtonSolver::solve）。
\end{itemize}
\paragraph{参数表}
\begin{paramtable}
\fld{enable\_lm\_trust\_region\_fallback} & — & — & on/off & true & 允许 LM 恢复步（仅在自动调度开启时生效）\\
\fld{lm\_lambda0}/\fld{lm\_lambda\_max} & $\lambda$ & — & — & $10^{-4}$/$10^{8}$ & LM 初始阻尼与上限（:NewtonSolver::solve 使用）\\
\fld{enable\_ptc\_ser} & — & — & on/off & true & 允许 PTC 恢复步（仅在自动调度开启时生效）\\
\fld{ptc\_dt0} 等四字段 & $\delta t$ 族 & — & — & 0.1 / $10^{-4}$ / $10^{6}$ / 0.7 & SER 控制器步长参数，已接线：PTC 分支以其构造 \texttt{PtcSerController}（\fld{ptc\_dt\_min}/\fld{\_max}/\fld{ptc\_gamma}；src/power\_flow/newton\_solver.cpp:NewtonSolver::solve）\\
\fld{trust\_region\_radius0} & $\Delta_0$ & — & 0.1--2 & 1.0 & 信赖域初始半径（\texttt{PowerFlowOptions}，:PowerFlowOptions）\\
\fld{trust\_region\_max} & $\Delta_{\max}$ & — & 5--50 & 10.0 & 信赖域半径上限\\
\fld{enable\_auto\_fallback\_scheduling} & — & — & on/off & true & 失败时自动追加 LM/PTC 恢复（src/power\_flow/newton\_solver.cpp:NewtonSolver::solve）\\
\end{paramtable}
\paragraph{验证与校核} \srcpath{tests/test\_unified\_pf\_upgrade.cpp}
「Globalization improves flat-start convergence」（:TEST\_CASE）：IEEE 118 混合
算例、$\pm5\%$ 初值噪声、30 次试验，断言 LineSearch 30/30 收敛、TrustRegion
与 PTC 各 $\ge50\%$；「Full comparison」（:TEST\_CASE）在 5 个压力算例
$\times$ 8 种配置（含 TR、PTC）打印迭代对照表。自动恢复调度（默认开）无
专门回归用例。

\subsection{结构保持 NCP 半光滑 Newton}\label{sec:robust:ncp}

\compmeta{powerflow::pv\_pq\_ncp、fischer\_burmeister 与 evaluate\_vsc\_limit\_ncp}{\srcpath{include/hacdcpf/power\_flow/ncp\_functions.hpp}；\srcpath{src/power\_flow/vsc\_limit\_ncp.cpp}}
\paragraph{功能说明} 生产 Newton 路径包含两类固定维度非光滑块：传统发电机
PV/PQ 无功限值使用中值 NCP；VSC 电流圆、P/Q priority、Vdc droop 饱和和
GFM Norton 端口使用每台设备六状态局部块。活动状态只改变数值，不改变行列
坐标，因此符号分解可跨限值激活复用。PV/PQ 块由
\fld{enable\_semi\_smooth\_newton} 开启；VSC 块由设备字段
\fld{enable\_limit\_ncp} 与正 \fld{i\_ac\_max\_pu} 独立开启。

\paragraph{精确目标方程} 标量互补关系采用严格 Fischer--Burmeister 映射
\[
\phi(a,b)=\sqrt{a^2+b^2}-a-b,
\qquad \phi(a,b)=0\Longleftrightarrow a\ge0,\ b\ge0,\ ab=0.
\]
实现不加入隐藏 $\varepsilon$，故双活跃交点满足 $\phi(0,0)=0$。原点处选择
Clarke 广义雅可比的对称元素
$\partial_a\phi=\partial_b\phi=1/\sqrt{2}-1$。PV/PQ 行使用
\[
F_Q=\operatorname{mid}(V_m-V_{set},\ Q_g-Q_{max},\ Q_g-Q_{min}),
\]
限内恢复 PV 电压控制，触及上下界时恢复 PQ 限值行为。Magnitude-priority
VSC 使用 $0\le\lambda\perp g(I,V)\ge0$；P-first/Q-first 使用嵌套中值投影。
网络行保持 $J=d(\mathrm{calc}-\mathrm{spec})/dx$，局部块存储 $dF/dx$，统一
线性系统右端分别为 $\mathrm{spec}-\mathrm{calc}$ 与 $-F$。

\paragraph{可选光滑延拓} \fld{enable\_smooth\_ncp} 启用后，magnitude 块采用
Kanzow 平滑半径
\[
\phi_\mu(a,b)=\sqrt{a^2+b^2+2\mu^2}-a-b,
\]
中值块采用 CHKS 光滑 min/max。VSC continuation 不要求同时打开 PV/PQ
半光滑开关；两类 NCP 共存时共享同一 $\mu$。$\mu$ 从
\fld{ncp\_mu0} 按粗/细两阶段因子降至 \fld{ncp\_mu\_min}，每次缩减后在
原状态重新装配残差与 Jacobian。公开 VSC 结果使用 $\mu=0$ 重新计算证书，
不把光滑代理根冒充精确互补解。

\paragraph{局部 VSC 块 Schur 消元与 BD-regular 条件}
将网络变量记为 $x$、每台 VSC 的六维局部变量合并为 $z$，则所选广义
Jacobian 元素具有
\[
J=\begin{bmatrix}A&B\\C&D\end{bmatrix},\qquad
D=\operatorname{blkdiag}(D_1,\ldots,D_m),\quad D_i\in\mathbb R^{6\times6}.
\]
生产实现 \srcpath{powerflow::VSCLocalSchurSolver} 不形成显式逆，而以每个
$D_i$ 的一次稠密分解同时求 $D_i^{-1}[C_i,r_i]$，构造
\[
S=A-BD^{-1}C,\qquad \widehat r_x=r_x-BD^{-1}r_z,
\]
解 $S\Delta x=\widehat r_x$ 后回代
$\Delta z=D^{-1}(r_z-C\Delta x)$。每块只连接本机端口
$V_a/V_m/V_{dc}$ 与 P/Q/DC 三行，故 Schur 更新最多九个固定坐标。局部阶数
$s=6$、端口宽度 $c\le3$ 时，附加工作为
$O(ms^3+ms^2(c+1))=O(m)$；总维度由 $n_x+6m$ 降至 $n_x$。完整固定模式
每块含 $36+18+3=57$ 个局部相关结构项，reduced 模式最多新增 9 项，因此
结构非零至少减少 $48m$。稀疏 LU 的 fill 仍取决于图与排序，本文不以母线数
直接套用未经测量的 $O(n^{1.5})$。

若全部 $D_i$ 与 $S$ 非奇异，则块 Gaussian 消元保证当前所选 $J$ 非奇异。
这只是可执行的单元素证书；BD-regular 要求解邻域内 B-subdifferential 的全部
相关元素一致可逆。在该条件与半光滑性成立时，Qi--Sun 半光滑 Newton 局部
超线性；强半光滑时得到标准二次局部率。程序报告
$\lVert F_{k+1}\rVert/\lVert F_k\rVert$ 与
$\lVert F_{k+1}\rVert/\lVert F_k\rVert^2$，但不把有限样本冒充 BD-regular
证明。

数值准入使用统一 \srcpath{NumericalConstants::kSqrtMachineEpsilon} 派生的
局部 reciprocal-condition 阈值，并对 reduced 与完整重构步执行 Higham
范数后向误差检查。任一检查失败即回退完整 sparse LU。生产默认仅在网络维度
不少于 \fld{Defaults::kVSCSchurMinNetworkDimension}$=1000$ 时启用；小系统可将
阈值置零做论文消融。两次 warmup、五次交替 A/B 的 Release/KLU 协议显示，
强制 case300 Schur 虽将维度 $640\to604$、结构非零 $4820\to4502$，但固定开销
使线性/总时长增加 $40.9\%/15.7\%$，故默认拒绝；ACTIVSg2000 将
$4054\to4006$、$29806\to29382$，线性/总时长降低 $14.0\%/12.5\%$，进入
生产准入。准入维度、局部 rcond 与后向误差门槛均可经 C++ 及生产 HTTP
配置并回显；机器 epsilon 保持只读。KLU 适配器当前不返回 factor nnz/work，
结果明确记为 unavailable。

\paragraph{参数表}
\begin{paramtable}
\fld{enable\_semi\_smooth\_newton} & — & — & on/off & false & 发电机 PV/PQ 中值 NCP 开关\\
\fld{enable\_smooth\_ncp} & — & — & on/off & false & PV/PQ 与 VSC 的 $\mu$ 光滑 FB/CHKS 延拓开关\\
\fld{ncp\_mu0} & $\mu_0$ & — & $10^{-3}$--$10^{-1}$ & $10^{-2}$ & 初始光滑参数\\
\fld{ncp\_mu\_min} & $\mu_{\min}$ & — & $10^{-14}$--$10^{-10}$ & $10^{-12}$ & 退火下限（趋近精确 NCP）\\
\fld{ncp\_mu\_factor(\_coarse)} & — & — & — & 0.1/0.5 & 细/粗两阶段退火因子；\fld{ncp\_mu\_phase\_transition}$=0.1$ 为切换残差比\\
\fld{enable\_activity\_hysteresis} 等 & — & — & — & true / 3 / $10^{-4}$ & 迟滞三字段（预留，未消费）\\
\fld{enable\_vsc\_local\_schur} & — & — & on/off & true & VSC 六维局部块精确 Schur 消元；证书失败回退 full LU\\
\fld{vsc\_schur\_min\_network\_dimension} & $n_{\min}$ & 维 & $\ge0$ & 1000 & 生产复杂度准入；0 仅用于小系统消融\\
\fld{vsc\_schur\_local\_rcond\_tolerance} & $\rho_D$ & — & $>0$ & $\sqrt{\epsilon_{mach}}$ & 当前所选 $D_i$ 的最小 reciprocal-condition 门槛\\
\fld{vsc\_schur\_backward\_error\_tolerance} & $\eta_b$ & — & $>0$ & $10^{-12}$ & reduced 与完整重构步最大范数后向误差\\
\end{paramtable}
\paragraph{验证与校核} \srcpath{tests/test\_vsc\_limit\_ncp.cpp} 覆盖精确 FB
原点、平滑导数有限差分、三种 power-port priority、三种限流 GFM priority、
多换流器同时激活、退化交点、弱网、不可行点以及 case300/ACTIVSg2000。
光滑算例均保持一次 Jacobian pattern build，$\mu$ 降至 $10^{-12}$ 后按精确
方程复核，证书残差不超过 $10^{-9}$。PV/PQ CHKS 数学恒等式与有限差分由
\srcpath{tests/test\_power\_flow\_math\_audit.cpp} 守卫。Schur/full-LU 小型代数
oracle 的步相对误差不超过 $10^{-10}$、完整后向误差不超过 $10^{-12}$；
case300/ACTIVSg2000 的 $V_m/V_a/V_{dc}$ 根差不超过 $10^{-8}$，并锁定
$6m$ 维度与至少 $48m$ 结构非零缩减。可复现实验入口为
\srcpath{tools/vsc\_schur\_benchmark.cpp}。

\subsection{HELM 全纯嵌入}\label{sec:robust:helm}

\compmeta{powerflow::HelmSolver}{\srcpath{include/hacdcpf/power\_flow/solvers/helm\_solver.hpp}}
\paragraph{功能说明} Holomorphic Embedding Load-flow Method：把潮流方程嵌入
复参数 $s$ 的全纯函数，在 $s=0$（平启动，精确解已知）展开幂级数，用 Padé
逼近解析延拓到 $s=1$。\textbf{不依赖初值}，Padé 发散可作为算例不可解的判证。
代码移植自 HELMpy 参考实现（AGPLv3，\srcpath{src/power\_flow/helm\_solver.cpp}），
增强：稀疏 $Y_{\mathrm{trans}}$、多 slack、K 因子分布式松弛与
DCPF/Gauss--Seidel 暖启动（)。纯 AC（只用 AC 子系统，头文件
:HelmSolver::solve），编译进 \texttt{hacdcpf} 库（\srcpath{CMakeLists.txt}）。
\paragraph{数学模型与算法}
\begin{itemize}
  \item \textbf{导纳分解}（\srcpath{build\_ytrans\_yshunt}，:build\_ytrans\_yshunt）：
        串联支路按 $\pi$ 模型生成 $Y_{\mathrm{trans}}$（$y_{ff}=y_s/|\tau|^2$ 等
        四元组，:build\_ytrans\_yshunt），对地支路 $jb/2$ 与母线 $G_s+jB_s$ 入
        $Y_{\mathrm{shunt}}$（:build\_ytrans\_yshunt）；
  \item \textbf{修正导纳阵} $Y_{\mathrm{mod}}\in\mathbb R^{2N\times2N}$
        （\srcpath{build\_ymod}，:build\_ymod）：slack 置单位行；PQ 行写
        $\Re/\Im$ 2$\times$2 块；PV 用 Model 2——P 行取实部、第 $2i{+}1$ 行写
        $|V|^2$ 约束（:build\_ymod）；
  \item \textbf{系数递推}（:HelmSolver::solve）：每阶解一次
        $Y_{\mathrm{mod}}\,c^{[n]}=\mathrm{rhs}$（共享一次 LU 分解，:HelmSolver::solve）。
        PQ 右端 $\overline{S_i}\,\overline{W_i^{[n-1]}}-Y_{\mathrm{sh},i}V_i^{[n-1]}$
        （:HelmSolver::solve）；PV 的 $|V|^2$ 行右端 $V\!V/2$（:HelmSolver::solve）、P 平衡行经
        $CC$ 缓存的卷积递推（:HelmSolver::solve）；$W=1/V$ 逆元递推
        $W^{[n]}=-\sum_{k<n}W^{[k]}V^{[n-k]}$（:HelmSolver::solve）；
  \item \textbf{Padé 求值与外环}（\srcpath{pade\_at\_one}，:pade\_at\_one）：$[L/L]$
        矩阵法，解 Toeplitz 方程得分母系数、卷积得分子系数，$s=1$ 处取
        $\sum a/\sum b$；每 2 阶做一次收敛检查（幅值/相角变化
        $<$\fld{mismatch}，:HelmSolver::solve），终值提取 :HelmSolver::solve；无功限值外环至多
        10 次 PV$\to$PQ 切换、越限 $10^{-6}$ 判据、Q 注入钉限
        （:verify\_physical\_residual（lambda）, :verify\_physical\_residual（lambda））；暖启动 DCPF 角度基线（:dcpf\_warmstart）/GS 预迭代
        （:gauss\_seidel\_warmstart）；K 因子在 PV+SLACK 间分摊不平衡（:apply\_distributed\_slack）。
\end{itemize}
\paragraph{参数表（\texttt{HelmOptions}，include/hacdcpf/power\_flow/solvers/helm\_solver.hpp）}
\begin{paramtable}
\fld{max\_coef} & $n_{\max}$ & 阶 & 50--200 & 100 & 幂级数最大系数数\\
\fld{mismatch} & $\varepsilon$ & pu/rad & $10^{-8}$--$10^{-4}$ & $10^{-6}$ & Padé 收敛容差（幅值与相角）\\
\fld{enforce\_q\_limits} & — & — & on/off & true & 无功限值 PV$\to$PQ 外环开关\\
\fld{warm\_start}/\fld{gauss\_iter} & — & — & None / DCPF / GS & None/5 & 暖启动策略与 GS 预迭代次数（）\\
\fld{allow\_multi\_slack} & — & — & on/off & true & 多 slack 母线（均置单位行）\\
\fld{participation\_factors} & $K$ & — & $\ge0$ & 空 & 分布式松弛参与因子（非空即启用）\\
\fld{sparse\_threshold} & — & 母线 & — & 200 & 稀疏阈值（当前实现恒为稀疏，该字段未消费）；\fld{pade\_trace\_out} 为收敛轨迹输出指针\\
\end{paramtable}
\paragraph{验证与校核} 如实说明：HELM 当前\textbf{无生产 API 自动挂接}
（\srcpath{solve\_power\_flow} 不分发该求解器，GUI 亦无对应路由），只能由
C++ 直接构造 \texttt{HelmSolver} 调用。回归覆盖：\srcpath{tests/test\_helm\_solver.cpp}
（注册于 \srcpath{tests/CMakeLists.txt}）「HELM matches Newton on a
well-conditioned AC case」——两节点良性算例上 HELM（\fld{max\_coef}$=120$、
关无功限值外环）与 Newton 逐母线 $|V|/\theta$ 对照，容差 $2\times10^{-5}$
（:make\_two\_bus\_hybrid\_dc\_case）。仍未覆盖：PV 无功限值外环、多 slack 与 K 因子分摊、暖启动路径、
不可解算例的 Padé 发散判证，覆盖仍薄；数值正确性同时依赖与 HELMpy 参考实现的
对口移植。

\subsection{同伦延拓}\label{sec:robust:homotopy}

\compmeta{powerflow::HomotopyContinuationSolver}{\srcpath{include/hacdcpf/power\_flow/continuation\_solver.hpp}}
\paragraph{功能说明} 把所有负荷与发电/换流器功率设定按同伦参数
$\lambda\in[0,1]$ 从平启动零注入基态缓步爬升到目标值，每步以上一步解暖启动，
作为难收敛算例的兜底路径（设计定位：鲁棒管线 Phase 4）。实现
\srcpath{src/power\_flow/homotopy\_continuation.cpp:HomotopyContinuationSolver::solve}；状态记录
\texttt{HomotopyState\{lambda, step, accepted\_steps, rejected\_steps, failed\}}
（\srcpath{include/hacdcpf/power\_flow/continuation\_solver.hpp:HomotopyState}）。
\paragraph{数学模型与算法} 嵌入构造
（\srcpath{scale\_system\_by\_lambda}，:scale\_system\_by\_lambda）：AC/DC 母线集总负荷、显式
负荷、发电机/储能/可再生/静态发电机有功无功、VSC $P_{set}$、DC/DC $P_{ref}$、
移动储能等全部注入 $\times\lambda$；电压设定与网络参数不变。路径跟踪
（:HomotopyContinuationSolver::solve）：子求解关闭同伦递归（:HomotopyContinuationSolver::solve），中间步迭代预算
$\max(10,\lfloor\,$\fld{max\_iter}$/2\rfloor)$（:HomotopyContinuationSolver::solve）；$\lambda=0$ 基态
必先可解（:HomotopyContinuationSolver::solve）；主循环 $\lambda_{\mathrm{next}}=\min(1,\lambda+
\Delta\lambda)$，以 \fld{initial\_state} 暖启动调 \srcpath{solve\_hybrid}
（:HomotopyContinuationSolver::solve）；成功则 $\Delta\lambda\leftarrow\min(2\Delta\lambda,
\Delta\lambda_{\max})$，失败则 $\Delta\lambda\leftarrow\Delta\lambda/2$
（:HomotopyContinuationSolver::solve）；步数超 \fld{homotopy\_max\_steps} 或步长低于
\fld{homotopy\_step\_min} 判失败（:HomotopyContinuationSolver::solve）。
\paragraph{参数表}
\begin{paramtable}
\fld{enable\_homotopy} & — & — & on/off & true & 同伦开关（当前仅用于子求解防递归，见下）\\
\fld{homotopy\_step0} & $\Delta\lambda_0$ & — & 0.05--0.5 & 0.2 & 初始步长\\
\fld{homotopy\_step\_min}/\fld{\_max} & — & — & — & $10^{-3}$/1.0 & 最小/最大步长（低于最小步长判失败）\\
\fld{homotopy\_max\_steps} & — & 步 & 20--100 & 30 & 接受+拒绝总步数上限\\
\end{paramtable}
\paragraph{验证与校核} 如实说明：同伦求解器\textbf{未被 Newton 失败路径自动
链入}——\fld{enable\_homotopy} 仅在 \srcpath{src/power\_flow/homotopy\_continuation.cpp:HomotopyContinuationSolver::solve}
用于防止子求解递归，GUI 仅透传该选项
（\srcpath{tests/run\_gui\_server.cpp:apply\_pf\_robust\_options}）；调用方需显式构造
\texttt{HomotopyContinuationSolver} 使用。回归覆盖：
\srcpath{tests/test\_homotopy\_continuation.cpp}（注册于
\srcpath{tests/CMakeLists.txt}）单用例——两节点电压稳定算例从 $\lambda=0$
爬升，断言收敛、\fld{state.failed}=false、终态 $\lambda=1.0$、
\fld{accepted\_steps}$\ge3$ 且负荷母线 $|V|<1.0$（:TEST\_CASE）；另
\srcpath{tests/test\_newton\_krylov.cpp:TEST\_CASE} 经
\srcpath{PowerFlowSolverFactory} 对含 \texttt{Continuation} 的六种方法做
可构造、可求解的接线冒烟。仍未覆盖：步长回退/失败路径
（$\Delta\lambda$ 折半、\fld{homotopy\_max\_steps} 判负）、混合 AC/DC 算例，
覆盖仍薄。

\subsection{Newton--Krylov（NK-GMRES 回退）}\label{sec:robust:nk}

\compmeta{powerflow::newton\_krylov\_step}{\srcpath{include/hacdcpf/power\_flow/solvers/newton\_krylov.hpp}}
\paragraph{功能说明} 以重启 GMRES($m$)（MGS Arnoldi + Givens 旋转）替代稀疏
LU 求解 Newton 线性系统，并配 Schur 补 AC/DC 块预条件器，面向病态雅可比
场景（\srcpath{src/power\_flow/newton\_krylov.cpp:newton\_krylov\_step}）。\textbf{默认关闭}
（\fld{enable\_newton\_krylov\_fallback}=false），仅在条件数代理越限时可选回退。
\paragraph{数学模型与算法}
\begin{itemize}
  \item \textbf{GMRES($m$)}（\srcpath{gmres\_solve}，:gmres\_solve）：左预条件
        重启型；MGS 正交化（:gmres\_solve）、Givens 旋转消 Hessenberg 次对角
        （:gmres\_solve）、估计残差 $|\,g_{j+1}\,|/\lVert b\rVert$（:gmres\_solve）；每轮
        重启后以真残差复判（:gmres\_solve）；
  \item \textbf{Schur 块预条件}（\srcpath{SchurBlockPreconditioner}，
        :SchurBlockPreconditioner）：利用 $J=\big[\begin{smallmatrix}A&B\\0&D\end{smallmatrix}\big]$
        块上三角结构（AC 块 $A$、DC 块 $D$ 各作 SparseLU 分解，:SchurBlockPreconditioner::build），
        $P^{-1}r$：$z_{dc}=D^{-1}r_{dc}$，$z_{ac}=A^{-1}(r_{ac}-Bz_{dc})$
        （:apply）；任一块分解失败则退化为单位预条件（:newton\_krylov\_step）；
  \item \textbf{接受准则}（:newton\_krylov\_step）：GMRES 收敛，或最终相对残差
        $<100\times$\fld{gmres\_tol} 且步有限、$\max|\Delta x|<10^{6}$；
  \item \textbf{触发与回退}（newton\_solver.cpp）：条件代理
        $>$\fld{nk\_condition\_trigger} 且开关开启时本轮用 NK（:NewtonSolver::solve）；
        三种全局化分支均可受益（:run\_line\_search（lambda）, :run\_line\_search（lambda）, :run\_line\_search（lambda）），
        默认分支 NK 失败或步不被接受时回落标准 LU（:NewtonSolver::solve）；状态写入
        \fld{profiling.linear\_solver\_status}（\texttt{nk\_gmres}/
        \texttt{nk\_gmres\_schur}/\texttt{nk\_gmres\_failed}，:solve\_linear\_nk（lambda））。
\end{itemize}
\paragraph{参数表}
\begin{paramtable}
\fld{enable\_newton\_krylov\_fallback} & — & — & on/off & false & NK-GMRES 回退总开关（默认关）\\
\fld{enable\_schur\_preconditioner} & — & — & on/off & false & Schur 补块预条件开关\\
\fld{nk\_condition\_trigger} & $\kappa_{\mathrm{trig}}$ & — & $10^{6}$--$10^{10}$ & $10^{8}$ & 触发阈值（对条件数代理）\\
\fld{gmres\_restart} & $m$ & — & 20--50 & 30 & Krylov 子空间重启维度\\
\fld{gmres\_max\_outer} & — & 轮 & 5--20 & 10 & 最大重启轮数（总 matvec 上界 $m\times$轮数）\\
\fld{gmres\_tol} & $\varepsilon$ & — & $10^{-12}$--$10^{-8}$ & $10^{-10}$ & 相对残差容差 $\lVert r\rVert/\lVert b\rVert$\\
\end{paramtable}
\paragraph{验证与校核} 如实说明：生产路径默认不经过 NK（默认关闭 + 触发
阈值 $10^{8}$），属覆盖较薄组件。回归覆盖：
\srcpath{tests/test\_newton\_krylov.cpp}（注册于
\srcpath{tests/CMakeLists.txt}）——(i) GMRES$+$Schur 预条件求解
$3\times3$ 耦合 AC/DC 块系统，断言 \fld{used\_schur} 且解误差
$<10^{-10}$（:TEST\_CASE）；(ii) \srcpath{PowerFlowSolverFactory} 对含
\texttt{NewtonKrylov} 的六种方法逐一构造并求解两节点算例均收敛
（:TEST\_CASE）；(iii) \texttt{SolverWorkspace} 布局保持的重置语义（:TEST\_CASE）。
仍未覆盖：Newton 主循环内经条件数代理触发的真实病态算例路径、GMRES
重启与失败回退路径。

\subsection{\texttt{RobustNonlinearOptions} 与五阶段鲁棒管线}\label{sec:robust:pipeline}

\compmeta{powerflow::RobustNonlinearOptions}{\srcpath{include/hacdcpf/power\_flow/globalization/robust\_nonlinear\_options.hpp}}
\paragraph{功能说明与接线状态} \texttt{RobustNonlinearOptions}
（\srcpath{include/hacdcpf/power\_flow/globalization/robust\_nonlinear\_options.hpp}，经
\fld{PowerFlowOptions::robust\_nonlinear} 挂载）按 Phase 1--5 组织全部鲁棒
组件开关与参数；\texttt{NonlinearSolveMode} 枚举与
\texttt{SolverModeSwitchRecord}（)为策略切换的诊断记录类型。
各 Phase 的\textbf{实际接线状态}（如实）：
\begin{enumerate}
  \item \textbf{Phase 1 缩放与诊断}：已接线，默认全开（\ref{sec:robust:scaling} 节）；
  \item \textbf{Phase 2 非单调线搜索与光滑 NCP}：GLL 已接线默认开
        （\ref{sec:robust:gll} 节）；光滑 NCP 已接线但默认关（\ref{sec:robust:ncp} 节）；
  \item \textbf{Phase 3 LM/PTC}：信赖域 dogleg 与 PTC 作为 \fld{globalization}
        三策略之二已接线；自动恢复调度默认关（\ref{sec:robust:lmptc} 节）；
  \item \textbf{Phase 4 同伦延拓}：求解器存在但\textbf{未自动链入}，需显式调用
        （\ref{sec:robust:homotopy} 节）；
  \item \textbf{Phase 5 高级项}：NK-GMRES 已接线默认关（\ref{sec:robust:nk} 节）；
        \fld{enable\_log\_voltage}/\fld{auto\_enable\_log\_voltage\_on\_failure}
        为保留字段，当前无实现消费（)。
\end{enumerate}
字段全集（含本表未列的诊断阈值 \fld{bad\_condition\_threshold}、
\fld{tiny\_pivot\_threshold} 等当前未消费字段）见第~\ref{sec:options}~章完整字段表。
\paragraph{验证与校核总览} 本章组件的回归矩阵汇总：缩放与条件监控——
\srcpath{tests/test\_scaling\_regression.cpp}（6 用例）；NCP 半光滑与信赖域/
PTC 全局化——\srcpath{tests/test\_unified\_pf\_upgrade.cpp}（四个用例：「Semi-smooth Newton
suppresses outer PV/PQ loops on IEEE 118」、「Globalization improves flat-start
convergence」、「MATPOWER cases: semi-smooth Newton vs heuristic PV/PQ」、
「Full comparison: all enhancements across stress cases」）；GLL 随默认配置被上述用例间接覆盖；HELM、同伦延拓、Newton--Krylov
现各有专项冒烟回归——\srcpath{tests/test\_helm\_solver.cpp}（HELM 对 Newton
两节点对照）、\srcpath{tests/test\_homotopy\_continuation.cpp}（端点到达）、
\srcpath{tests/test\_newton\_krylov.cpp}（GMRES/Schur 精度与求解器工厂接线），
注册于 \srcpath{tests/CMakeLists.txt}。如实说明：三者仍
\textbf{无生产 API 自动挂接}，回归均为单算例/单路径冒烟，失败与回退路径
未覆盖，覆盖仍薄（与第~\ref{sec:validation}~章现状声明一致）；CPF 的相关
诚实口径见连续潮流章。

% 第 6 章：直流网络与混合交直流潮流
% 本章文件由 \input 引入，不含导言区。
\section{直流网络与混合交直流潮流}\label{sec:hybrid}

本章给出 DC 网络方程、纯 DC 求解器、混合 AC/DC 联立求解的耦合机制、
VSC 与 DC/DC 换流器稳态模型、DC 岛电压参考规划，以及换流器协调校核体系。
统一 Newton 主循环（状态变量排序、全局化、收敛判定）见第~\ref{sec:newton}~章，
本章只展开 DC 侧与换流器相关的部分。设计背景文档为
\srcpath{docs/archive/theory/multiple\_converter.md}（多变换器统一模型）与
\srcpath{docs/archive/theory/multiple\_converter\_r1.md}（AC 侧构网与七类控制模式补充）。

\subsection{功能定位与入口}

混合 AC/DC 潮流的公共入口族（薄包装，均委托统一主链，见第~\ref{sec:overview}~章）：

\begin{center}
\footnotesize
\begin{tabular}{@{}L{4.2cm}L{4.6cm}L{6.2cm}@{}}
\toprule
入口 & 实现 & 说明\\
\midrule
\srcpath{powerflow::solve\_dc} & \srcpath{src/power\_flow/dc.cpp:7--10} & 委托 \srcpath{solve\_dc\_power\_flow}（\srcpath{src/api/hacdcpf.cpp:1022}）\\
\srcpath{powerflow::solve\_hybrid} & \srcpath{src/power\_flow/hybrid.cpp:7--10} & 委托 \srcpath{solve\_power\_flow}（统一 Newton）\\
\srcpath{powerflow::solve\_hybrid\_adaptive} & \srcpath{src/power\_flow/hybrid.cpp:12--15} & 委托分岛自适应求解\\
\bottomrule
\end{tabular}
\end{center}

DC 母线类型枚举 \texttt{DCBusType}（\srcpath{include/hacdcpf/model/enums/bus\_types.hpp:14--18}）：
\texttt{DC\_P}=1（功率母线）、\texttt{DC\_V}=3（电压参考母线）、
\texttt{DC\_ISOLATED}=4（失电母线，剔除出方程组）。
\textbf{如实说明}：\srcpath{include/hacdcpf/power\_flow/dc.hpp:9} 的注释称
"linear approximation"，但实现（见下节）是 $V\odot GV=P$ 的\textbf{非线性} Newton，
与线性化角度直流潮流（$B\theta=P$，纯 AC）不是一回事。

\subsection{DC 网络方程与纯 DC 求解器}

\compmeta{powerflow::DCSolver}{\srcpath{include/hacdcpf/power\_flow/solvers/dc\_solver.hpp}}

\subsubsection*{功能说明}
\texttt{DCSolver} 是独立纯 DC 求解器（实现 \srcpath{src/power\_flow/dc\_solver.cpp:24--241}）：
只解 DC 电压向量 $\mathbf v_{dc}$，AC 侧仅以其初始 $\fld{vm\_pu}$ 参与换流器注入求值
（\texttt{va} 取 0，dc\_solver.cpp:145--153），不回代交流网络。输出
\texttt{DCPowerFlowResult\{vdc, converged, iterations, residual\}}
（\srcpath{include/hacdcpf/power\_flow/power\_flow\_result.hpp:228--233}）。

\subsubsection*{数学模型与算法}
DC 节点功率方程（注入发电为正）与失配（dc\_solver.cpp:160--170）：
\[
p_{dc,i}^{\mathrm{calc}}=v_{dc,i}\sum_{j}g_{ij}v_{dc,j},\qquad
F_i=p_{dc,i}^{\mathrm{spec}}-p_{dc,i}^{\mathrm{calc}}
\]
其中 $G_{dc}=[g_{ij}]$ 为第~\ref{sec:assembly}~章的正对角 Laplacian 电导阵
（$g=1/r_{pu}$）；\fld{pdc\_spec} 由 \srcpath{assemble\_dc\_injections}
（\srcpath{src/power\_flow/pf\_injection\_assembly.cpp:12--101}）装配：DC 负荷为负
（Load 表空时回退母线 \fld{pd\_mw}，:20--37）、DC 储能放电为正（:39--46）、
静态发电机与 PV 阵列（:48--64）、VSC 经 \srcpath{converter\_dc\_injection}（:66--74）、
DC/DC 入端 $-P_{in}$、出端 $+P_{out}$（:76--88）、能量路由器 DC 端口（:90--100）。

求解流程（Newton + 回溯线搜索，稀疏代数）：
\begin{enumerate}
  \item \textbf{参考与方程集}：\srcpath{first\_dc\_slack\_or\_default}
        （\srcpath{include/hacdcpf/power\_flow/pf\_utils.hpp:34--46}）优先首个
        \texttt{DC\_V} 母线，否则首个非孤立母线；\texttt{DC\_ISOLATED} 母线钉住初值、
        剔除出方程组（dc\_solver.cpp:113--129，避免零行/零列奇异）；方程数为 0 时
        直接按 \fld{vm\_pu} 返回（:131--140）；
  \item \textbf{雅可比}（\srcpath{build\_dc\_newton\_jacobian}，:24--93）：
        $J_{kk}=p_{dc,k}^{\mathrm{linear}}+g_{kk}v_{dc,k}$（$p^{\mathrm{linear}}=G_{dc}\mathbf v_{dc}$），
        $J_{kl}=g_{kl}v_{dc,k}\ (l\neq k)$；
  \item \textbf{线性求解}：稀疏 \texttt{Eigen::SparseLU}（\texttt{\#include <Eigen/SparseLU>}，
        :9），分解不可逆或解含 NaN/Inf 即判失败
        （:181--195）；
  \item \textbf{回溯线搜索}（:197--227）：$\alpha$ 初值 1、最多 8 次减半，试验步满足
        $v_{dc}\ge\texttt{kMinVdc}=0.05$（:20）；有改善试探则采纳最优试探，
        无改善时取全步（:220--227）；
  \item \textbf{收敛}：$\lVert F\rVert_\infty<\fld{opt.tol}$；结束后以末态残差重新判定
        \fld{converged}（:230--234，与统一 Newton 的诚实收敛口径一致）。
\end{enumerate}

\subsubsection*{参数表}
\begin{paramtable}
\fld{DCBus::bus\_type} & --- & --- & DC\_P / DC\_V / DC\_ISOLATED & DC\_P & 母线角色，决定参考与剔除（\srcpath{include/hacdcpf/model/dc\_components.hpp:22}）\\
\fld{DCBus::vm\_pu} & $v_{dc,0}$ & pu & $0.9$--$1.1$ & 1.0 & DC 电压初值/钉住值（dc\_components.hpp:23）\\
\fld{DCBus::vmax\_pu} / \fld{vmin\_pu} & --- & pu & --- & 1.1 / 0.9 & 仅事后越限告警（newton\_solver.cpp:1839--1857）\\
\fld{DCBranch::r\_pu} & $r$ & pu & $10^{-4}$--$10^{-1}$ & 0.0 & 支路电阻，$g=1/r$（dc\_components.hpp:52）\\
\fld{opt.tol} / \fld{opt.max\_iter} & $\varepsilon$ & pu/次 & --- & $10^{-8}$ / 50 & 收敛容差/迭代上限（同第~\ref{sec:newton}~章默认值）\\
\texttt{kMinVdc} & --- & pu & --- & 0.05 & 线搜索电压下限（dc\_solver.cpp:20）\\
\end{paramtable}

\subsubsection*{验证与校核}
纯 DC 路径由 \srcpath{tests/test\_all\_components\_pf.cpp} 的 DC 组件族
（DCBus、DCBranch、DCLoad、StaticGeneratorDC、PVArrayDC，文件头注释 :15--16）与
\srcpath{tools/validate\_case\_capabilities.py} 的 NR/FDPF/DC 三法互验覆盖；
DC 岛参考行为见第~\ref{sec:hybrid-dcref}~节。

\subsection{混合 AC/DC 联立求解耦合机制}

\subsubsection*{功能说明}
统一 Newton 把 AC 与 DC 未知量放在\textbf{同一}方程组、同一稀疏雅可比中联立求解
（非 AC/DC 主从交替）。状态排序 $x=[\Delta\theta;\Delta V_m;\Delta V_{dc}]$ 与
DC 行列映射（偏移 $n_p+n_q$）见第~\ref{sec:newton}~章；本小节展开换流器注入
进 spec 的方式与换流器导数的\textbf{三档注入}机制。

\subsubsection*{换流器注入与三档雅可比注入}
残差与雅可比统一装配于 \srcpath{evaluate\_residual\_impl}（:181--367，
\srcpath{src/power\_flow/jacobian\_builder.cpp}）。换流器对指定注入的贡献在
\srcpath{build\_power\_spec}（同文件 :25--178）：AC 侧
$P^{\mathrm{spec}} \mathrel{+}=p_{ac}$、$Q^{\mathrm{spec}}\mathrel{+}=q_{ac}$
（:143--147）；DC 侧 \fld{pdc\_spec} 加 \srcpath{converter\_dc\_injection}（:172--176）。
DC 网络块与其雅可比（:285--293）：
\[
p_{dc}^{\mathrm{calc}}=\mathbf v_{dc}\odot(G_{dc}\mathbf v_{dc}),\qquad
J_{ii}=p_{dc,i}^{\mathrm{linear}}+g_{ii}v_{dc,i},\qquad
J_{ij}=g_{ij}v_{dc,i}\ (i\neq j)
\]
DC 失配行 $=p_{dc}^{\mathrm{spec}}-p_{dc}^{\mathrm{calc}}$（:364--367）。

换流器导数分\textbf{三档}注入（符号约定 $J\mathrel{+}=-(\partial\mathrm{spec}/\partial x)$，
jacobian\_builder.cpp:327）：
\begin{enumerate}
  \item \textbf{AC 构网能量导管（无条件）}：\texttt{AC\_GRID\_FORMING} 换流器的
        AC 母线是 slack，其 AC 有功为自由平衡量；DC 注入按能量守恒跟踪
        $P_{dc}=-\big(P_{ac}+\mathrm{loss}(P_{ac})\big)$，
        $P_{ac}=p^{\mathrm{calc}}-p^{\mathrm{spec}}$（:150--172，
        可选叠加 $r\cdot P_{ac}^2/V_m^2$ 交流传导损耗）；
  \item \textbf{DC 自洽导数（始终注入）}：VSC 的
        $\partial p_{dc}^{\mathrm{spec}}/\partial v_{dc}$ 与 DC/DC droop/Voltage 模式的
        DC 块内导数（:322--353），保证耦合关闭时 DC 方程仍有精确导数；
  \item \textbf{AC$\leftrightarrow$DC 交叉块（仅 \fld{enable\_coupled\_jacobian} 开启）}：
        AC P/Q 行 $\times$ $V_{dc}$ 列注入 $-\partial p_{ac}/\partial v_{dc}$、
        $-\partial q_{ac}/\partial v_{dc}$（:296--303）；DC 行 $\times$ AC $V_m$ 列
        （仅 \fld{r\_conv\_ac\_pu}$>0$）注入 $-\partial p_{dc}/\partial V_m$（:304--314）。
\end{enumerate}
稀疏模式侧：DC 对角槽 \fld{dc\_vdc\_diag\_nz} 与 DC/DC droop 槽
\fld{dcdc\_droop\_entries} \textbf{始终}构建
（\srcpath{src/power\_flow/jacobian\_pattern.cpp:204--242}）；交叉耦合槽
\fld{coupling\_entries} 与 \fld{dcdc\_coupling\_entries}（后者仅 \fld{r\_eq\_pu}$>0$）
仅在耦合开启时构建（:312--363）。DC/DC 导数注入：输出行$\times$输入电压列
$-\partial p_{out}/\partial v_{dc,in}$（jacobian\_builder.cpp:339--343），
droop/Voltage 项 $-\partial p_{out}/\partial v_{dc,out}$ 与
$+\partial p_{in}/\partial v_{dc,out}$（:344--353）。

\subsubsection*{参数表}
\begin{paramtable}
\fld{enable\_coupled\_jacobian} & --- & --- & on/off & \textbf{false} & AC$\leftrightarrow$DC 交叉雅可比块开关（\srcpath{include/hacdcpf/power\_flow/power\_flow\_options.hpp:149}）\\
\fld{enable\_augmented\_equations} & --- & --- & on/off & false & VDC\_VAC 母线增广方程（Direction 2，见第~\ref{sec:newton}~章）\\
\fld{r\_conv\_ac\_pu} & $r_{ac}$ & pu & 0--0.05 & 0.0 & AC 传导损耗耦合（opt-in，0 保持 legacy 行为）\\
\fld{r\_eq\_pu} & $r_{eq}$ & pu & 0--0.05 & 0.0 & DC/DC 等值电阻（$I^2R$ 损耗，opt-in）\\
\end{paramtable}

\subsubsection*{验证与校核}
\srcpath{tests/test\_unified\_pf\_upgrade.cpp:232--341} 在强化耦合的混合算例上
（\fld{enable\_coupled\_jacobian}=true、\fld{loss\_model}=CurrentBased）
对 $V_{dc}$ 列（耦合项）及前 10 个 $\theta/V_m$ 列做中心差分对拍：
$\varepsilon=10^{-6}$，max\_rel\_error $<10^{-3}$、violations==0；
同文件 :349 起的第二个用例证明耦合雅可比降低 $V_{dc}$ 扰动的 Taylor 模型误差。
\fld{r\_conv\_ac\_pu} 注入与导数由 \srcpath{tests/test\_converter\_coordination.cpp:610--650}
逐项精确断言（容差 $10^{-12}$）。

\subsection{VSC 换流器稳态模型}\label{sec:hybrid-vsc}

\compmeta{VSCConverter}{\srcpath{include/hacdcpf/model/converter\_components.hpp}}

\subsubsection*{功能说明}
VSC 是 AC/DC 双端口耦合设备，稳态模型在 \srcpath{src/power\_flow/converter\_model.cpp}。
控制模式枚举 \texttt{ConverterMode} 共\textbf{六种}
（\srcpath{include/hacdcpf/model/enums/converter\_enums.hpp:6--23}），
对应七类典型控制分类 \texttt{ACDCControlMode}（:25--43，
\srcpath{docs/archive/theory/multiple\_converter\_r1.md} §4）的 Mode 1--7
（$\delta_s{+}V_s$ / $P_s{+}V_s$ / $P_s{+}Q_s$ / $U_{dc}{+}Q_s$ / $U_{dc}{+}V_s$ /
droop $U_{dc}{+}V_s$ / droop $U_{dc}{+}Q_s$）；控制角色统一由
\srcpath{resolve\_device\_control\_role}（\srcpath{include/hacdcpf/model/device\_control\_role.hpp:62--118}）判定。

\subsubsection*{数学模型：注入与损耗}
\textbf{AC 侧注入} \srcpath{converter\_ac\_injection}（converter\_model.cpp:52--97）：
\texttt{AC\_GRID\_FORMING} 返回 $(0,0)$（AC 母线为 slack，P/Q 行整体剔除，:66--73）；
\texttt{PQ\_MODE}/\texttt{AC\_PV} 返回 $(p_{set},q_{set})$（AC\_PV 的 Q 行为占位，
其 AC 母线经装配层提升为 PV，:75--83）；\texttt{VDC\_*} 族先算下垂转移功率
（:85--87）
\[
P_{tr}=k_{vdc}\big(v_{dc}^{2}-v_{dc,set}^{2}\big),\qquad P_{ac}=P_{tr}-P_{loss}(P_{tr},v_{dc})
\]
\texttt{VDC\_Q} 返回 $(P_{ac},q_{set})$；\texttt{VDC\_VAC} 与
\texttt{DC\_V\_DROOP\_AC\_V} 返回 $(P_{ac},0)$（:89--96）。

\textbf{DC 侧注入} \srcpath{converter\_dc\_injection}（:99--146）：
PQ/AC\_PV/AC\_GRID\_FORMING 族 $P_{dc}=-(p_{set}+P_{loss}(p_{set}))$（:113--124）；
\texttt{VDC\_*} 族 $P_{dc}=-P_{tr}$（:125--131）；opt-in 交流传导损耗
（\fld{r\_conv\_ac\_pu}$>0$）使 DC 侧多抽
$P_{loss}^{AC}=r_{ac}\,P_{ac}^{2}/\max(V_m,0.1)^{2}$（:133--143）。

\textbf{损耗模型} \srcpath{converter\_loss}（:8--23）与导数
\srcpath{converter\_loss\_jacobian}（:25--50），按 \fld{loss\_model} 二选一：
\[
P_{loss}=\begin{cases}
(1-\eta)\,|p| & \text{Linear（默认）}\\[2pt]
a+bI+cI^{2},\quad I=\dfrac{|p|}{\max(v_{dc},0.1)} & \text{CurrentBased}
\end{cases}
\]
其中 $a=\fld{loss\_mw}/S_{base}$（固定损耗）、$b=\fld{loss\_percent}/100$（$\propto I$）、
$c=1-\eta$（$\propto I^2$）。

\textbf{交叉导数}（耦合雅可比用，converter\_model.cpp:152--211, 330--344）：
\texttt{VDC\_*} 族
$\partial P_{ac}/\partial v_{dc}=2k_{vdc}v_{dc}\big(1-\partial P_{loss}/\partial p\big)-\partial P_{loss}/\partial v_{dc}$
（:181）；DC 侧 PQ 族 $\partial P_{dc}/\partial v_{dc}=-\partial P_{loss}/\partial v_{dc}$、
\texttt{VDC\_*} 族 $-2k_{vdc}v_{dc}$（:200--209）；
$\partial P_{dc}/\partial V_m=2r_{ac}P_{ac}^{2}/V_m^{3}$（:341--343）。

\subsubsection*{PQ$\leftrightarrow$VDC\_Q 模式切换迟滞}
\srcpath{apply\_converter\_mode\_switching}（\srcpath{src/power\_flow/newton\_solver.cpp:797--861}）：
误差 $e=|v_{dc}-v_{dc,set}|$；PQ 下 $e>$ 高阈值且持续 \fld{mode\_hysteresis\_iters}
次迭代则升 \texttt{VDC\_Q}（:826--838）；\texttt{VDC\_Q} 下 $e<$ 低阈值且持续则降回 PQ，
但被 DC 岛参考规划提升/刚性构网锁定的换流器禁止降级（:839--854，
\fld{promotion\_lock} 见 :953--961）。切换仅在迭代残差 $<10^{-6}$ 后评估
（:1444--1452），切换计数入 \fld{profiling.converter\_mode\_switches}。
\texttt{VDC\_VAC} 母线每次试探后由 \srcpath{enforce\_vac\_setpoints}（:446--457）
钉住 $V_m=v_{ac,set}$（增广方程模式下跳过）。

\subsubsection*{参数表（converter\_components.hpp:19--128）}
\begin{paramtable}
\fld{control\_mode} & --- & --- & 六种模式 & PQ\_MODE & 控制模式（:24）\\
\fld{p\_set\_mw} / \fld{q\_set\_mvar} & $p_{set},q_{set}$ & MW/Mvar & --- & 0.0 & AC 侧功率整定（:27,43）\\
\fld{p\_is\_hard\_constraint} / \fld{p\_schedule\_mw} / \fld{p\_initial\_mw} & --- & --- & --- & false / 0 / 0 & P 硬约束与释放后的调度/初值（:40--42）\\
\fld{v\_dc\_set\_pu} & $v_{dc,set}$ & pu & 0.95--1.05 & 1.0 & DC 电压整定（:44）\\
\fld{v\_ac\_set\_pu} / \fld{v\_ac\_angle\_set\_deg} & $v_{ac,set}$ & pu/deg & --- & 1.0 / 0.0 & AC 构网/AC\_PV 整定（:45,49）\\
\fld{eta} & $\eta$ & --- & 0.95--0.99 & 0.99 & 效率（Linear 损耗，:51）\\
\fld{loss\_mw} / \fld{loss\_percent} & $a,b$ & MW/\% & --- & 0.0 & CurrentBased 损耗系数（:52--53）\\
\fld{k\_vdc} & $k_{vdc}$ & pu & 0.1--25 & 0.1 & $v_{dc}^2$ 下垂增益（:54）\\
\fld{r\_conv\_ac\_pu} & $r_{ac}$ & pu & 0--0.05 & 0.0 & AC 传导电阻（opt-in，:89）\\
\fld{i\_ac\_max\_pu} / \fld{i\_dc\_max\_pu} & --- & pu & --- & 0.0 & 电流限（0=不校核，:70--71）\\
\fld{k\_m\_modulation} / \fld{m\_min} / \fld{m\_max} & $K_m,m$ & --- & --- & 0.0 & 调制比校核参数（0=不校核，:72--74）\\
\fld{coordination\_group\_id} / \fld{is\_master} / \fld{participation\_factor} & $\alpha$ & --- & --- & 空/false/0.0 & 多源 DC 电压协调（:115--117）\\
\fld{converter\_vdc\_switch\_high\_pu} / \fld{\_low\_pu} & --- & pu & --- & 0.03 / 0.01 & PQ$\leftrightarrow$VDC\_Q 切换阈值（power\_flow\_options.hpp:120--121）\\
\fld{mode\_hysteresis\_iters} & --- & 次 & --- & 2 & 切换持续次数（:122）\\
\fld{enable\_converter\_mode\_switching} & --- & --- & on/off & true & 模式切换总开关（:126）\\
\end{paramtable}

\subsubsection*{验证与校核}
AC\_PV 保持交流电压（:1174）、\texttt{DC\_V\_DROOP\_AC\_V} 下垂构网（:1577）、
\texttt{AC\_GRID\_FORMING} 孤岛成参考（:1670）、七模式映射（:1299）等用例见
\srcpath{tests/test\_converter\_coordination.cpp}（括号为行号）；
换流器接入全组件回归见 \srcpath{tests/test\_all\_components\_pf.cpp:486--502}。

\subsection{DC/DC 变换器}\label{sec:hybrid-dcdc}

\compmeta{DCDCConverter}{\srcpath{include/hacdcpf/model/converter\_components.hpp}}

\subsubsection*{功能说明}
DC/DC 是两 DC 母线间的电力电子接口：耦合两侧功率但\textbf{不}在电气上合并岛
（岛识别注释 newton\_solver.cpp:86--90）。控制模式 \texttt{DCDCControlMode}：
\texttt{Voltage}=0、\texttt{Power}=1、\texttt{Droop}=2
（converter\_enums.hpp:52--56）。统一实现在
\srcpath{dcdc\_power\_transfer}（converter\_model.cpp:227--282）。

\subsubsection*{数学模型}
符号约定：\fld{p\_ref\_mw} 为\textbf{输出侧}参考功率，正值表示 bus\_in$\to$bus\_out
正向传输（\srcpath{src/power\_flow/pf\_injection\_assembly.cpp:76--88}）。
效率 $\eta$ 钳制 $[0.01,1]$（converter\_model.cpp:245）。输出侧参考指令：
\[
P_{out}^{ref}=\begin{cases}
p_{ref}/S_{base} & \text{Power（默认路径）}\\
p_{ref}/S_{base}-k_{droop}\big(v_{dc,out}-v_{ref}\big) & \text{Droop（:247--250）}\\
-k_{form}\big(v_{dc,out}-v_{ref}\big) & \text{Voltage（:251--269）}
\end{cases}
\]
Voltage 模式为真电压成形：刚性负反馈律，整定项 \fld{p\_ref\_mw} 故意不加；
增益 $k_{form}=\max\big(p_{rated,pu}\cdot 100,\ 25\big)$，\fld{k\_droop}$\neq0$ 时再与
$|k_{droop}|$ 取大（:260--267），$p_{rated,pu}$ 取 \fld{sn\_mva} 或
$\max(|\fld{pmax\_mw}|,|\fld{pmin\_mw}|)$ 除以基准。功率反射与 $I^2R$ 损耗（:272--283）：
\[
P_{in}=P_{out}^{ref}\cdot\big(1/\eta\ \text{若}\ P_{out}^{ref}\ge 0\ \text{否则}\ \eta\big),\qquad
P_{loss}^{I^2R}=r_{eq}\frac{P_{in}^{2}}{\max(v_{dc,in},0.1)^{2}},\qquad
P_{out}=P_{out}^{ref}-P_{loss}^{I^2R}
\]
DC 注入：入端 $-P_{in}$、出端 $+P_{out}$（pf\_injection\_assembly.cpp:86--87）。
雅可比由 \srcpath{dcdc\_jacobian\_vdc}（converter\_model.cpp:219--231）共享同一
transfer 结构给出 $\partial P_{in}/\partial v_{dc,out}$、$\partial P_{out}/\partial v_{dc,in}$、
$\partial P_{out}/\partial v_{dc,out}$。

\textbf{占空比反演} \srcpath{dcdc\_duty\_ratio}（:296--334）：端口 pu 电压乘
\fld{vn\_in\_kv}/\fld{vn\_out\_kv}（缺省按 1.0）得名值后按理想 CCM 律反演：
\[
D=\begin{cases}
V_{out}/V_{in} & \text{Buck（:309--312）}\\
1-V_{in}/V_{out} & \text{Boost（:313--316）}\\
V_{out}/(V_{in}+V_{out}) & \text{BuckBoost（:317--320）}\\
V_{out}/(n\,V_{in}) & \text{Isolated（调制增益 }M\text{，:321--326）}
\end{cases}
\]
\texttt{Generic} 拓扑不评估（\fld{defined}=false，:299）；
可行性窗口 $D\in[\fld{d\_min},\fld{d\_max}]$，容差 $10^{-9}$（:331--332）。

\subsubsection*{参数表（converter\_components.hpp:133--173）}
\begin{paramtable}
\fld{control\_mode} & --- & --- & Voltage / Power / Droop & Voltage & 控制模式（:140）\\
\fld{p\_ref\_mw} & $p_{ref}$ & MW & --- & 0.0 & 输出侧参考功率（:142）\\
\fld{v\_ref\_pu} & $v_{ref}$ & pu & 0.95--1.05 & 1.0 & 输出电压整定（:143）\\
\fld{sn\_mva}、\fld{pmax/pmin\_mw} & --- & MVA/MW & --- & 0.0 & 容量与功率限（:144,150--151）\\
\fld{eta} & $\eta$ & --- & 0.95--0.99 & 0.98 & 效率（钳制 $[0.01,1]$，:147）\\
\fld{r\_eq\_pu} & $r_{eq}$ & pu & 0--0.05 & 0.0 & 等值电阻（$I^2R$ 损耗，:148）\\
\fld{k\_droop} & $k_{droop}$ & pu & --- & 0.0 & 下垂增益（:153）\\
\fld{topology} & --- & --- & 五种拓扑 & Generic & 占空比律选择（:160）\\
\fld{d\_min} / \fld{d\_max} & $D$ & --- & 0.05--0.95 & 0.05 / 0.95 & 占空比可行性窗（:161--162）\\
\fld{n\_ratio} & $n$ & --- & --- & 1.0 & 隔离型变比（:163）\\
\end{paramtable}

\subsubsection*{验证与校核}
\srcpath{tests/test\_converter\_coordination.cpp}：
DC/DC 不合并电压岛（:279、:295 对比闭合金属路径）、
Buck 占空比反演与可行性（:493）、droop 输出成岛参考（:713）、
Voltage 模式输出成形与输入侧参考要求（:394、:419）；
OPF 三模式建模见同文件 :1493；全组件回归
\srcpath{tests/test\_all\_components\_pf.cpp:516--528}。

\subsection{LCC 换流站稳态模型}\label{sec:hybrid-lcc}

\compmeta{LCCConverter}{\srcpath{include/hacdcpf/model/converter\_components.hpp:162}}

\subsubsection*{功能说明}
LCC（电网换相换流器）站以准稳态模型接入统一 Newton 潮流（BPA/DSP dat 手册第 4 章
BD/LD 卡），实现在 \srcpath{src/power\_flow/lcc\_model.cpp/.hpp}。与 VSC 注入助手一样，
站在\textbf{当前 Newton 迭代点}求值；存在换流变时刻意区分两个交流电压量
（lcc\_model.hpp:9--15）：$E_t$ 为带载阀侧母线线电压（用于基波电流无功关系），
$E_0$ 为换流变漏抗后的理想空载阀侧电压（由一次侧母线电压与分接头求得，用于
换相特性）。站数据经 \fld{SolverData::lcc\_converters}
（solver\_data.hpp:31）进入残差与雅可比装配。

\subsubsection*{数学模型：外特性与注入}
$X_c$ 为每桥每相换相电抗，$n_b$ 为每极串联 6 脉动桥数（lcc\_model.hpp:16--21）：
\[
U_{d0}=\tfrac{3\sqrt{2}}{\pi}\,n_b\,E_0,\qquad
U_d=\begin{cases}
U_{d0}\cos\alpha-\tfrac{3}{\pi}n_b X_c I_d-n_b\,dU_v & \text{整流}\\
U_{d0}\cos\gamma-\tfrac{3}{\pi}n_b X_c I_d+n_b\,dU_v & \text{逆变}
\end{cases}
\]
无功 $Q=P\tan\varphi$，$\cos\varphi\approx U_d/U_{d0,t}$，
$U_{d0,t}=\tfrac{3\sqrt{2}}{\pi}n_b E_t$。\srcpath{lcc\_operating\_point}
（lcc\_model.hpp:80--91）按 \fld{station\_role}+\fld{control\_mode} 分派：
整流 ConstantPower/ConstantCurrent/ConstantAlpha，逆变
ConstantGamma/ConstantPower/ConstantCurrent；不支持组合返回 \fld{valid}=false，
该站不注入（:75--79）。直流电流钳制 $[0,\fld{rated\_current\_a}]$：LCC 电流不可反号，
钳制生效时 CEA/定 $\alpha$ 整定不再保持，站退化为恒流特性（$\gamma$ 保持在
$\gamma_{\min}$ 以上），经 \fld{LCCTransfer::id\_at\_limit} 与
\texttt{[LCC-PHYS-03]} 诊断如实报告（lcc\_model.hpp:30--36）。

残差侧：AC 阀侧母线 P/Q 注入与 DC 母线注入由
\srcpath{lcc\_ac\_injection}/\srcpath{lcc\_dc\_injection} 在当前迭代点求值
（jacobian\_builder.cpp:217--235）。雅可比侧：CEA（ConstantGamma）逆变站的 DC 注入
经外特性依赖端电压，$\partial p_{dc}/\partial v_{dc}$ 必须精确以使 DC Newton 方程收敛
（jacobian\_builder.cpp:451--465）；AC$\leftrightarrow$DC 耦合块恒常装配
（\srcpath{lcc\_ac\_dc\_jacobian}，jacobian\_builder.cpp:466--488），纯 DC 求解器
同样取其对角元（dc\_solver.cpp:72--76）。

\subsubsection*{分接头与 DC 岛参考}
Newton 方程在固定分接头下求值；分接头不是 Newton 状态量，公共 PF API 可用
R 卡外环包络：求解潮流→反推换相角→更新绑定 T 卡分接头→重建 Ybus→再求解
（lcc\_model.hpp:23--28）。特征控制站（CEA 逆变/定 $\alpha$ 整流）经外特性
锚定所在 DC 岛电压，作用等同下垂 VSC，无需钉母线或提升
（\srcpath{lcc\_forms\_dc\_voltage}，lcc\_model.hpp:171--175；
newton\_solver.cpp:406--428）。

\subsubsection*{结果与诚实口径}
求解后各站在收敛点重估并输出完整运行点（$\alpha$/$\gamma$/$U_{d0}$/$U_d$/$I_d$、
AC/DC 侧功率、有效分接头与 R 卡外环状态）至 \fld{PowerFlowResult::lcc\_transfers}
（\texttt{LCCTransfer}，power\_flow\_result.hpp:282；装配 hacdcpf.cpp:1525--1576）。
$\alpha$/$\gamma$ 越限不静默：\texttt{[LCC-PHYS-01]} 等事后诊断写入
\fld{diagnostics.warnings}（hacdcpf.cpp:1578 起）。

\subsection{DC 岛参考规划}\label{sec:hybrid-dcref}

\subsubsection*{功能说明}
每个 DC 电压岛必须恰有一个电压参考，否则 DC 子块奇异。规划器
（\srcpath{src/power\_flow/newton\_solver.cpp:238--444}，求解开始时调用一次 :944--945）
在求解器的换流器\textbf{可变副本}上工作：必要时把 PQ 换流器提升为电压成形，
并产出 \texttt{DcSlackPlan\{dc\_slacks, promoted\_converter\_idx, rigid\_pins, warnings\}}
（:187--193）。岛识别 \srcpath{compute\_dc\_components}（:71--111）：
邻接仅来自在运 DC 支路（闭合 DC 断路器已在投影层折叠为支路），
VSC 与 DC/DC \textbf{不}构成连边（:86--90），与协调校核器的岛定义一致。

\subsubsection*{算法：逐岛优先级}
对每岛按以下顺序确定参考（:247--441）：
\begin{enumerate}
  \item 岛内有 \texttt{DC\_V} 母线 $\to$ 钉住该母线（:300--304）；
  \item 已有 \texttt{VDC\_Q}/\texttt{VDC\_VAC}/\texttt{DC\_V\_DROOP\_AC\_V} 换流器
        下垂持压（:264--275），或 droop/Voltage 模式 DC/DC 锚定其输出母线
        （:283--298），或特征控制 LCC 站（CEA 逆变/定 $\alpha$ 整流）经外特性
        锚定（:406--428）$\to$ 无需钉母线；
  \item 否则在岛内\textbf{最大额定}的在运 PQ 换流器中选候选
        （\fld{is\_master} 声明者优先，:310--327），分两路：
        \begin{itemize}
          \item \textbf{刚性 DC-slack 构网}（opt-in \fld{enable\_rigid\_vdc\_former}，
                默认 false，\srcpath{power\_flow\_options.hpp:135--140}）：仅当岛为
                单非孤立母线、唯一 VSC 且无 DC/DC/ER 耦合时适用（:343--374）——
                母线钉于 $v_{set}$（初值种子 :950--952），换流器保持 PQ 并按线性损耗
                口径整定 $p_{set}$ 以恰好抵消岛固定注入（:376--398）：
                $\mathrm{net}\ge0$ 时 $p_{set}=\mathrm{net}/(2-\eta)$，
                $\mathrm{net}<0$ 时 $p_{set}=\mathrm{net}/\eta$；
                $\mathrm{net}$ 由 \srcpath{net\_fixed\_dc\_injection\_mw}（:198--218）计算；
          \item \textbf{刚性下垂提升（默认）}：提升为 \texttt{VDC\_Q}，增益强制
                $k_{vdc}\leftarrow\max\big(|k_{vdc}|,\ \max(p_{rated,pu}\cdot100,\ 25)\big)$，
                \fld{v\_dc\_set\_pu}$\le0.1$ 时置 1.0（:400--429），避免弱下垂下
                $V_{dc}$ 漂移越限；
        \end{itemize}
  \item 岛内无任何 VSC $\to$ 钉首个非孤立母线并告警"非物理参考"（:431--440）。
\end{enumerate}
配套机制：\srcpath{apply\_participation\_factor\_sharing}（:153--184，调用 :940）
对成员 $\ge2$、$\alpha_i\ge0$ 且 $\sum\alpha\approx1$（容差 $10^{-6}$）的协调组
重分配总下垂增益 $k_{vdc,i}\leftarrow k_{total}\cdot\alpha_i$
（$k_{total}=\sum|k_{vdc}|$，退化时取 $0.1\times$ 成员数）；被提升/刚性锁定的换流器
记入 \fld{promotion\_lock}（:953--961）防止外环降级；提升名单与告警写入
\fld{diagnostics.promoted\_vsc\_indices} 与 \fld{diagnostics.warnings}（:962--963），
终态换流器列表经 \fld{diagnostics.effective\_converters} 输出（:975, :1862）。

\subsubsection*{验证与校核}
\srcpath{tests/test\_dc\_island\_reference.cpp}：无参考 DC 岛自动提升 PQ 换流器后
收敛、$|v_{dc}-1|<0.05$ 且无越限告警（:33--70）；真 \texttt{DC\_V} 母线钉住、
不提升（:72--90）；opt-in 刚性构网把母线\textbf{精确}钉在整定值（$10^{-9}$），
且换流器恰好吸收 0.05 MW 光伏盈余（:92--120）；默认仍为下垂提升（:122--133）。
主从指定与参与因子重分配见 \srcpath{tests/test\_converter\_coordination.cpp:756--813}。

\subsection{换流器协调校核与诚实口径}\label{sec:hybrid-coord}

\subsubsection*{功能说明}
两层校核体系：求解\textbf{前}的可行性预校核
\srcpath{powerflow::evaluate\_converter\_coordination}
（\srcpath{src/power\_flow/converter\_coordination.cpp:809--1368}，头文件
\srcpath{include/hacdcpf/power\_flow/converter\_coordination.hpp}），
与求解\textbf{后}的物理不等式校核（\srcpath{src/api/hacdcpf.cpp}，:661--718 与 :941--951）。
预校核为 opt-in（开关 \fld{enable\_converter\_coordination\_check}，默认 false，
见 \srcpath{power\_flow\_options.hpp:127}）：Fatal/Error 级问题直接阻断求解
（\srcpath{src/api/hacdcpf.cpp:894--903}），Warning/Info 附入诊断（:928--936）。
岛识别与求解器一致（converter\_coordination.cpp:56--109，DC/DC 不合并岛）。

\subsubsection*{预校核规则目录}
严重度四级 \texttt{Info/Warning/Error/Fatal}（converter\_coordination.hpp:10--15）。
主要规则（实现行号括注）：
\begin{center}
\footnotesize
\begin{tabular}{@{}llp{8.6cm}@{}}
\toprule
规则 ID & 级别 & 触发条件\\
\midrule
\texttt{ACDC-REF-01} & Fatal & VSC 引用缺失 DC 母线（:323--333）\\
\texttt{ACDC-GFM-01} & Error & DC 构网换流器又声明 AC P 硬约束（:353--366）\\
\texttt{ACDC-03} & Warning & VDC 模式下 \fld{p\_set\_mw} 仅作调度值（:367--382）\\
\texttt{ACDC-DROOP-01} & Error & VDC 模式但 \fld{k\_vdc}$=0$，无法成参考（:384--398）\\
\texttt{ACDC-CTRL-03/04} & Error/Warn & AC\_PV 无有效 $v_{ac,set}$ / \fld{q\_set\_mvar} 被忽略（:403--428）\\
\texttt{ACDC-GFM-03} & Fatal & 普通两端口同时 AC/DC 双侧构网（:437--452）\\
\texttt{ACDC-GFM-04/05/06} & Error & AC 构网又钉 P / 无 DC 侧能量支撑 / DC 构网无 AC 侧平衡源（:456--515）\\
\texttt{DCDC-REF-01} & Fatal & DC/DC 引用缺失母线（:534--544）\\
\texttt{DCDC-DROOP-01} & Error & Droop 模式但 \fld{k\_droop}$=0$（:546--556）\\
\texttt{DCDC-TOPO-01} & Warning & DC/DC 两端落在同一金属 DC 岛（:926--937）\\
\texttt{DCDC-CTRL-03} & Error & 输出成形但输入岛无电压参考（:1330--1346）\\
\texttt{DCDC-CTRL-05} & Error & Power 模式两侧缺一参考（:1347--1364）\\
\texttt{DCBUS-01} & Warning & \texttt{DC\_V} 仅为声明类型、非物理电压源（:847--858）\\
\texttt{DCISLAND-01} & Fatal & 岛无电压成形设备且无可提升 PQ VSC（:1104--1121）\\
\texttt{DCISLAND-PROMOTE-01} & Info & 将由求解器自动提升形成参考（:1122--1134）\\
\texttt{DCISLAND-05} & Fatal & 多电压源整定冲突（$>10^{-4}$ pu，:1143--1165）\\
\texttt{DCISLAND-04} & Error & 多刚性电压源且无下垂分摊（:1173--1187）\\
\texttt{DCISLAND-DROOP-01} & Error & 下垂源总增益为零（:1192--1203）\\
\texttt{DCISLAND-MS-01} & Error & 主从组 master 数 $\neq1$（:1224--1238）\\
\texttt{DCISLAND-PARTICIPATION-01/02} & Error & 参与因子和 $\neq1$ / 刚性源与参与因子并存（:1240--1267）\\
\texttt{DCISLAND-02/03} & Fatal/Warn & 纯固定注入无平衡源 / 灵活性裕度不足（:1269--1310）\\
\texttt{ACISLAND-REF-01/02} & Warning & AC 岛缺/多相角参考（:702--731）\\
\texttt{ACISLAND-BALANCE-01} & Warning & AC 岛无有功平衡源（:739--751）\\
\bottomrule
\end{tabular}
\end{center}

\subsubsection*{事后物理不等式校核与诚实口径（重要）}
\textbf{换流器物理不等式约束未嵌入 Newton 方程}：容量圆、AC/DC 电流限、
调制比窗口、占空比窗口\textbf{不}作为约束参与求解。默认（仅告警）模式下，
收敛后逐台校核并以\textbf{告警}输出，不阻断收敛
（\srcpath{src/api/hacdcpf.cpp}，:661--718 与 :941--951，
裕度容差 $10^{-6}$；各检查 opt-in，限值字段为 0 即不校核）：
\begin{gather*}
\underbrace{\sqrt{p_{ac}^{2}+q_{ac}^{2}}>S_{rated}}_{\text{ACDC-PHYS-04 容量圆}},\qquad
\underbrace{S_{ac}/V_m>\fld{i\_ac\_max\_pu}}_{\text{PHYS-02 AC 电流}},\qquad
\underbrace{|p_{dc}|/v_{dc}>\fld{i\_dc\_max\_pu}}_{\text{PHYS-03 DC 电流}}\\
\underbrace{m=\frac{V_m\cdot\mathrm{vn}_{ac}}{K_m\cdot v_{dc}\cdot\mathrm{vn}_{dc}}\notin[m_{min},m_{max}]}_{\text{PHYS-05/01 调制比}},\qquad
\underbrace{D\notin[\fld{d\_min},\fld{d\_max}]}_{\text{DCDC-PHYS-01 占空比}}
\end{gather*}
opt-in 开关 \fld{enforce\_converter\_physical\_limits}（默认 false，
\srcpath{power\_flow\_options.hpp:128--134}）把同一组不等式转为\textbf{硬可行性门卫}：
确定性潮流没有再调度自由度去修复不相容的整定，因此收敛后若存在任一
ACDC-PHYS 告警或占空比越窗，该 Newton 根被\textbf{拒绝为物理不可行}——
\fld{converged} 置 false、\fld{diagnostics.termination\_reason} 记
"Converter physical-limit feasibility check failed" 并追加
[CONVERTER-PHYS-HARD] 告警（hacdcpf.cpp:957--978）；需要再调度寻找可行
运行点时应改用 OPF。
求解后另有 DC 母线电压越限告警
（$v_{dc}\notin[\fld{vmin\_pu},\fld{vmax\_pu}]$，\srcpath{src/power\_flow/newton\_solver.cpp:1839--1857}）。

结果对象的 \fld{converter\_model\_scope}（结构
\srcpath{include/hacdcpf/model/converter\_model\_scope.hpp:21--53}，
填充 \srcpath{src/api/hacdcpf.cpp:1002--1018}）如实声明本次快照求解的模型覆盖：
\fld{model\_scope}=\texttt{"steady-state-newton:vsc-3mode+dcdc-power-transfer"}；
\texttt{true} 项：VSC 损耗、AC 传导损耗耦合、$V_{dc}$ 控制、DC/DC 损耗、
多源协调、方程闭合检查；四个 \texttt{enforced} 标志
\fld{vsc\_capacity\_circle\_enforced}、\fld{vsc\_current\_limits\_enforced}、
\fld{vsc\_modulation\_limits\_enforced}、\fld{dcdc\_duty\_ratio\_enforced}
回填 \fld{opt.enforce\_converter\_physical\_limits}（:1012--1016）——
默认仅告警模式下为 \texttt{false}（不等式未在方程内强制），
开启硬约束门卫后为 \texttt{true}（违反限值的根被拒绝而非报告收敛）。

\subsubsection*{验证与校核}
规则覆盖见 \srcpath{tests/test\_converter\_coordination.cpp}：GFM-01/03/04/05/06
（:433, 1142, 1162, 1766, 1788）、DCISLAND-05 整定冲突判 Fatal（:379）、
MS/PARTICIPATION 系列（:527--586）、ACISLAND 系列（:1106, 1830）；
\fld{converter\_model\_scope} 声明与 PHYS-03 触发分别由 :652--669、:671--693 断言
（后者证明限值越限时\textbf{仍收敛}、仅告警——默认模式行为）；方程闭合检查见 :695--711。

\subsection{验证与校核汇总}
本章各节的专项验证之外，统一交叉验证包括：\srcpath{tests/test\_unified\_pf\_upgrade.cpp}
的耦合雅可比中心差分对拍（:232--341，$\varepsilon=10^{-6}$，max\_rel\_error $<10^{-3}$）、
\srcpath{tests/test\_all\_components\_pf.cpp} 的 29 类组件逐类禁用回归
（混合族含 VSCConverter、DCDCConverter、EnergyRouter，文件头 :15--19）。
\textbf{已知边界（如实）}：换流器不等式默认仅事后告警，硬可行性门卫需显式打开
\fld{enforce\_converter\_physical\_limits}（第~\ref{sec:hybrid-coord}~节）；
\fld{enable\_coupled\_jacobian} 默认关闭，交叉块导数需显式开启；
DC/DC \texttt{Generic} 拓扑无占空比约束；纯 DC 求解器的 AC 侧固定为初值、
不做 AC/DC 联立。

% 第 7 章：配电网潮流（前推回代 / 调压器外环 / PV 功率曲线 / 支路回算）
% 本章文件由 \input 引入，不含导言区。
\section{配电网潮流}\label{sec:dist}\label{sec:distribution}

\subsection{功能定位与入口}

\compmeta{analysis::solve\_distribution\_pf}{\srcpath{include/hacdcpf/power\_flow/distribution\_power\_flow.hpp}}

配电网潮流族以辐射状馈线的\textbf{前推回代}（Backward-Forward Sweep, BFS；
Shirmohammadi et al.\ 1988）为主路径，面向单相（正序）等值模型
（include/hacdcpf/power\_flow/distribution\_power\_flow.hpp）。两个重载入口（:solve\_distribution\_pf）：
\begin{quote}\footnotesize
\texttt{DPFResult solve\_distribution\_pf(const ACSystem\& ac\_sys, const DPFOptions\& opt = \{\});}\\
\texttt{DPFResult solve\_distribution\_pf(const HybridPowerSystem\& sys, const DPFOptions\& opt = \{\});}
\end{quote}
实现位于 \srcpath{src/power\_flow/distribution\_power\_flow.cpp}（单相内核
:390--540；入口 :1041--1387）。结果索引空间与 \fld{ac\_sys.buses}/\fld{branches}
\textbf{完全同序}（hpp:75--76, 432--433）：\fld{vm\_pu[i]} 对应
\texttt{buses[i]}，\fld{p\_branch\_mw[b]} 对应 \texttt{branches[b]}。
\texttt{HybridPowerSystem} 重载先经 \srcpath{projection::RichToCanonicalOperator::apply}
投影到 canonical 模型（cpp:1067--1068），求解后把 \fld{vm\_pu}/\fld{va\_deg}
按 \fld{bus\_merge\_map} 以强度量语义广播回原始母线编号空间
（:1048--1057，回投影契约见第~\ref{sec:overview}~章；注意
\fld{p\_branch\_mw} 等支路量\textbf{不}回投影，仍按 canonical 支路序索引）。

如实说明定位边界：\srcpath{solve\_distribution\_pf} 是库级 API，
\textbf{未}接入公共门面 \srcpath{include/hacdcpf/api/hacdcpf.hpp}，也未接入 GUI
后端 \srcpath{tests/run\_gui\_server.cpp} 的 \texttt{method} 分发；环网与多电源
三相场景应改用三相 NR（第~\ref{sec:threephase}~章）或统一 Newton
（第~\ref{sec:newton}~章）。同文件还提供 abc 域三相 BFS
\srcpath{solve\_three\_phase\_distribution\_pf}（hpp:456--464，
实现 cpp:1632--1804），其公式与 NR 对照细节见第~\ref{sec:threephase}~章，
本章仅在支路回算一处引用其损耗口径。

\subsection{辐射拓扑处理与 BFS 树}
\srcpath{build\_bfs\_tree}（cpp:74--117）从首个 SLACK 母线出发对在役支路
邻接表做广度优先遍历，生成三个并行数组：BFS 访问序 \fld{order}、父母线
\fld{parent\_bus}、父支路 \fld{parent\_branch} 及父侧方向标志
\fld{parent\_is\_from\_side}（:build\_bfs\_tree）。无 SLACK 直接抛异常（:solve\_distribution\_pf\_ac\_snapshot）。
诚实口径：
\begin{itemize}
  \item \textbf{辐射假设不强制校验}：头文件仅要求调用方先用
        \srcpath{topology\_analysis.hpp} 的 \texttt{is\_radial()} 自查
        （hpp:429--430）；环网送入时 BFS 树按先访问者确定父支路，环路支路
        不参与扫描但仍出现在支路潮流回算中，结果\textbf{不报错也不保证正确}；
  \item \textbf{与根不连通的孤岛}：未进入 \fld{order} 的母线电压保持初值
        （\fld{vm\_pu}、相角 $0$），不参与残差计算，结果中静默保留（:solve\_distribution\_pf\_ac\_snapshot）；
  \item \textbf{多 SLACK}：只取首个 SLACK 为根（:solve\_distribution\_pf\_ac\_snapshot），其余 SLACK 母线
        当作普通节点处理，其出力由 \fld{compute\_net\_injection} 按普通发电机
        扣减（:260--267 仅跳过根）。
\end{itemize}

\subsection{前推回代算法}
净注入 \fld{compute\_net\_injection}（cpp:199--312）按
$S^{\mathrm{net}}_i=(S^{\mathrm{load}}_i-S^{\mathrm{gen}}_i)/S_{\mathrm{base}}$
汇总：优先取 \fld{loads} 表（为空时回退母线级 \fld{pd\_mw}/\fld{qd\_mvar}，
:208--225）、充电站（kW/kvar 换算，:228--235）、可选母线与可分组并联补偿
（\fld{include\_shunts}，:238--257）、常规/静态/可再生发电机、PV（经功率曲线，
:290--299，见本节下文）与储能（放电为正，:301--309），发电机侧一律跳过根节点。
基准容量 \fld{base\_mva}$\le 0$ 时回退 $100$~MVA（:solve\_distribution\_pf\_ac\_snapshot）。注意 ZIP 负荷在此
被\textbf{恒功率近似}（注释 :214--216），不随电压迭代更新。

内层扫描 \srcpath{run\_bfs\_sweep}（:run\_bfs\_sweep）。记复变比
$t=|\mathrm{tap}|\angle\delta$（$|\mathrm{tap}|<10^{-12}$ 视为 $1$，:121--125），
支路串联电流
\[
I_s=\big(V_f/t - V_t\big)/z,\qquad z=r+jx,\quad |z|<10^{-20}\Rightarrow I_s=0
\quad(:127\text{--}135).
\]
\textbf{回推}（:run\_bfs\_sweep）：由 BFS 序逐点计算注入电流
$I_b=\overline{S^{\mathrm{net}}_j/V_j}$（$|V_j|\le10^{-12}$ 时取 $0$），再逆序
把子支路电流累加进父支路，跨越变比时按父侧方向共轭折算：父在 from 侧
$I_{\mathrm{up}}=I_{\mathrm{child}}/\overline{t}$，父在 to 侧
$I_{\mathrm{up}}=\overline{t}\,I_{\mathrm{child}}$（:branch\_subtree\_current\_seen\_from\_parent）。
\textbf{前推}（:branch\_child\_voltage\_from\_parent）沿 BFS 序更新电压：
\[
V_j=\begin{cases}
V_i/t - z\,I_b & \text{父在 from 侧}\\[2pt]
t\,\big(V_i - z\,\overline{t}\,I_b\big) & \text{父在 to 侧}
\end{cases}
\qquad(:158\text{--}168).
\]
收敛判据为全母线 $\max|\Delta V|<\texttt{opt.tol}$（:run\_bfs\_sweep）；初值取各母线
\fld{vm\_pu}、相角 $0$ 的平启动（:solve\_distribution\_pf\_ac\_snapshot）。收敛后由
\srcpath{compute\_single\_phase\_branch\_flow}（:compute\_single\_phase\_branch\_flow）回算两端功率
$S_f=V_f\overline{I_f}\,S_{\mathrm{base}}$、$S_t$ 与损耗
$S_{\mathrm{loss}}=S_f+S_t$，写入 \fld{p\_branch\_mw}/\fld{q\_branch\_mvar}
与总损耗（:solve\_distribution\_pf\_ac\_snapshot）。

可选\textbf{无功限值外环}（\fld{enforce\_q\_limits}，:438--507）：内层收敛后
逐台发电机复算母线所需无功 $Q_{\mathrm{req}}$（由其相连支路潮流求和再扣回本地
负荷，:457--476），越 \fld{qmax\_mvar}/\fld{qmin\_mvar} 则把 \fld{qg} 钉到限值、
标记已钳位并重启平启动；外环上限 \fld{q\_limit\_max\_outer\_iter}=30 次，
\fld{qmax}=\fld{qmin}=0 视为无限值跳过（:solve\_distribution\_pf\_ac\_snapshot）。钳位不可逆、无迟滞恢复，
与统一 Newton 的 PV$\leftrightarrow$PQ 外环（第~\ref{sec:newton}~章）相比更简化。

\begin{paramtable}
\fld{max\_iter} & $N_{\mathrm{it}}$ & 次 & $50\sim500$ & 100 & BFS 内层迭代上限（hpp:65）\\
\fld{max\_control\_iter} & $N_{\mathrm{ctl}}$ & 次 & $10\sim100$ & 100 & 调压器离散控制外环上限（:DPFOptions）\\
\fld{tol} & $\varepsilon$ & pu & $10^{-10}\sim10^{-6}$ & $10^{-6}$ & 收敛容差，$\max|\Delta V|$（:DPFOptions）\\
\fld{verbose} & --- & bool & --- & false & 打印逐迭代残差（:DPFOptions）\\
\fld{include\_shunts} & --- & bool & --- & true & 计入母线/补偿器 $g_s,b_s$（:DPFOptions）\\
\fld{enforce\_q\_limits} & --- & bool & --- & false & 启用发电机 Q 限值外环（:DPFOptions）\\
\fld{q\_limit\_max\_outer\_iter} & --- & 次 & $10\sim30$ & 30 & Q 限值外环上限（:DPFOptions）\\
\end{paramtable}

\subsection{调压器（OLTC/稳压器）离散控制外环}
\texttt{HybridPowerSystem} 重载检测到任一 \fld{regulator\_controls} 启用即进入
控制外环（cpp:1058--1066, 1106--1380）：每轮先投影、跑一趟 BFS，再逐个调压器
评估一次档位动作，全部 \texttt{within\_band} 才判 \fld{control\_converged}。
控制律为 OpenDSS RegControl 风格（伏特域，非标幺）：

\begin{enumerate}
  \item \textbf{监控电压}：默认监控 \fld{winding} 侧端母线，\fld{monitored\_bus}$\neq0$
        时取远端母线（:solve\_distribution\_pf）；电压幅值乘 $U_{\mathrm{base}}=$
        \fld{base\_kv}$\times1000$（:629--635, 1196--1202，缺 \fld{base\_kv} 抛异常）；
        PT 变比远端时优先 \fld{remote\_ptratio}（:solve\_distribution\_pf）；
  \item \textbf{线路压降补偿（LDC）}：\fld{r\_volts}/\fld{x\_volts} 非零时
        要求 \fld{ct\_primary\_amps}$>0$，取调压绕组侧支路电流
        $I_w$（安培，基值 $I_{\mathrm{base}}=S_{\mathrm{base}}\!\times\!10^{3}/U_{\mathrm{base,kV}}$，
        :1238--1241），补偿相量
        $\Delta U_c=(I_w/I_{\mathrm{ct}})\,(R+jX)$（:solve\_distribution\_pf）；
  \item \textbf{控制电压与死区}：
        $U_{\mathrm{ctl}}=\big|V_{\mathrm{mon}}/\mathrm{ptratio}+\Delta U_c\big|$
        （:solve\_distribution\_pf）；$U_{\mathrm{ctl}}<U_{\mathrm{reg}}-B/2-10^{-6}$ 升压、
        $>U_{\mathrm{reg}}+B/2+10^{-6}$ 降压，否则 \texttt{within\_band} 停
        （:1248--1250, 1270--1299；$10^{-6}$~V 数值迟滞）；
  \item \textbf{档位动作}：连续变比
        $\mathrm{tap}_{\mathrm{pu}}=1+(\mathrm{tap\_pos}-\mathrm{tap\_neutral})\cdot s\%/100$
        （:transformer\_tap\_pu）；单步动作
        $\Delta\mathrm{pos}=\pm\min(\texttt{max\_tap\_change},|\mathrm{tap\_max}-\mathrm{tap\_min}|)$，
        方向由 \fld{tap\_side} 决定（升压时 $=1$ 取 $+1$、$=0$ 取 $-1$，:567--569），
        并钳位到 $[\mathrm{tap\_min},\mathrm{tap\_max}]$（:solve\_distribution\_pf）；触限即
        fail-close（\texttt{tap\_limit\_reached\_*}）；
  \item \textbf{振荡防护与诚实停机}：以全体启用调压器的 \fld{tap\_pos} 向量
        做历史状态集，重复即 \texttt{cycle\_detected}；耗尽
        \fld{max\_control\_iter} 记 \texttt{max\_control\_iter\_reached}；两种情形
        均置 \fld{control\_converged}=\fld{converged}=false（:solve\_distribution\_pf）。
\end{enumerate}
不支持的模式\textbf{显式 fail-close} 而非静默近似（:1164--1187, 1309--1317）：
\fld{reversible}、\fld{monitored\_node}$\neq1$、\fld{tap\_winding} 与
\fld{tap\_side} 不匹配、远端监控与 LDC 同用、缺失离散档位网格等均抛出原因串
写入 \fld{stop\_reason}。逐轮决策轨迹见 \fld{regulator\_trace}
（hpp:78--96），终态汇总 \fld{regulator\_states}（hpp:98--128）。

\begin{paramtable}
\fld{vreg\_volts} & $U_{\mathrm{reg}}$ & V & $115\sim125$ & 0 & 目标控制电压（include/hacdcpf/model/ac\_components.hpp:RegulatorControl）\\
\fld{band\_volts} & $B$ & V & $1\sim4$ & 0 & 全死区宽度，半宽 $B/2$（:RegulatorControl）\\
\fld{ptratio} & $n_{\mathrm{PT}}$ & --- & $20\sim120$ & 0 & PT 变比；$\le10^{-12}$ 抛错（:RegulatorControl）\\
\fld{remote\_ptratio} & --- & --- & 同左 & 0 & 远端 PT 变比；$\le0$ 回退 \fld{ptratio}（:RegulatorControl）\\
\fld{ct\_primary\_amps} & $I_{\mathrm{ct}}$ & A & $100\sim1000$ & 0 & CT 一次电流，LDC 必填（:RegulatorControl）\\
\fld{r\_volts}/\fld{x\_volts} & $R+jX$ & V & $0\sim30$ & 0 & LDC 补偿阻抗（伏特域，:1127--1128）\\
\fld{max\_tap\_change} & --- & 档 & $1\sim4$ & 1 & 单轮最大动作档数；$<1$ 抛错（:RegulatorControl）\\
\fld{winding}/\fld{tap\_winding} & --- & 1|2 & --- & 0 & 监控/调档绕组（:RegulatorControl）\\
\fld{monitored\_bus} & --- & id & --- & 0 & 0=监控绕组端母线（:RegulatorControl）\\
\fld{reversible} & --- & bool & --- & false & 逆功率模式；当前不支持（:RegulatorControl）\\
\end{paramtable}

\subsection{PV 功率曲线}
\compmeta{powerflow::compute\_pv\_power\_mw}{\srcpath{include/hacdcpf/power\_flow/pv\_power\_curve.hpp}}

经验单二极管族曲线（纯头文件实现，include/hacdcpf/power\_flow/pv\_power\_curve.hpp:pv\_module\_current）：
\[
I(V)=I_{\mathrm{sc}}\Big(1-\big(V/V_{\mathrm{oc}}\big)^{a}\Big)^{b}
      \Big(1-c\big(V/V_{\mathrm{mpp}}\big)^{2}\Big),
\qquad a=10.0,\ b=0.547596,\ c=0.023812 .
\]
阵列出力（:compute\_pv\_power\_mw）在 MPPT 点取值（$V=V_{\mathrm{mpp}}$）：
\[
P_{\mathrm{MW}}=I(V_{\mathrm{mpp}})\,V_{\mathrm{mpp}}\,
  N_{\mathrm{series}}N_{\mathrm{parallel}}\,
  \frac{G}{1000}\,\eta_{\mathrm{inv}}\times10^{-6}.
\]
\fld{voc}/\fld{isc}/\fld{vmpp} 任一 $\le0$ 时回退静态 \fld{p\_mw}（:compute\_pv\_power\_mw）。
相关字段默认值（\srcpath{include/hacdcpf/model/ac\_components.hpp:PVSystem}）：
\fld{inverter\_eff}=0.97、\fld{irradiance}=1000~W/m$^2$、串并联数
\fld{num\_series}/\fld{num\_parallel} 默认 0（触发回退）。BFS 在净注入阶段调用
（src/power\_flow/distribution\_power\_flow.cpp:compute\_net\_injection），GUI 时序仿真亦复用同一函数
（\srcpath{tests/run\_gui\_server.cpp:has\_profile（lambda）}）。诚实口径：该模型是
\textbf{MPPT 单点估值}而非 $P(V)$ 电压依赖注入——潮流迭代中 PV 出力恒定，
曲线自变量并不取端电压；\fld{alpha\_isc}/\fld{beta\_voc} 温度系数字段存在
但此函数未使用。

\subsection{支路功率回算与结果结构}
两套支路回算并存，勿混淆：
\begin{itemize}
  \item \textbf{BFS 内部回算} \srcpath{compute\_single\_phase\_branch\_flow}
        （src/power\_flow/distribution\_power\_flow.cpp:compute\_single\_phase\_branch\_flow）：仅串联阻抗 + 复变比，
        $S_f=V_f\overline{I_f}S_{\mathrm{base}}$，损耗 $=S_f+S_t$；
        不含对地支路充电 $b/2$——BFS 结果中的 \fld{q\_branch\_mvar} 与统一
        Newton 口径在高充电线路上存在系统性差异；
  \item \textbf{统一 Newton 回算} \srcpath{powerflow::compute\_branch\_flows}
        （\srcpath{src/power\_flow/branch\_flow.cpp:compute\_branch\_flows}，规范头
        \srcpath{include/hacdcpf/power\_flow/assembly/branch\_flow.hpp}，
        \srcpath{power\_flow/branch\_flow.hpp} 为转发头）：MATPOWER $\pi$ 模型
        $y_{ff}=(y_s+jb/2)/|t|^2$、$y_{ft}=-y_s/\overline{t}$、
        $y_{tf}=-y_s/t$（:compute\_branch\_flows），$|z|=0$ 与停运支路跳过，
        供 \srcpath{solve\_power\_flow} 派生结果使用（见第~\ref{sec:newton}~章）。
\end{itemize}
输出 \fld{DPFResult}（include/hacdcpf/power\_flow/distribution\_power\_flow.hpp:DPFResult）：
\fld{vm\_pu}/\fld{va\_deg}（相角单位\textbf{度}）、\fld{p\_branch\_mw}/
\fld{q\_branch\_mvar}（送端 from 侧）、总损耗、\fld{converged}/
\fld{control\_converged}、\fld{iterations}、\fld{residual}（末轮
$\max|\Delta V|$），以及调压器 \fld{regulator\_states}/\fld{regulator\_trace}。
\fld{specialized\_path\_id} 记录三绕组三角等专属致密 NR 试验路径的触发与
fail-close 原因（cpp:637--648, 995--996）。三相 BFS 的损耗回算保留
$Z_{abc}$ 互耦交叉项（注释明确反对仅用自阻抗 $|I_p|^2$ 近似，
src/power\_flow/distribution\_power\_flow.cpp:solve\_three\_phase\_distribution\_pf）。

\subsection{验证与校核}
全部 BFS/调压器对照集中在 \srcpath{tests/test\_opendss\_pf\_compare.cpp}
（Catch2 + OpenDSS 桥，受 \srcpath{HACDCPF\_ENABLE\_OPENDSS\_COMPARE} 门控，
\srcpath{tests/CMakeLists.txt:254--257}，OpenDSS 缺失时整组 skip）：
\begin{itemize}
  \item \textbf{3 母线 1ph 基线}（:TEST\_CASE）：vs OpenDSS 快照，
        $|V|$ 容差 $10^{-3}$、相角 $0.1^\circ$、送端 $P/Q$ 各 $10^{-3}$、
        总损耗 $10^{-3}$；并联补偿变体（:TEST\_CASE）与物理电容器 vs BFS 恒功率
        并联近似的探索性对照（:TEST\_CASE）；
  \item \textbf{case33bw 全馈线}（:TEST\_CASE）：Baran \& Wu 33 节点平衡三相
        OpenDSS 模型 vs BFS 单相等值，33 母线 $|V|$ 容差 $2\times10^{-3}$、
        相角 $0.2^\circ$，32 支路三相合计 $P/Q$ 各 $10^{-3}$，总损耗
        $5\times10^{-3}$；同测试含 OpenDSS 自相一致性回归（相间幅值差
        $<10^{-4}$、$\pm120^\circ$ 偏差 $<10^{-3}$~deg）；
  \item \textbf{调压器闭环}（:TEST\_CASE）：本地 PT、本地 PT+LDC、
        RemotePTRatio 三算例，终态档位数与 OpenDSS \texttt{regcontrol}
        精确一致（容差 $10^{-12}$），控制电压与按 DSS 数据独立复算值差
        $<0.2$~V（:verify\_regulator\_case\_matches\_repository），终止原因 \texttt{within\_band}；另有运行时
        语义矩阵（:1970，remote/local/LDC 组合逐场景核对 \fld{stop\_reason}
        与 DSS 桥一致）与带母线合并的失败路径索引回归（:2490, :2519）；
  \item \textbf{固定档变压器}：1ph 直通变压器（:TEST\_CASE）、固定档 vs 变比感知
        BFS（:TEST\_CASE）、LV 侧固定档端到端（:TEST\_CASE）。
\end{itemize}
诚实边界：单相 BFS \textbf{没有}独立单元测试文件，全部覆盖经由上述
OpenDSS 门控集成测试；环网/多 SLACK/死岛行为无测试；case33bw 对照为
\textbf{平衡}系统（单相等值严格成立的特例）；IEEE 13 节点正式交叉校核
当前状态为 \texttt{deferred}、原因 \texttt{model\_limitations}（:3235--3264,
:3438），即不平衡全馈线的 BFS 精度尚未验收——不平衡场景以三相 NR 的
回归（见第~\ref{sec:threephase}~章）为准。

% 第 8 章：三相相域潮流与三相混合潮流
% 本章文件由 \input 引入，不含导言区。
\section{三相相域潮流与三相混合潮流}\label{sec:threephase}

\subsection{功能定位与入口}

\compmeta{analysis::solve\_three\_phase\_nr}{\srcpath{include/hacdcpf/power\_flow/distribution\_power\_flow.hpp:ThreePhaseNROptions}}

相域（abc domain）潮流是平台处理\textbf{不平衡配电系统}的生产路径：不再把网络
压缩为正序单相等值，而是直接在相节点上建立紧凑导纳矩阵并联立求解。与
第~\ref{sec:newton}~章的统一 Newton 不同，本章求解器只消费
\texttt{ThreePhaseACSystem} 子系统（\fld{HybridPowerSystem::three\_phase\_ac}），
支持辐射状与环网拓扑、多电源、全相域/序分量混合建模的线路、任意接线组变压器
与外部电网 Norton 等值。入口两个重载
（实现 \srcpath{src/power\_flow/three\_phase\_nr.cpp:solve\_three\_phase\_nr\_snapshot} 与 :5915）：
\begin{quote}\footnotesize
\texttt{ThreePhaseDPFResult solve\_three\_phase\_nr(const ThreePhaseACSystem\& sys,\\
\hspace*{2em}const ThreePhaseNROptions\& opt = \{\});}\\
\texttt{ThreePhaseDPFResult solve\_three\_phase\_nr(const HybridPowerSystem\& sys,\\
\hspace*{2em}const ThreePhaseNROptions\& opt = \{\});}
\end{quote}
\texttt{HybridPowerSystem} 重载仅转发到 \fld{three\_phase\_ac} 子系统，缺失时抛异常
（src/power\_flow/three\_phase\_nr.cpp:capture\_tap\_state（lambda））。便捷包装
\srcpath{powerflow::solve\_three\_phase}（\srcpath{include/hacdcpf/power\_flow/three\_phase.hpp:solve\_three\_phase}）
把 \fld{ThreePhaseFlowOptions} 转写为 NR 选项并回传简化结果；该头文件明确
\textbf{不再暴露}前推回代式 legacy 接口（include/hacdcpf/power\_flow/three\_phase.hpp:ThreePhaseFlowOptions，配网 BFS 见
第~\ref{sec:dist}~章）。求解前提是\textbf{至少一个 SLACK 母线或一台在运
\texttt{external\_grid}}（src/power\_flow/three\_phase\_nr.cpp:finalize\_three\_phase\_result\_from\_state）。

\subsection{相域网络模型}

\paragraph{紧凑相节点索引} \srcpath{PhaseNodeIndexer::build}
（src/power\_flow/three\_phase\_nr.cpp:build）按母线 \fld{phase\_mask} 只为\textbf{实际存在的相}
分配节点号：三相母线占 3 个节点，单相支线占 1 个，空 mask 直接抛错（:build）。
全 abc 网络的节点数为 $3N$，含单相/两相支线时严格小于 $3N$，不产生幻影节点。

\paragraph{线路 3$\times$3 阻抗块} 线路有两种参数化
（src/power\_flow/three\_phase\_nr.cpp:validate\_generator\_dispatch\_semantics）：
\begin{itemize}
  \item \textbf{序分量参数}（\fld{r1\_pu}/\fld{x1\_pu}/\fld{r0\_pu}/\fld{x0\_pu}）：
        由对称分量变换合成循环对称的相域阻抗矩阵
        \[
        Z_{abc}=T\,\mathrm{diag}(z_0,z_1,z_1)\,T^{-1},\qquad
        T=\begin{bmatrix}1&1&1\\1&a^2&a\\1&a&a^2\end{bmatrix},\quad a=e^{j2\pi/3}
        \]
        （\srcpath{build\_sequence\_zabc}，:570--587）；$|z_0|<10^{-20}$ 时取 $z_0=z_1$
        （\srcpath{build\_line\_zabc}，:589--594），即退化为平衡线路；
  \item \textbf{全相域矩阵}（\fld{use\_phase\_matrix}，\fld{r\_matrix\_pu}/\fld{x\_matrix\_pu}/%
        \fld{b\_matrix\_pu}）：直接读入 $3\times3$ 矩阵，并校验 mask 外元素必须为零，
        否则抛错（\srcpath{validate\_line\_phase\_matrix\_against\_mask}，:610--628）。
\end{itemize}
按线路 \fld{phase\_mask} 提取子矩阵后求逆得串联导纳，$\pi$ 型半并纳各半
（\srcpath{build\_line\_primitive}，:1080--1121）。

\paragraph{变压器与电源} 变压器以\textbf{绕组域原语}进入导纳矩阵：接线种类
（GroundedWye/Wye/Delta，src/power\_flow/three\_phase\_nr.cpp:ThreePhaseSolveSnapshot）经连接嵌入矩阵把绕组量
映射到相节点（\srcpath{build\_transformer\_connection\_embedding}，:1581；
\srcpath{delta\_connection\_matrix}，:1377），支持钟点相移与显式
\fld{\_winding\_topology}；接地星--接地星路径含零序磁化 T 模型
（\srcpath{build\_transformer\_zero\_sequence\_port\_admittance}，:1193，按
\fld{si0\_hv\_partial} 拆分两侧；合成入口 \srcpath{build\_transformer\_primitive}，:2485）。
外部电网建模为序阻抗后的 Norton 源
（\srcpath{build\_source\_norton\_primitive}，:2313；序阻抗推导 :2121--2176），
其注入电流进入 \fld{fixed\_current} 向量。

\paragraph{每相基准} 相域 NR 采用\textbf{每相功率基准}
$S_{base,1\phi}=S_{base}/3$：有名值 MW/MVAr 乘 $3/S_{base}$ 化为每相标幺
（\srcpath{phase\_domain\_power\_base\_scale}，:3108--3114；
\srcpath{stamp\_injection}，:2476--2482），支路功率回算时再乘回 $S_{base}/3$
（:solve\_three\_phase\_nr\_snapshot）。\fld{base\_mva} 未设置时取 100（:finalize\_three\_phase\_result\_from\_state）。

\subsection{abc 极坐标 Newton：失配方程与雅可比}

节点角色（\srcpath{build\_tp\_nr\_context}，src/power\_flow/three\_phase\_nr.cpp:compute\_phase\_spec）：
SLACK 母线全部相的 $\theta$、$V_m$ 固定；PV 母线仅 $\theta$ 为未知量；
PQ 母线 $\theta$、$V_m$ 皆未知。状态向量
$x=[\theta_{n_p};\ V_{m,\,n_q}]$，$n_{\mathrm{var}}=n_p+n_q$。
指定注入 $P^{spec}=P_{gen}-P_{load}$（注入为正，\srcpath{compute\_phase\_spec}，
:3842--3916；发电机支持每相调度或总量按相均分 :3882--3912）。
失配约定 $\mathrm{mismatch}=F^{spec}-F^{calc}$，每个相节点 $r$ 的注入
（:add\_delta\_load\_contributions）：
\begin{align*}
P_r^{calc}&=\sum_{s}V_rV_s\big(G_{rs}\cos\theta_{rs}+B_{rs}\sin\theta_{rs}\big)\\
Q_r^{calc}&=\sum_{s}V_rV_s\big(G_{rs}\sin\theta_{rs}-B_{rs}\cos\theta_{rs}\big)
\end{align*}
其中 $r,s$ 遍历紧凑相节点，$G_{rs}+jB_{rs}$ 是含相间互耦的导纳块元素。
Norton 源电流单独累加（:evaluate\_tp\_mismatch\_and\_jacobian）。

雅可比非对角元（$r\neq s$，:4301--4326）与第~\ref{sec:newton}~章极坐标形式
逐相同构：
\begin{align*}
\partial P_r/\partial\theta_s&=V_rV_s\big(G_{rs}\sin\theta_{rs}-B_{rs}\cos\theta_{rs}\big) &
\partial P_r/\partial V_s&=V_r\big(G_{rs}\cos\theta_{rs}+B_{rs}\sin\theta_{rs}\big)\\
\partial Q_r/\partial\theta_s&=-V_rV_s\big(G_{rs}\cos\theta_{rs}+B_{rs}\sin\theta_{rs}\big) &
\partial Q_r/\partial V_s&=V_r\big(G_{rs}\sin\theta_{rs}-B_{rs}\cos\theta_{rs}\big)
\end{align*}
对角元（:evaluate\_tp\_mismatch\_and\_jacobian）：
\begin{align*}
\partial P_r/\partial\theta_r&=-Q_r^{calc}-B_{rr}V_r^2 &
\partial P_r/\partial V_r&=P_r^{calc}/V_r+G_{rr}V_r\\
\partial Q_r/\partial\theta_r&=P_r^{calc}-G_{rr}V_r^2 &
\partial Q_r/\partial V_r&=Q_r^{calc}/V_r-B_{rr}V_r
\end{align*}
恒功率/恒电流负荷的失配与导数以注入功率形式叠加到上述方程
（\srcpath{add\_wye\_load\_contributions}，:4019；
\srcpath{add\_delta\_load\_contributions}，:4119，含中性点漂移的链式求导）。

\paragraph{求解循环与全局化} Newton 内核复用引擎层
\srcpath{engine::NativeNewtonAdapter}（:find\_transformer\_position\_by\_index）：对角正则初值
$\lambda_0=10^{-10}$、回溯线搜索步退 0.5、最多 20 步；投影算子把 $V_m$ 钳在
下限 0.05（:try\_factorize（lambda））。$n_{\mathrm{var}}=0$（全 SLACK）时直接判收敛（:finalize\_three\_phase\_result\_from\_state）。
收敛后以末态残差与迭代数写入结果（:find\_transformer\_observation\_position），并计算电压不平衡度
\[
V_1=\tfrac{1}{3}\big(V_a+aV_b+a^2V_c\big),\qquad
V_2=\tfrac{1}{3}\big(V_a+a^2V_b+aV_c\big),\qquad
\mathrm{VUF}=|V_2|/|V_1|\times100\%
\]
取全 abc 母线上的最大值 \fld{max\_vuf\_percent}
（\srcpath{compute\_max\_vuf}，:4433--4455）。结果结构
\fld{ThreePhaseDPFResult}（include/hacdcpf/power\_flow/distribution\_power\_flow.hpp:ThreePhaseDPFResult）含分相母线电压、
分相支路功率、分相损耗与调压器轨迹，并带诚实口径字段
\fld{primary\_solver}/\fld{solver\_used}/\fld{result\_source}。

\paragraph{调压器外环} 存在启用的 \fld{regulator\_controls} 时，NR 解被包进
离散档位外环（:5576--5910，上限 \fld{max\_control\_iter}）：每次内层收敛后按
控制电压 $V_{ctl}=|V_{mon}|/\mathrm{PTRatio}+(I_{tap}/I_{CT})(R_{ldc}+jX_{ldc})$
与死区 $\texttt{vreg}\pm\texttt{band}/2$ 决定升/降一档（:runpf\_phase），
档位组合重复即判振荡失败（:runpf\_phase）。覆盖边界如实说明：仅支持
\textbf{单相}调压器、接地星--接地星零相移变压器、\fld{max\_tap\_change}=1，
其余契约一律 fail-close 并记入 \fld{stop\_reason}（:runpf\_phase）。

\begin{paramtable}
\fld{max\_iter} & $k_{\max}$ & 次 & 30--100 & 50 & 内层 NR 迭代上限（include/hacdcpf/power\_flow/distribution\_power\_flow.hpp:solve\_three\_phase\_distribution\_pf）\\
\fld{max\_control\_iter} & --- & 次 & 20--100 & 100 & 调压器档位外环上限（:solve\_three\_phase\_distribution\_pf）\\
\fld{tol} & $\varepsilon$ & pu & $10^{-8}\sim10^{-4}$ & $10^{-6}$ & $\max|\Delta P|,|\Delta Q|$ 失配容差（:solve\_three\_phase\_distribution\_pf）\\
\fld{verbose} & --- & bool & --- & false & 逐迭代打印残差（:ThreePhaseNROptions）\\
\fld{include\_shunts} & --- & bool & --- & true & 是否在 $Y_{bus}$ 中含母线并纳（:ThreePhaseNROptions）\\
\end{paramtable}

\subsection{ZIP 与 Y/$\Delta$ 负荷相域建模}

\paragraph{ZIP 权重} 权重顺序为\textbf{恒阻抗、恒电流、恒功率}（Z、I、P）：
内部结构 \texttt{LoadZipWeights\{z,i,p\}}（src/power\_flow/three\_phase\_nr.cpp:LoadNeutralKind），
legacy 共享字段 \fld{const\_z\_percent}/\fld{const\_i\_percent}/\fld{const\_p\_percent}
默认 $0/0/100$（纯恒功率，\srcpath{include/hacdcpf/model/ac\_components.hpp:ThreePhaseRegulatorControl}）；
另有不与 legacy 混用的 P/Q 分离六字段（默认 $-1$ 表示未设，
include/hacdcpf/model/ac\_components.hpp:ThreePhaseLoad）。解析规则
（\srcpath{resolve\_active\_zip\_weights}/\srcpath{resolve\_reactive\_zip\_weights}，
src/power\_flow/three\_phase\_nr.cpp:build\_phase\_network\_model）：显式六字段必须全设或全不设，且不得与非默认
legacy 权重混用，否则抛错；任一权重为负或三者之和 $\neq100\%$ 抛错
（:build\_phase\_network\_model）。OpenDSS \texttt{ZIPV} 导入同样按 Z-I-P 顺序映射
（src/power\_flow/distribution\_power\_flow.cpp:load\_three\_phase\_system\_from\_opendss）。

\paragraph{电压窗与 cutoff} 恒 Z 部分折算为导纳在装配期\textbf{直接并入}
$Y_{bus}$（\srcpath{stamp\_prepared\_constant\_impedance\_loads}，:3398--3418）；
恒 I、恒 P 部分为动态项，按窗 $[\fld{vmin\_pu},\fld{vmax\_pu}]$ 分段
（\srcpath{evaluate\_dynamic\_load\_power}，:3444--3485，默认窗 $0.95\sim1.05$，
include/hacdcpf/model/ac\_components.hpp:ThreePhaseRegulatorControl）：
\[
S_{dyn}(v)=
\begin{cases}
0, & v<\fld{zipv\_cutoff\_pu}\ (\text{且 cutoff}>0)\\[2pt]
S_I^{nom}\dfrac{v^2}{v_{min}}+S_P^{nom}\dfrac{v^2}{v_{min}^2}, & v<v_{min}\\[6pt]
S_I^{nom}\,v+S_P^{nom}, & v_{min}\le v\le v_{max}\\[4pt]
S_I^{nom}\dfrac{v^2}{v_{max}}+S_P^{nom}\dfrac{v^2}{v_{max}^2}, & v>v_{max}
\end{cases}
\]
即窗外恒 I/恒 P 均退化为二次电压依赖（防止低压下电流发散），
\fld{zipv\_cutoff\_pu}（默认 0，即不启用）以下负荷整体退出。

\paragraph{Y 接与 $\Delta$ 接} Y 接负荷区分实接地/浮中性点/阻抗接地
（\srcpath{resolve\_load\_neutral\_kind}，:3078--3088；中性阻抗按母线基准折算，
:3138--3155）：浮中性点与阻抗接地时，每次求值以局部 $2\times2$ Newton 解
中性点电位（\srcpath{solve\_local\_neutral\_state}，:3739），雅可比经中性点
线性化链式求导（:IslandInfo）。$\Delta$ 接负荷拆为相间支路（abc 三相时
ab/bc/ca 三条，两相时合并，:3375--3392），其恒阻抗部分按线电压基准折算
$y_{branch}=S_Z^{*}/3$（:IslandInfo），动态项以线电压除以 $\sqrt3$ 进窗
（:IslandInfo），电流 $I=(S_{dyn}/V_{LL})^{*}$ 注入两端节点（:IslandInfo）。

\begin{paramtable}
\fld{connection} & --- & --- & wye/delta & wye & 接线方式，其余取值抛错（include/hacdcpf/model/ac\_components.hpp:ThreePhaseRegulatorControl；:3329--3333）\\
\fld{grounded}/\fld{r\_neut\_ohm}/\fld{x\_neut\_ohm} & $Z_n$ & $\Omega$ & --- & true/0/0 & Y 接中性点接地方式与阻抗（:ThreePhaseRegulatorControl）\\
\fld{vmin\_pu}/\fld{vmax\_pu} & $v_{min},v_{max}$ & pu & 0.9--1.1 & 0.95/1.05 & ZIP 电压窗（:ThreePhaseRegulatorControl）\\
\fld{zipv\_cutoff\_pu} & $v_{cut}$ & pu & 0--0.9 & 0.0 & 低于该值动态负荷退出（:1175；OpenDSS 导入取 0.65，src/io/case\_builders.cpp:add\_pv（lambda））\\
\fld{const\_z/i/p\_percent} & $w_z,w_i,w_p$ & \% & 0--100 & 0/0/100 & legacy 共享 ZIP 权重（Z-I-P 顺序，:1177--1179）\\
\fld{p\_/q\_const\_z/i/p\_percent} & --- & \% & 0--100 & $-1$ & P/Q 分离显式权重，须六字段全设（:add\_ac\_storage（lambda））\\
\end{paramtable}

\subsection{紧凑法与定点（fixed-point）法}

\compmeta{analysis::solve\_three\_phase\_compact\_pf / solve\_three\_phase\_fixed\_point}{\srcpath{include/hacdcpf/power\_flow/distribution\_power\_flow.hpp:build\_compact\_pf\_data}}

相域家族另有两个\textbf{电流失配不动点}求解器（实现同在
\srcpath{src/power\_flow/three\_phase\_nr.cpp}），由
\fld{PhaseDomainSolverAlgorithm} 枚举选择（Newton/FixedPoint/Compact/OpenDSS，
include/hacdcpf/power\_flow/distribution\_power\_flow.hpp:ThreePhaseDPFResult），经 \srcpath{runpf\_phase} 统一分发
（src/power\_flow/distribution\_power\_flow.cpp:find\_bus（lambda））。把节点分为固定电压集 $f$
（SLACK 与电源内节点）与变量集 $v$，不动点迭代为
\[
V_v^{(k+1)}\ \longleftarrow\ Y_{vv}^{-1}\Big(I\big(V^{(k)}\big)-Y_{vf}V_f\Big),
\qquad I_i=\Big(\frac{S_i^{spec}}{V_i}\Big)^{*}+I_i^{load}(V)+I_i^{fixed}
\]
电流注入由 \srcpath{evaluate\_fixed\_point\_current\_injections} 求值
（:4842--4895，恒 PQ、Y/$\Delta$ 动态负荷与 Norton 源同式）。初值为一次
稀疏 LU 解出的\textbf{空载电压}（compact :6014--6024；fixed-point :6176--6184），
显式电压猜测可覆盖（:runpf\_phase）。收敛判据为 $\max_i|\Delta V_i|<\texttt{tol}$。

两变体差异（three\_phase\_nr.cpp）：
\begin{itemize}
  \item \textbf{紧凑法} \srcpath{solve\_three\_phase\_compact\_pf}（:capture\_tap\_state（lambda））：
        在紧凑相节点空间求解，欠阻尼更新
        $V\leftarrow\alpha\hat V+(1-\alpha)V$，$\alpha$ 按迭代数 0.3/0.5/0.8 分级
        （:capture\_tap\_state（lambda））；要求 $k>1$ 才允许判收敛（:capture\_tap\_state（lambda））；
  \item \textbf{全维定点法} \srcpath{solve\_three\_phase\_fixed\_point}（:capture\_tap\_state（lambda））：
        在全 $3N$ 节点空间求解（\srcpath{build\_full\_phase\_matrix\_model}，:4997），
        带\textbf{负荷延拓}：按 $\{0.05,0.1,0.2,0.3,0.5,0.7,1.0\}$ 逐级加载
        （:build\_full\_ybus\_phase\_entries），中间级迭代上限 $\min(500,\texttt{max\_iter})$（:build\_full\_ybus\_phase\_entries），
        阻尼 $\alpha$ 分四级 0.2/0.3/0.5/0.7（:validate\_fixed\_point\_contract），电压幅值钳制
        $v_{min}$ 随延拓级 0.3/0.2/0.1 放松、$v_{max}=1.5$（:validate\_fixed\_point\_contract）；
        任一级不收敛即整体失败并报告最后收敛级（:apply\_fixed\_point\_voltage\_validity）。
\end{itemize}
$Y_{vv}$ 分解失败时先试对角平移 $10^{-8}$ 再试 $10^{-4}$
（\srcpath{factorize\_compact\_matrix}，:5270--5285，Eigen SparseLU）。

\begin{paramtable}
\fld{max\_iter} & $k_{\max}$ & 次 & 50--500 & 100 & 不动点迭代上限（include/hacdcpf/power\_flow/distribution\_power\_flow.hpp:ThreePhaseCompactPFData）\\
\fld{tol} & $\varepsilon$ & pu & $10^{-8}\sim10^{-4}$ & $10^{-6}$ & $\max|\Delta V|$ 收敛容差（:ThreePhaseCompactPFData）\\
\fld{verbose} & --- & bool & --- & false & 打印迭代信息（:ThreePhaseCompactPFData）\\
\fld{include\_shunts} & --- & bool & --- & true & 是否含母线并纳（:ThreePhaseFixedPointOptions）\\
\fld{vmin\_pu} & --- & pu & --- & 0.95 & \textbf{注意}：由 \srcpath{runpf\_phase} 赋值（src/power\_flow/distribution\_power\_flow.cpp:find\_bus（lambda）），但两求解器内均未读取，实际无效（诚实口径）\\
\end{paramtable}

\subsection{三相混合潮流（相域 AC + DC 联立）}

\compmeta{powerflow::solve\_three\_phase\_hybrid\_pf}{\srcpath{include/hacdcpf/power\_flow/three\_phase\_hybrid.hpp:solve\_three\_phase\_hybrid\_pf}}

这是与第~\ref{sec:hybrid}~章正序混合 Newton 不同的另一条路径：AC 侧为
\textbf{相域复数电压}（非极坐标），DC 侧为电压幅值，二者与换流器控制变量在
同一方程组、同一雅可比中联立（实现
\srcpath{src/power\_flow/three\_phase\_hybrid.cpp:solve\_three\_phase\_hybrid\_pf}）。接口为\textbf{矩阵级 API}：
调用方自备 \fld{y\_ac}、\fld{i\_ac\_fixed}、\fld{p\_load\_pu}/\fld{q\_load\_pu}、
\fld{g\_dc} 与换流器列表（\fld{ThreePhaseHybridPFCase}，include/hacdcpf/power\_flow/three\_phase\_hybrid.hpp:ThreePhaseHybridPFCase）。
状态布局 $x=[e;\ f;\ u;\ c]$（AC 电压实部、虚部、DC 电压、控制变量，
:135--139）。残差（\srcpath{evaluate}，:278--331）：
\begin{align*}
\text{AC 节点：}\quad & \Re/\Im\{\,v_i\,\mathrm{conj}[(Y_{ac}v+i_{fixed})_i]\,\}
  +p_i^{load}+jq_i^{load}\quad(\text{参考节点改为 }v_i=v_i^{ref})\\
\text{DC 节点：}\quad & u_i\,(G_{dc}u)_i+p_{dc,i}^{load}+\eta\sum_{\phi}P_{cv,\phi}\\
\text{控制方程：}\quad & u_{dc}-v_{dc}^{set}=0
\end{align*}
换流器六种控制模式（include/hacdcpf/power\_flow/three\_phase\_hybrid.hpp:ThreePhaseHybridPFControlMode，校验 :46--60, 96--134）：
\begin{itemize}
  \item \textbf{等分 PQ / 等分 Vdc-Q}：每相 $S=(P_{set}/3,\ Q_{set}/3)$，
        Vdc-Q 时 $P$ 由控制变量解出（:solve\_three\_phase\_hybrid\_pf）；
  \item \textbf{GFL-PQ / GFL-VdcQ}：跟网型正序电流注入
        $V_1=\frac13\sum_\phi a^{\phi}v_\phi$，$I_a^{*}=S_{set}/(3V_1)$，
        各相电流随正序旋转对称分配（:solve\_three\_phase\_hybrid\_pf）；$|V_1|$ 低于
        \fld{minimum\_positive\_sequence\_voltage\_pu} 即拒绝求值；
  \item \textbf{GFM-Voltage / GFM-Vdc}：构网型经虚拟阻抗
        $z_v=\fld{virtual\_r\_pu}+j\fld{virtual\_x\_pu}$ 注入
        $s_\phi=v_\phi\,\mathrm{conj}(e_\phi-v_\phi)/z_v^{*}$，
        $e_\phi$ 为内电势 $E_1$ 的正序旋转（:solve\_three\_phase\_hybrid\_pf）；GFM-Vdc 的内电势角为
        控制变量。
\end{itemize}
约束：每台换流器恰好三相；每个 DC 端子至多一台电压控制换流器；存在 DC 网络
时至少一台（:solve\_three\_phase\_hybrid\_pf）。求解为阻尼 Newton：Eigen SparseLU（COLAMD）分解
雅可比（:solve\_three\_phase\_hybrid\_pf），Armijo 回溯线搜索（merit $\tfrac12\lVert r\rVert^2$，
接受条件 $\le(1-c_1\alpha)\cdot$merit，$\alpha$ 逐次减半，:480--495）。
\textbf{诚实审计}：收敛后独立复算末态残差并分类报告
\fld{ac\_active}/\fld{ac\_reactive}/\fld{dc\_balance}/\fld{converter\_coupling}，
末态残差仍超容差则翻转为不收敛（:solve\_three\_phase\_hybrid\_pf）。

\begin{paramtable}
\fld{max\_iterations} & $k_{\max}$ & 次 & 50--200 & 80 & Newton 迭代上限（include/hacdcpf/power\_flow/three\_phase\_hybrid.hpp:ThreePhaseHybridPFOptions）\\
\fld{max\_line\_search\_steps} & --- & 步 & 8--30 & 18 & Armijo 回溯上限（:ThreePhaseHybridPFOptions）\\
\fld{tolerance} & $\varepsilon$ & pu & $10^{-10}\sim10^{-6}$ & $10^{-9}$ & $\lVert r\rVert_\infty$ 收敛容差（:ThreePhaseHybridPFOptions）\\
\fld{minimum\_dc\_voltage\_pu} & $\underline u$ & pu & --- & 0.05 & DC 电压下限，触及即中止（:60；:300）\\
\fld{minimum\_positive\_sequence\_voltage\_pu} & $\underline V_1$ & pu & --- & $10^{-6}$ & GFL 正序电压门槛（:61；:209）\\
\fld{armijo} & $c_1$ & --- & $10^{-4}$ & $10^{-4}$ & Armijo 充分下降系数（:ThreePhaseHybridPFOptions）\\
\fld{efficiency} & $\eta$ & --- & 0.9--1.0 & 0.98 & 换流器效率，耦合项 $\eta P_{ac}$（:25；:326）\\
\fld{virtual\_r\_pu}/\fld{virtual\_x\_pu} & $z_v$ & pu & --- & 0.01/0.10 & GFM 虚拟阻抗（:ThreePhaseHybridPFConverter）\\
\end{paramtable}

\subsection{与正序模型的关系与回退}

\begin{itemize}
  \item \textbf{正序投影（向上回退）}：
        \srcpath{project\_three\_phase\_if\_needed}
        （\srcpath{src/model/network\_utils.cpp:add\_equivalent\_storage\_from\_mobile\_storage}，投影管线调用点 :2177）在
        \fld{ac.buses} 为空时把 \fld{three\_phase\_ac} 折叠为正序单相等值：
        负荷按相求和（:project\_chargers\_into\_stations）、电压取三相平均（:project\_chargers\_into\_stations）、线路仅取正序
        $r_1,x_1,b_1$（:676--678，零序量保留在 \fld{r0\_pu} 等字段供短路分析）、
        变压器折叠为 \texttt{Transformer2W}（:project\_three\_phase\_if\_needed）。折叠后即可走
        第~\ref{sec:newton}~章统一 Newton——这也是仅有 \fld{three\_phase\_ac} 的
        算例调用 \srcpath{solve\_power\_flow} 时的隐式路径；相间耦合与不平衡
        信息在折叠中\textbf{丢弃}，属近似口径。
  \item \textbf{GUI 分阶段混合（staged hybrid）}：\texttt{three\_phase} 路由（\srcpath{tests/run\_gui\_server.cpp}，:16904--16990）
        先以统一 Newton 解
        正序混合边界，再把 VSC 传输功率\textbf{三相平衡}注入相域 AC 网后求解
        （:cell\_range（lambda））；返回中如实声明
        \fld{coupling=one\_way\_vsc\_boundary\_injection} 与
        \fld{monolithic\_abc\_dc\_jacobian=false}（:Group）——即单向边界、
        非联立。真正联立的相域 AC+DC 仅由上节矩阵级 API（GUI
        \texttt{three\_phase\_hybrid} 路由组装调用，:16820 附近）提供。
\end{itemize}

\subsection{验证与校核}

\begin{itemize}
  \item \textbf{OpenDSS IEEE13 对照}：\srcpath{tests/test\_phase\_domain\_ieee13.cpp:main}
        载入官方 \texttt{IEEE13Nodeckt.dss}，紧凑法与 OpenDSS 官方快照逐相节点
        比对（$\ge30$ 个比对点）：$\max|\Delta V_m|<0.002$~pu、
        $\max|\Delta\theta|<0.2^\circ$；
  \item \textbf{OpenDSS IEEE123 对照}：\srcpath{tests/test\_phase\_domain\_ieee123.cpp:main}：
        $\max|\Delta V_m|<0.001$~pu、$\max|\Delta\theta|<0.1^\circ$、$\ge270$ 点；
        另以 step8 组合算例交叉验证定点法与 OpenDSS YMatrix 路径（对紧凑法
        容差均为 2\%，:97--132）。两对照测试需 \texttt{HACDCPF\_HAVE\_OPENDSS}
        且数据目录在位，否则不注册（\srcpath{tests/CMakeLists.txt:236--252}）——
        默认构建下可能整体跳过，属覆盖薄弱点；
  \item \textbf{单元/集成测试}：\srcpath{tests/test\_three\_phase\_nr.cpp}（平衡
        三相电压相等至 $10^{-5}$、环网、多源 PV、平衡算例 VUF 为 0、分相损耗和容差
        $10^{-6}$、两节点精确功率平衡 :518）；\srcpath{test\_three\_phase\_nr\_load\_models.cpp}
        （ZIP 顺序语义、电压窗、cutoff、浮中性点/阻抗接地、$\Delta$ 接线、
        非法契约 fail-close）；\srcpath{test\_three\_phase\_nr\_transformer\_source.cpp}
        （YNyn/Dyn/YNd/Yd/Dd 等 20 余接线组、零序磁化、外部电网 Norton）；
        \srcpath{test\_three\_phase\_nr\_regulator.cpp}（单相档位闭环、三相分调
        银行、远程 PTRatio+LDC、振荡与限值 fail-close）；
  \item \textbf{混合潮流}：\srcpath{tests/test\_three\_phase\_hybrid\_pf.cpp:options}
        两算例（GFL-VdcQ、GFM-Vdc）：残差与分类平衡残差 $<10^{-9}$，
        $V_{dc}$ 钉在设定值至 $10^{-10}$，GFM 下电流不平衡度 $>10^{-5}$
        如实体现；
  \item \textbf{正序投影回归}：\srcpath{tests/test\_three\_phase\_pf.cpp:TEST\_CASE}
        投影后经 \srcpath{solve\_power\_flow} 收敛且 slack 电压保持；
  \item 已知边界：相域 NR 无 PV 无功限值切换与分布式松弛（PV 仅固定 $V_m$）；
        调压器仅单相、接地星--接地星零相移、单档步进；紧凑/定点法的基准
        覆盖弱于 NR（step8 容差仅 2\%）；混合潮流负荷仅恒定 PQ，无 ZIP。
\end{itemize}

% 第 9 章：电压稳定与连续潮流
% 本章文件由 \input 引入，不含导言区。
\section{电压稳定与连续潮流}\label{sec:cpf}

\subsection{功能定位与入口}

\compmeta{powerflow::CpfSolver}{\srcpath{include/hacdcpf/power\_flow/voltage\_stability.hpp}}

连续潮流（Continuation Power Flow，CPF）从基态运行点（$\lambda=0$）出发，
沿给定的负荷/发电增长方向逐步加大系统负荷水平，追踪 P--V 鼻曲线。新版实现
默认采用\textbf{弧长（arc-length）延拓}：延拓参数 $\lambda$ 在第一个延拓点
之后成为增广方程组的未知量，求解器因而可以\textbf{翻越鼻点}并采样 P--V 曲线
下支，而不仅仅是从上支逼近崩溃点。\textbf{回退分支（诚实口径）}：弧长增广
校正要求方程布局固定；只要存在任一在运 PV 母线，求解器即自动回退自然参数化
（src/power\_flow/voltage\_stability.cpp:solve），此时只能从上支逼近鼻点，结果以
\fld{model\_scope}=\texttt{cpf:natural-parameterization+pv-pq-active-set}、
\fld{model\_limitations} 与 \fld{warnings} 如实声明（:solve），并以
\fld{arc\_length\_used}=false 可机读地区分（:328--329；另两种口径为
\texttt{cpf:arc-length-fixed-active-set} 与 \texttt{cpf:natural-parameterization}，
:338--342）。入口签名
（include/hacdcpf/power\_flow/voltage\_stability.hpp:CpfSolver；实现
\srcpath{src/power\_flow/voltage\_stability.cpp:solve}）：
\begin{quote}\footnotesize
\texttt{CpfResult solve(const SolverData\& base\_data,\\
\hspace*{2em}const CpfDirection\& direction) const;}
\end{quote}
求解器选项挂在公开成员 \fld{CpfSolver::opts}（\fld{CpfOptions}，
include/hacdcpf/power\_flow/voltage\_stability.hpp:CpfSolver）上。\textbf{接入口径（诚实声明）}：
\fld{CpfSolver} 位于 \texttt{hacdcpf::powerflow} 命名空间，
\textbf{未挂接}公共门面 \srcpath{include/hacdcpf/api/hacdcpf.hpp}，
GUI 后端 \srcpath{tests/run\_gui\_server.cpp} 亦无对应 \texttt{method} 分发；
调用方需自行构造 \fld{SolverData}（经 \srcpath{make\_solver\_data}，
见第~\ref{sec:overview}~章）后以库内 API 方式使用。能力表条目见
\srcpath{README.md:204}；求解器工厂枚举 \fld{PowerFlowMethod::Continuation}
虽已声明（\srcpath{include/hacdcpf/power\_flow/solver\_factory.hpp:PowerFlowMethod}），
但与第~\ref{sec:overview}~章所述抽象层一样未接线到生产路径。

\subsection{参数化模型与负荷方向}
CPF 把注入写为延拓参数 $\lambda$ 的仿射函数
（头文件注释 include/hacdcpf/power\_flow/voltage\_stability.hpp:CpfDirection；实现
\srcpath{apply\_direction}，src/power\_flow/voltage\_stability.cpp:apply\_direction）：
\begin{align*}
P_{L,i}(\lambda)&=P_{L,i}^{0}+\lambda\,\Delta P_{L,i},&
Q_{L,i}(\lambda)&=Q_{L,i}^{0}+\lambda\,\Delta Q_{L,i},&
P_{G,i}(\lambda)&=P_{G,i}^{0}+\lambda\,\Delta P_{G,i}\ \ (\text{非 slack 机组})
\end{align*}
其中 $\lambda=0$ 对应基态；采用比例方向时 $\lambda=1$ 即全部负荷翻倍。
增长方向由 \fld{CpfDirection}（include/hacdcpf/power\_flow/voltage\_stability.hpp:CpfDirection）携带，
三个向量 \fld{dp\_load}/\fld{dq\_load}/\fld{dp\_gen} 均按 AC 母线 0-based
位置索引，单位为 MW/Mvar（与 SolverData 物理单位一致，非标幺）；
长度必须等于 AC 母线数，否则抛 \texttt{std::invalid\_argument}
（src/power\_flow/voltage\_stability.cpp:solve）。两个工厂方法：
\begin{itemize}
  \item \srcpath{CpfDirection::proportional}（cpp:25--63）：各母线负荷增量取
        基态值（:proportional）；非 slack 在役机组按基态出力份额分摊
        \textbf{总负荷}增量，即
        $\Delta P_{G,b(g)}=\dfrac{P_{G,g}^{0}}{\sum_{g'}P_{G,g'}^{0}}\sum_i P_{L,i}^{0}$
        （:proportional）， slack 机组吸收不平衡量；
  \item \srcpath{CpfDirection::at\_buses}（:at\_buses）：仅列表内母线参与，
        每母线均匀 $1$~MW 与 $1$~Mvar 增量，\fld{dp\_gen} 全零
        （发电增量全部由 slack 承担）。
\end{itemize}
\textbf{组件聚合负荷口径}：基态负荷 $P_{L,i}^{0}$/$Q_{L,i}^{0}$ 的取值尊重
SolverData 的组件聚合结果——\fld{has\_component\_loads} 为真时取
\fld{pd\_pu}/\fld{qd\_pu}$\times$\fld{base\_mva}（即经组件负荷
\texttt{scaling} 聚合后的标幺值折算），否则取 \fld{ac\_buses} 的
\fld{pd\_mw}/\fld{qd\_mvar}（proportional 于 cpp:34--39，
\srcpath{total\_load} 于 :139--144 同口径）。相应地，
\srcpath{apply\_direction} 在施加方向时同步更新 \fld{pd\_pu}/\fld{qd\_pu}
（:apply\_direction），保证后续残差求值与聚合路径看到一致的负荷水平。
每个校正步把方向施加到基态副本上（:apply\_direction），随后调用
\srcpath{aggregate\_generation}（:105；实现
\srcpath{src/power\_flow/solver\_data.cpp:aggregate\_generation}）把机组出力重新聚合进
\fld{pg}/\fld{qg} 注入向量——\textbf{因此 CPF 每个 $\lambda$ 取值都是一次
完整的、独立的 SolverData 潮流问题}。

\subsection{弧长延拓预测--校正算法与自适应步长}
默认模式（\fld{enable\_arc\_length}=true 且无在运 PV 母线；有在运 PV 母线时
自动回退自然参数化，见第~\ref{sec:cpf}~章功能定位节）的主流程
（src/power\_flow/voltage\_stability.cpp:solve）如下：
\begin{enumerate}
  \item \textbf{基态求解}：先在 $\lambda=0$ 跑一次统一 NR（:run\_nr（lambda）），
        不收敛即以 \texttt{"base case (lambda=0) did not converge"} 终止
        （:solve）；
  \item \textbf{首个延拓点（自然参数化）}：以 $\lambda=\texttt{step\_init}$
        施加方向后用普通 NR 求第二个点，失败则按 \fld{step\_shrink} 收缩
        步长重试直至 \fld{step\_min}（:solve）；仍不收敛则以
        \texttt{"no continuation step converged from base case"} 终止
        （:397--398，\textbf{不置} \fld{nose\_found}）。该点与基态共同确定
        第一条割线切向量；
  \item \textbf{增广状态空间}：此后所有步在增广状态
        \[
        y=\bigl[\theta_{\text{非 slack}},\;V_{m,\text{PQ}},\;
        V_{dc,\text{非 slack}},\;\lambda\bigr]\in\mathbb{R}^{n+1}
        \]
        中工作，$\lambda$ 成为未知量。状态布局由
        \srcpath{make\_state\_layout}（:make\_state\_layout）经
        \srcpath{classify\_ac\_buses} 与 DC slack 判定生成；
        打包/解包见 \srcpath{pack\_augmented\_state}（:pack\_augmented\_state）与
        \srcpath{unpack\_augmented\_state}（:190--213，解包时
        $V_m$、$V_{dc}$ 钳制下限 $0.02$~pu，:211--212；相角以弧度计，
        :203--204）；
  \item \textbf{预测步（割线切向量）}：取最近两个被接受增广解点的归一化
        割线作为切向量 $t=(y_{\mathrm{curr}}-y_{\mathrm{prev}})/
        \|y_{\mathrm{curr}}-y_{\mathrm{prev}}\|$，并与上一步切向量做符号
        一致性校正（$t\cdot t_{\mathrm{prev}}<0$ 时反号，:484）；
        预测点 $y_{\mathrm{pred}}=y_{\mathrm{curr}}+s\,t$（:solve）。
        弧长步长 $s$ 初值取首段割线长度并钳入
        $[\texttt{step\_min},\texttt{step\_max}]$（:464；首段割线零范数或
        非有限时以 \texttt{"arc-length initial secant is zero"} 终止，
        :460--463）；
        \fld{use\_secant\_predictor}=false 时复用上一步切向量
        （定方向预测，:474）；割线零范数或非有限（坍缩）时以
        \texttt{"arc-length secant collapsed to zero"} 终止（:solve）。
        尚未进入下支时若预测 $\lambda$ 超过
        \fld{lambda\_max} 即以 \texttt{"lambda\_max reached"} 终止
        （:solve）；
  \item \textbf{校正步（弧长超平面）}：\srcpath{arc\_length\_correct}
        （:arc\_length\_correct）在预测点处求解增广系统
        \[
        F(y)=0,\qquad t^{\top}(y-y_{\mathrm{pred}})=0,
        \]
        其中 $F$ 由残差求值器 \srcpath{evaluate\_power\_flow\_residual}
        在给定 $(v_m,v_a,v_{dc})$ 与参数化 SolverData 下求值
        （:evaluate\_cpf\_residual）。Newton 迭代用\textbf{有限差分}构造稠密
        $(n+1)\times(n+1)$ 雅可比（扰动步
        $h=10^{-6}\max(1,|y_j|)$，:243--253；最后一行为切向量 $t^{\top}$，
        :253），以 \texttt{Eigen::FullPivLU} 求解（:arc\_length\_correct），并带
        \textbf{回溯线搜索}（最多 10 次二分，要求增广残差 max 范数下降，
        :260--278）。收敛判据为增广残差 max 范数 $<\,$\fld{corrector\_tol}
        （:arc\_length\_correct），迭代上限 \fld{corrector\_max\_iter}（:arc\_length\_correct）；
  \item \textbf{接受、鼻点检测与下支停止}：校正收敛后先查电压下限
        （任一母线 $v_m<\,$\fld{vm\_min\_pu} 即以
        \texttt{"voltage below vm\_min\_pu"} 终止，:505--510；
        \textbf{注意此分支不置 \fld{nose\_found}}）；随后入 trace
        （:record\_state（lambda））。当接受点的 $\lambda$ 首次低于上一点（$\lambda$ 切向
        符号反转）时判定已翻越鼻点，置 \fld{nose\_found}=true 并进入
        下支计数（:solve）；再接受 \fld{lower\_branch\_steps} 个下支
        点后以
        \texttt{"nose point passed by arc-length continuation"} 停止
        （:solve）。校正失败则 $s\leftarrow\texttt{step\_shrink}\cdot s$
        重试，$s<\texttt{step\_min}$ 时以
        \texttt{"arc-length corrector step below step\_min"} 终止
        （:solve）；成功则
        $s\leftarrow\min(\texttt{step\_grow}\cdot s,\texttt{step\_max})$
        （:solve）。步数耗尽时报 \texttt{"max\_steps reached"}（:solve）。
\end{enumerate}
\textbf{legacy 模式}：\fld{enable\_arc\_length}=false 时保留旧的自然参数化
实现（:solve）——$\lambda$ 固定递增、上一解点暖启动 NR 校正、成败倍率
步长；步长跌破 \fld{step\_min} 时以
\texttt{"natural parameterization stalled; no lambda fold was observed"}
终止（:443--449，\textbf{不置} \fld{nose\_found}），该模式只能
从上支逼近鼻点，供需要严格递增 $\lambda$ 采样的调用方使用。

终止原因（\fld{CpfResult::termination\_reason}）全集：
\texttt{"empty system"}（:solve）、
\texttt{"base case (lambda=0) did not converge"}（:solve）、
\texttt{"no continuation step converged from base case"}（:solve）、
\texttt{"lambda\_max reached"}（:489，legacy 模式
:457--458）、\texttt{"arc-length initial secant is zero"}（:solve）、
\texttt{"arc-length secant collapsed to zero"}（:solve）、
\texttt{"arc-length corrector step below step\_min"}（:solve）、
\texttt{"voltage below vm\_min\_pu"}（:solve）、
\texttt{"nose point passed by arc-length continuation"}（:solve）、
legacy 模式另有
\texttt{"natural parameterization stalled; no lambda fold was observed"}
（:solve）、兜底 \texttt{"max\_steps reached"}（:solve）。
全集之中仅弧长模式翻越鼻点一分支置 \fld{nose\_found}=true（:solve），
其余终止分支均不置位。

\textbf{能力边界（诚实口径）}：
\begin{itemize}
  \item 弧长模式下求解器\textbf{可以翻越鼻点}，但下支仅保留
        \fld{lower\_branch\_steps} 个采样点（默认 2）即主动停止——这是
        刻意的工程取舍（头文件注释 include/hacdcpf/power\_flow/voltage\_stability.hpp:CpfOptions：
        小正值使鼻点可观测而不深入低电压支），\textbf{不是}完整的下支
        追踪；需要更多下支点可手动调大该字段；
  \item 弧长校正器使用\textbf{有限差分稠密雅可比}+\texttt{FullPivLU}，
        单次 Newton 迭代的代价为 $n+2$ 次残差求值加一次稠密 LU，复杂度
        约 $O(n^3)$ 且不复用统一 NR 的稀疏解析雅可比与 PatternCache——
        因此该模式面向中小规模系统的电压稳定分析，大系统下校正步耗时
        会显著高于一次普通 NR；
  \item \fld{nose\_found} 仅在弧长模式\textbf{已观察到} $\lambda$ 符号反转
        （确实越过鼻点）时置位（:solve）；首步失败、legacy 停滞、
        电压下限停止（\texttt{"voltage below vm\_min\_pu"}）等其余终止分支
        一律\textbf{不}置 \fld{nose\_found}——对这些分支调用方应读
        \fld{termination\_reason} 而非 \fld{nose\_found}；
  \item 鼻点定位精度仍受 \fld{step\_min}/\fld{step\_max} 与
        \fld{corrector\_tol} 控制；弧长模式报告的是轨迹上 $\lambda$
        最大的采样点（见下节汇总口径），相邻采样点之间不做鼻点加密
        二分。
\end{itemize}
监控母线 \fld{monitor\_bus} 缺省自动选为基态有功负荷最大的母线
（基态负荷尊重组件聚合口径，:307--320）。

\subsection{参数选项}
\fld{CpfOptions}（include/hacdcpf/power\_flow/voltage\_stability.hpp:CpfOptions）全部字段与默认值如下
（行号见“说明”列）：

\begin{paramtable}
\fld{lambda\_max} & $\lambda_{\max}$ & --- & $1\sim10$ & 5.0 & 追踪的负荷因子上限，上支预测越过即停（:CpfOptions）\\
\fld{step\_init} & $\Delta\lambda_0$ & --- & $0.01\sim0.2$ & 0.05 & 初始延拓步长（首点自然参数化步长）（:CpfOptions）\\
\fld{step\_min} & $s_{\min}$ & --- & $10^{-6}\sim10^{-4}$ & $10^{-5}$ & 最小步长；跌破即停（弧长模式下为弧长步长）（:CpfOptions）\\
\fld{step\_max} & $s_{\max}$ & --- & $0.1\sim0.5$ & 0.20 & 步长上限，防止跨过鼻点特征（:CpfOptions）\\
\fld{step\_grow} & --- & --- & $1.2\sim2$ & 2.0 & 校正成功后的步长放大因子（:CpfOptions）\\
\fld{step\_shrink} & --- & --- & $0.3\sim0.7$ & 0.5 & 校正失败后的步长收缩因子（:CpfOptions）\\
\fld{corrector\_max\_iter} & --- & 次 & $20\sim100$ & 50 & 校正器迭代上限（弧长 Newton 与 legacy NR 复用同一字段）（:CpfOptions）\\
\fld{corrector\_tol} & --- & pu & $10^{-8}\sim10^{-6}$ & $10^{-8}$ & 校正收敛容差（弧长模式为增广残差 max 范数）（:CpfOptions）\\
\fld{monitor\_bus} & --- & 0-based & --- & $-1$ & P--V 曲线监控母线；$-1$=自动选最大有功负荷母线（:CpfOptions）\\
\fld{trace\_all\_buses} & --- & bool & --- & false & 每步是否保存全母线 \fld{vm}/\fld{va}/\fld{vdc}（:CpfOptions）\\
\fld{max\_steps} & --- & 步 & $500\sim5000$ & 2000 & 预测--校正总步数安全上限（:CpfOptions）\\
\fld{vm\_min\_pu} & --- & pu & $0.2\sim0.6$ & 0.25 & 电压下限；任一母线跌破即停（不置 \fld{nose\_found}）（:CpfOptions）\\
\fld{use\_secant\_predictor} & --- & bool & --- & true & 割线切向量预测；false=复用上步切向量（:CpfOptions）\\
\fld{enable\_arc\_length} & --- & bool & --- & true & 弧长增广校正（$\lambda$ 作为未知量，可翻越鼻点）；false=自然参数化 legacy 模式（:CpfOptions）\\
\fld{lower\_branch\_steps} & --- & 点 & $1\sim5$ & 2 & 翻越鼻点后保留的下支采样点数，达到即停（:CpfOptions）\\
\end{paramtable}

\subsection{结果结构与电压稳定指标}
每个被接受点记为 \fld{CpfPoint}（include/hacdcpf/power\_flow/voltage\_stability.hpp:CpfPoint）：
\fld{lambda}、监控母线电压 \fld{vm\_monitor}、该 $\lambda$ 下的系统总负荷
\fld{p\_total\_mw}/\fld{q\_total\_mvar}（由 \srcpath{total\_load} 按
$P^0+\lambda\Delta P$ 复算，src/power\_flow/voltage\_stability.cpp:total\_load），以及可选的
全母线 \fld{vm}/\fld{va}/\fld{vdc}（\fld{trace\_all\_buses} 开启时填充）。
\textbf{单位注意}：\fld{va} 为\textbf{弧度}（头文件注释
include/hacdcpf/power\_flow/voltage\_stability.hpp:CpfPoint；统一 NR 的 \fld{PowerFlowResult::va}
同样为弧度，\srcpath{src/power\_flow/newton\_solver.cpp:solve}），
\fld{vdc} 为新增的全 DC 母线电压（pu，:123--125；弧长模式下经
\srcpath{record\_state} 填充，cpp:410--432）。
汇总结果 \fld{CpfResult}（hpp:129--166）含 \fld{trace}、鼻点摘要
\fld{lambda\_max}/\fld{p\_max\_mw}/\fld{vm\_at\_nose}、\fld{nose\_found}、
\fld{termination\_reason}、监控母线 \fld{monitor\_bus}、诚实口径字段
\fld{q\_limits\_enforced}/\fld{arc\_length\_used}/\fld{model\_scope}/%
\fld{model\_limitations}/\fld{warnings}（hpp:148--159）与校正求解计数
\fld{total\_pf\_solves}。
\textbf{汇总口径}：三个鼻点摘要字段取 trace 中 $\lambda$ \textbf{最大}的
点（\texttt{std::max\_element}，cpp:537--544）而非末点——弧长模式翻越鼻点
后末点位于下支，此口径保证 $\lambda_{\max}$ 始终对应鼻点附近的最高负荷
水平。

指标族由自由函数 \srcpath{compute\_vsi}（声明 hpp:218--223；实现
cpp:553--579）从一次完成的 \fld{CpfResult} 派生，输出
\fld{VoltageStabilityIndex}（hpp:169--182）：
\begin{itemize}
  \item \textbf{最大负荷因子} \fld{lambda\_max} 与鼻点总负荷
        \fld{p\_max\_mw}（直接转录，:556--559）；
  \item \textbf{负荷裕度}
        $\fld{p\_margin\_mw}=P_{\max}-\sum_i P_{L,i}^{0}$（MW，:561--568；
        基态负荷尊重组件聚合口径：\fld{has\_component\_loads} 为真时取
        \fld{pd\_pu} 求和$\times$\fld{base\_mva}（:compute\_vsi），否则按
        \fld{ac\_buses} 的 \fld{pd\_mw} 求和（:compute\_vsi））；
  \item 监控母线鼻点电压 \fld{vm\_at\_nose} 与基态电压 \fld{vm\_base}
        （:570--576；trace 非空时以首点监控电压为准）。
\end{itemize}
\textbf{指标边界（诚实口径）}：当前仅提供\textbf{负荷裕度族}指标
（$\lambda_{\max}$、$P_{\max}$、$P_{\mathrm{margin}}$ 与监控母线电压），
\textbf{未实现} L 指标、模态分析/参与因子、Q--V 曲线等其他经典电压稳定
指标；且指标全部是标量化的系统级/单母线量，无逐母线弱节点排序。

\subsection{与统一 Newton 求解器及残差求值器的复用关系}
新版 CPF 的复用关系按模式分两层：
\begin{itemize}
  \item \textbf{NR 复用层}：基态（$\lambda=0$）、首个延拓点与 legacy
        模式的全部校正步\textbf{原样复用}第~\ref{sec:newton}~章的统一
        Newton 求解器：头文件 \srcpath{power\_flow/newton\_solver.hpp}
        为转发头，实际类为 \srcpath{engine::NewtonSolver}；
        \fld{CpfSolver::solve} 在栈上构造单个 \fld{NewtonSolver} 实例并
        在整个求解过程复用（src/power\_flow/voltage\_stability.cpp:solve），每次调用
        经 \srcpath{run\_nr}（:run\_nr（lambda））并累计 \fld{total\_pf\_solves}。
        校正选项由 \srcpath{make\_corrector\_opts}（:make\_corrector\_opts）生成：
        仅覆写 \fld{max\_iter}$\leftarrow$\fld{corrector\_max\_iter}、
        \fld{tol}$\leftarrow$\fld{corrector\_tol}，并强制
        \fld{enable\_pv\_pq\_conversion}=true；其余取
        \fld{PowerFlowOptions} 默认。初值
        \fld{InitialState\{vm, va, vdc\}}（定义
        \srcpath{include/hacdcpf/power\_flow/power\_flow\_options.hpp:InitialState}）
        由 \srcpath{make\_init} 自上一解点暖启动（:119--129，混合 AC/DC
        系统 DC 电压空时补 $1.0$~pu）；
  \item \textbf{残差求值层}：弧长校正器\textbf{不经过}统一 NR，而是直接
        调用独立残差求值器 \srcpath{evaluate\_power\_flow\_residual}
        （\srcpath{include/hacdcpf/power\_flow/residual\_evaluator.hpp}；
        调用点 cpp:215--223）在给定状态与参数化 SolverData 下复算全残差，
        再以有限差分稠密雅可比自洽求解（见上节）。每步校正经
        \srcpath{apply\_direction} 在栈上新建 \fld{SolverData} 副本
        （:apply\_direction）。
\end{itemize}
\textbf{性能提示}：
\begin{itemize}
  \item 弧长校正器单次 Newton 迭代需 $n+2$ 次全残差求值（构造差分列）
        加一次 $(n+1)\times(n+1)$ 稠密 \texttt{FullPivLU}，不利用稀疏性
        与解析雅可比；这是弧长模式相对 legacy 模式的主要额外开销；
  \item NR 复用层一侧，\fld{NewtonSolver} 的 PatternCache 以数据指针
        \fld{data\_ptr}、\fld{data\_build\_id} 与 Ybus/Gdc 存储指针为键
        （\srcpath{include/hacdcpf/engine/solver/native/nle/newton\_solver.hpp:NewtonSolver}；
        命中判定 \srcpath{src/power\_flow/newton\_solver.cpp:ensure\_pattern（lambda）}），
        而 CPF 每步在栈上新建 \fld{SolverData} 副本，故符号模式缓存逐步
        失效、雅可比上下文重建。NEW 中 \fld{NewtonSolver} 另持有每实例
        \fld{SolverWorkspace}（engine 头 :47--48；
        \srcpath{include/hacdcpf/power\_flow/workspace.hpp:SolverWorkspace} 的
        \fld{prepare\_state}/\fld{prepare\_equations}，复用逻辑
        src/power\_flow/newton\_solver.cpp:solve），可在维度不变时保留 Eigen 分配，
        但这不改变每步模式重建的事实；
  \item \fld{total\_pf\_solves} 只统计经 \srcpath{run\_nr} 的 NR 求解
        （基态、首点与 legacy 校正步；弧长校正步在 cpp:494 处单独累加），
        不能用它折算弧长模式的总耗时。
\end{itemize}
由于基态与首点就是标准 NR，第~\ref{sec:newton}~章的全部性质（混合 AC/DC
联立、多 slack 岛、PV$\leftrightarrow$PQ 切换）对 CPF 的起点同样成立；
弧长校正的残差求值器同样覆盖换流器耦合等 SolverData 特性。

\subsection{验证与校核}
\begin{itemize}
  \item \textbf{专项回归}：NEW 新增
        \srcpath{tests/test\_voltage\_stability.cpp}，已注册于
        \srcpath{tests/CMakeLists.txt:124}，含两个 Catch2 用例：
        \begin{enumerate}
          \item \textbf{弧长翻越鼻点}（:TEST\_CASE）：两母线电压稳定算例
                （\srcpath{make\_two\_bus\_voltage\_stability\_case}，
                \srcpath{tests/power\_flow\_solver\_test\_utils.hpp:make\_two\_bus\_voltage\_stability\_case}），
                设 \fld{lower\_branch\_steps}=3、\fld{trace\_all\_buses}=true，
                断言 \fld{nose\_found}、trace 不少于 5 点、
                $\lambda_{\max}\in(0.25,\,0.50)$、trace 中 $\lambda$ 最大点
                之后存在 $\lambda$ 回落的下支点、且
                \fld{vm\_at\_nose} 低于基态监控电压——直接验证了弧长模式
                “翻越鼻点并保留下支采样”的核心能力；
          \item \textbf{组件聚合负荷方向}（:make\_two\_bus\_voltage\_stability\_case）：向两母线系统追加一个
                带 \texttt{scaling}=0.5 的显式 \fld{Load} 组件，断言
                \fld{has\_component\_loads} 为真且
                \srcpath{CpfDirection::proportional} 给出的
                \fld{dp\_load}/\fld{dq\_load} 等于聚合后数值
                （$55$~MW/$22$~Mvar）——覆盖上节所述组件聚合口径。
        \end{enumerate}
        同批新增的高级求解器回归还有
        \srcpath{tests/test\_homotopy\_continuation.cpp} 与
        \srcpath{tests/test\_newton\_krylov.cpp}
        （\srcpath{tests/CMakeLists.txt:123,125}），第~\ref{sec:robust}~章
        所述同伦求解器与 Newton--Krylov 的“无专项回归”状态同步解除；
  \item \textbf{仍未覆盖（诚实口径）}：混合 AC/DC 系统的 CPF（\fld{vdc}
        轨迹与 DC slack 布局）、legacy 自然参数化模式、在运 PV 母线触发的
        弧长$\to$自然参数化回退分支（src/power\_flow/voltage\_stability.cpp:solve）、
        \fld{vm\_min\_pu} 电压下限停止分支、\srcpath{at\_buses} 方向、
        \srcpath{compute\_vsi} 指标派生以及多 slack 系统均无专项断言；
        纯 NR 求解器之外的 MATPOWER 大算例 CPF 基准也不存在。GUI 后端无
        CPF 分发，E2E 自然不覆盖；
  \item \textbf{间接保障}：基态与首点校正依赖统一 Newton 的既有验证
        （MATPOWER 85 算例回归、雅可比对差分与缩放回归，见
        第~\ref{sec:validation}~章）；弧长校正所依赖的残差求值器
        \srcpath{src/power\_flow/residual\_evaluator.cpp} 本身是独立致密
        残差复算通道，与 NR 内核互为对照。CPF 自身逻辑层约
        560 行（voltage\_stability.cpp 全文件 581 行）；
  \item \textbf{文档对照}：技术笔记 Algorithm A13
        （\srcpath{docs/archive/reference/technical\_notebook/sections/15\_algorithm\_index.tex:609--672}）
        的伪代码仍是\textbf{自然参数化}版本（割线预测+固定 $\lambda$ 的
        NR 校正、\fld{step\_min} 判鼻），与当前默认弧长实现\textbf{不一致}，
        仅对应 \fld{enable\_arc\_length}=false 的 legacy 模式；阅读时应以
        本章与源码为准。另注意技术笔记 05 节把 CPF 实现位置误记为
        \srcpath{src/power\_flow/homotopy\_continuation.cpp}
        （\srcpath{docs/archive/reference/technical\_notebook/sections/05\_advanced\_pf.tex:21}），
        实际实现为 \srcpath{src/power\_flow/voltage\_stability.cpp}，
        同伦求解器（收敛兜底用，见第~\ref{sec:robust}~章）与 CPF 是两套
        独立代码；
  \item \textbf{使用建议}：$\lambda_{\max}$ 结果对 \fld{step\_min}、
        \fld{step\_max} 与 \fld{vm\_min\_pu} 敏感；弧长模式下
        $\lambda_{\max}$ 取自轨迹中 $\lambda$ 最大的采样点而非末点，
        用于工程结论时建议收紧步长上下限重跑以确认鼻点定位的数值稳定性。
        需要完整下支曲线时应调大 \fld{lower\_branch\_steps} 并放宽
        \fld{vm\_min\_pu}，同时注意下支低电压解的物理解释需另行论证。
\end{itemize}

% 第 10 章：选项、结果与诊断参考
% 本章文件由 \input 引入，不含导言区。
\section{选项、结果与诊断参考}\label{sec:options}

\subsection{概述与阅读约定}

\compmeta{hacdcpf::PowerFlowOptions}{\srcpath{include/hacdcpf/power\_flow/power\_flow\_options.hpp:PowerFlowOptions}}

本章给出潮流计算全部对外选项与结果结构的\textbf{逐字段参考}，每个字段均与头文件
逐一核对，说明列注明源码读取点；默认值以 \srcpath{include/hacdcpf/model/defaults.hpp}
与头内初始化器为准，\texttt{src/} 中无读取点的字段一律标注“\textbf{当前未接线}”。

\texttt{PowerFlowOptions} 采用平铺布局（向后兼容），另提供 7 个分组子结构与静态
工厂 \srcpath{from\_parts}（src/api/hacdcpf.cpp:PowerFlowOptions::from\_parts）按组装配。主入口
\srcpath{solve\_power\_flow}（第~\ref{sec:overview}~章）把 ZIP 权重与三开关抄入
\texttt{SolverData}（src/api/hacdcpf.cpp:apply\_zip\_weights）。内部计算一律标幺值
（$S_{\mathrm{base}}=100$~MVA）；选项/结果的验证体系见第~\ref{sec:validation}~章。

\subsection{收敛控制}
\begin{paramtable}
\fld{max\_iter} & $K_{\max}$ & 次 & 10--100 & 50 & Newton 内层迭代上限；\texttt{Defaults::kPFMaxIter}（include/hacdcpf/model/defaults.hpp:Defaults）；读取点\srcpath{src/power\_flow/newton\_solver.cpp:NewtonSolver::solve}\\
\fld{tol} & $\varepsilon$ & pu & $10^{-10}$--$10^{-6}$ & $10^{-8}$ &缩放后 $\lVert F\rVert_\infty$ 收敛容差；\texttt{Defaults::kPFTol}（include/hacdcpf/model/defaults.hpp:Defaults）；判定点 src/power\_flow/newton\_solver.cpp:NewtonSolver::solve 与最终诚实重判 :1827--1829\\
\fld{fdpf\_max\_iter} & --- & 次 & 100--5000 & 1000 & FDPF 专用迭代上限，仅 \srcpath{src/power\_flow/fdpf\_solver.cpp:solve} 读取（Newton 不用）\\
\end{paramtable}

\subsection{步长截断与线性回退}
\begin{paramtable}
\fld{max\_delta\_va\_rad} & $\Delta\theta_{\max}$ & rad & 0.5--3.0 & 1.50 &单步相角截断；\srcpath{clip\_step}（src/power\_flow/newton\_solver.cpp:clip\_step）\\
\fld{max\_delta\_vm\_pu} & $\Delta V_{\max}$ & pu & 0.1--1.0 & 0.50 &单步 $V_m$ 截断（src/power\_flow/newton\_solver.cpp:clip\_step）\\
\fld{max\_delta\_vdc\_pu} & $\Delta V_{dc,\max}$ & pu & 0.1--1.0 & 0.50 &单步 $V_{dc}$ 截断（src/power\_flow/newton\_solver.cpp:clip\_step）\\
\fld{max\_line\_search\_steps} & --- & 次 & 4--20 & 10 & 回溯线搜索最大步数（src/power\_flow/newton\_solver.cpp:run\_line\_search（lambda），见第~\ref{sec:newton}~章回退链）\\
\fld{max\_regularization\_steps} & --- & 次 & 0--16 & 8 & 对角正则化$J+\lambda I$ 最大步数（src/power\_flow/newton\_solver.cpp:NewtonSolver::solve）\\
\fld{regularization\_lambda0} & $\lambda_0$ & --- & $10^{-8}$--$10^{-4}$ & $10^{-6}$ &正则化初值（src/power\_flow/newton\_solver.cpp:NewtonSolver::solve）\\
\fld{regularization\_growth} & --- & --- & 2--100 & 10.0 & 正则化增长因子（src/power\_flow/newton\_solver.cpp:NewtonSolver::solve）\\
\end{paramtable}

\subsection{全局化策略}
\begin{paramtable}
\fld{globalization} & --- & --- & 三选一 & LineSearch & 全局化策略枚举\texttt{GlobalizationStrategy\allowbreak\{LineSearch,\allowbreak TrustRegion,\allowbreak PseudoTransient\}}（include/hacdcpf/power\_flow/power\_flow\_options.hpp:GlobalizationStrategy）；分派点 src/power\_flow/newton\_solver.cpp:NewtonSolver::solve\\
\fld{trust\_region\_radius0} & $\Delta_0$ & --- & 0.1--5 & 1.0 & 信赖域初始半径（src/power\_flow/newton\_solver.cpp:NewtonSolver::solve）\\
\fld{trust\_region\_max} & $\Delta_{\max}$ & --- & 1--100 & 10.0 & 信赖域半径上限（src/power\_flow/newton\_solver.cpp:NewtonSolver::solve）\\
\fld{ptc\_delta0} & $\delta t_0$ & --- & 0.01--10 & 1.0 & 伪瞬态延拓初始步长；\textbf{保留字段，当前无消费点}（PTC 实际由 \fld{RobustNonlinearOptions::ptc\_dt0} 族驱动，见下表）\\
\fld{ptc\_growth} & --- & --- & 1.2--5 & 2.0 & PTC 步长增长因子；\textbf{保留字段，当前无消费点}（SER 近零残差改乘 $2$ 为 \srcpath{include/hacdcpf/power\_flow/globalization/lm\_trust\_region.hpp:PtcSerController} 内硬编码）\\
\end{paramtable}

\subsection{PV/PQ 切换与参考母线}
\begin{paramtable}
\fld{enable\_pv\_pq\_conversion} & --- & --- & true/false & true &MATPOWER 风格无功限值切换外环总开关（src/power\_flow/newton\_solver.cpp:NewtonSolver::solve）\\
\fld{pv\_q\_hysteresis\_pu} & --- & pu & 0--0.05 & 0.01 & PV$\to$PQ 无功越限迟滞带（src/power\_flow/newton\_solver.cpp:check\_q\_limits\_and\_switch（lambda））\\
\fld{pv\_recover\_vm\_tol\_pu} & --- & pu & 0--0.05 & 0.01 & PQ$\to$PV 恢复判据的$V_m$ 偏差容差（src/power\_flow/newton\_solver.cpp:check\_q\_limits\_and\_switch（lambda））\\
\fld{enable\_auto\_swing\_selection} & --- & --- & true/false & true &孤岛缺 slack 时自动遴选参考；读取点\srcpath{src/power\_flow/adaptive\_solver.cpp:AdaptiveSolver::solve}（Newton 主路径不读）\\
\end{paramtable}

\subsection{换流器、DC 与方程模式}
\begin{paramtable}
\fld{enable\_converter\_mode\_switching} & --- & --- & true/false & true &换流器 PQ$\leftrightarrow$VDC\_Q 模式切换（src/power\_flow/newton\_solver.cpp:NewtonSolver::solve，逻辑见第~\ref{sec:hybrid}~章）\\
\fld{converter\_vdc\_switch\_high\_pu} & --- & pu & 0.01--0.1 & 0.03 &切换评估高阈值（src/power\_flow/newton\_solver.cpp:apply\_converter\_mode\_switching）\\
\fld{converter\_vdc\_switch\_low\_pu} & --- & pu & 0--0.05 & 0.01 &切换评估低阈值（src/power\_flow/newton\_solver.cpp:apply\_converter\_mode\_switching）\\
\fld{mode\_hysteresis\_iters} & --- & 次 & 1--5 & 2 &模式切换迟滞迭代数（src/power\_flow/newton\_solver.cpp:apply\_converter\_mode\_switching）\\
\fld{enable\_rigid\_vdc\_former} & --- & --- & true/false & false &单换流器 DC 岛刚性电压源模式（DC-slack，多换流器模型 \S 6.3）；默认关，droop 提升已使 $V_{dc}$ 偏差约 0.1\%；读取点 src/power\_flow/newton\_solver.cpp:plan\_dc\_island\_references\\
\fld{enable\_converter\_coordination\_check} & --- & --- & true/false & false &预求解换流器协调可行性校核（src/api/hacdcpf.cpp:solve\_power\_flow；src/power\_flow/distributed\_slack\_solver.cpp:DistributedSlackSolver::solve\_simplified），Error/Fatal 阻断求解\\
\fld{enforce\_converter\_physical\_limits} & --- & --- & true/false & false &换流器容量圆/电流/调制比/占空比限值转为硬可行性门卫（定义 include/hacdcpf/power\_flow/power\_flow\_options.hpp:PowerFlowOptions）：违反限值的 Newton 根被拒绝为物理不可行而非报告收敛（src/api/hacdcpf.cpp:solve\_power\_flow），并回填四个 \texttt{enforced} validity 标志（:populate\_newton\_converter\_scope）；默认 false 保持仅告警行为；需再调度找可行点时用 OPF。仅平铺字段，不经 \srcpath{from\_parts} 装配\\
\fld{loss\_model} & --- & --- & Linear/\allowbreak CurrentBased & Linear &换流器损耗模型 \texttt{LossModelType}（include/hacdcpf/model/enums/converter\_enums.hpp:LossModelType）；作为 SolverData 缓存键（src/api/hacdcpf.cpp:hash\_system\_signature）并参与派生结果回算（:populate\_derived\_results）\\
\fld{enable\_coupled\_jacobian} & --- & --- & true/false & false &注入 AC$\leftrightarrow$DC 交叉耦合雅可比块；经 SolverData 转交（src/api/hacdcpf.cpp:apply\_zip\_weights），消费点 src/power\_flow/jacobian\_builder.cpp:evaluate\_residual\_impl、src/power\_flow/jacobian\_pattern.cpp:build\_jacobian\_pattern\\
\fld{enable\_augmented\_equations} & --- & --- & true/false & false &增强方程模式（VDC\_VAC 母线改 PQ + $V_m{=}V_{set}$ 行替换）；转交 src/api/hacdcpf.cpp:apply\_zip\_weights，消费点 src/power\_flow/newton\_solver.cpp:build\_jacobian\_context\\
\fld{enable\_semi\_smooth\_newton} & --- & --- & true/false & false &半光滑 Newton（PV 全转 PQ 记 NCP 限值）；转交 src/api/hacdcpf.cpp:apply\_zip\_weights，消费点 src/power\_flow/newton\_solver.cpp:build\_jacobian\_context\\
\end{paramtable}

\subsection{ZIP 负荷权重}
\begin{paramtable}
\fld{zip\_pw[3]} & $w_{p}$ & --- & $[0,1]$，$\sum=1$ & $(1,0,0)$ &有功 ZIP 权重（恒功率/恒电流/恒阻抗）；抄入 SolverData（src/api/hacdcpf.cpp:apply\_zip\_weights），装配消费 src/power\_flow/jacobian\_builder.cpp:build\_power\_spec\\
\fld{zip\_qw[3]} & $w_{q}$ & --- & $[0,1]$，$\sum=1$ & $(1,0,0)$ &无功 ZIP 权重（src/api/hacdcpf.cpp:apply\_zip\_weights；src/power\_flow/jacobian\_builder.cpp:build\_power\_spec）\\
\end{paramtable}

\subsection{运行时与诊断开关}
\begin{paramtable}
\fld{ac\_eval\_threads} & --- & 线程 & 0=自动 & 0 & AC 内核并行线程数，$<1024$ 条目或 1 线程回退串行（src/power\_flow/newton\_solver.cpp:effective\_ac\_eval\_threads）\\
\fld{enable\_solver\_profiling} & --- & --- & true/false & false &计时统计开关，填充 \fld{profiling} 各 ms 计数（src/power\_flow/newton\_solver.cpp:NewtonSolver::solve）\\
\fld{enable\_iteration\_log} & --- & --- & true/false & false &逐迭代日志开关；服务器层读取（\srcpath{src/server/runtime\_api\_v1.cpp:power\_flow\_options}），Newton 主路径\textbf{当前未接线}（\fld{iteration\_log} 不被填充）\\
\fld{verbose} & --- & --- & true/false & false & 详细日志；三相 NR 转发使用（include/hacdcpf/power\_flow/three\_phase.hpp:solve\_three\_phase），统一 Newton 主路径\textbf{当前未接线}\\
\end{paramtable}

\subsection{初始状态、嵌套选项与分组构造}
结构体成员 \fld{initial\_state}（\texttt{PowerFlowOptions} 内，\srcpath{include/hacdcpf/power\_flow/power\_flow\_options.hpp:PowerFlowOptions}）类型为
\texttt{std::optional<\allowbreak InitialState>}；API 层取指针传给 Newton（src/api/hacdcpf.cpp:solve\_power\_flow），
DC 初值由 \srcpath{load\_dc\_initial\_state} 应用（src/power\_flow/newton\_solver.cpp:NewtonSolver::solve）；
\fld{robust\_nonlinear}（:PowerFlowOptions）为五阶段鲁棒管线选项（第~\ref{subsec:robustopts}~节）。

\compmeta{hacdcpf::InitialState}{\srcpath{include/hacdcpf/power\_flow/power\_flow\_options.hpp:InitialState}}
\begin{paramtable}
\fld{vm} & $V_m^{(0)}$ & pu & $[0.9,1.1]$ & 空=平启动 & AC 母线电压幅值初值；索引空间为 canonical AC 母线位置（0-based）\\
\fld{va} & $\theta^{(0)}$ & rad & $(-\pi,\pi]$ & 空=0 & AC 母线相角初值，同上索引空间\\
\fld{vdc} & $V_{dc}^{(0)}$ & pu & $[0.9,1.1]$ & 空=1.0 & DC 母线电压初值；索引空间为 canonical DC 母线位置（0-based）\\
\end{paramtable}

同头文件另有两个辅助结构（如实列出，本章不展开）：\fld{ReactiveLimit}
（\fld{qmin}$=-0.5$、\fld{qmax}$=0.5$，:25--28）与 \fld{DistributedSlack}
（参与母线/因子/参考母线/参与上限，:30--35）；以及 7 个分组子结构
（\fld{ConvergenceOptions} 等，:41--95），字段与默认值同平铺字段一一对应，
由 \srcpath{from\_parts} 展开装配（例外：\fld{enforce\_converter\_physical\_limits}
仅存在于平铺层，分组子结构未收录）。

\subsection{RobustNonlinearOptions：五阶段鲁棒管线选项}\label{subsec:robustopts}

\compmeta{powerflow::RobustNonlinearOptions}{\srcpath{include/hacdcpf/power\_flow/globalization/robust\_nonlinear\_options.hpp:RobustNonlinearOptions}}
（\srcpath{power\_flow/robust\_nonlinear\_options.hpp} 为转发头，对应 .cpp 仅含
include。）算法语义见第~\ref{sec:robust}~章，此处仅列字段契约；Newton 内以
\texttt{ropts}（\texttt{opt.robust\_nonlinear}）引用。

\subsubsection*{Phase 1：缩放与诊断（默认全开）}
\begin{paramtable}
\fld{enable\_residual\_scaling} & --- & --- & true/false & true &残差按 $\max(|F_i^{spec}|,1)$ 行缩放；收敛判定用缩放范数（src/power\_flow/newton\_solver.cpp:NewtonSolver::solve；实现 src/power\_flow/nonlinear\_scaling.cpp:build\_nonlinear\_scaling）\\
\fld{enable\_variable\_scaling} & --- & --- & true/false & true &变量缩放 $D_x=\mathrm{diag}(\pi,V_i,V_{dc,i})$（src/power\_flow/nonlinear\_scaling.cpp:build\_nonlinear\_scaling）\\
\fld{enable\_jacobian\_row\_col\_equilibration} & --- & --- & true/false & true &$\hat J=D_FJD_x^{-1}$ 均衡，需至少一个缩放开关打开（src/power\_flow/newton\_solver.cpp:NewtonSolver::solve）\\
\fld{enable\_condition\_monitor} & --- & --- & true/false & true &逐迭代条件数代理，写入 \fld{profiling.condition\_estimate}（src/power\_flow/newton\_solver.cpp:NewtonSolver::solve）\\
\fld{min\_scale} & --- & --- & $[10^{-9},10^{-3}]$ & $10^{-6}$ &缩放因子下限（src/power\_flow/nonlinear\_scaling.cpp:build\_nonlinear\_scaling）\\
\fld{max\_scale} & --- & --- & $[10^{3},10^{9}]$ & $10^{6}$ &缩放因子上限（src/power\_flow/nonlinear\_scaling.cpp:build\_nonlinear\_scaling）\\
\fld{min\_vm\_pu} & $\underline V$ & pu & $10^{-9}$--$10^{-4}$ & $10^{-8}$ &雅可比 $\partial/\partial V$ 项分母电压下限；抄入雅可比上下文（src/power\_flow/newton\_solver.cpp:NewtonSolver::solve），消费 src/power\_flow/jacobian\_builder.cpp:build\_power\_spec\\
\fld{min\_branch\_x\_pu} & --- & pu & $10^{-20}$--$10^{-6}$ & $10^{-20}$ &FDPF $B'$ 矩阵零电抗支路剔除阈值（src/power\_flow/fdpf\_solver.cpp:solve）\\
\fld{bad\_condition\_threshold} & --- & --- & $10^{6}$--$10^{12}$ & $10^{10}$ &“坏条件数”阈值；\textbf{当前未接线}（仅 GUI/服务器层暴露）\\
\fld{tiny\_pivot\_threshold} & --- & --- & $10^{-15}$--$10^{-9}$ & $10^{-12}$ &微小主元日志阈值；\textbf{当前未接线}（仅 GUI/服务器层暴露）\\
\end{paramtable}

\subsubsection*{Phase 2：非单调线搜索与光滑 NCP}
\begin{paramtable}
\fld{enable\_nonmonotone\_linesearch} & --- & --- & true/false & true &GLL 非单调 Armijo 接受准则（src/power\_flow/newton\_solver.cpp:run\_line\_search（lambda））\\
\fld{nonmonotone\_window} & $M$ & 次 & 3--10 & 5 & 非单调参考 Merit 滑窗（src/power\_flow/newton\_solver.cpp:NewtonSolver::solve）\\
\fld{armijo\_c} & $c_1$ & --- & $(0,1)$ & $10^{-4}$ & Armijo 充分下降常数（src/power\_flow/newton\_solver.cpp:NewtonSolver::solve）\\
\fld{line\_search\_beta} & $\beta$ & --- & $(0,1)$ & 0.5 & 回溯缩减因子（src/power\_flow/newton\_solver.cpp:NewtonSolver::solve）\\
\fld{max\_line\_search\_trials} & --- & 次 & 5--30 & 12 & 非单调搜索最大试探数；\textbf{当前未接线}（实际试探数由 \fld{max\_line\_search\_steps} 控制）\\
\fld{enable\_smooth\_ncp} & --- & --- & true/false & false &PV/PQ 与 VSC 限流块的光滑 FB/CHKS 延拓；VSC 局部块可独立启用，不要求\fld{enable\_semi\_smooth\_newton}（newton\_solver.cpp）\\
\fld{ncp\_mu0} & $\mu_0$ & --- & $10^{-4}$--$10^{-1}$ & $10^{-2}$ &NCP 平滑参数初值（src/power\_flow/newton\_solver.cpp:NewtonSolver::solve）\\
\fld{ncp\_mu\_min} & $\mu_{\min}$ & --- & $10^{-14}$--$10^{-8}$ & $10^{-12}$ &平滑参数下限（src/power\_flow/newton\_solver.cpp:NewtonSolver::solve）\\
\fld{ncp\_mu\_factor} / \fld{ncp\_mu\_factor\_coarse} & --- & --- & 0.05--0.9 & 0.1 / 0.5 & 精细/粗糙阶段退火因子（src/power\_flow/newton\_solver.cpp:NewtonSolver::solve）\\
\fld{ncp\_mu\_phase\_transition} & --- & --- & 0.01--0.5 & 0.1 &粗/细阶段切换的相对残差阈值（src/power\_flow/newton\_solver.cpp:NewtonSolver::solve）\\
\fld{enable\_activity\_hysteresis} & --- & --- & true/false & true &PV/PQ 活跃集保持防抖；\textbf{当前未接线}（仅 GUI/服务器层暴露）\\
\fld{min\_active\_set\_hold\_iters} & --- & 次 & 1--10 & 3 &活跃集最小保持迭代数；\textbf{当前未接线}\\
\fld{q\_limit\_hysteresis} & --- & pu & $10^{-5}$--$10^{-2}$ & $10^{-4}$ &Q 限值邻近迟滞带；\textbf{当前未接线}（实际迟滞用\fld{pv\_q\_hysteresis\_pu}）\\
\fld{enable\_vsc\_local\_schur} & --- & --- & true/false & true &VSC 六状态局部块精确 Schur 消元；每个候选步通过局部条件与完整后向误差守卫，失败回退 full sparse LU\\
\fld{vsc\_schur\_min\_network\_dimension} & $n_{\min}$ & 维 & $\ge0$ & 1000 &生产复杂度准入；小于阈值保留 full LU，0 仅用于研究消融\\
\fld{vsc\_schur\_local\_rcond\_tolerance} & $\rho_D$ & --- & $>0$ & $\sqrt{\epsilon_{mach}}$ &所选 $6\times6$ 局部广义 Jacobian 块最小 reciprocal-condition 门槛；默认来自全局 \fld{NumericalConstants}\\
\fld{vsc\_schur\_backward\_error\_tolerance} & $\eta_b$ & --- & $>0$ & $10^{-12}$ &reduced solve 与完整重构 Newton 步的范数后向误差上限\\
\end{paramtable}

\subsubsection*{Phase 3：LM 信赖域与 PTC-SER}
\begin{paramtable}
\fld{enable\_lm\_trust\_region\_fallback} & --- & --- & true/false & true &LM 恢复步总开关（src/power\_flow/newton\_solver.cpp:NewtonSolver::solve）；经\fld{enable\_auto\_fallback\_scheduling} 挂入主线搜索回退链（:NewtonSolver::solve）\\
\fld{lm\_lambda0} & $\lambda_0$ & --- & $10^{-6}$--$10^{-2}$ & $10^{-4}$ &LM 阻尼初值（src/power\_flow/newton\_solver.cpp:NewtonSolver::solve）\\
\fld{lm\_lambda\_min} & $\lambda_{\min}$ & --- & $10^{-14}$--$10^{-8}$ & $10^{-12}$ &LM 阻尼下限；\textbf{当前未接线}（仅 GUI/服务器层暴露）\\
\fld{lm\_lambda\_max} & $\lambda_{\max}$ & --- & $10^{4}$--$10^{10}$ & $10^{8}$ &LM 阻尼上限（src/power\_flow/newton\_solver.cpp:NewtonSolver::solve）\\
\fld{trust\_region\_radius0} & $\Delta_0$ & --- & 0.1--5 & 1.0 &LM 信赖域初始半径（归一化）；\textbf{当前未接线}（信赖域全局化实际用\fld{PowerFlowOptions::trust\_region\_radius0}）\\
\fld{trust\_region\_radius\_min} & $\Delta_{\min}$ & --- & $10^{-8}$--$10^{-4}$ & $10^{-6}$ &信赖域半径下限；\textbf{当前未接线}\\
\fld{trust\_region\_radius\_max} & $\Delta_{\max}$ & --- & 10--$10^{3}$ & 100.0 &信赖域半径上限；\textbf{当前未接线}\\
\fld{enable\_ptc\_ser} & --- & --- & true/false & true &PTC 恢复步开关（src/power\_flow/newton\_solver.cpp:NewtonSolver::solve）；PTC 步长由 \fld{ptc\_dt0} 族经 \texttt{PtcSerController} 的 SER 更新驱动\\
\fld{ptc\_dt0} & $\delta t_0$ & --- & $10^{-3}$--1 & 0.1 & SER 初始伪时间步（经 \texttt{PtcSerController} 接线，src/power\_flow/newton\_solver.cpp:NewtonSolver::solve）\\
\fld{ptc\_dt\_min} / \fld{ptc\_dt\_max} & --- & --- & $10^{-6}$--$10^{8}$ & $10^{-4}$ / $10^{6}$ &伪时间步下/上限（同上接线）\\
\fld{ptc\_gamma} & $\gamma$ & --- & 0.1--2 & 0.7 &SER 指数 $\delta t_{k+1}=\delta t_k(r_{k-1}/r_k)^\gamma$（同上接线）\\
\end{paramtable}

\subsubsection*{Phase 4：同伦延拓}
\begin{paramtable}
\fld{enable\_homotopy} & --- & --- & true/false & true &同伦求解器使能；\texttt{HomotopyContinuationSolver::\allowbreak solve} 内对子求解关闭以防递归（src/power\_flow/homotopy\_continuation.cpp:HomotopyContinuationSolver::solve）；该求解器当前由测试/GUI 直接驱动，不在 Newton 失败链中自动调用（见第~\ref{sec:robust}~章）\\
\fld{homotopy\_step0} & $\Delta\lambda_0$ & --- & 0.05--0.5 & 0.2 &负荷/设定值爬坡初始步长（src/power\_flow/homotopy\_continuation.cpp:HomotopyContinuationSolver::solve）\\
\fld{homotopy\_step\_min} & --- & --- & $10^{-4}$--$10^{-2}$ & $10^{-3}$ &最小步长，低于则判同伦失败（src/power\_flow/homotopy\_continuation.cpp:solve）\\
\fld{homotopy\_step\_max} & --- & --- & 0.5--1.0 & 1.0 & 步长上限（src/power\_flow/homotopy\_continuation.cpp:solve）\\
\fld{homotopy\_max\_steps} & --- & 步 & 10--100 & 30 & $\lambda$ 路径最大步数（src/power\_flow/homotopy\_continuation.cpp:solve）；每步以 \fld{initial\_state} 热启动（:solve）\\
\end{paramtable}

\subsubsection*{Phase 5：高级与 Newton--Krylov（默认全关）}
\begin{paramtable}
\fld{enable\_log\_voltage} & --- & --- & true/false & false &以 $\ln V$、$\ln V_{dc}$ 为未知量；\textbf{当前未接线}（仅 GUI/服务器层暴露）\\
\fld{auto\_enable\_log\_voltage\_on\_failure} & --- & --- & true/false & false &失败后自动重试对数电压；\textbf{当前未接线}\\
\fld{enable\_newton\_krylov\_fallback} & --- & --- & true/false & false &病态雅可比时 GMRES 内迭代回退（src/power\_flow/newton\_solver.cpp:NewtonSolver::solve）\\
\fld{enable\_schur\_preconditioner} & --- & --- & true/false & false &AC/DC 块 Schur 补预条件子（src/power\_flow/newton\_solver.cpp:solve\_linear\_nk（lambda））\\
\fld{enable\_auto\_fallback\_scheduling} & --- & --- & true/false & true &线搜索耗尽后自动挂接 LM/PTC 恢复步（src/power\_flow/newton\_solver.cpp:NewtonSolver::solve）；需 Phase 3 对应开关同时打开\\
\fld{nk\_condition\_trigger} & --- & --- & $10^{4}$--$10^{10}$ & $10^{8}$ &NK-GMRES 触发的条件数代理阈值（src/power\_flow/newton\_solver.cpp:NewtonSolver::solve）\\
\fld{gmres\_restart} & $m$ & 维 & 10--100 & 30 & GMRES($m$) 重启维数（src/power\_flow/newton\_solver.cpp:solve\_linear\_nk（lambda））\\
\fld{gmres\_max\_outer} & --- & 轮 & 5--50 & 10 & GMRES 外重启轮数上限（src/power\_flow/newton\_solver.cpp:solve\_linear\_nk（lambda））\\
\fld{gmres\_tol} & --- & --- & $10^{-12}$--$10^{-6}$ & $10^{-10}$ &GMRES 相对残差容差（src/power\_flow/newton\_solver.cpp:solve\_linear\_nk（lambda））\\
\end{paramtable}

同头文件还定义策略切换日志类型：枚举 \texttt{NonlinearSolveMode}（ScaledNewton /
DampedNewton / NonmonotoneLineSearch / LMTrustRegion / PtcSer / Homotopy /
NewtonKrylovFallback，:213--221）与记录结构 \fld{SolverModeSwitchRecord}
\{\fld{iteration}, \fld{from}, \fld{to}, \fld{reason}\}（:SolverModeSwitchRecord）。

\subsection{PowerFlowProblem：求解任务描述符}

\compmeta{hacdcpf::PowerFlowProblem}{\srcpath{include/hacdcpf/power\_flow/power\_flow\_problem.hpp:PowerFlowProblem}}

由投影层工厂 \srcpath{build\_power\_flow\_problem}（声明
\srcpath{include/hacdcpf/power\_flow/power\_flow\_problem.hpp:build\_power\_flow\_problem}，实现
\srcpath{src/power\_flow/solver\_interface.cpp:build\_power\_flow\_problem}）装配，
供求解器抽象层 \texttt{IPowerFlowSolver} 消费；生产主路径（统一 Newton）
不经过该描述符（见第~\ref{sec:overview}~章诚实口径）。
\begin{paramtable}
\fld{network} & --- & --- & --- & --- & canonical 网络装配数据\texttt{powerflow::SolverData}（含 $Y_{bus}$、$G_{dc}$ 与映射，见第~\ref{sec:assembly}~章）\\
\fld{options} & --- & --- & --- & --- & 本次求解的 \texttt{PowerFlowOptions}\\
\fld{name} & --- & --- & --- & 空 & 可选算例名（人类可读）\\
\fld{n\_ac\_buses} & $n_{ac}$ & 个 & --- & 0 & canonical AC 母线数（规模摘要）\\
\fld{n\_dc\_buses} & $n_{dc}$ & 个 & --- & 0 & canonical DC 母线数\\
\fld{n\_converters} & --- & 个 & --- & 0 & 换流器台数\\
\end{paramtable}

\subsection{PowerFlowResult：顶层字段与索引空间}

\compmeta{hacdcpf::PowerFlowResult}{\srcpath{include/hacdcpf/power\_flow/power\_flow\_result.hpp:PowerFlowResult}}

索引空间总约定：求解器内部向量为 canonical 位置（0-based）；API 出口经
\srcpath{unproject\_pf\_result}（src/api/hacdcpf.cpp:unproject\_pf\_result，调用点 :solve\_power\_flow）广播回
\textbf{原始母线编号空间}，内部辅助 DC 母线经
\srcpath{trim\_internal\_dc\_bus\_results}（:819--826，调用点 :985）裁剪。\fld{va} 以\textbf{弧度}
存储，展示层转换为度（tests/run\_gui\_server.cpp:publish\_balanced\_pf（lambda））。
\begin{paramtable}
\fld{vm} & $V_m$ & pu & $[0.5,1.5]$ & 1.0 & AC 母线电压幅值；索引=原始 AC母线向量位置（0-based），长度=原始 AC 母线数\\
\fld{va} & $\theta$ & rad & $\pm\pi$ & 0 & AC 母线相角（弧度；展示层为度）；索引空间同 \fld{vm}\\
\fld{vdc} & $V_{dc}$ & pu & $[0.5,1.5]$ & 1.0 & DC 母线电压；索引=对外 DC母线向量位置（内部辅助 DC 母线已裁剪）\\
\fld{converged} & --- & --- & true/false & false & 末态残差重新判定的诚实收敛标志（src/power\_flow/newton\_solver.cpp:NewtonSolver::solve）；\fld{enforce\_converter\_physical\_limits} 门卫拒绝时置 false（src/api/hacdcpf.cpp:solve\_power\_flow）\\
\fld{iterations} & --- & 次 & --- & 0 & 外层+内层+1（src/power\_flow/newton\_solver.cpp:NewtonSolver::solve）\\
\fld{residual} & $\lVert F\rVert_\infty$ & pu & --- & 0 &最终（缩放）失配无穷范数\\
\fld{profiling} & --- & --- & --- & --- & 计数/计时统计\texttt{SolverProfiling}，见第~\ref{subsec:profiling}~节\\
\fld{diagnostics} & --- & --- & --- & --- & 高层诊断\texttt{SolverDiagnostics}，见第~\ref{subsec:diagnostics}~节\\
\fld{branch\_flows} & --- & MW/Mvar & --- & 空 & AC 支路两端潮流\texttt{BranchFlow}；索引=原始 AC 支路 .index 顺序（\texttt{sys.ac.branches}）\\
\fld{vsc\_transfers} & --- & MW/Mvar & --- & 空 & VSC 换流器传输\texttt{VSCTransfer}；\fld{index} 为组件稳定 .index\\
\fld{dcdc\_transfers} & --- & MW & --- & 空 & DC/DC 传输\texttt{DCDCTransfer}，含占空比事后校核\\
\fld{trafo3w\_flows} & --- & MW/Mvar & --- & 空 & 三绕组变压器三侧潮流\texttt{Trafo3WFlow}\\
\fld{er\_port\_transfers} & --- & MW/Mvar & --- & 空 & 能量路由器端口潮流\texttt{ERPortTransfer}\\
\fld{ac\_switch\_flows} & --- & MW/Mvar & --- & 空 & 被合并开关的端子潮流恢复\texttt{DeviceTerminalFlow}（src/api/hacdcpf.cpp:solve\_power\_flow）；无开关时为空\\
\fld{ac\_circuit\_breaker\_flows} & --- & MW/Mvar & --- & 空 &被合并断路器的端子潮流恢复，同上\\
\fld{converter\_model\_scope} & --- & --- & --- & unset &换流器模型覆盖声明 \texttt{ConverterModelScope}，见第~\ref{subsec:scope}~节\\
\end{paramtable}

\subsection{组件潮流结果结构}
以下结构同头文件（include/hacdcpf/power\_flow/power\_flow\_result.hpp:SolverProfiling）：功率 MW/Mvar、容量 MVA、
负载率 \%、电压 pu；\fld{index}/\fld{bus\_*} 均为对外稳定 ID，非向量位置。

\subsubsection*{BranchFlow（:BranchFlow）——AC 支路两端潮流}
\begin{paramtable}
\fld{pf\_mw} & $P_f$ & MW & --- & 0 & 首端（from）有功，流出母线为正\\
\fld{qf\_mvar} & $Q_f$ & Mvar & --- & 0 & 首端无功\\
\fld{pt\_mw} & $P_t$ & MW & --- & 0 & 末端（to）有功\\
\fld{qt\_mvar} & $Q_t$ & Mvar & --- & 0 & 末端无功\\
\end{paramtable}

\subsubsection*{VSCTransfer（:VSCTransfer）——VSC 换流器传输}
\begin{paramtable}
\fld{index} & --- & --- & --- & 0 & 换流器组件 .index（稳定 ID）\\
\fld{bus\_ac} & --- & --- & --- & 0 & 交流侧母线编号（原始编号空间，经 \srcpath{restore\_original\_vsc\_bus\_ac} 恢复，src/api/hacdcpf.cpp:solve\_power\_flow）\\
\fld{bus\_dc} & --- & --- & --- & 0 & 直流侧母线编号\\
\fld{p\_ac\_mw} & $P_{ac}$ & MW & --- & 0 & 交流侧有功（注入约定见阅读指南）\\
\fld{q\_ac\_mvar} & $Q_{ac}$ & Mvar & --- & 0 & 交流侧无功\\
\fld{p\_dc\_mw} & $P_{dc}$ & MW & --- & 0 & 直流侧有功；$P_{ac}+P_{dc}+P_{loss}=0$\\
\fld{loss\_mw} & $P_{loss}$ & MW & $\ge 0$ & 0 & 换流器损耗（依\fld{loss\_model}）\\
\end{paramtable}

\subsubsection*{DCDCTransfer（:DCDCTransfer）——DC/DC 传输与占空比校核}
\begin{paramtable}
\fld{index} & --- & --- & --- & 0 & DC/DC 组件 .index\\
\fld{bus\_in} / \fld{bus\_out} & --- & --- & --- & 0 & 输入/输出侧 DC 母线编号\\
\fld{p\_in\_mw} & $P_{in}$ & MW & --- & 0 & 输入侧有功（解后状态回算，非计划值，src/api/hacdcpf.cpp:populate\_derived\_results）\\
\fld{p\_out\_mw} & $P_{out}$ & MW & --- & 0 & 输出侧有功\\
\fld{loss\_mw} & $P_{loss}$ & MW & $\ge 0$ & 0 & 效率+$I^2R$ 损耗\\
\fld{duty} & $D$ & --- & $[d_{\min},d_{\max}]$ & 0 & 由解出端口电压反推的理想 CCM 占空比（Isolated 拓扑为调制增益）\\
\fld{voltage\_ratio} & $V_{out}/V_{in}$ & --- & --- & 0 & 端口电压比\\
\fld{duty\_defined} & --- & --- & true/false & false & Generic 拓扑为 false\\
\fld{duty\_feasible} & --- & --- & true/false & true & 越窗时触发[DCDC-PHYS-01] 告警（src/api/hacdcpf.cpp:solve\_power\_flow）\\
\end{paramtable}

\subsubsection*{Trafo3WFlow（:Trafo3WFlow）——三绕组变压器}
\begin{paramtable}
\fld{index} & --- & --- & --- & 0 & 组件 .index\\
\fld{hv\_bus} / \fld{mv\_bus} / \fld{lv\_bus} & --- & --- & --- & 0 &高/中/低压侧母线编号\\
\fld{p\_hv\_mw}，\fld{q\_hv\_mvar} 等六量 & $P,Q$ & MW/Mvar & --- & 0 &三侧各 $P$、$Q$（\fld{p\_mv\_mw}、\fld{q\_mv\_mvar}、\fld{p\_lv\_mw}、\fld{q\_lv\_mvar} 同构）\\
\fld{loss\_mw} & $P_{loss}$ & MW & $\ge 0$ & 0 & 三侧功率不平衡（损耗）\\
\fld{loading\_pct} & --- & \% & --- & 0 & 负载率\\
\fld{rate\_mva} & $S_{rate}$ & MVA & --- & 0 & 额定容量\\
\end{paramtable}

\subsubsection*{ERPortTransfer（:ERPortTransfer）——能量路由器端口}
\begin{paramtable}
\fld{router\_index} / \fld{port\_index} & --- & --- & --- & 0 &路由器与端口组件 .index\\
\fld{bus} & --- & --- & --- & 0 & 端口所接母线编号\\
\fld{is\_ac} & --- & --- & true/false & true & 端口域（AC/DC）\\
\fld{p\_mw} / \fld{q\_mvar} & $P,Q$ & MW/Mvar & --- & 0 & 端口有功/无功\\
\fld{v\_pu} & $V$ & pu & --- & 1.0 & 端口母线电压幅值\\
\end{paramtable}

\subsubsection*{DeviceTerminalFlow（:DeviceTerminalFlow）——开关/断路器端子潮流}
投影层把闭合开关/断路器折叠进零阻抗合并（第~\ref{sec:assembly}~章），
求解后由 \srcpath{compute\_device\_terminal\_flows}（src/api/hacdcpf.cpp:solve\_power\_flow）恢复其端子潮流：
\begin{paramtable}
\fld{index} & --- & --- & --- & 0 & 设备组件 .index\\
\fld{bus\_from} / \fld{bus\_to} & --- & --- & --- & 0 & 两端母线编号（原始编号空间）\\
\fld{closed} & --- & --- & true/false & true & 合/分状态\\
\fld{pf\_mw}，\fld{pt\_mw}，\fld{qf\_mvar}，\fld{qt\_mvar} & $P,Q$ & MW/Mvar & --- & 0 &两端 $P$、$Q$（四字段同行声明）\\
\fld{rate\_mva} & $S_{rate}$ & MVA & --- & 0 & 容量（0=不限）\\
\fld{loading\_pct} & --- & \% & --- & 0 & 负载率\\
\end{paramtable}

\subsection{SolverProfiling：计数与计时}\label{subsec:profiling}

\compmeta{hacdcpf::SolverProfiling}{\srcpath{include/hacdcpf/power\_flow/power\_flow\_result.hpp:SolverProfiling}}

逐迭代向量索引空间为 Newton 迭代序号（0-based）；计时字段单位毫秒，仅在
\fld{enable\_solver\_profiling} 打开时填充（src/power\_flow/newton\_solver.cpp:NewtonSolver::solve）。
\begin{paramtable}
\fld{linear\_solver\_backend} & --- & --- & KLU/UMFPACK/\dots & 空 &实际选用的稀疏线性后端名（优先级见第~\ref{sec:newton}~章）\\
\fld{ac\_eval\_threads} & --- & 线程 & --- & 1 & 实际生效的 AC 内核线程数\\
\fld{jacobian\_pattern\_rebuilds} & --- & 次 & --- & 0 & 稀疏模式重建次数（缓存失效计数）\\
\fld{jacobian\_analyze\_calls} & --- & 次 & --- & 0 &\texttt{analyzePattern} 调用次数\\
\fld{factorization\_calls} & --- & 次 & --- & 0 & 数值分解次数\\
\fld{linear\_solve\_calls} & --- & 次 & --- & 0 & 三角回代次数\\
\fld{regularization\_attempts} & --- & 次 & --- & 0 & 对角正则化尝试次数\\
\fld{line\_search\_evaluations} & --- & 次 & --- & 0 & 线搜索残差评估总次数\\
\fld{rejected\_steps} & --- & 次 & --- & 0 & 被拒绝步数\\
\fld{pv\_to\_pq\_switches} / \fld{pq\_to\_pv\_switches} & --- & 次 & --- & 0 &PV/PQ 双向切换计数\\
\fld{converter\_mode\_switches} & --- & 次 & --- & 0 & 换流器模式切换计数\\
\fld{residual\_by\_iter} & --- & pu & --- & 空 & 逐迭代（缩放）残差；索引=迭代序号\\
\fld{line\_search\_evals\_by\_iter} & --- & 次 & --- & 空 & 逐迭代线搜索评估数\\
\fld{eval\_jacobian\_ms\_by\_iter} / \fld{linear\_solve\_ms\_by\_iter} / \fld{line\_search\_ms\_by\_iter} & --- & ms & --- & 空 & 逐迭代三段耗时（装配求值/线性求解/线搜索）\\
\fld{eval\_jacobian\_ms\_total} / \fld{linear\_solve\_ms\_total} / \fld{line\_search\_ms\_total} & --- & ms & --- & 0 & 三段累计耗时\\
\fld{raw\_residual\_norm} / \fld{scaled\_residual\_norm} & --- & pu & --- & 0 & 未缩放/缩放后残差范数（后者为收敛判定所用）\\
\fld{condition\_estimate} & $\kappa$ 代理 & --- & --- & 0 &条件数 $\ell_1$ 行/列范数比代理（src/power\_flow/newton\_solver.cpp:NewtonSolver::solve 起；\fld{enable\_condition\_monitor} 控制，写入 :1429）\\
\fld{regularization\_count} / \fld{linear\_solver\_status} & --- & 次/串 & --- & 0 / "not\_run" & 正则化累计次数；线性后端状态串\\
\end{paramtable}

\subsection{SolverDiagnostics：诊断与诚实口径}\label{subsec:diagnostics}

\compmeta{hacdcpf::SolverDiagnostics}{\srcpath{include/hacdcpf/power\_flow/power\_flow\_result.hpp:SolverDiagnostics}}

\begin{paramtable}
\fld{converged} & --- & --- & true/false & false & 诊断层收敛标志；\textbf{当前未接线}（以顶层 \fld{PowerFlowResult::converged} 为准；\fld{enforce\_converter\_physical\_limits} 门卫拒绝时同步置 false，src/api/hacdcpf.cpp:solve\_power\_flow）\\
\fld{termination\_reason} & --- & --- & --- & 空 & 终止原因；仅阻断/拒绝时写入：参考母线资格失败（src/api/hacdcpf.cpp:solve\_power\_flow）、协调校核失败（:solve\_power\_flow）、无合法 slack AC 岛（:solve\_power\_flow）、硬约束门卫判物理不可行（:solve\_power\_flow）\\
\fld{iterations} & --- & 次 & --- & 0 & 诊断层迭代数；\textbf{当前未接线}（以顶层 \fld{iterations} 为准）\\
\fld{initial\_mismatch\_norm} & --- & pu & --- & 0 & 初始失配范数；\textbf{当前未接线}（逐迭代残差见 \fld{profiling.residual\_by\_iter}）\\
\fld{final\_mismatch\_norm} & --- & pu & --- & 0 & 末态失配范数；\textbf{当前未接线}（以顶层 \fld{residual} 为准）\\
\fld{max\_p\_mismatch\_pu} & --- & pu & --- & 0 & 最大有功失配；\textbf{当前未接线}（市场层读取，src/market/market\_simulation.cpp:summarize\_ac\_validation）\\
\fld{max\_q\_mismatch\_pu} & --- & pu & --- & 0 & 最大无功失配；\textbf{当前未接线}\\
\fld{max\_i\_violation\_pu} & --- & pu & --- & 0 & 最大电流越限；\textbf{当前未接线}\\
\fld{condition\_estimate} & --- & --- & --- & 0 & 诊断层条件数；\textbf{当前未接线}（以 \fld{profiling.condition\_estimate} 为准）\\
\fld{equation\_closure\_checked} & --- & --- & true/false & false &结构闭合扫描已执行（首外层一次，src/power\_flow/newton\_solver.cpp:NewtonSolver::solve）\\
\fld{equation\_closure\_ok} & --- & --- & true/false & true &雅可比无全零行/列（:NewtonSolver::solve）；失败时追加告警（:NewtonSolver::solve）\\
\fld{n\_equations} / \fld{n\_variables} & $n_{eq},n_{var}$ & --- & --- & 0 & 方程数/未知量数（闭合校核 $n_{eq}=n_{var}$）\\
\fld{empty\_jacobian\_rows} / \fld{empty\_jacobian\_cols} & --- & 行/列 & --- & 0 & 全零行/列计数（\srcpath{scan\_empty\_rows\_cols}）\\
\fld{final\_breakdown} & --- & pu & --- & --- & 分类残差\texttt{ResidualBreakdown}（见下）\\
\fld{iteration\_log} & --- & --- & --- & 空 & 逐迭代日志\texttt{IterationLogEntry}；\textbf{当前未接线}（Newton 不填充）\\
\fld{warnings} & --- & --- & --- & 空 & 告警字符串（含规则码，见下）\\
\fld{converter\_coordination} & --- & --- & --- & --- & 预求解协调报告\texttt{ConverterCoordinationReport}（src/api/hacdcpf.cpp:solve\_power\_flow 回填）\\
\fld{promoted\_vsc\_indices} & --- & --- & --- & 空 & 因 DC 岛无电压参考而被自动提升为调压模式的 VSC 下标（求解器换流器列表 0-based，src/power\_flow/newton\_solver.cpp:plan\_dc\_island\_references）\\
\fld{effective\_converters} & --- & --- & --- & 空 & 求解终态换流器列表（含自动提升/刚性增益/模式切换），派生结果回算以它为准（src/power\_flow/newton\_solver.cpp:NewtonSolver::solve）\\
\end{paramtable}

分类残差 \texttt{ResidualBreakdown}（include/hacdcpf/power\_flow/power\_flow\_result.hpp:ResidualBreakdown）：
\fld{ac\_p\_norm}、\fld{ac\_q\_norm}、\fld{dc\_p\_norm}、\fld{converter\_p\_norm}、
\fld{converter\_q\_norm}、\fld{control\_eq\_norm} 六个分类无穷范数（pu），
\texttt{total\_norm()} 取六者最大。逐迭代日志条目 \texttt{IterationLogEntry}
（:IterationLogEntry）：\fld{iteration}、\fld{residual\_norm}（pu）、\fld{step\_norm}、
\fld{damping\_factor}、\fld{jacobian\_refactorized}、\fld{lm\_fallback\_used}、\fld{breakdown}。

告警码（写入 \fld{warnings}；默认模式下均为\textbf{事后校核}而非硬约束，
\fld{enforce\_converter\_physical\_limits} 开启时转为收敛拒绝依据，
读取点 src/api/hacdcpf.cpp:populate\_derived\_results）：
\begin{itemize}
  \item [ACDC-PHYS-01] VSC 调制比 $m<m_{\min}$；
  \item [ACDC-PHYS-02] VSC 交流电流 $I_{ac}>$ \fld{i\_ac\_max\_pu}；
  \item [ACDC-PHYS-03] VSC 直流电流 $I_{dc}>$ \fld{i\_dc\_max\_pu}；
  \item [ACDC-PHYS-04] VSC 视在功率越额定 \fld{p\_rated\_mw}；
  \item [ACDC-PHYS-05] VSC 调制比 $m>m_{\max}$（直流电压不足以支撑交流电压）；
  \item [DCDC-PHYS-01] DC/DC 占空比越 $[d_{\min},d_{\max}]$ 窗口；
  \item [CONVERTER-PHYS-HARD] 硬约束门卫：Newton 已收敛但运行点违反已声明的
        换流器不等式，被判物理不可行（src/api/hacdcpf.cpp:solve\_power\_flow）；
  \item 换流器协调规则码（\fld{rule\_id}，Warning/Info 级汇入，
        src/api/hacdcpf.cpp:solve\_power\_flow；规则目录见第~\ref{sec:hybrid}~章）。
\end{itemize}
\texttt{ConverterCoordinationReport}（include/hacdcpf/power\_flow/converter\_coordination.hpp:ConverterCoordinationReport）：
\fld{enabled}、\fld{feasible}、\fld{issues}（严重度 Info/Warning/Error/Fatal、
\fld{rule\_id}、组件 .index、岛号、消息）与逐 DC 岛摘要 \fld{dc\_islands}
（\texttt{DCIslandCoordinationSummary}，:40--63）。

\subsection{ConverterModelScope：模型覆盖声明}\label{subsec:scope}

\compmeta{hacdcpf::ConverterModelScope}{\srcpath{include/hacdcpf/model/converter\_model\_scope.hpp:ConverterModelScope}}

潮流快照求解的填充点：src/api/hacdcpf.cpp:populate\_newton\_converter\_scope。\texttt{ValidityFlags} 十标志逐一
声明该次求解\textbf{实际}建模/执行的换流器特性，false 项即诚实口径边界：
\begin{paramtable}
\fld{model\_scope} & --- & --- & --- & "unset" & 作用域标签；快照潮流填 \texttt{steady-state-newton:\allowbreak vsc-3mode+\allowbreak dcdc-power-transfer}\\
\fld{vsc\_loss\_modelled} & --- & --- & --- & true & VSC 损耗（线性/电流基）\\
\fld{vsc\_ac\_conduction\_loss\_modelled} & --- & --- & --- & true &交流侧传导损耗（\fld{r\_conv\_ac\_pu} 耦合，opt-in）\\
\fld{vsc\_capacity\_circle\_enforced} & --- & --- & --- & 随选项 &容量圆 $P^2+Q^2\le S^2$ 未嵌入方程；回填 \fld{enforce\_converter\_physical\_limits}（src/api/hacdcpf.cpp:populate\_newton\_converter\_scope）——false 时仅 [ACDC-PHYS-04] 事后告警，true 时违反即拒绝收敛\\
\fld{vsc\_current\_limits\_enforced} & --- & --- & --- & 随选项 &电流限，同上回填（:populate\_newton\_converter\_scope）；false 时仅事后告警（02/03）\\
\fld{vsc\_modulation\_limits\_enforced} & --- & --- & --- & 随选项 &调制比限，同上回填（:populate\_newton\_converter\_scope）；false 时仅事后告警（01/05）\\
\fld{vsc\_vdc\_control\_modelled} & --- & --- & --- & true &VDC\_Q/VDC\_VAC 控制与刚性 droop 成压\\
\fld{dc\_multisource\_coordination\_modelled} & --- & --- & --- & true &多源协调（\fld{is\_master}+参与因子）影响解点\\
\fld{dcdc\_loss\_modelled} & --- & --- & --- & true & DC/DC 效率+$I^2R$ 损耗\\
\fld{dcdc\_duty\_ratio\_enforced} & --- & --- & --- & 随选项 &占空比，同上回填（:populate\_newton\_converter\_scope）；false 时仅事后校核（[DCDC-PHYS-01]）\\
\fld{equation\_closure\_checked} & --- & --- & --- & 随求解 &回填 \fld{diagnostics.equation\_closure\_checked}\\
\end{paramtable}

\subsection{其他求解器结果类型}
同头文件另有四个轻量结果结构（include/hacdcpf/power\_flow/power\_flow\_result.hpp:DeviceTerminalFlow），字段含义与
\texttt{PowerFlowResult} 同名项一致，\fld{vm}/\fld{va}/\fld{vdc} 索引空间与
单位约定亦同（弧度相角、回投影后原始编号空间）：
\begin{itemize}
  \item \texttt{DCPowerFlowResult}（:DCPowerFlowResult）：\fld{vdc}、\fld{converged}、
        \fld{iterations}、\fld{residual}——纯 DC 网络求解（第~\ref{sec:hybrid}~章）；
  \item \texttt{AdaptiveSolveResult}（:AdaptiveSolveResult）/ \texttt{IslandedSolveResult}
        （:IslandedSolveResult）：增 \fld{islands}（\texttt{IslandInfo} 逐岛信息），前者另含 \fld{branch\_flows}；
  \item \texttt{DistributedSlackResult}（:DistributedSlackResult）：增 \fld{diagnostics}、
        \fld{distributed\_slack\_p}（母线 .index $\to$ 有功修正，系统
        \fld{base\_mva} 基准上的 pu；单位字符串见
        \fld{distributed\_slack\_p\_unit}）、
        \fld{hit\_limits}（触及参与上限的母线 .index）。
\end{itemize}

% 第 11 章：验证与校核总览
% 本章文件由 \input 引入，不含导言区。
\section{验证与校核总览}\label{sec:validation}

本章集中汇总潮流计算全家族的验证资产：测试矩阵、算例目录、外部对照通道、
已知边界与运行方法。各求解器节末的“验证与校核”小节（如第~\ref{sec:newton}~章）
给出局部依据，本章给出全局地图与诚实口径。

\subsection{验证体系分层}
平台潮流验证分四层，层层判据不同、互为补充：
\begin{enumerate}
  \item \textbf{单元与回归层}（Catch2，100+ 目标、约 1250 用例）：自建小算例的
        解析断言、收敛性回归、选项开/关对照；
  \item \textbf{外部对照层}：MATPOWER 三档回归、OpenDSS dss\_capi 桥、
        GridLAB-D 快照、Julia/JPC 相域基线；
  \item \textbf{解析与差分层}：雅可比中心差分校验、独立致密残差复算、
        两节点解析解；
  \item \textbf{能力合同与端到端层}：HTTP 能力覆盖矩阵与 GUI 全链路冒烟。
\end{enumerate}
铁律沿用第~\ref{sec:overview}~章的诚实结果口径：近似、fallback、门控跳过、
模型覆盖不足必须显式声明，验证报告不得把 skip 计为 pass。

\subsection{潮流测试矩阵总表}
\srcpath{tests/} 下潮流相关测试文件（仅列实际存在者，判据以文件内
\texttt{SECTION}/\texttt{REQUIRE}/\texttt{CHECK} 实际断言为准）：

\begin{footnotesize}
\begin{longtable}{@{}L{3.9cm}L{3.2cm}L{3.5cm}L{4.4cm}@{}}
\toprule
测试文件 & 覆盖求解器 & 算例/对照 & 容差/判据\\
\midrule
\endhead
\srcpath{test\_matpower\_cases.cpp} & 统一 Newton（+FDPF 一致性、DC OPF） &
  MATPOWER 三档（见 \ref{subsec:val-cases} 节）&
  收敛 + $|V|\in[0.5,1.5]$、$|\delta|<60^\circ$、残差 $<10^{-4}$；Tier2 批收敛率 $\ge 85\%$\\
\srcpath{test\_powerflow\_performance.cpp} & Newton（\texttt{[slow]} 标签）&
  MATPOWER 全目录（排除 \texttt{contab}/\texttt{scenarios}）&
  独立致密残差复算 $<10^{-6}$；电压窗 $(0.5,1.5)$\\
\srcpath{test\_unified\_pf\_upgrade.cpp} & 耦合雅可比/增强方程/半光滑/全局化 &
  解析雅可比 vs 中心差分（$\varepsilon=10^{-6}$）&
  $\max$ 相对误差 $<10^{-3}$；VDC\_VAC 定点 $<10^{-6}$\\
\srcpath{test\_scaling\_regression.cpp} & Newton 缩放开/关 &
  case9/14/30/57/118 &
  均收敛；缩放/原始残差与条件数估计有限\\
\srcpath{test\_helm\_solver.cpp} & HELM（+FDPF 稀疏注入、组件负荷聚合）&
  两母线电压稳定算例 vs Newton &
  HELM 与 NR 的 $|V|,|\theta|$ 差 $<2\times10^{-5}$；FDPF 与 NR 差 $<2\sim3\times10^{-4}$\\
\srcpath{test\_voltage\_stability.cpp} & CPF 弧长延拓 &
  两母线 P--V 鼻点 + 组件负荷方向 &
  翻越鼻点、$\lambda_{\max}\in(0.25,0.5)$、下支 $\lambda$ 回落；方向量精确到 $10^{-12}$\\
\srcpath{test\_homotopy\_continuation.cpp} & 同伦延拓 &
  两母线重负荷 &
  $\lambda=1.0$ 抵达端点、接受步 $\ge3$、收敛\\
\srcpath{test\_newton\_krylov.cpp} & NK-GMRES / 求解器工厂 / SolverWorkspace &
  $3\times3$ 耦合块系统 + 两母线 &
  GMRES 步误差 $<10^{-10}$；六方法工厂均可解；工作区复位保布局\\
\srcpath{test\_hacdcpf.cpp} & 公共 API 全族（AC/DC/FDPF/自适应/孤岛/线性化）&
  自建 2 母线 + IEEE14 AC/DC case builder &
  收敛、结果尺寸、有限性\\
\srcpath{test\_advanced\_pf.cpp} & 孤岛检测/自适应/分布式松弛/PV--PQ 策略 &
  9 母线、case14/118/2383wp &
  岛数精确、参与因子归一、双路径收敛一致\\
\srcpath{test\_dc\_island\_reference.cpp} & DC 岛参考规划/自动提升 &
  自建混合微网 &
  $V_{dc}$ 钉住、提升记录并告警、真参考不提升\\
\srcpath{test\_converter\_coordination.cpp} & 换流器协调预校核（46 个用例）&
  规则目录 \srcpath{ACDC-*/DCDC-*/DCISLAND-*} &
  阻断/告警语义、\fld{converter\_model\_scope} 声明\\
\srcpath{test\_all\_components\_pf.cpp} & 29 类组件全链路投影+潮流 &
  all29 集成系统 &
  收敛 + 每类组件对结果有可测影响\\
\srcpath{test\_distribution\_pipeline.cpp} & 配网 BFS 管线 &
  CB/开关/VSC rich$\to$graph$\to$PF$\to$OPF &
  管线端到端收敛\\
\srcpath{test\_three\_phase\_pf.cpp} & 三相投影与平衡化 &
  自建 3 母线（平衡/不平衡/环网/多源） &
  $|V|\in[0.8,1.1]$、迭代 $<15$、投影功率守恒\\
\srcpath{test\_three\_phase\_nr.cpp} & 相域 NR（16 个用例）&
  平衡/不平衡/支线/环网/多源 &
  残差 $<$ tol、VUF $\approx 0$、逐相功率平衡\\
\srcpath{test\_three\_phase\_nr\_load\_models.cpp} & 相域 NR 负荷模型 &
  wye/delta ZIP、中性点位移 &
  阻抗--电流--功率语义次序、不支持合同 fail-close\\
\srcpath{test\_three\_phase\_nr\_regulator.cpp} & 相域 NR 调压器闭环 &
  分相独立、远方 PT+LDC &
  档位单向提升、限值/振荡 fail-close\\
\srcpath{test\_three\_phase\_nr\_phase\_indexing.cpp} & 紧凑节点索引 &
  单相/两相支线 &
  无幻影相节点、phase\_mask 越界拒绝\\
\srcpath{test\_three\_phase\_nr\_transformer\_source.cpp} & 相域变压器源 &
  接地 wye 抽头、YNyn6 移相 &
  abc 域戳记精确\\
\srcpath{test\_three\_phase\_hybrid\_pf.cpp} & 相域 AC+DC 联立 &
  GFL Vdc--Q / GFM Vdc--内电势控制 &
  联立收敛（见第~\ref{sec:threephase}~章）\\
\srcpath{test\_opendss\_pf\_compare.cpp} & OpenDSS 桥对照（门控）&
  IEEE13、case33bw 全三相、调压器闭环 &
  分层 evidence/accepted 门槛，回归时退出码 2\\
\srcpath{test\_gridlabd\_compare.cpp} & GridLAB-D 快照对照 &
  两/三母线、元件阶梯、馈线阶梯 &
  运行时探测，缺失自动 skip\\
\srcpath{test\_phase\_domain\_ieee13.cpp} 等 plain 测试 & 相域求解器 vs JPC/OpenDSS 基线 &
  IEEE13/IEEE123 &
  \srcpath{HACDCPF\_HAVE\_OPENDSS} + 数据目录双重门控\\
\srcpath{test\_hacdcdss.cpp} & \srcpath{NetworkModel/PowerModels} 基元 &
  2 母线、IEEE14 AC/DC &
  $Y_{bus}$ 对称、注入自洽、收敛尺寸\\
\bottomrule
\end{longtable}
\end{footnotesize}

\subsection{算例目录}\label{subsec:val-cases}
\srcpath{external\_data/} 下潮流相关数据（均经 \texttt{ls} 核实）：
\begin{itemize}
  \item \textbf{\srcpath{matpower/}}：85 个 \texttt{.m} 文件，其中
        \texttt{contab\_*}（4 个）与 \texttt{scenarios\_*}（2 个）为扰动/场景
        辅助文件，潮流回归只取其余 79 个标准算例，分三档
        （\srcpath{tests/test\_matpower\_cases.cpp}）：Tier1 九算例
        （case5/9/14/30/57/118/300/39/24\_ieee\_rts）逐案断言收敛与电压合理性；
        Tier2 批量 54 算例（case533mt\_hi/lo 因解析器不支持算术表达式允许
        skip 不计失败）批收敛率 $\ge 85\%$；Tier3 大系统 14 算例至
        case\_SyntheticUSA（约 8.2 万节点，先关 PV--PQ 切换求解、失败再回退
        全开，:504--515）；
  \item \textbf{\srcpath{opendss\_ieee\_pes/}}：IEEE PES 测试馈线，
        \texttt{extracted/} 含 13/34/123 节点原始资料，
        \texttt{opendss\_reference/13\_node/official\_full/IEEE13Nodeckt.dss}
        是相域对照的官方基准（\srcpath{tests/CMakeLists.txt:250}）；
  \item \textbf{\srcpath{profiles/}}：\texttt{case33bw} 与
        \texttt{nansha\_network} 两套 8760 小时负荷/辐照/电价时序
        （时序生产模拟用）；
  \item \textbf{\srcpath{classical\_example/}}：4 个 JSON 算例（含
        \srcpath{hybrid\_ac\_dc\_mg\_case\_01.json}、
        \srcpath{nansha\_full\_network.json}）；
  \item \textbf{\srcpath{opf\_hand\_cases/}}：4 个手算 OPF 小算例（含 README
        给出解析期望；该目录当前不在仓库检出中）。
\end{itemize}

\subsection{外部对照通道细节}
\begin{itemize}
  \item \textbf{MATPOWER 回归}：Catch2 侧为上述三档测试；脚本侧
        \srcpath{tools/matpower\_benchmark.py}（217 行）驱动 MATLAB 批量
        \texttt{runpf} 与本仓库 \srcpath{tools/matpower\_pf\_compare.cpp}
        （输出收敛标志与逐母线电压 JSON，支持 \texttt{--no-pv-pq}/
        \texttt{--flows}/\texttt{--repeat N}），按匹配母线号比较
        $\max/\mathrm{mean}\,|\Delta V_m|$（pu）与 $|\Delta\theta|$（度），
        并对比两侧耗时；产物落在 \texttt{output/benchmarks/matpower\_pf/}。
        注意脚本内 MATLAB 路径与 \texttt{build/macos-release} 二进制路径为
        硬编码，跨机使用需改路径或走 Catch2 通道；
  \item \textbf{OpenDSS（dss\_capi 桥）}：三级门控——CMake 选项
        \srcpath{HACDCPF\_ENABLE\_OPENDSS}（默认 OFF）与
        \srcpath{HACDCPF\_ENABLE\_OPENDSS\_COMPARE} 先生成内部变量
        \srcpath{HACDCPF\_HAVE\_OPENDSS}（CMakeLists.txt:265--334），再派生
        \srcpath{HACDCPF\_HAVE\_OPENDSS\_COMPARE}（:solve\_once（lambda））；
        \srcpath{tests/test\_opendss\_pf\_compare.cpp}（3717 行）仅在后者为真时
        编译（\srcpath{tests/CMakeLists.txt:254--257}），覆盖单相/三相桥对照、
        case33bw 全三相 vs BFS 单相等值、调压器闭环与变压器分接头投影，
        接受层回归时对照门返回退出码 2；
  \item \textbf{GridLAB-D}：\srcpath{tests/test\_gridlabd\_compare.cpp} 目标
        始终构建，外部对照在运行时按 \texttt{HACDCPF\_GRIDLABD\_BIN}/
        \texttt{GRIDLABD\_BIN}/PATH/兄弟构建目录顺序探测可执行文件，缺失则
        干净 skip（\srcpath{tests/CMakeLists.txt:229--234}）；PV 母线电压调节
        不精确是已记录的 diagnostic 用例而非 pass；
  \item \textbf{解析解与雅可比中心差分}：
        \srcpath{src/power\_flow/jacobian\_check.cpp}（75 行）对任意
        残差/雅可比函数对做中心差分（默认 $h=10^{-6}$），按相对误差容差判定；
        \srcpath{src/power\_flow/residual\_evaluator.cpp}（157 行）从解向量
        以稀疏复数形式 $S=V\cdot\overline{Y_{bus}V}$ 独立重算注入残差
        （不再转致密矩阵，仅用于校核），被
        \srcpath{test\_powerflow\_performance.cpp} 以 $<10^{-6}$ 门槛复算
        每个收敛解（PV--PQ 切换或多岛算例因指定量不同显式跳过复算，:163--193）；
  \item \textbf{能力合同与 GUI E2E}：
        \srcpath{tools/validate\_case\_capabilities.py}（698 行）对 19 个能力域
        逐一发 HTTP 断言，输出 \texttt{output/capability\_coverage.md}
        覆盖矩阵（PASS/FAIL/SKIP 分列，skip 不计 pass）；
        \srcpath{tools/gui\_api\_e2e.py}（1050 行）启动
        \texttt{run\_gui\_server} 走 REST 全链路（装载算例$\to$ETAP 导出/再导入
        $\to$潮流$\to$短路），注册为 ctest 用例 \texttt{gui\_api\_e2e}
        （标签 \texttt{gui;e2e}、超时 120~s，由 \texttt{HACDCPF\_ENABLE\_ETAP}
        与 Python3 双重门控，\srcpath{tests/CMakeLists.txt:423--433}）。
\end{itemize}

\subsection{已知边界与覆盖缺口}
对照主文档“现状与诚实口径声明”逐条展开，每条附代码或测试证据：
\begin{enumerate}
  \item \textbf{HELM、同伦延拓、CPF、Newton--Krylov 已有专属回归（NEW 新增），
        但仍无生产 API 自动挂接}：实现分别在
        \srcpath{src/power\_flow/helm\_solver.cpp}、
        \srcpath{homotopy\_continuation.cpp}、\srcpath{voltage\_stability.cpp}、
        \srcpath{newton\_krylov.cpp}；NEW 起四者均有专属 Catch2 目标
        （\srcpath{tests/CMakeLists.txt:122--125}）：
        \srcpath{test\_helm\_solver.cpp}（HELM 与 Newton 在两母线算例上
        $|V|$/$|\theta|$ 一致至 $2\times10^{-5}$，另含 FDPF 稀疏注入与组件负荷
        聚合两条回归）、\srcpath{test\_homotopy\_continuation.cpp}
        （$\lambda$ 抵达 $1.0$ 端点）、\srcpath{test\_voltage\_stability.cpp}
        （弧长 CPF 翻越两母线 P--V 鼻点、比例方向计入组件负荷）、
        \srcpath{test\_newton\_krylov.cpp}（GMRES 块系统、求解器工厂、
        SolverWorkspace 复位）。覆盖深度仍属小算例 smoke 级，大系统回归
        仍缺；四者依旧由测试/GUI 直接驱动或作为选项层能力暴露。
        Newton--Krylov 回退默认关闭
        （\fld{enable\_newton\_krylov\_fallback\{false\}}，
        \srcpath{include/hacdcpf/power\_flow/globalization/robust\_nonlinear\_options.hpp:RobustNonlinearOptions}），
        仅在条件数代理越限时可选启用（src/power\_flow/newton\_solver.cpp:WorkspaceGuard）；
  \item \textbf{换流器物理不等式默认仅事后校核}：容量圆、交/直流电流限、调制比、
        占空比不嵌入 Newton 方程，求解后由
        \srcpath{src/api/hacdcpf.cpp:merge\_solver\_profiling} 发出
        \texttt{ACDC-PHYS-01..05}、:945 发出 \texttt{DCDC-PHYS-01} 告警；
        NEW 新增可选硬约束模式
        \fld{enforce\_converter\_physical\_limits\{false\}}
        （\srcpath{include/hacdcpf/power\_flow/power\_flow\_options.hpp:PowerFlowOptions}）：
        置 true 时带物理告警的 Newton 根被判不可行而非报收敛
        （\srcpath{src/api/hacdcpf.cpp:solve\_newton\_with\_lcc\_tap\_control}），并在结果 \fld{validity} 中
        如实标记各限值已强制（:solve\_newton\_with\_lcc\_tap\_control）；默认仍为仅告警模式。
        覆盖边界经 \fld{converter\_model\_scope} 声明（回归：
        \srcpath{test\_converter\_coordination.cpp} 的
        “Power-flow result declares its converter model scope”，:652）；
  \item \textbf{FDPF 注入已稀疏化}：\srcpath{src/power\_flow/fdpf\_solver.cpp}
        的注入计算改为稀疏复数形式 $S=V\cdot\overline{Y_{bus}V}$
        （\srcpath{compute\_power\_injections}，:19--35），不再转致密矩阵，
        原稠密 $O(n^2)$ 双重循环性能项随之消除；指定注入在
        \fld{has\_component\_loads} 时改用逐母线 ZIP 系数
        （\srcpath{compute\_specified\_injections}，:38--71）。回归：
        \srcpath{tests/test\_helm\_solver.cpp} 的 “FDPF sparse injection
        path remains numerically consistent” 与 “FDPF uses aggregated
        component loads”（与 NR 结果差 $<2\times10^{-4}$/$<3\times10^{-4}$）；
  \item \textbf{CPF 默认弧长延拓}：NEW 起 \srcpath{CpfSolver} 默认以弧长增广
        求解（\fld{enable\_arc\_length\{true\}}，
        \srcpath{include/hacdcpf/power\_flow/voltage\_stability.hpp:CpfOptions}），
        $\lambda$ 作为增广未知量，可翻越鼻点并保留下支
        \fld{lower\_branch\_steps\{2\}}（:CpfOptions）个点；置 false 时保留自然
        参数化旧模式（\srcpath{src/power\_flow/voltage\_stability.cpp:record（lambda）}）。
        \fld{CpfPoint} 新增 \fld{vdc}、\fld{va} 单位改为弧度（hpp:118--124）；
        鼻点摘要取迹线 $\lambda$ 最大点而非末点（cpp:495--500）。回归：
        \srcpath{tests/test\_voltage\_stability.cpp}（模型细节见
        第~\ref{sec:cpf}~章）；
  \item \textbf{\srcpath{workspace.hpp} 已落地为真实工作区}：
        \srcpath{include/hacdcpf/power\_flow/workspace.hpp} 新增
        \srcpath{prepare\_state}/\srcpath{prepare\_equations}/%
        \srcpath{reset\_iteration}（:35--42，实现
        \srcpath{src/power\_flow/workspace.cpp:prepare\_state}）；
        \srcpath{NewtonSolver} 复用每实例 \texttt{SolverWorkspace}，并以
        \fld{workspace\_in\_use\_} 递归守卫防止嵌套求解互踩
        （\srcpath{src/power\_flow/newton\_solver.cpp:build\_jacobian\_context}）。回归：
        \srcpath{tests/test\_newton\_krylov.cpp} 的 “SolverWorkspace resets
        values while retaining its layout”；
  \item \textbf{\srcpath{solver\_interface.hpp} 已接线}：NEW 新增
        \srcpath{src/power\_flow/solver\_interface.cpp}（137 行），实现
        \texttt{NewtonHybridSolver::solve} 等接口求解器与
        \srcpath{PowerFlowSolverFactory::create}
        （\srcpath{include/hacdcpf/power\_flow/solver\_factory.hpp:PowerFlowSolverFactory}）；
        工厂对 Newton/FDPF/DC/Adaptive/Continuation/NewtonKrylov 六方法
        均可产出可求解实例（回归 “Unified solver factory returns linked,
        callable solvers”，\srcpath{tests/test\_newton\_krylov.cpp:TEST\_CASE}）。
        生产 API 路径仍统一走 \srcpath{engine::NewtonSolver}
        （见第~\ref{sec:overview}~章求解器抽象层一节）。
\end{enumerate}

\subsection{如何运行验证}
configure/build/test 三段 preset（\srcpath{CMakePresets.json}，与根
\srcpath{AGENTS.md} 一致）：configure 侧 \texttt{macos-release}/
\texttt{linux-release}/\texttt{windows-msvc-release}/
\texttt{windows-vcpkg-release}/\texttt{full-dev}；test 侧同名 preset
仅前四个（\texttt{full-dev} 无 testPreset）。
\begin{quote}\footnotesize
\texttt{cmake --preset linux-release}\\
\texttt{cmake --build --preset linux-release}\\
\texttt{ctest --preset linux-release}
\end{quote}
常用定向运行（Catch2 标签与 ctest 正则）：
\begin{itemize}
  \item MATPOWER 回归：\texttt{ctest -R matpower\_cases}；慢速性能复算
        \texttt{test\_powerflow\_performance} 带 \texttt{[slow]} 标签，
        默认 Catch2 运行不含，需显式 \texttt{[slow]} 选择；
  \item 缩放/雅可比升级回归：\texttt{ctest -R}
        \srcpath{'scaling\_regression|unified\_pf\_upgrade'}；
  \item 非默认算法与工作区/工厂契约回归：\texttt{ctest -R}
        \srcpath{'helm\_solver|voltage\_stability|homotopy\_continuation|newton\_krylov'}；
  \item GUI 全链路：\texttt{ctest -R gui\_api\_e2e}（需
        \texttt{HACDCPF\_ENABLE\_ETAP=ON}，默认 ON，与 Python3）；
  \item OpenDSS 对照：configure 时加
        \texttt{-DHACDCPF\_ENABLE\_OPENDSS\_COMPARE=ON} 并提供 dss\_capi
        运行时；GridLAB-D 对照无需 configure 开关，运行时探测、缺失自动 skip；
  \item 能力合同矩阵：\texttt{python3}
        \srcpath{tools/validate\_case\_capabilities.py}
        \texttt{--server <run\_gui\_server 路径>}。
\end{itemize}
外部依赖（MATLAB、gridlabd、Julia、OpenDSS、chromium）缺失时对应测试自动
skip；判读报告时须区分“通过”与“跳过”，与主文档诚实口径一致。

\section{跨模块工业级技术评价基线}

本章规定本手册所述模块进入工程算例、批量研究或生产服务前必须共同满足的评价基线。
具体模型公式、变量、约束和专用算法以正文相应章节为准；本章不扩大源码已经实现的范围。

\subsection{输入、输出、边界条件与数据结构审计}
输入审计应逐项检查有限性、单位、上下界、枚举值和引用完整性。系统对象必须先按所需层级完成
校验；AC 与 DC 母线采用域限定映射，组件稳定 \texttt{index}、向量位置和临时图索引不得混用。
输出审计必须记录单位、索引空间、模型范围、收敛或可行状态、后端、容差、迭代/节点计数和告警。
空系统、空时域、孤岛、重复 ID、零容量、退化约束与不可用可选后端均属于边界测试集。

\subsection{工程假设与数值稳定性分析}
模型使用的平衡、线性化、稳态、独立性、概率分布、时间离散或网络等值假设必须与应用场景匹配。
数值评价至少包含变量缩放、绝对/相对残差、可行性违反、条件数或枢轴质量、步长/阻尼、迭代停滞
和随机种子复现性。若采用正则化、平滑、回退、松弛或近似，应同时报告触发条件、参数值、对主指标
的敏感性以及是否改变模型口径；不得用“已返回结果”替代收敛或有效性证明。

\subsection{复杂度与资源分析}
令 $n_x$、$n_c$、$n_t$、$n_s$、$|V|$、$|E|$ 分别表示变量、约束、时段、场景、节点和边数。
报告应分别给出建模、求解、后处理的时间与内存。图遍历通常为 $O(|V|+|E|)$；逐时段/逐场景
流程至少线性增长为 $O(n_t n_s)$ 次子问题调用；稠密线性代数最坏可达 $O(n_x^3)$，稀疏方法则应
报告结构非零数、填充率和因子分解峰值内存；MILP/NLP 不给出虚假的多项式最坏界，而报告节点数、
间隙、时限和实测扩展曲线。复杂度结论必须基于固定硬件、固定后端和可复现算例族。

\subsection{异常工况处理}
非法输入应在进入求解器前返回结构化错误；不可行、无界、不收敛、时限、内存不足和后端缺失必须
相互区分。部分结果只有在对应 \fld{ValidityFlags}、\fld{model\_scope}、
\fld{model\_limitations} 或等价字段允许时才可使用。异常路径不得静默丢弃 AC/DC 资产、制造平衡
节点、把 NaN 改写为零或把近似解标为全局最优；系统原始对象应保持可恢复和可重试。

\subsection{与其他模块的接口关系}
接口链路遵循“工程语义模型 $\to$ 校验 $\to$ 投影/装配 $\to$ 求解或分析 $\to$ 结果回投影”。
跨模块传递时保留稳定 ID、单位、符号、时间基准和场景身份；消费潮流、OPF、动态或可靠性结果前，
先检查其收敛、覆盖范围和有效性标志。HTTP/JSON/Python 层只做语义保持适配，不重新解释求解结果。

\subsection{测试与验证建议}
最低验证矩阵包括：手算小例、标准算例、边界/异常输入、随机性质测试、算法或后端交叉校核、
序列化往返、稳定 ID 回投影、AC/DC 同号母线、重复运行确定性和规模扩展测试。数值算法还应加入
残差独立重算、雅可比有限差分或自动微分校核；近似模型应与高保真模型比较并给出误差分布，而非
只报告单个平均值。回归报告必须列明构建配置、算例版本、容差、通过/失败/跳过数及未验证范围。

\subsection{工业级评估指标}
\begin{itemize}
  \item \textbf{正确性}：守恒残差、约束最大违反、解析/基准误差、结果回投影一致性；
  \item \textbf{鲁棒性}：收敛率、不可行识别率、异常分类准确率、扰动敏感性与随机复现率；
  \item \textbf{性能}：端到端时延、吞吐量、峰值内存、并行效率与规模增长斜率；
  \item \textbf{可追溯性}：输入模式版本、稳定 ID 覆盖率、模型范围/限制字段完整率、告警可定位率；
  \item \textbf{工程有效性}：与标准、外部引擎或现场数据的偏差，以及误判/漏判对业务指标的影响。
\end{itemize}
指标应预先给出验收阈值和失败处置；未达到阈值时只能标记为实验性或受限适用。

\subsection{适用范围、局限性与后续改进}
本手册只评价当前源码覆盖的模型、后端和数据结构，不对未建模物理过程、未安装求解器或未执行的
外部验证作保证。后续改进应按“缺口 $\to$ 可量化指标 $\to$ 源码/测试变更 $\to$ 独立验证”闭环，
优先消除会造成错误结论的语义缺口，其次提升数值鲁棒性与性能，最后扩展模型范围。


\end{document}
